\documentclass[11pt]{article}
\usepackage[utf8]{inputenc}
\usepackage[T1]{fontenc}
\usepackage{amsmath}

\usepackage{libertinus-otf}
\providecommand{\square}{\mdlgwhtsquare}
\ifdefined\mathscr\else\usepackage{mathrsfs}\fi
\usepackage[margin=1in]{geometry}
\usepackage{graphicx,booktabs,array,longtable,pdflscape,rotating}
\usepackage{enumitem}
\usepackage{microtype}
\usepackage{titlesec}
\usepackage[hidelinks]{hyperref}          % links are clickable but printed in black
\urlstyle{same}

\linespread{1.04}
\allowdisplaybreaks
\setlength{\parskip}{0.25em}
\setlength{\parindent}{1.2em}
\setlength{\emergencystretch}{3em}
\tolerance=1000
\clubpenalty=10000 \widowpenalty=10000 \displaywidowpenalty=10000
\renewcommand{\arraystretch}{1.15}
\setlist{itemsep=0.2em,topsep=0.4em,leftmargin=*}
\newcolumntype{P}[1]{>{\raggedright\arraybackslash}p{#1}}

\titleformat{\section}{\Large\bfseries}{\thesection.}{0.6em}{}
\titleformat{\subsection}{\large\bfseries}{\thesubsection}{0.6em}{}
\titleformat{\subsubsection}{\normalsize\bfseries\itshape}{\thesubsubsection}{0.6em}{}
\titlespacing*{\section}{0pt}{3ex plus 1ex minus .2ex}{1.5ex plus .2ex}
\titlespacing*{\subsection}{0pt}{2.5ex plus 1ex minus .2ex}{1ex plus .2ex}
\titlespacing*{\subsubsection}{0pt}{2ex plus 1ex minus .2ex}{0.8ex plus .2ex}

% Appendix headings carry the word "Appendix" and reset the subsection counter.
\newcommand{\appsection}[1]{\refstepcounter{section}%
  \section*{Appendix \thesection. #1}%
  \addcontentsline{toc}{section}{Appendix \thesection. #1}}

\pagestyle{plain}                         % page number centred at the foot, no running head

\hypersetup{pdftitle={Consistent forward-curve families with scheduled jumps}}

\usepackage[numbers]{natbib}
\begin{document}
\begin{center}
{\LARGE\bfseries Consistent forward-curve families with scheduled jumps\par}
\vspace{1.2em}
% {\large Author name\par}\vspace{0.6em}   % add author line here
{\normalsize\scshape Draft of 2026-09-24\par}
\end{center}
\vspace{0.4em}

\begin{abstract}
\noindent We study which finite-dimensional forward-curve families remain consistent with arbitrage-free dynamics when the curve jumps at scheduled policy dates, as the SOFR curve does at central-bank meetings. Under the drift and jump martingale conditions for term structures with predetermined jump dates, the meeting-date step family admits only the trivial dynamics, and its extension by ramps is consistent under deterministic step volatility exactly when each ramp slope is the squared step volatility and each level jump is centered Gaussian under the exponential tilt accumulated over earlier intervals; joint Gaussianity is not forced. Consistency alone imposes no restriction on meeting variances across meetings. In a Gaussian model with ideal payoffs, contracts settled after a date see the earlier variances only through two aggregates, and on the listed calendar of three-month SOFR futures options every meeting variance of 2026 and 2027 is identified in principle, though not to tick precision for distant meetings. For realizations in time to maturity that serve every schedule, one factor with a quasi-exponential shape and linearly recurrent meeting loadings is realized for every bounded Lipschitz scale, and with a constant scale the recurrence is also necessary. With deterministic or bounded predictable volatilities, a smooth exponential-polynomial block can be attached to a scheduled-step front end only by an independence splice or by indicator-times-exponential terms, up to a stated counting convention, and in the uncorrelated splice the front end frees only the level direction the two families share. We also correct a published sufficient condition for the number of correlated factors of a biexponential-polynomial family. The results are machine-checked in Lean, with the coverage in Appendix \hyperref[sec:A]{A}.
\end{abstract}
\vspace{0.6em}

\section{Introduction}\label{sec:1}

This paper proves four things about forward curves that jump at scheduled policy dates. The meeting-date step family admits only the trivial dynamics, and its extension by ramps is consistent exactly under the differentiated drift and tilted centered Gaussian level jumps (Theorems \hyperref[res:Theorem-1]{1} and \hyperref[res:Theorem-11]{11}). Consistency alone gives no cross-meeting restriction on event variances, and in a scalar Gaussian model ideal option panels identify them exactly when a rank condition holds, which the listed SOFR option calendar satisfies in principle (Remark 12, Propositions \hyperref[res:Proposition-19]{19}, \hyperref[res:Proposition-20]{20} and \hyperref[res:Proposition-24]{24}). Realizations in time to maturity that serve every schedule exist for one factor when the shape is quasi-exponential and the meeting loadings are linearly recurrent, and with a constant scale the recurrence is also necessary (Propositions \hyperref[res:Proposition-30]{30}, \hyperref[res:Proposition-31]{31} and \hyperref[res:Proposition-33]{33}). And attaching a smooth block to a scheduled-step front end is possible only through the two known constructions, up to a stated counting convention (Propositions \hyperref[res:Proposition-34]{34}, \hyperref[res:Proposition-35]{35}, \hyperref[res:Proposition-41]{41}, \hyperref[res:Proposition-43]{43}, \hyperref[res:Proposition-44]{44} and \hyperref[res:Proposition-47]{47}).

Which finite-dimensional forward-curve families are compatible with arbitrage-free dynamics when the curve jumps at scheduled policy dates, and what does compatibility force about the sizes of the jumps across meetings? The classical program (\cite{bjork1999interest}, \cite{filipovic1999note}, \cite{filipovic2000exponential}, \cite{filipovic2001consistency}) asks the first half of this question for parametric families such as Nelson--Siegel and the exponential-polynomial families, under Wiener-driven or stochastically continuous jump-driven Heath--Jarrow--Morton dynamics. Every jump in that literature is stochastically continuous. We found no paper that poses or solves the consistency or finite-dimensional-realization question for curve families with breakpoints at scheduled dates, or for term-structure models with predictable jump times; Section \hyperref[sec:4]{4} records what the literature does contain.

The motivation is the front end of the SOFR curve. The policy rate moves at the scheduled meetings of the central bank, so the forward curve has steps at known calendar dates; desks build it as a meeting-date step function spliced onto a smooth long-end family, and \cite{heitfield2019inferring} is the static version of that construction fitted to futures. The constructive precedents are the SOFR models of \cite{gellert2021short} and \cite{brace2024sofr}. The results are meant to be usable in three ways: as a consistent front-end curve for carry and roll-down (Sections \hyperref[sec:3.1]{3.1} to 3.9), as a test of whether consistency restricts option-implied event variances across meetings (Sections \hyperref[sec:3.10]{3.10} to 3.23), and as a classification of how the scheduled front end can carry stochastic volatility and be attached to a smooth long-end block (Sections \hyperref[sec:3.24]{3.24} to 3.41).

\emph{The front end.} Theorem \hyperref[res:Theorem-1]{1} says that the pure maturity-step family supports only the trivial dynamics: under the two no-arbitrage conditions of Section \hyperref[sec:2]{2} its drift, volatility and curve jumps vanish, so \(f(t,T)=f(0,T)\) (Corollary \hyperref[res:Corollary-2]{2}), the scheduled-jump analogue of the negative result for Nelson--Siegel (\cite{filipovic1999note}). Its diffusive leg is not new in mechanism: \cite{gellert2021short} §3.1.5 and §5 state informally that piecewise-constant short-rate paths constrain the initial curve, and the ramp drift behind the argument is their eq. 28 and \cite{brace2024sofr} eq. 18; Theorem \hyperref[res:Theorem-1]{1} gives the statement for the curve family with a proof and adds the jump leg, which neither paper treats. Theorem \hyperref[res:Theorem-11]{11} classifies the extension by ramps, the step-plus-ramp family (2.5), under deterministic step volatility: consistency is equivalent to the differentiated HJM drift (Proposition \hyperref[res:Proposition-8]{8}) and to tilted centered Gaussian level jumps (Proposition \hyperref[res:Proposition-9]{9}), and Proposition \hyperref[res:Proposition-10]{10} is an explicit multi-date instance. Sufficiency is the content of the models of \cite{gellert2021short}, \cite{brace2024sofr}, \cite{fontana2022term}, \cite{kellerressel2018affine} and \cite{kim2014jumps}; the necessity directions, including the equivalence of the integrated and differentiated drift conditions without a differentiability hypothesis, and the non-Gaussian example of Proposition \hyperref[res:Proposition-9]{9}(d), are new. Proposition \hyperref[res:Proposition-48]{48} shows that on the interval after a scheduled date the compensating shape of the curve jump is the derivative of the conditional cumulant generating function of the level jump, so that a straight ramp forces a Gaussian level jump, and proves the pre-meeting forward-rate formula (5.1). Section \hyperref[sec:3.9]{3.9} shows that the forced ramps are two orders of magnitude below the diffusive convexity and below quote precision, by a factor of about fifteen to thirty per meeting and less in aggregate over a year of meetings.

\emph{Event variances.} Remark 12 separates state dimension from parameter dimension: scalar-state models fill the nonnegative orthant of meeting variances as their parameters vary, so consistency alone gives no cross-meeting restriction. Propositions \hyperref[res:Proposition-13]{13}, \hyperref[res:Proposition-14]{14}, \hyperref[res:Proposition-16]{16} and \hyperref[res:Proposition-26]{26} give the restrictions of the fixed Gaussian and stochastic-volatility specifications of \cite{brace2024sofr}, and Proposition \hyperref[res:Proposition-25]{25} those that survive every recalibration of the scales. In the scalar Gaussian model, Propositions \hyperref[res:Proposition-15]{15}, \hyperref[res:Proposition-17]{17}, \hyperref[res:Proposition-19]{19}, \hyperref[res:Proposition-20]{20} and \hyperref[res:Proposition-23]{23}, Corollaries \hyperref[res:Corollary-21]{21} and \hyperref[res:Corollary-22]{22} and Lemma \hyperref[res:Lemma-18]{18} determine which meeting variances ideal bond calls, compounded-rate futures and European, American and futures-style calls identify; Proposition \hyperref[res:Proposition-20]{20} isolates the mechanism, two aggregates of the earlier variances, and Proposition \hyperref[res:Proposition-24]{24} evaluates the rank criterion on the listed SOFR option calendar. Proposition \hyperref[res:Proposition-49]{49} adds a diffusion between meetings and an arbitrary initial curve, which the surfaces do not see: the diffusion variance is read off the expiry gaps without a meeting, so the listed calendar identifies a constant diffusion together with every meeting variance and at most eight diffusion cells, and the diffusion's error, entering each meeting in proportion to its gap, degrades practical identification. Proposition \hyperref[res:Proposition-50]{50} lets the diffusion depend on maturity through a known shape: the rank criterion keeps its form with window weights, and on the listed calendar the answers survive for a loading that starts at the next meeting and, for exponential decay, at every rate except one critical value.

\emph{Stochastic volatility and realizations.} Proposition \hyperref[res:Proposition-27]{27} gives a bounded three-exponent family with one scheduled breakpoint whose variance state enters the curve. Proposition \hyperref[res:Proposition-28]{28} gives the finite representation of the curve under separable volatilities with meeting-dependent loadings, and Proposition \hyperref[res:Proposition-29]{29} a time-inhomogeneous Markov realization with a Lipschitz scale state; both extend known calendar-maturity results (\cite{bjork2001existence} §3.3 and Proposition 6.2, \cite{fontana2022term} Example 3.10) to meeting restarts. For realizations in time to maturity that serve every schedule, Propositions \hyperref[res:Proposition-30]{30} and \hyperref[res:Proposition-33]{33} show that the ordinal meeting loadings must satisfy a linear recurrence, and Proposition \hyperref[res:Proposition-31]{31} gives the converse for one factor and every bounded Lipschitz scale, adding meeting restarts to the quasi-exponential answer of \cite{bjork2001existence} and \cite{bjork2004finite}. On a finite horizon, Proposition \hyperref[res:Proposition-51]{51} bounds the forward-curve and bond-price errors of a recurrent approximation, each model with its own drift, by the loading error, with constants free of the number and spacing of meetings, and Proposition \hyperref[res:Proposition-31]{31} realizes the approximating model.

\emph{The splice.} Propositions \hyperref[res:Proposition-34]{34}, \hyperref[res:Proposition-35]{35}, \hyperref[res:Proposition-41]{41}, \hyperref[res:Proposition-43]{43} and \hyperref[res:Proposition-44]{44} show that a step factor with maturity steps can be attached to a fixed exponential-polynomial block only through the independence splice of \cite{gellert2021short} §3 or the indicator-times-exponential terms of \cite{brace2024sofr} eq. 36: with deterministic volatilities, with bounded predictable scales, with a block volatility driven by the block's own state, and with several factors, where the condition is orthogonality of the block's noise to the aggregate maturity-step noise. Propositions \hyperref[res:Proposition-37]{37}, \hyperref[res:Proposition-38]{38}, \hyperref[res:Proposition-40]{40}, \hyperref[res:Proposition-45]{45}, \hyperref[res:Proposition-46]{46} and \hyperref[res:Proposition-47]{47} treat the block's own restrictions: in the uncorrelated splice the front end frees only the level direction that the two families share, and in the correlated splice with one step factor only that level may correlate with the step noise. Proposition \hyperref[res:Proposition-52]{52} combines several step factors with a full polynomial part and states the answer without a counting convention: only the level may covary with the aggregate maturity-step noise, the current front-end noise acts as extra level noise, and consistency is the block-alone condition for that effective level. Proposition \hyperref[res:Proposition-53]{53} extends it, without a stopping argument, to locally square-integrable noises, which admit Heston-type scales. Two of these statements rest on counting conventions, stated where they are used: a block uncorrelated with the aggregate step noise counts as the independence splice, and a correlated level counts with the front end. Lemma \hyperref[res:Lemma-36]{36} and Proposition \hyperref[res:Proposition-39]{39} settle, at a parameter point, the number of correlated factors of a biexponential-polynomial block and show that a published sufficient condition for it (\cite{sharef2004conditions} §4.2) is false as stated.

\emph{What is conditional.} Assumptions 2.1 and 2.2 are hypotheses throughout. Every result that uses stochastic integration, Itô's formula, solutions of stochastic differential equations or their Markov property takes these from cited sources rather than proving them, and several take solutions of stated equations as given; each such result is marked conditional, and its hypotheses are stated with it. The corresponding hypothesis structures of the formalization have only degenerate or vacuous instances, which show that their fields are consistent but give no nontrivial model of them. Section \hyperref[sec:5.1]{5.1} lists what the results rest on, and Section \hyperref[sec:5]{5} as a whole what is not done.

Every result in Section \hyperref[sec:3]{3} is proved in the text, conditional on the cited inputs stated with it. Appendix \hyperref[sec:A]{A} lists the exact Lean coverage: Theorem \hyperref[res:Theorem-1]{1} to Remark 12 are formalized, except Corollary \hyperref[res:Corollary-2]{2}; Propositions \hyperref[res:Proposition-13]{13} and \hyperref[res:Proposition-14]{14} have conditional formalizations, Proposition \hyperref[res:Proposition-15]{15} is fully formalized in its stated ideal model, Proposition \hyperref[res:Proposition-16]{16} is formalized using the cited stochastic-calculus fields, and Proposition \hyperref[res:Proposition-26]{26} also uses the cited predictability field; both assume solutions of their variance equations. Proposition \hyperref[res:Proposition-27]{27} is formalized conditional on the two supplied calculus structures, Proposition \hyperref[res:Proposition-28]{28} conditional on the supplied calculus, except the word "semimartingale", Proposition \hyperref[res:Proposition-29]{29} conditional on the supplied calculus, strong-existence, predictability and Markov inputs, except its descriptive paragraph, Proposition \hyperref[res:Proposition-30]{30} conditional on the supplied calculus, Proposition \hyperref[res:Proposition-31]{31} conditional on the supplied calculus and predictability input, Lemma \hyperref[res:Lemma-32]{32} and Proposition \hyperref[res:Proposition-33]{33} conditional on the supplied calculus, Propositions \hyperref[res:Proposition-34]{34} and \hyperref[res:Proposition-35]{35} conditional on the supplied calculus, Lemma \hyperref[res:Lemma-36]{36} without conditions, Propositions \hyperref[res:Proposition-37]{37} and \hyperref[res:Proposition-38]{38} conditional on the supplied calculus, Propositions \hyperref[res:Proposition-39]{39} and \hyperref[res:Proposition-40]{40} without conditions, Proposition \hyperref[res:Proposition-41]{41}, Lemma \hyperref[res:Lemma-42]{42} and Propositions \hyperref[res:Proposition-43]{43} and \hyperref[res:Proposition-44]{44} conditional on the supplied calculus, Proposition \hyperref[res:Proposition-45]{45} as a statement at a parameter point, and Proposition \hyperref[res:Proposition-46]{46} conditional on the supplied calculus and, for (b), on the cited results of \cite{filipovic2000exponential}, Proposition \hyperref[res:Proposition-47]{47} conditional on the supplied calculus, Proposition \hyperref[res:Proposition-48]{48} without conditions, Proposition \hyperref[res:Proposition-49]{49} with (a) and (b) conditional on a stated Gaussian law of the Wiener integrals, satisfied in the formalization only without diffusion, and (c) to (e) without conditions, Proposition \hyperref[res:Proposition-50]{50} likewise, with (a) and (b) conditional and (c) and (d) without conditions, Proposition \hyperref[res:Proposition-51]{51} with its deterministic estimates without conditions, (a) under the supplied calculus, and (b) and (c) also under a Brownian driver, Proposition \hyperref[res:Proposition-52]{52} conditional on Assumption 2.1 for the representation, stated as an explicit hypothesis, and Proposition \hyperref[res:Proposition-53]{53} likewise, with its integrability class entering only the converse. Lemma \hyperref[res:Lemma-18]{18}, Propositions \hyperref[res:Proposition-17]{17}, \hyperref[res:Proposition-19]{19}, \hyperref[res:Proposition-20]{20} and \hyperref[res:Proposition-23]{23} to 25 and Corollaries \hyperref[res:Corollary-21]{21} and \hyperref[res:Corollary-22]{22} are fully formalized for their stated ideal prices, exercise problems, calendar and cones. Propositions \hyperref[res:Proposition-27]{27}, \hyperref[res:Proposition-29]{29} and \hyperref[res:Proposition-46]{46}(b) rest on hypothesis structures whose instances are vacuous, with separate checks that their premises and conclusions are inhabited; the other calculus structures have degenerate instances, and the structure for the standing assumptions a nonzero one; no calculus structure is given a nontrivial model of its fields.

Section \hyperref[sec:2]{2} fixes the setting and the standing assumptions, Section \hyperref[sec:3]{3} states and proves the results, Section \hyperref[sec:4]{4} surveys related work, Section \hyperref[sec:5]{5} states what is not done, and Appendix \hyperref[sec:A]{A} records the machine verification.

\section{Setting and standing assumptions}\label{sec:2}

Throughout, \((\Omega, \mathcal F, (\mathcal F_t)_{t \ge 0}, Q)\) is a filtered probability space, \(Q\) the pricing measure, and the scheduled dates are

\begin{equation*}
0 = T_0 < T_1 < \dots < T_N, \qquad I_k := [T_k, T_{k+1}) \ (k < N), \qquad I_N := [T_N, \infty). \tag{2.1}\label{eq:2.1}
\end{equation*}

The intervals \(I_0, \dots, I_N\) partition \([0, \infty)\); for \(t \ge 0\), \(j(t)\) denotes the index with \(t \in I_{j(t)}\). The forward rate \(f(t, T)\), \(T \ge t\), has the short rate \(r_t = f(t, t)\) and the bank account \(B_t = \exp(\int_0^t r_s\,ds)\), and is assumed to satisfy, for \(T \ge t\),

\begin{equation*}
f(t, T) = f(0, T) + \int_0^t \alpha(s, T)\,ds + \int_0^t \langle \sigma(s, T), dW_s \rangle + \sum_{n : T_n \le t} \xi_n(T), \tag{2.2}\label{eq:2.2}
\end{equation*}

with \(W\) a \(d\)-dimensional Brownian motion, \(\alpha\) real, \(\sigma\) \(\mathbb R^d\)-valued, and \(\xi_n(\cdot)\) an \(\mathcal F_{T_n}\)-measurable random function, the jump of the whole curve at the scheduled date \(T_n\). Between scheduled dates the curve is continuous in \(t\); at \(T_n\) it jumps by \(\xi_n\). Because the discontinuities sit at fixed calendar dates, families are parametrized in \((t, T)\) rather than in time to maturity alone, and are time-inhomogeneous; the same parametrization appears in \cite{kim2014jumps} Proposition 1 and \cite{brace2024sofr} eqs. 2 to 5.

\textbf{Regularity.} The coefficients vanish before the current date, \(\alpha(t, T) = 0\) and \(\sigma(t, T) = 0\) for \(T < t\) and \(\xi_n(u) = 0\) for \(u < T_n\); \(\alpha\) and \(\sigma\) are progressively measurable in \((t, T)\), \(\xi_n\) is \(\mathcal F_{T_n} \otimes \mathcal B(\mathbb R_+)\)-measurable, and all three are locally integrable in the maturity variable almost surely. These are the standing assumptions of \cite{fontana2022term} (their Assumption 3.3) in the special case used here.

The two no-arbitrage conditions are taken from \cite{fontana2022term} Theorem 3.5, which characterizes the risk-neutral measures of a term structure with an overnight numéraire, roll-over dates at which the numéraire may jump, and expected jump dates at which the curve may jump. We use the case with no roll-over dates, so the term-structure measure has no atoms, the bank account is absolutely continuous, the policy rate jumps but nothing pays a lump sum at \(T_n\); with no random measure, so no totally inaccessible jumps; with finitely many fixed jump dates; and with the overnight rate equal to the short end, \(r_t = f(t, t)\). This is the case the source itself exercises in its Example 3.10. The characterization is an equivalence; only its two conditions are used, as hypotheses, and bond prices, the numéraire and local martingales do not appear below.

\textbf{Assumption 2.1 (drift condition between scheduled dates, integrated form).} For every \(T > 0\), for \(Q \otimes dt\)-almost every \((\omega, t) \in \Omega \times [0, T]\),

\begin{equation*}
\int_t^T \alpha(t, u)\,du = \tfrac12 \Bigl| \int_t^T \sigma(t, u)\,du \Bigr|^2, \tag{2.3}\label{eq:2.3}
\end{equation*}

with \(|\cdot|\) the Euclidean norm on \(\mathbb R^d\) and both integrals Lebesgue integrals in \(u\). Source: \cite{fontana2022term} Theorem 3.5(ii), whose integrated coefficients reduce to the two integrals in (2.3) once the roll-over set is empty and the random measure is absent; the original drift condition is \cite{heath1992bond}, and an equivalent statement in affine-semimartingale form, allowing a deterministic clock with atoms, is \cite{kellerressel2018affine} Proposition 6.11 eq. 74. The pointwise form \(\alpha(t, T) = \langle \sigma(t, T), \int_t^T \sigma(t, u)\,du \rangle\) is the \(T\)-derivative of (2.3); the two agree for almost every \(T\) at each fixed \((\omega, t)\), and exactly when \(\alpha(t, \cdot)\) and \(\sigma(t, \cdot)\) are piecewise constant on the same finite partition. Only the integrated form is assumed.

\textbf{Assumption 2.2 (jump martingale condition at scheduled dates).} For every \(n = 1, \dots, N\) and every real \(T \ge T_n\), \(Q\)-almost surely,

\begin{equation*}
E\Bigl[ \exp\Bigl( -\int_{T_n}^T \xi_n(u)\,du \Bigr) \Bigm| \mathcal F_{T_n-} \Bigr] = 1, \tag{2.4}\label{eq:2.4}
\end{equation*}

where \(\mathcal F_{T_n-}\) is the sigma-algebra generated by the \(\mathcal F_t\), \(t < T_n\), and the conditional expectation of the nonnegative random variable is taken in \([0, \infty]\); (2.4) forces that variable to be integrable. \(T_0 = 0\) is an interval endpoint, not a jump date. Source: \cite{fontana2022term} Theorem 3.5(iv) together with the sigma-integrability of their (3.6), simplified by their Assumption 3.6; with no atoms and the overnight rate equal to \(f(t, t)\), the display of their Remark 3.9 reads exactly (2.4) with \(\xi_n = \Delta V(T_n, \cdot)\) in their notation, and their Example 3.10 deduces its compensating drift from this condition. Its content is that the size and direction of the jump at a scheduled date cannot be predicted, as the source says after its Theorem 3.5.

Assumptions 2.1 and 2.2 are hypotheses, not results of this paper. Each result below says in words which of them it uses; the lemmas of Section \hyperref[sec:3]{3} use neither, and Theorem \hyperref[res:Theorem-1]{1} uses both.

\textbf{The step-plus-ramp family.} The curves \(S^+\) consist of a step and a ramp attached to each scheduled date: at time \(t\in I_j\), for every maturity \(T\in I_k\) with \(k\ge j\),

\begin{equation*}
f(t,T)=\sum_{m\le k}\bigl[M_m(t)+C_m(t)\,(T-T_m)\bigr], \tag{2.5}\label{eq:2.5}
\end{equation*}

with \(M_m\) the level and \(C_m\) the ramp coefficient attached to \(T_m\), so that on the current interval \(f(t,T)=r_t+\sum_{m\le j}C_m(t)(T-t)\). The case \(C_m\equiv0\) is the maturity-step family \(S\) of Section \hyperref[sec:3.1]{3.1}. In the bookkeeping of \cite{gellert2021short} eq. 28 the ramps are indexed by pairs of scheduled dates, with pairwise coefficients \(C_{nm}\) entered with a minus sign; only their sums attached to a date are identifiable (Section \hyperref[sec:3.8]{3.8}). We call \(S^+\) the extension of \(S\) by ramps; no minimality is claimed.

\textbf{Scope.} Every result below is a statement about the model just described, with prices generated exactly by it. Sections \hyperref[sec:3.10]{3.10} to 3.23 use two fixed specifications: a scalar Gaussian model with zero initial curve whose only parameters are the meeting variances, and the Gaussian and stochastic-volatility specifications of \cite{brace2024sofr} with loadings, decay and scales fixed. The payoffs priced there are ideal: European bond calls, backward-looking futures with continuously compounded or exact discrete bond-implied accrual, and cash-settled European or American calls on those futures with stated exercise windows. Section \hyperref[sec:3.21]{3.21} takes the listed contracts' expiry calendar and tick size from the exchange rules as inputs to that ideal model, not their payoffs. Sections \hyperref[sec:3.43]{3.43} and \hyperref[sec:3.44]{3.44} add to the scalar Gaussian model a Gaussian diffusion with piecewise-constant variance, parallel or with a known maturity shape, and an arbitrary bounded initial curve. Not modelled: listed SOFR contracts with their settlement, margin and exercise-into-futures rules; accrual conventions other than those stated; quoted or noisy prices; an unknown or fitted initial curve outside Sections \hyperref[sec:3.43]{3.43} and \hyperref[sec:3.44]{3.44}; additional diffusion or stochastic volatility where none is stated; and calibration to market data. An identification result therefore says which meeting variances the stated ideal observations determine within the stated specification, and a non-identification pair says that they do not determine them there; neither is a statement about listed instruments, and an infeasible panel is a rejection of the specification, not an arbitrage certificate. Section \hyperref[sec:5.3]{5.3} states what a transfer to listed instruments would require.

\section{Main results}\label{sec:3}

\subsection{The maturity-step family}\label{sec:3.1}

\textbf{Definition 3.1.} Coefficients \((\alpha, \sigma, (\xi_n))\) as in Section \hyperref[sec:2]{2} belong to the \emph{maturity-step family} \(S\) if there are functions \(\mu_k(t, \omega) \in \mathbb R\) and \(s_k(t, \omega) \in \mathbb R^d\) (\(k = 0, \dots, N\)) and random variables \(X_{n,k}\) (\(1 \le n \le k \le N\)) such that, for every \(\omega\),

\begin{equation*}
\alpha(t, T, \omega) = \mu_k(t, \omega) \quad \text{and} \quad \sigma(t, T, \omega) = s_k(t, \omega) \qquad \text{whenever } t \le T \text{ and } T \in I_k, \tag{3.1}\label{eq:3.1}
\end{equation*}

\begin{equation*}
\xi_n(u, \omega) = X_{n,k}(\omega) \qquad \text{whenever } u \in I_k \text{ and } n \le k. \tag{3.2}\label{eq:3.2}
\end{equation*}

That is, at each time the drift and the volatility, and at each scheduled date the curve jump, are step functions of maturity on the scheduled intervals. No measurability or integrability of \(\mu_k\), \(s_k\), \(X_{n,k}\) is assumed beyond what (3.1), (3.2) and the regularity of Section \hyperref[sec:2]{2} imply. At the level of the curve, \(S\) is the family \(f(t, T) = b_k(t)\) for \(T \in I_k\), \(T \ge t\), with each level an Itô process \(db_k = \mu_k\,dt + \langle s_k, dW \rangle\) between scheduled dates and a jump \(X_{n,k}\) of \(b_k\) at \(T_n\); (3.1) and (3.2) are the coefficient form of that description, and we take them as the definition. Desks build the front-end SOFR curve as such a step function.

\phantomsection\label{res:Theorem-1}
\textbf{Theorem 1.} Let the coefficients belong to \(S\) and satisfy the regularity conditions of Section \hyperref[sec:2]{2} together with Assumptions 2.1 and 2.2. Then:

(i) (jump leg) for every \(1 \le n \le k \le N\), \(X_{n,k} = 0\) \(Q\)-almost surely; hence for every \(n\) there is an almost-sure event on which \(\xi_n(u) = 0\) for all \(u \in \mathbb R\);

(ii) (drift and volatility leg) for \(Q \otimes dt\)-almost every \((\omega, t) \in \Omega \times [0, \infty)\), \(\mu_k(t, \omega) = 0\) and \(s_k(t, \omega) = 0\) for every \(k \ge j(t)\); hence \(\alpha(t, T, \omega) = 0\) and \(\sigma(t, T, \omega) = 0\) for all \(T \in \mathbb R\).

The theorem uses both standing assumptions: (i) uses Assumption 2.2 only and (ii) uses Assumption 2.1 only. The values \(\mu_k(t)\), \(s_k(t)\) for \(k < j(t)\) concern maturities before \(t\) and are not constrained.

\phantomsection\label{res:Corollary-2}
\textbf{Corollary 2.} If in addition \(W\) is a \(d\)-dimensional Brownian motion for \((\mathcal F_t)\) and \(f\) is given by (2.2), then for every \(T \ge 0\) and \(t \in [0, T]\), \(f(t, T) = f(0, T)\) \(Q\)-almost surely: the maturity-step family supports no randomness at all.

\emph{Proof of Corollary \hyperref[res:Corollary-2]{2} from Theorem \hyperref[res:Theorem-1]{1}.} Fix \(T \ge 0\) and \(t \le T\). By (ii) and Fubini's theorem, for \(Q\)-almost every \(\omega\), \(\alpha(s, T, \omega) = 0\) for \(ds\)-almost every \(s \in [0, t]\), so the drift integral in (2.2) vanishes almost surely; the integrand \(s \mapsto \sigma(s, T)\) vanishes \(Q \otimes ds\)-almost everywhere on \([0, t]\), so its stochastic integral vanishes almost surely; and by (i) the jump sum vanishes almost surely.\hfill$\square$

The corollary uses two standard facts, Fubini's theorem and the vanishing of the Itô integral of an integrand that vanishes almost everywhere; it is the one statement of this section not among the machine-checked ones (Appendix \hyperref[sec:A]{A}).

\subsection{Four lemmas}\label{sec:3.2}

The proof of Theorem \hyperref[res:Theorem-1]{1} rests on four lemmas about a fixed step curve: a closed form for the integral of a step function on the partition (Lemma \hyperref[res:Lemma-3]{3}), a two-point moment lemma (Lemma \hyperref[res:Lemma-4]{4}), the jump leg for one scheduled date (Lemma \hyperref[res:Lemma-5]{5}) and the drift leg for one time (Lemma \hyperref[res:Lemma-6]{6}). None of them uses Assumption 2.1 or 2.2.

\phantomsection\label{res:Lemma-3}
\textbf{Lemma 3 (integral of a step function).} Let the dates be as in (2.1), let \(s_0, \dots, s_N \in \mathbb R^d\) be fixed vectors, and let \(\sigma : [0, \infty) \to \mathbb R^d\) be the step function \(\sigma(u) := s_k\) for \(u \in I_k\). For \(0 \le t \le T < \infty\) with \(t \in I_j\) and \(T \in I_k\) one has \(j \le k\) and

\begin{equation*}
\int_t^T \sigma(u)\,du = c_{jk}(t) + s_k \bigl( T - \max(T_k, t) \bigr), \tag{3.3}\label{eq:3.3}
\end{equation*}

where

\begin{equation*}
c_{jj}(t) := 0, \qquad c_{jk}(t) := s_j (T_{j+1} - t) + \sum_{m=j+1}^{k-1} s_m (T_{m+1} - T_m) \quad (k > j); \tag{3.4}\label{eq:3.4}
\end{equation*}

in particular \(c_{jk}(t)\) does not depend on \(T\). Consequently \(T \mapsto \langle s_k, \int_t^T \sigma(u)\,du \rangle\) is affine on \(I_k \cap [t, \infty) = [\max(T_k, t), T_{k+1})\) with slope \(|s_k|^2\):

\begin{equation*}
\Bigl\langle s_k, \int_t^T \sigma(u)\,du \Bigr\rangle = a_{jk}(t) + |s_k|^2\, T, \qquad a_{jk}(t) := \langle s_k, c_{jk}(t) \rangle - |s_k|^2 \max(T_k, t). \tag{3.5}\label{eq:3.5}
\end{equation*}

\emph{Proof.} Split \([t, T]\) at the scheduled dates it contains, \(t < T_{j+1} \le \dots \le T_k \le T\). On each piece \(\sigma\) is constant except at one endpoint, a null set, so each piece integrates to the constant times the length; adding the pieces gives (3.3), and the \(T\)-dependence sits only in the last piece \(s_k (T - T_k)\). When \(k = j\) there is a single piece and (3.3) reads \(s_j (T - t)\). Linearity of the inner product turns (3.3) into (3.5).\hfill$\square$

The identity is displayed but not stated in \cite{gellert2021short}, whose eq. 28 contains the ramp terms and whose appendix eqs. 65 to 67 evaluate the integral of a step indicator case by case; \cite{fontana2018term} Theorem 3.7 and Remark 3.8 give the general drift condition without specializing to step volatilities.

\phantomsection\label{res:Lemma-4}
\textbf{Lemma 4 (two exponential moments).} Let \(X\) be a real random variable on a probability space \((\Omega, \mathcal F, P)\) and let \(\lambda \in \mathbb R\), \(\lambda \ne 0\). If

\begin{equation*}
E[\exp(-\lambda X)] \text{ is finite and equals } 1, \tag{3.6}\label{eq:3.6}
\end{equation*}

\begin{equation*}
E[\exp(-\lambda X / 2)] \text{ is finite and equals } 1, \tag{3.7}\label{eq:3.7}
\end{equation*}

then \(X = 0\) \(P\)-almost surely.

\emph{Proof.} Put \(Y := \exp(-\lambda X / 2)\), so \(Y > 0\) and \(Y^2 = \exp(-\lambda X)\). By (3.6), \(E[Y^2] = 1\), so \(Y\) is square-integrable; by (3.7), \(E[Y] = 1\). Then \(E[(Y - 1)^2] = E[Y^2] - (E[Y])^2 = 0\), so \(Y = 1\) almost surely, and injectivity of \(\exp\) with \(\lambda \ne 0\) gives \(X = 0\) almost surely.\hfill$\square$

This is the equality case of the Cauchy--Schwarz inequality \((E[Y])^2 \le E[Y^2]\), or equivalently the strict convexity of the cumulant generating function of \(X\), which vanishes at the three collinear points \(0\), \(-\lambda/2\), \(-\lambda\). Both conditions are needed: a Gaussian \(X\) with mean \(\lambda \sigma^2 / 2\) and variance \(\sigma^2\) satisfies (3.6) alone with \(X \ne 0\). The finiteness half of (3.7) follows from (3.6).

\phantomsection\label{res:Lemma-5}
\textbf{Lemma 5 (a maturity-step jump at a scheduled date vanishes).} Let \((\Omega, \mathcal F, P)\) be a probability space, let the dates be as in (2.1) with \(N \ge 1\), fix \(n \in \{1, \dots, N\}\) and real random variables \(X_n, \dots, X_N\), and let \(\xi\) be the jump that is a step in maturity with these values,

\begin{equation*}
\xi(u, \omega) := X_k(\omega) \ \text{for } u \in I_k,\ k \ge n, \qquad \xi(u, \omega) := 0 \ \text{for } u < T_n. \tag{3.8}\label{eq:3.8}
\end{equation*}

For real \(T \ge T_n\) put \(Y_T := \exp\bigl(-\int_{T_n}^T \xi(u, \cdot)\,du\bigr)\), a nonnegative random variable. If

\begin{equation*}
Y_T \text{ is integrable and } E[Y_T] = 1 \quad \text{for every real } T \ge T_n, \tag{3.9}\label{eq:3.9}
\end{equation*}

then \(X_k = 0\) \(P\)-almost surely for every \(k = n, \dots, N\). The conclusion already follows when (3.9) is assumed only at the \(2(N - n + 1)\) maturities

\begin{equation*}
T_k + h_k \ \text{ and } \ T_k + h_k / 2, \qquad k = n, \dots, N, \qquad h_k := (T_{k+1} - T_k)/2 \ (k < N), \quad h_N := 1. \tag{3.10}\label{eq:3.10}
\end{equation*}

Conditional form: if \(\mathcal G \subseteq \mathcal F\) is a sub-sigma-algebra and

\begin{equation*}
E[Y_T \mid \mathcal G] = 1 \quad P\text{-almost surely, for every real } T \ge T_n, \tag{3.11}\label{eq:3.11}
\end{equation*}

with the conditional expectation of the nonnegative \(Y_T\) taken in \([0, \infty]\), then (3.9) holds and the same conclusion follows.

\emph{Proof.} Fix \(\omega\). Then \(\xi(\cdot, \omega)\) is the step function of Lemma \hyperref[res:Lemma-3]{3} with \(d = 1\), values \(X_k(\omega)\) for \(k \ge n\) and \(0\) before \(T_n\), so Lemma \hyperref[res:Lemma-3]{3} with \(t = T_n\) gives, for \(T \in I_k\) with \(k \ge n\), \(\int_{T_n}^T \xi(u, \omega)\,du = c_k(\omega) + X_k(\omega)(T - T_k)\), where \(c_n = 0\) and \(c_k\) is a fixed linear combination of \(X_n, \dots, X_{k-1}\) with the interval lengths as coefficients. Induct on \(k\) from \(n\) to \(N\). If \(X_m = 0\) almost surely for \(n \le m < k\), then \(c_k = 0\) almost surely, and (3.9) at the two maturities (3.10), which lie in \(I_k\), reads \(E[\exp(-h_k X_k)] = 1\) and \(E[\exp(-h_k X_k / 2)] = 1\), both finite. Lemma \hyperref[res:Lemma-4]{4} with \(\lambda = h_k > 0\) gives \(X_k = 0\) almost surely. For the conditional form, \(Y_T \ge 0\) and the tower property in \([0, \infty]\) give \(E[Y_T] = E[E[Y_T \mid \mathcal G]] = 1 < \infty\), which is (3.9).\hfill$\square$

Both maturities in each interval are needed, and no interval can be skipped: a Gaussian \(X_k\) with mean \(h_k s^2 / 2\) and variance \(s^2\) satisfies (3.9) at \(T_k + h_k\) and fails it at \(T_k + h_k / 2\), and if \(X_{k_0}\) is Gaussian with the mean tuned to the length of \(I_{k_0}\) while all other \(X_k\) vanish, (3.9) holds at every maturity outside \(I_{k_0}\) with \(X_{k_0} \ne 0\); any two distinct maturities in \(I_k\) would serve in place of (3.10). With \(\mathcal G = \mathcal F_{T_n-}\) and \(\xi = \xi_n\), (3.11) is Assumption 2.2 for the jump at \(T_n\), specialized to a jump that is a step in maturity; that specialization is the definition of \(S\), not a consequence of the assumption. The finite form (3.10) is what lets Theorem \hyperref[res:Theorem-1]{1} pass from "for every \(T\), almost surely" in (2.4) to a single almost-sure event.

\cite{fontana2022term} Example 3.10 is the boundary case: a Gaussian jump of the curve at a scheduled date is consistent only when accompanied by a ramp in maturity whose slope is the jump variance, so a jump that is a pure step in maturity is forced to vanish; the direct computation in the Gaussian case gives the same. \cite{kim2014jumps} Proposition 1 and \cite{kellerressel2018affine} Examples 6.15 and 6.16 are Gaussian scheduled jumps of an affine state with compensating ramps, and \cite{backwell2021expected} computes its curve by pricing from a short-rate model with discretely distributed expected jumps rather than assuming a step family; none contains the statement.

\phantomsection\label{res:Lemma-6}
\textbf{Lemma 6 (a maturity-step drift and volatility under the drift condition vanish).} Let the dates be as in (2.1), fix \(j \in \{0, \dots, N\}\) and \(t \in I_j\), real numbers \(\mu_0, \dots, \mu_N\) and vectors \(s_0, \dots, s_N \in \mathbb R^d\), and let \(\alpha : [0, \infty) \to \mathbb R\) and \(\sigma : [0, \infty) \to \mathbb R^d\) be the step functions \(\alpha(u) := \mu_k\), \(\sigma(u) := s_k\) for \(u \in I_k\). Put \(m_k := \max(T_k, t)\), so \(m_j = t\) and \(m_k = T_k\) for \(k > j\). If

\begin{equation*}
\int_t^T \alpha(u)\,du = \tfrac12 \Bigl| \int_t^T \sigma(u)\,du \Bigr|^2 \quad \text{for every real } T \ge t, \tag{3.12}\label{eq:3.12}
\end{equation*}

then \(\mu_k = 0\) and \(s_k = 0\) for every \(k = j, \dots, N\); the values for \(k < j\) never enter (3.12) and are unconstrained. The conclusion already follows when (3.12) is assumed only for \(T\) in a set \(M \subset [t, \infty)\) such that

\begin{equation*}
\text{for every } k = j, \dots, N,\ M \text{ contains two distinct points of the open interval } (m_k, T_{k+1}); \tag{3.13}\label{eq:3.13}
\end{equation*}

for instance \(M\) the rationals in \((t, \infty)\), or the \(2(N - j + 1)\) points \(m_k + h_k\), \(m_k + h_k / 2\) with \(h_k := (T_{k+1} - m_k)/2\) for \(k < N\) and \(h_N := 1\).

\emph{Proof.} Fix \(k \ge j\) and write \(T = m_k + \tau\) with \(0 \le \tau < T_{k+1} - m_k\), so that \(T \in I_k\) and \(T \ge t\). Lemma \hyperref[res:Lemma-3]{3}, applied to \(\sigma\) and, with \(d = 1\), to \(\alpha\), gives \(\int_t^T \sigma = C_k + s_k \tau\) and \(\int_t^T \alpha = A_k + \mu_k \tau\), where \(C_k\) and \(A_k\) are the constants (3.4) for \(\sigma\) and for \(\alpha\): zero for \(k = j\), and for \(k > j\) fixed linear combinations of \(s_j, \dots, s_{k-1}\) and of \(\mu_j, \dots, \mu_{k-1}\). Substituting into (3.12) and expanding the square,

\begin{equation*}
A_k + \mu_k \tau = \tfrac12 |C_k|^2 + \langle C_k, s_k \rangle\, \tau + \tfrac12 |s_k|^2\, \tau^2 \tag{3.14}\label{eq:3.14}
\end{equation*}

at every \(\tau\) for which (3.12) is assumed. Induct on \(k\) from \(j\) to \(N\). If \(s_m = 0\) and \(\mu_m = 0\) for \(j \le m < k\), then \(C_k = 0\) and \(A_k = 0\) exactly, and (3.14) at two distinct \(\tau_1, \tau_2 > 0\), supplied by (3.13) or by the two points \(h_k, h_k/2\), reads \(\mu_k \tau_i = \tfrac12 |s_k|^2 \tau_i^2\). Dividing by \(\tau_i\) and subtracting, \(\tfrac12 |s_k|^2 (\tau_1 - \tau_2) = 0\), so \(s_k = 0\), and then \(\mu_k = 0\).\hfill$\square$

One maturity per interval does not suffice, since \(\mu_k := \tfrac12 |s_k|^2 \tau_1\) with \(s_k \ne 0\) satisfies (3.14) at \(\tau_1\) alone; skipping an interval does not suffice either, since with the volatility nonzero on one interval only and the drift tuned to its length both sides of (3.12) agree at every maturity outside that interval; and the point \(m_k\) itself is uninformative, which is why (3.13) asks for points of the open interval. For fixed \((\omega, t)\), condition (3.12) is Assumption 2.1 with \(\alpha(u) = \alpha(t, u, \omega)\) and \(\sigma(u) = \sigma(t, u, \omega)\); that assumption holds for every \(T > 0\) outside a null set of \((\omega, t)\) that may depend on \(T\), and the version of the lemma with \(M\) the rationals is what lets Theorem \hyperref[res:Theorem-1]{1} use it at countably many fixed maturities, so that a single null set suffices.

\cite{gellert2021short} §3.1.5, eqs. 29 and 30 and the sentences after them, is the short-rate shadow of this lemma: with step volatilities the short rate acquires a deterministic term depending on \(t\) which must cancel against the initial curve, so piecewise-constant short-rate paths preclude free interpolation of the initial term structure, and their footnote 18 identifies the term as the steepening of a Gaussian model without mean reversion; no statement about the curve family and no proof are given. \cite{brace2024sofr} eqs. 18 and 19 carry the same drift with the integral of the step function left unevaluated. The classical analogues are \cite{filipovic1999note} Proposition 3.2 and Theorem 4.1, where a consistent Itô process for the Nelson--Siegel family is quasi-deterministic, and \cite{filipovic2000exponential} Theorem 3.2, which forces only the exponent coordinates of an exponential-polynomial family to be locally constant; as Remark 4.3 of \cite{filipovic1999note} says, wider exponential-polynomial families do admit restricted nontrivial consistent processes, so total vanishing is specific to Nelson--Siegel there and to the maturity-step family here.

\subsection{Proof of Theorem 1}\label{sec:3.3}

\emph{Proof of (i).} Fix \(n\) with \(1 \le n \le N\) and \(\omega\). By (3.2) and the vanishing of \(\xi_n\) before \(T_n\), the function \(u \mapsto \xi_n(u, \omega)\) is the step function (3.8) with \(X_k := X_{n,k}\) for \(k = n, \dots, N\). Assumption 2.2 says that for every real \(T \ge T_n\), (3.11) holds with \(\mathcal G := \mathcal F_{T_n-}\). The conditional form of Lemma \hyperref[res:Lemma-5]{5} gives \(X_{n,k} = 0\) almost surely for every \(k = n, \dots, N\). There are finitely many pairs \((n, k)\), so the event on which all \(X_{n,k}\) vanish has probability one, and on it \(\xi_n(u) = 0\) for every \(u\): for \(u < T_n\) by regularity, for \(u \ge T_n\) by (3.2).\hfill$\square$

\emph{Proof of (ii).} Step 1, one null set. Let \(M := \mathbb Q \cap (0, \infty)\). For each \(T \in M\), Assumption 2.1 provides a \(Q \otimes dt\)-null set \(N_T \subset \Omega \times [0, T]\) outside which (2.3) holds at \(T\). The countable union \(N := \bigcup_{T \in M} N_T\) is null. Fix \((\omega, t) \notin N\) with \(t \ge 0\) and put \(j := j(t)\); for every \(T \in M\) with \(T \ge t\), \((\omega, t) \in \Omega \times [0, T]\) and \((\omega, t) \notin N_T\), so (2.3) holds at \(T\).

Step 2, Lemma \hyperref[res:Lemma-6]{6} applies. Define step functions on \([0, \infty)\) by \(\tilde\alpha(u) := \mu_k(t, \omega)\) and \(\tilde\sigma(u) := s_k(t, \omega)\) for \(u \in I_k\). By (3.1), \(\tilde\alpha(u) = \alpha(t, u, \omega)\) and \(\tilde\sigma(u) = \sigma(t, u, \omega)\) for every \(u \ge t\), so the integrals over \([t, T]\) agree and, by Step 1, (3.12) holds for \(\tilde\alpha, \tilde\sigma\) at every \(T \in M\) with \(T \ge t\). For each \(k = j, \dots, N\) the open interval \((m_k, T_{k+1})\) is nonempty and contained in \((t, \infty)\), so it contains two distinct rationals; thus \(M \cap [t, \infty)\) satisfies (3.13), and Lemma \hyperref[res:Lemma-6]{6} gives \(\mu_k(t, \omega) = 0\) and \(s_k(t, \omega) = 0\) for every \(k = j, \dots, N\).

Step 3. Step 2 holds for every \((\omega, t)\) outside the null set \(N\). On the same set, \(\alpha(t, T, \omega) = 0\) for all \(T\): for \(T < t\) by regularity, for \(T \ge t\) by (3.1) and Step 2; likewise \(\sigma\).\hfill$\square$

The proof uses Assumption 2.1 only at rational maturities and Assumption 2.2 only at two maturities per scheduled interval; the finite and countable forms of Lemmas \hyperref[res:Lemma-5]{5} and \hyperref[res:Lemma-6]{6} exist for this bookkeeping. The hypotheses of Theorem \hyperref[res:Theorem-1]{1} are consistent: the zero coefficients satisfy everything, and a nonzero example lying outside \(S\) is given in Appendix \hyperref[sec:A]{A}.

\subsection{Piecewise-affine curves}\label{sec:3.4}

The step-plus-ramp family (2.5) replaces steps by steps plus ramps. The lemma below is the analogue of Lemma \hyperref[res:Lemma-3]{3} one degree up, with the identification statement that reads ramp slopes and levels off an identity between integrals; it is deterministic and uses no standing assumption.

\phantomsection\label{res:Lemma-7}
\textbf{Lemma 7 (integral of a piecewise-affine function, and identification).} Let the dates be as in (2.1), fix vectors \(a_0, \dots, a_N\) and \(b_0, \dots, b_N\) in \(\mathbb R^d\), and define the piecewise-affine function

\begin{equation*}
g(u) := a_k + b_k (u - T_k) \qquad \text{for } u \in I_k. \tag{3.15}\label{eq:3.15}
\end{equation*}

(a) For \(0 \le t \le T < \infty\) with \(t \in I_j\), \(T \in I_k\), and \(m_k := \max(T_k, t)\), one has \(j \le k\) and

\begin{equation*}
\int_t^T g(u)\,du = e_{jk}(t) + a_k (T - m_k) + \tfrac12 b_k \bigl[ (T - T_k)^2 - (m_k - T_k)^2 \bigr], \tag{3.16}\label{eq:3.16}
\end{equation*}

where \(e_{jj}(t) := 0\) and, for \(k > j\),

\begin{equation*}
e_{jk}(t) := a_j (T_{j+1} - t) + \tfrac12 b_j \bigl[ (T_{j+1} - T_j)^2 - (t - T_j)^2 \bigr] + \sum_{m=j+1}^{k-1} \bigl[ a_m (T_{m+1} - T_m) + \tfrac12 b_m (T_{m+1} - T_m)^2 \bigr]; \tag{3.17}\label{eq:3.17}
\end{equation*}

\(e_{jk}(t)\) does not depend on \(T\). Equivalently, with \(\tau := T - m_k \in [0, T_{k+1} - m_k)\),

\begin{equation*}
\int_t^T g(u)\,du = e_{jk}(t) + g(m_k)\, \tau + \tfrac12 b_k\, \tau^2, \tag{3.18}\label{eq:3.18}
\end{equation*}

so on \([m_k, T_{k+1})\) the integral is a polynomial of degree at most two in \(\tau\), with \(\tau^2\)-coefficient \(\tfrac12 b_k\), \(\tau\)-coefficient \(g(m_k)\) and constant term \(e_{jk}(t)\).

(b) If \(P(\tau) := c_0 + c_1 \tau + c_2 \tau^2\) with \(c_0, c_1, c_2 \in \mathbb R^d\) vanishes at three distinct real numbers, then \(c_0 = c_1 = c_2 = 0\); if \(c_2 = 0\) and \(P\) vanishes at two distinct real numbers, then \(c_0 = c_1 = 0\). Consequently two functions of the form \(p + q\tau + r\tau^2\) that agree at three distinct points have the same coefficients, so the three coefficients in (3.18) are determined by the values of the integral at any three distinct maturities of \([m_k, T_{k+1})\); and two functions of the form (3.15) that agree at two distinct points of every \(I_k\) have the same \(a_k\) and \(b_k\), so the representation (3.15) on the fixed partition is unique.

\emph{Proof.} (a) On a piece \([p, q) \subset I_k\) the integrand is \(a_k + b_k(u - T_k)\) except at the null point \(q\), so \(\int_p^q g = a_k (q - p) + \tfrac12 b_k [(q - T_k)^2 - (p - T_k)^2]\). For \(k = j\) the single piece \([t, T)\) gives (3.16) with \(m_j = t\) and \(e_{jj} = 0\). For \(k > j\), additivity along \(t < T_{j+1} \le \dots \le T_k \le T\) gives the sum of the piece integrals; the first \(k - j\) pieces are \(e_{jk}(t)\) by (3.17) and the last is \(a_k (T - T_k) + \tfrac12 b_k (T - T_k)^2\), which is the rest of (3.16) since \(m_k = T_k\). For (3.18), write \(T - T_k = \tau + (m_k - T_k)\), so the bracket in (3.16) is \(\tau^2 + 2\tau(m_k - T_k)\), and collect terms using \(g(m_k) = a_k + b_k (m_k - T_k)\), valid because \(m_k \in I_k\). (b) With \(\tau_1, \tau_2, \tau_3\) distinct, \(0 = P(\tau_i) - P(\tau_1) = (\tau_i - \tau_1)[c_1 + c_2(\tau_1 + \tau_i)]\) for \(i = 2, 3\); dividing by \(\tau_i - \tau_1 \ne 0\) and subtracting gives \(c_2 (\tau_2 - \tau_3) = 0\), so \(c_2 = 0\), then \(c_1 = 0\), then \(c_0 = P(\tau_1) = 0\). The affine case is the first half of the same computation. The two consequences follow by applying this to differences.\hfill$\square$

With all \(b_k = 0\), Lemma \hyperref[res:Lemma-7]{7} reduces to Lemma \hyperref[res:Lemma-3]{3}. Three points are needed in (b), since \((\tau - \tau_1)(\tau - \tau_2)\) is a nonzero quadratic vanishing at two. The step-plus-ramp curves (2.5) are of the form (3.15) on each \(I_k\) with \(a_k = \sum_{m \le k} [M_m + C_m (T_k - T_m)]\) and \(b_k = \sum_{m \le k} C_m\). The partition bookkeeping is stated nowhere in the literature; \cite{gellert2021short} eqs. 65 to 67 and \cite{brace2024sofr} eqs. 18 and 19 integrate or leave unevaluated the step case only.

\subsection{The drift of a step-plus-ramp curve}\label{sec:3.5}

In the step-plus-ramp family the volatility at a fixed time is a step function of maturity and the drift is piecewise affine. The next result identifies the drift from the volatility: the drift condition holds if and only if the drift is the differentiated HJM drift, so that on each scheduled interval the ramp slope equals the squared step volatility. Like Lemmas \hyperref[res:Lemma-6]{6} and \hyperref[res:Lemma-7]{7} it is deterministic and pointwise in \((\omega, t)\), and it uses no standing assumption; the passage from Assumption 2.1 to its hypothesis for almost every \((\omega, t)\) is the bookkeeping of Section \hyperref[sec:3.3]{3.3} with three rational maturities per interval in place of two.

Let the dates be as in (2.1), fix \(j \in \{0, \dots, N\}\) and \(t \in I_j\), and put \(m_k := \max(T_k, t)\) for \(k \ge j\). Let \(\sigma\) be the step function with values \(s_0, \dots, s_N \in \mathbb R^d\) as in Lemma \hyperref[res:Lemma-3]{3}, let \(C_k := c_{jk}(t)\) be its constants (3.4), so that \(\int_t^u \sigma = C_k + s_k (u - m_k)\) for \(u \in [m_k, T_{k+1})\), and let \(\alpha\) be a piecewise-affine function with data \(a_0, \dots, a_N, b_0, \dots, b_N \in \mathbb R\) as in (3.15). The \emph{differentiated HJM drift} of \(\sigma\) from time \(t\) is

\begin{equation*}
\alpha^*(u) := \Bigl\langle \sigma(u), \int_t^u \sigma(v)\,dv \Bigr\rangle, \qquad u \ge t, \tag{3.19}\label{eq:3.19}
\end{equation*}

which is piecewise affine on \([t, \infty)\) with level and slope

\begin{equation*}
b^*_k := |s_k|^2, \qquad a^*_k := \langle s_k, C_k \rangle - |s_k|^2 (m_k - T_k), \qquad k \ge j, \tag{3.20}\label{eq:3.20}
\end{equation*}

on the pieces \([m_k, T_{k+1})\): \(\alpha^*(u) = \langle s_k, C_k \rangle + |s_k|^2 (u - m_k) = a^*_k + b^*_k (u - T_k)\) there.

\phantomsection\label{res:Proposition-8}
\textbf{Proposition 8.} The following are equivalent.

(A) The drift condition holds for every real \(T \ge t\):

\begin{equation*}
\int_t^T \alpha(u)\,du = \tfrac12 \Bigl| \int_t^T \sigma(u)\,du \Bigr|^2. \tag{3.21}\label{eq:3.21}
\end{equation*}

(A\('\)) Condition (3.21) holds for every \(T\) in some set \(M \subset [t, \infty)\) such that

\begin{equation*}
\text{for every } k = j, \dots, N,\ M \text{ contains three distinct points of the open interval } (m_k, T_{k+1}). \tag{3.22}\label{eq:3.22}
\end{equation*}

(B) \(a_k = a^*_k\) and \(b_k = b^*_k\) for every \(k = j, \dots, N\).

(B\('\)) \(\alpha(u) = \alpha^*(u)\) for every \(u \ge t\).

The values \(a_k, b_k\) for \(k < j\) enter none of these and are unconstrained.

\emph{Proof.} Fix \(k \ge j\), let \(L_k := T_{k+1} - m_k \in (0, \infty]\) be the length of the piece \([m_k, T_{k+1})\), and for \(T\) in the piece write \(\tau := T - m_k\). By Lemma \hyperref[res:Lemma-7]{7}(a) in the form (3.18) with \(d = 1\), and by (3.4) with bilinearity of the inner product,

\begin{equation*}
\int_t^T \alpha(u)\,du = e_{jk}(t) + \alpha(m_k)\,\tau + \tfrac12 b_k \tau^2, \qquad \tfrac12 \Bigl| \int_t^T \sigma(u)\,du \Bigr|^2 = \tfrac12 |C_k|^2 + \langle C_k, s_k \rangle\, \tau + \tfrac12 |s_k|^2 \tau^2, \tag{3.23}\label{eq:3.23}
\end{equation*}

with \(\alpha(m_k) = a_k + b_k (m_k - T_k)\) and \(e_{jk}(t)\) the constant (3.17). The constants obey, for \(k < N\), the recursions \(e_{j,k+1}(t) = e_{jk}(t) + \alpha(m_k) L_k + \tfrac12 b_k L_k^2\), the increment being the integral of \(\alpha\) over the piece, and \(C_{k+1} = C_k + s_k L_k\), hence \(\tfrac12 |C_{k+1}|^2 = \tfrac12 |C_k|^2 + \langle C_k, s_k \rangle L_k + \tfrac12 |s_k|^2 L_k^2\).

(B) implies (A): under (B), \(\alpha(m_k) = \langle s_k, C_k \rangle\) and \(\tfrac12 b_k = \tfrac12 |s_k|^2\), so the \(\tau\)- and \(\tau^2\)-coefficients in (3.23) agree on every piece; the constants agree by induction on \(k\), since \(e_{jj}(t) = 0 = \tfrac12 |C_j|^2\) and the two recursions then carry \(e_{jk}(t) = \tfrac12 |C_k|^2\) to \(k + 1\). So the two sides of (3.21) coincide as polynomials in \(\tau\) on every piece, which is (A). (A) implies (A\('\)) with \(M = [t, \infty)\). (A\('\)) implies (B): on each piece the two quadratics in (3.23) agree at three distinct points of \((0, L_k)\), so by Lemma \hyperref[res:Lemma-7]{7}(b) their coefficients agree; the \(\tau^2\)-coefficients give \(b_k = |s_k|^2\), and the \(\tau\)-coefficients, with \(\alpha(m_k) = a_k + b_k(m_k - T_k)\), give \(a_k = a^*_k\). This step is interval-local and uses no induction. (B) and (B\('\)) are equivalent because the pieces \([m_k, T_{k+1})\), \(k \ge j\), partition \([t, \infty)\) and on each piece \(\alpha\) and \(\alpha^*\) are affine in \(u - T_k\): equal data give equal functions, and equal functions agree at two distinct points of a piece of positive length, so by the affine case of Lemma \hyperref[res:Lemma-7]{7}(b) the data agree.\hfill$\square$

Three maturities per interval cannot be reduced to two: a quadratic with the wrong \(\tau^2\)-coefficient, hence the wrong slope \(b_k\), can match the right side of (3.23) at two points of an interval, and Lemma \hyperref[res:Lemma-6]{6} needed only two because there the slope was known to vanish. With all \(b_k = 0\) imposed, Proposition \hyperref[res:Proposition-8]{8} recovers Lemma \hyperref[res:Lemma-6]{6}. The induction in the sufficiency direction is not vacuous: a deviation from (B) on one interval changes the constants \(e_{jk}(t)\) of every later interval. The proposition determines the drift from the volatility with no differentiability hypothesis on \(\alpha\) in the maturity beyond membership in the family, which is what makes (A) and (B\('\)) equivalent: \cite{heath1992bond} state the drift condition in differentiated form under differentiability in maturity, and \cite{fontana2022term} Theorem 3.5(ii) in integrated form; within this family the two forms are the same statement. Sufficiency, (B) implies (A), is the content of the explicit models of \cite{gellert2021short} eqs. 28 to 30, \cite{brace2024sofr} eqs. 18 and 19, \cite{fontana2022term} Example 3.10 eq. 3.14 with constant jump shape, \cite{kellerressel2018affine} Examples 6.15 and 6.16 and \cite{kim2014jumps} Proposition 1, all of which construct members of the family and verify the drift condition; none states the converse. The classical analogue is the consistency calculus of \cite{bjork1999interest} §2 and §3 and \cite{filipovic2000exponential} §2 and §3, which match the drift generated by the volatility against the tangent space of the family, and which do not treat a piecewise-affine family with breakpoints at scheduled dates. In the parametrization (2.5) with drifts \(\dot M_m, \dot C_m\), the slope of the drift on \(I_k\) is \(\sum_{m \le k} \dot C_m(t)\), so (B) reads \(\sum_{m \le k} \dot C_m(t) = |s_k(t)|^2\) for every \(k \ge j(t)\) (a computation, not part of the proposition).

\subsection{The jump of a step-plus-ramp curve at a scheduled date}\label{sec:3.6}

At a scheduled date the curve of the step-plus-ramp family jumps by a function that is affine on each later scheduled interval. The next result characterizes the level jumps that the jump martingale condition allows, given the slope jumps. It is stated for one scheduled date and an arbitrary sub-sigma-algebra in place of \(\mathcal F_{T_n-}\), like the conditional form of Lemma \hyperref[res:Lemma-5]{5}; the assembly over all dates and with the drift half is left to the classification of Section \hyperref[sec:5.2]{5.2}.

Let \((\Omega, \mathcal F, P)\) be a probability space and \(\mathcal G \subseteq \mathcal F\) a sub-sigma-algebra. Let the dates be as in (2.1) with \(N \ge 1\), fix \(n \in \{1, \dots, N\}\), real random variables \(X_n, \dots, X_N\) (the level jumps) and \(\mathcal G\)-measurable real random variables \(y_n, \dots, y_N\) (the slope jumps), put \(L_k := T_{k+1} - T_k\) for \(k < N\), and let \(\xi\) be the piecewise-affine jump

\begin{equation*}
\xi(u, \omega) := X_k(\omega) + y_k(\omega)\,(u - T_k) \ \text{for } u \in I_k,\ k \ge n, \qquad \xi(u, \omega) := 0 \ \text{for } u < T_n. \tag{3.24}\label{eq:3.24}
\end{equation*}

The hypothesis is the jump martingale condition for this jump: for every real \(T \ge T_n\),

\begin{equation*}
E\Bigl[ \exp\Bigl( -\int_{T_n}^T \xi(u, \cdot)\,du \Bigr) \Bigm| \mathcal G \Bigr] = 1 \quad P\text{-almost surely}, \tag{3.25}\label{eq:3.25}
\end{equation*}

the conditional expectation of the nonnegative variable taken in \([0, \infty]\). By Lemma \hyperref[res:Lemma-7]{7}(a) with \(d = 1\) and \(t = T_n\), for \(T \in I_k\), \(k \ge n\), and \(\tau := T - T_k\),

\begin{equation*}
\int_{T_n}^T \xi(u, \cdot)\,du = S_k + X_k \tau + \tfrac12 y_k \tau^2, \qquad S_n := 0, \quad S_k := \sum_{m=n}^{k-1} \bigl( X_m L_m + \tfrac12 y_m L_m^2 \bigr) \ (k > n). \tag{3.26}\label{eq:3.26}
\end{equation*}

Put \(R_k := \exp(-\sum_{m=n}^{k-1} L_m X_m)\) and \(Z_k := \exp(\sum_{m=n}^{k-1} \tfrac12 y_m L_m^2)\), so that \(\exp(-S_k) = R_k / Z_k\), with \(R_k > 0\) and \(Z_k\) positive and \(\mathcal G\)-measurable; \(R_n = Z_n = 1\).

\phantomsection\label{res:Proposition-9}
\textbf{Proposition 9.} (a) If (3.25) holds, then for every \(k \in \{n, \dots, N\}\):

\begin{equation*}
E[R_k \mid \mathcal G] = Z_k \quad P\text{-almost surely}, \tag{3.27}\label{eq:3.27}
\end{equation*}

so that \(P_k(A) := E[(R_k / Z_k) \mathbf 1_A]\) is a probability measure on \((\Omega, \mathcal F)\) agreeing with \(P\) on \(\mathcal G\); \(y_k \ge 0\) almost surely; and the regular conditional distribution \(\kappa_k\) of \(X_k\) given \(\mathcal G\) under \(P_k\) satisfies

\begin{equation*}
\kappa_k(\omega, \cdot) = N(0, y_k(\omega)) \quad \text{for } P\text{-almost every } \omega, \tag{3.28}\label{eq:3.28}
\end{equation*}

with \(N(0, 0) := \delta_0\). In words: given \(\mathcal G\), the level jump on the date's own interval is centered Gaussian with variance the slope jump, and on each later interval the level jump is centered Gaussian with variance the slope jump under the tilt by the exponent accumulated over the earlier intervals.

(b) Conversely, if \(y_k \ge 0\) almost surely and (3.27) and (3.28) hold for every \(k\), then (3.25) holds.

(c) If conditionally on \(\mathcal G\) the vector \((X_n, \dots, X_N)\) is Gaussian with \(\mathcal G\)-measurable conditional mean \(\mu\) and covariance \(\Sigma\), then (3.25) holds if and only if, almost surely, for every \(k\),

\begin{equation*}
\mu_k = \sum_{m=n}^{k-1} L_m\, \Sigma_{km} \qquad \text{and} \qquad y_k = \Sigma_{kk}. \tag{3.29}\label{eq:3.29}
\end{equation*}

(d) There are a probability space, \(N = 2\), \(n = 1\), trivial \(\mathcal G\), deterministic \(y_1, y_2 > 0\) and \((X_1, X_2)\) satisfying (3.25) with \(X_1 \sim N(0, y_1)\) and \(X_2\) not Gaussian: under a measure \(Q\) carrying independent \(X_1 \sim N(-L_1 y_1, y_1)\) and \(Z \sim N(0, y_2)\), put \(s := \operatorname{sign}(X_1 + L_1 y_1)\), \(X_2 := s |Z|\) and \(dP/dQ := \exp(L_1 X_1 + \tfrac12 L_1^2 y_1)\). Under \(P\), \(X_1 \sim N(0, y_1)\), (3.25) holds for every \(T \ge T_1\), and \(X_2\) has the density \(x \mapsto 2 \varphi_{y_2}(x)\,[p \mathbf 1\{x > 0\} + (1 - p) \mathbf 1\{x < 0\}]\) with \(p := \Phi(L_1 \sqrt{y_1}) > \tfrac12\), which is discontinuous at \(0\); under the tilt \(P_2 = Q\) the law of \(X_2\) is \(N(0, y_2)\), as (3.28) requires.

The proposition uses no standing assumption. With \(\mathcal G = \mathcal F_{T_n-}\) and \(\xi = \xi_n\), (3.25) is Assumption 2.2 for the jump at \(T_n\), specialized to a jump of the form (3.24) with slopes measurable before the date; that specialization is a hypothesis on the family. Condition (3.29) is the "Gaussian jump laws with matching covariances" of the literature: the level jump on a later interval has conditional mean equal to its covariances with the earlier level jumps weighted by the earlier interval lengths, and each slope jump equals the corresponding variance. Part (d) shows that (3.25) does not force joint Gaussianity; what it forces is exactly the per-interval structure (3.28), of which (3.29) is the jointly Gaussian case. With zero drift and volatility and no further jumps the example of (d) is a consistent member of the step-plus-ramp family (a statement, not proved here); Theorem \hyperref[res:Theorem-11]{11} draws the consequence for the classification, and Proposition \hyperref[res:Proposition-48]{48} treats shapes that are not affine.

\emph{Proof.} (a) Step 1, a change of measure. If \(D \ge 0\) is \(\mathcal F\)-measurable with \(E[D \mid \mathcal G] = 1\), then \(P_D(A) := E[D \mathbf 1_A]\) is a probability measure agreeing with \(P\) on \(\mathcal G\), and for every real random variable \(Y\) and measurable \(h \ge 0\), \(E_{P_D}[h(Y) \mid \mathcal G] = E_P[D\, h(Y) \mid \mathcal G]\) almost surely, both sides taken in \([0, \infty]\); hence a regular conditional distribution \(\kappa\) of \(Y\) given \(\mathcal G\) under \(P_D\) satisfies \(\int h\,d\kappa(\omega, \cdot) = E_P[D\, h(Y) \mid \mathcal G](\omega)\) for almost every \(\omega\). Step 2. Condition (3.25) at \(T = T_k\), where \(\tau = 0\), reads \(E[R_k / Z_k \mid \mathcal G] = 1\); pulling out the \(\mathcal G\)-measurable \(Z_k\) gives (3.27), so \(D_k := R_k / Z_k\) qualifies for Step 1 and defines \(P_k\). Step 3. For \(\tau \in [0, L_k)\), (3.25) at \(T = T_k + \tau\) with (3.26) reads, after pulling out the \(\mathcal G\)-measurable factors, \(E[D_k\, e^{-\tau X_k} \mid \mathcal G] = \exp(\tfrac12 y_k \tau^2)\); by Step 1 the left side is the Laplace transform \(M_k(\omega, \tau) := \int e^{-\tau x}\,\kappa_k(\omega, dx)\) of the regular conditional distribution of \(X_k\) under \(P_k\), so \(M_k(\omega, \tau) = \exp(\tfrac12 y_k(\omega) \tau^2)\) for almost every \(\omega\), for each fixed \(\tau\). Step 4. Taking the union of the exceptional sets over rational \(\tau\) gives one null set; off it, \(M_k(\omega, \cdot)\) is the Laplace transform of a probability measure, hence convex with values in \((0, \infty]\), finite at the rationals of \([0, L_k)\), therefore finite and continuous on \((0, L_k)\), and equal to \(\exp(\tfrac12 y_k \tau^2)\) there by density. Step 5. Its logarithm is convex (Hölder), while \(\tfrac12 y_k \tau^2\) is strictly concave if \(y_k < 0\); so \(y_k \ge 0\) off the null set. Step 6. Fix \(\omega\) off the null set, \(y := y_k(\omega)\), and \(\tau_0 \in (0, L_k)\), and tilt: \(\nu(dx) := e^{-\tau_0 x}\,\kappa_k(\omega, dx) / M_k(\omega, \tau_0)\) has Laplace transform \(\exp(y \tau_0 s + \tfrac12 y s^2)\) for \(s\) in the open interval \((-\tau_0, L_k - \tau_0)\) around \(0\), which is the Laplace transform of \(N(-y\tau_0, y)\). Two probability measures on \(\mathbb R\) whose Laplace transforms are finite and agree on an open interval containing \(0\) coincide (\cite{curtiss1942note}; the statement is standard), so \(\nu = N(-y\tau_0, y)\), and tilting back by \(e^{\tau_0 x}\) gives \(\kappa_k(\omega, \cdot) = N(0, y)\), with \(\delta_0\) when \(y = 0\). This is (3.28).

(b) Given (3.27), Step 1 applies to \(D_k\), and (3.28) gives \(E[D_k e^{-\tau X_k} \mid \mathcal G] = \exp(\tfrac12 y_k \tau^2)\); multiplying by \(\exp(-\tfrac12 y_k \tau^2)\) and using (3.26) returns (3.25) at \(T = T_k + \tau\), and every \(T \ge T_n\) is of this form for one \(k\).

(c) Conditionally on \(\mathcal G\), tilting the Gaussian law \(N(\mu, \Sigma)\) by \(\exp(-\langle \ell, x \rangle)\) with \(\ell_m = L_m\) for \(m < k\) and \(0\) otherwise gives \(N(\mu - \Sigma \ell, \Sigma)\), whose \(k\)-th marginal, \(N(\mu_k - \sum_{m<k} L_m \Sigma_{km}, \Sigma_{kk})\), is \(\kappa_k\). By (a) and (b), (3.25) holds if and only if this marginal is \(N(0, y_k)\) for every \(k\) and (3.27) holds; the first is (3.29), and (3.29) implies (3.27) because \(-\langle \ell, \mu \rangle + \tfrac12 \langle \ell, \Sigma \ell \rangle\) reduces, after the cross terms cancel against (3.29) for the earlier indices, to \(\tfrac12 \sum_{m<k} L_m^2 y_m = \log Z_k\).

(d) Under \(Q\), \(X_1 + L_1 y_1\) is symmetric and independent of \(Z\), so \(s\) is a fair sign independent of \(|Z|\) and \(X_2 \sim N(0, y_2)\); \(dP/dQ\) integrates to one. Under \(P\), completing the square gives \(X_1 \sim N(0, y_1)\); (3.25) on \(I_1\) is the Gaussian Laplace transform, and on \(I_2\) the density \(dP/dQ\) cancels \(\exp(-S_2)\) exactly, leaving \(E_Q[e^{-\tau X_2}]\,e^{-\frac12 y_2 \tau^2} = 1\). Since \(dP/dQ\) depends on \(X_1\) alone, \((s, |Z|)\) stay independent under \(P\) with \(P(s = 1) = \Phi(L_1 \sqrt{y_1}) = p\), which gives the stated density of \(X_2\); its jump at \(0\) rules out a Gaussian law. Finally \(D_2 = (dP/dQ)^{-1}\), so \(P_2 = Q\).\hfill$\square$

The one analytic input is the uniqueness of a law from its Laplace transform on an interval; everything else is the change of measure and the quadratic (3.26). The mechanism of (a) is the one the literature anticipates for a single Gaussian jump, "a linear ramp forces a quadratic cumulant generating function", made precise for a whole sequence of intervals, and (d) shows that it yields Gaussianity on later intervals only under the accumulated tilt. \cite{fontana2022term} Example 3.10 is the sufficiency instance for one date with a Gaussian jump, whose compensator slope for a constant jump shape is the variance, which is (3.29) for \(k = n\); their Proposition 4.5 carries the general-law tilted-mean mechanism for the overnight rate. \cite{kellerressel2018affine} Examples 6.15 and 6.16 and \cite{kim2014jumps} Proposition 1 are the jointly Gaussian multi-date instances, sufficiency in case (c). \cite{silva2024discretely} and \cite{calvia2025short} have non-Gaussian scheduled jump laws in short-rate models whose curves are not piecewise affine, so they do not bear on (d). The characterization and the example are stated nowhere.

\subsection{An explicit step-plus-ramp model}\label{sec:3.7}

The one-date version of the following construction is \cite{fontana2022term} Example 3.10 with constant jump shape and centered Gaussian jump; \cite{kellerressel2018affine} Example 6.16 and \cite{kim2014jumps} Proposition 1 give multi-date Gaussian affine constructions. Here the coefficients are given directly on the scheduled partition and both standing assumptions are verified for every date and maturity.

Fix deterministic vectors \(s_0, \dots, s_N \in \mathbb R^d\) and numbers \(y_{n,k} \ge 0\), \(1 \le n \le k \le N\). Let

\begin{equation*}
\mathcal P := \{(n,k): 1 \le n \le k \le N\}, \qquad \Omega := \mathbb R^{\mathcal P}, \qquad Q := \bigotimes_{(n,k) \in \mathcal P} N(0,y_{n,k}), \tag{3.30}\label{eq:3.30}
\end{equation*}

where \(N(0,0) := \delta_0\), and write \(X_{n,k}\) for the coordinate maps. Thus the \(X_{n,k}\) are independent and \(X_{n,k} \sim N(0,y_{n,k})\). Take the filtration

\begin{equation*}
\mathcal F_t := \sigma\bigl(X_{m,k}: (m,k) \in \mathcal P,\ T_m \le t\bigr). \tag{3.31}\label{eq:3.31}
\end{equation*}

For \(T \ge t\), \(T \in I_k\), and \(1 \le n \le k\), define

\begin{equation*}
\sigma(t,T) := s_k, \qquad \alpha(t,T) := \Bigl\langle s_k, \int_t^T \sigma(t,u)\,du \Bigr\rangle, \tag{3.32}\label{eq:3.32}
\end{equation*}

\begin{equation*}
\xi_n(T) := X_{n,k} + y_{n,k}(T-T_k), \tag{3.33}\label{eq:3.33}
\end{equation*}

and set each coefficient to zero before its current date as in Section \hyperref[sec:2]{2}. Let \(f(0,\cdot)\) be Borel measurable and bounded on bounded intervals.

\phantomsection\label{res:Proposition-10}
\textbf{Proposition 10.} The coefficient system (3.30)--(3.33) satisfies the regularity conditions of Section \hyperref[sec:2]{2}, Assumption 2.1 pointwise for every \(0 \le t \le T\), and Assumption 2.2 for every scheduled date and every \(T \ge T_n\).

\emph{Proof.} The initial curve is locally integrable. The volatility is deterministic, Borel measurable and bounded by \(S := \max_k |s_k|\). Lemma \hyperref[res:Lemma-3]{3} makes \((t,T) \mapsto \int_t^T \sigma(t,u)\,du\), hence \(\alpha\), Borel measurable on each of the finitely many regions determined by the scheduled intervals. The bounds \(|\sigma| \le S\) and \(|\alpha(t,T)| \le S^2(T-t)\) give the required local integrability. Each \(\xi_n\) is a finite sum of indicator-times-affine functions with \(\mathcal F_{T_n}\)-measurable coefficients and is locally integrable pathwise. This proves the regularity conditions.

Equation (3.32) is the differentiated HJM drift. Proposition \hyperref[res:Proposition-8]{8}, (B\('\)) implies (A), therefore gives (2.3) for every \(0 \le t \le T\).

Fix \(n\) and let \(T \in I_k\), \(k \ge n\), with \(\tau := T-T_k\). Lemma \hyperref[res:Lemma-7]{7} gives

\begin{equation*}
\int_{T_n}^T \xi_n(u)\,du = \sum_{m=n}^{k-1} \bigl(L_m X_{n,m} + \tfrac12 y_{n,m}L_m^2\bigr) + \tau X_{n,k} + \tfrac12 y_{n,k}\tau^2. \tag{3.34}\label{eq:3.34}
\end{equation*}

The left limit of (3.31) at \(T_n\) is generated by the coordinates \(X_{m,k}\) with \(m<n\), while (3.34) uses only coordinates whose first index is \(n\). The two sets of coordinates are independent. The conditional expectation in (2.4) is therefore the unconditional expectation, and independence together with the Gaussian Laplace transform gives

\begin{equation*}
\prod_{m=n}^{k-1} \exp\bigl(\tfrac12 y_{n,m}L_m^2\bigr)\exp\bigl(-\tfrac12 y_{n,m}L_m^2\bigr)\,
\exp\bigl(\tfrac12 y_{n,k}\tau^2\bigr)\exp\bigl(-\tfrac12 y_{n,k}\tau^2\bigr) = 1. \tag{3.35}\label{eq:3.35}
\end{equation*}

This is (2.4).\hfill$\square$

Choosing some \(s_k\) or \(y_{n,k}\) nonzero makes the coefficient system nonzero. On an extension carrying a Brownian motion independent of the coordinates, insert (3.32) and (3.33) into (2.2). If the initial curve is affine on each \(I_k\), then the curve remains affine there. Its slope on \(I_k \cap [t,\infty)\) is the initial slope plus \(|s_k|^2t + \sum_{n:T_n \le t} y_{n,k}\): the drift contributes \(|s_k|^2t\), the diffusion is constant in \(T\) on the interval, and each realized jump contributes \(y_{n,k}\). Thus the construction is a member of the step-plus-ramp family, not merely a coefficient system satisfying the two standing conditions.

\subsection{Classification of the step-plus-ramp family}\label{sec:3.8}

We now assemble Propositions \hyperref[res:Proposition-8]{8} and \hyperref[res:Proposition-9]{9} under the quantifiers of the standing assumptions. Coefficients belong to the \emph{step-plus-ramp family} \(S^+\) with deterministic step volatility if there are deterministic functions \(s_k : [0,\infty) \to \mathbb R^d\), coefficient processes \(a_k,b_k\), level jumps \(X_{n,k}\), and \(\mathcal F_{T_n-}\)-measurable slope jumps \(y_{n,k}\) such that, for \(T \ge t\), \(T \in I_k\), and \(u \in I_k\), \(k \ge n\),

\begin{equation*}
\sigma(t,T) = s_k(t), \qquad \alpha(t,T) = a_k(t) + b_k(t)(T-T_k), \tag{3.36}\label{eq:3.36}
\end{equation*}

\begin{equation*}
\xi_n(u) = X_{n,k} + y_{n,k}(u-T_k). \tag{3.37}\label{eq:3.37}
\end{equation*}

The coefficients also satisfy the regularity conditions of Section \hyperref[sec:2]{2}. Put \(j := j(t)\), \(m_k := \max(T_k,t)\), and \(C_k(t) := c_{j k}(t)\), with \(c_{jk}(t)\) defined by (3.4) for the step values \(s_0(t),\dots,s_N(t)\). Thus, for \(u \in [m_k,T_{k+1})\),

\begin{equation*}
\int_t^u \sigma(t,v)\,dv = C_k(t) + s_k(t)(u-m_k). \tag{3.38}\label{eq:3.38}
\end{equation*}

For \(1 \le n \le k \le N\), set

\begin{equation*}
R_{n,k} := \exp\Bigl(-\sum_{m=n}^{k-1}L_m X_{n,m}\Bigr), \qquad Z_{n,k} := \exp\Bigl(\sum_{m=n}^{k-1}\tfrac12 y_{n,m}L_m^2\Bigr). \tag{3.39}\label{eq:3.39}
\end{equation*}

Sufficiency and explicit Gaussian instances occur in \cite{gellert2021short} eqs. 28 to 30, \cite{brace2024sofr} eqs. 18 and 19, \cite{fontana2022term} Example 3.10, \cite{kellerressel2018affine} Examples 6.15 and 6.16, and \cite{kim2014jumps} Proposition 1. The theorem adds necessity and states the exact jump law outside the jointly Gaussian class.

\phantomsection\label{res:Theorem-11}
\textbf{Theorem 11.} For coefficients in \(S^+\) with deterministic step volatility, Assumptions 2.1 and 2.2 hold if and only if both of the following conditions hold.

(i) Outside one \(Q \otimes dt\)-null set, for every \(t \ge 0\) and every \(k \ge j(t)\),

\begin{equation*}
b_k(t) = |s_k(t)|^2, \qquad a_k(t) = \langle s_k(t),C_k(t)\rangle - |s_k(t)|^2(m_k-T_k). \tag{3.40}\label{eq:3.40}
\end{equation*}

Equivalently, outside that null set,

\begin{equation*}
\alpha(t,T) = \Bigl\langle \sigma(t,T),\int_t^T \sigma(t,u)\,du\Bigr\rangle \qquad \text{for every } T \ge t. \tag{3.41}\label{eq:3.41}
\end{equation*}

In particular, the drift is determined by the deterministic volatility and is deterministic outside the null set; its slope on \(I_k\) is \(|s_k(t)|^2\).

(ii) For every \(1 \le n \le k \le N\), \(y_{n,k} \ge 0\) almost surely,

\begin{equation*}
E[R_{n,k}\mid\mathcal F_{T_n-}] = Z_{n,k} \quad Q\text{-almost surely}, \tag{3.42}\label{eq:3.42}
\end{equation*}

and, under the probability measure \(Q_{n,k}(A) := E[(R_{n,k}/Z_{n,k})\mathbf 1_A]\), the conditional distribution of \(X_{n,k}\) given \(\mathcal F_{T_n-}\) is \(N(0,y_{n,k})\). Thus the level jump is centered Gaussian with variance the slope jump on the date's own interval and under the accumulated tilt on every later interval.

If, conditionally on \(\mathcal F_{T_n-}\), \((X_{n,n},\dots,X_{n,N})\) is jointly Gaussian with mean \(\mu^{(n)}\) and covariance \(\Sigma^{(n)}\), condition (ii) is equivalent to

\begin{equation*}
\mu^{(n)}_k = \sum_{m=n}^{k-1}L_m\Sigma^{(n)}_{km}, \qquad y_{n,k} = \Sigma^{(n)}_{kk}, \qquad k=n,\dots,N. \tag{3.43}\label{eq:3.43}
\end{equation*}

\emph{Proof.} For necessity in (i), let \(M\) be the positive rationals. Assumption 2.1 at each \(T \in M\) fails only on a \(Q \otimes dt\)-null set in \(\Omega \times [0,T]\). Their countable union is one null set in \(\Omega \times [0,\infty)\). Fix \((\omega,t)\) outside it. Every open piece \((m_k,T_{k+1})\), including the unbounded last piece, contains three distinct points of \(M\). Equations (3.36) and (3.38) therefore put the fixed functions \(T \mapsto \alpha(t,T,\omega)\) and \(T \mapsto \sigma(t,T)\) under condition (A\('\)) of Proposition \hyperref[res:Proposition-8]{8}. Its implication (A\('\)) to (B) gives (3.40), and its equivalence of (B) and (B\('\)) gives (3.41).

For necessity in (ii), fix \(n\). Equations (3.37) and (2.4) are (3.24) and (3.25) with \(\mathcal G = \mathcal F_{T_n-}\). Proposition \hyperref[res:Proposition-9]{9}(a) gives nonnegativity, (3.42), and the conditional Gaussian law for every \(k \ge n\).

Conversely, (3.40) is condition (B) of Proposition \hyperref[res:Proposition-8]{8} for almost every \((\omega,t)\). The implication (B) to (A) gives (2.3) outside that null set for every maturity. For each \(n\), condition (ii) is the hypothesis of Proposition \hyperref[res:Proposition-9]{9}(b), which gives (2.4) for every \(T \ge T_n\). Hence both standing assumptions hold. The jointly Gaussian statement is Proposition \hyperref[res:Proposition-9]{9}(c) at each date.\hfill$\square$

In the curve parametrization (2.5), the ramp coefficient \(C_m\) attached to \(T_m\) is, in the pairwise bookkeeping of \cite{gellert2021short} eq. 28 with pairwise coefficients \(C_{nm}\) entered with a minus sign, the identifiable combination \(C_m = -\sum_n C_{nm}\) over the pairs whose later date is \(T_m\); suppose these coefficients have drifts \(\dot C_m\). The slope of the curve drift on \(I_k\) is \(\sum_{m \le k}\dot C_m\), so (3.40) gives \(\sum_{m \le k}\dot C_m = |s_k|^2\), or \(\dot C_k = |s_k|^2-|s_{k-1}|^2\) for \(k \ge 1\). With volatility increments \(s_k=\sum_{m \le k}v_m\), this is the corresponding sum of the pairwise relations \(dC_{nm}/dt=-\langle v_n,v_m\rangle\). This last identification assumes uniqueness of the coefficient semimartingale decomposition and is not part of Theorem \hyperref[res:Theorem-11]{11}.

\subsection{Magnitudes}\label{sec:3.9}

The classification of Theorem \hyperref[res:Theorem-11]{11} is structural, not a pricing correction. Take a meeting at which the policy rate moves by \(25\) basis points or not at all, with equal probability. The level jump has variance

\begin{equation*}
v = \tfrac14\,(0.0025)^2 \approx 1.6 \times 10^{-6} \tag{3.44}\label{eq:3.44}
\end{equation*}

in units of rate squared. A consistent model gives the jump a Gaussian law of that variance and puts on the interval after the meeting a ramp of slope \(v\) (Theorem \hyperref[res:Theorem-11]{11}(ii)), so the forward rate at time to maturity \(\tau\) past the meeting is corrected by \(v\tau\). Six weeks past the meeting, \(\tau \approx 0.115\), the correction is \(1.8 \times 10^{-7}\), about \(0.002\) basis points; one year past it, \(1.6 \times 10^{-6}\), about \(0.016\) basis points, and eight such meetings spread over a year contribute in total less than \(0.1\) basis point at the one-year point. A year of diffusion with volatility \(100\) basis points per year, \(\sigma = 0.01\), gives by the same theorem the drift correction \(\sigma^2\tau = 10^{-4}\) at \(\tau = 1\), about one basis point. The ramps that consistency forces at meetings are two orders of magnitude below the diffusive convexity at the one-year point and below quoting precision at the front end. What Theorem \hyperref[res:Theorem-11]{11} fixes is the form of the family and the law of its jumps; the bound (3.185) of Proposition \hyperref[res:Proposition-48]{48} shows that a two-point jump law and its Gaussian ramp are indistinguishable at pricing precision within one meeting interval.

\subsection{The scalar Gaussian model}\label{sec:3.10}

The following construction is the constant-shape, zero-diffusion specialization of \cite{fontana2022term} Example 3.10, eq. 3.14. Fix \(N \ge 3\), the dates (2.1), and deterministic \(v_1, \dots, v_N \ge 0\). On the product space carrying independent coordinates \(Z_n \sim N(0, v_n)\), with \(N(0, 0)\) the point mass at zero, let \(\mathcal F_t\) reveal \(Z_n\) at \(T_n\), and put

\begin{equation*}
X_t := \sum_{T_n \le t} Z_n, \qquad f(t, T) := X_t + \sum_{T_n \le t} v_n (T - T_n), \qquad r_t := f(t, t). \tag{3.45}\label{eq:3.45}
\end{equation*}

The initial curve is zero, the drift and volatility vanish, and the curve jump at \(T_n\) is

\begin{equation*}
\xi_n(u) := \mathbf 1\{u \ge T_n\}\,[Z_n + v_n (u - T_n)]. \tag{3.46}\label{eq:3.46}
\end{equation*}

Write \(B_t := \exp(\int_0^t r_s\,ds)\) and \(P(t, T) := \exp(-\int_t^T f(t, u)\,du)\). For \(h = T - T_n \ge 0\), \(\int_{T_n}^T \xi_n = h Z_n + \tfrac12 v_n h^2\), whose exponential has conditional expectation one given the earlier coordinates, so Assumption 2.2 holds, including at \(v_n = 0\); the zero drift satisfies Assumption 2.1; the regularity conditions are immediate; and

\begin{equation*}
\frac{P(t, T)}{B_t} = \prod_{T_n \le t} \exp\bigl( -(T - T_n) Z_n - \tfrac12 v_n (T - T_n)^2 \bigr) \tag{3.47}\label{eq:3.47}
\end{equation*}

is a finite product of independent unit-mean factors, so the discounted bond is a true martingale. Every current curve lies in \(S^+\); on a later interval the level and slope coefficients are \(Z_n + v_n(T_k - T_n)\) and \(v_n\), whose Gaussian mean and covariance satisfy (3.43), so the model is in the jointly Gaussian class of Theorem \hyperref[res:Theorem-11]{11}. This is the model of Sections \hyperref[sec:3.13]{3.13} and \hyperref[sec:3.15]{3.15} to 3.21.

\textbf{Remark 12 (state dimension and parameter dimension).} The scalar \(X_t\) is a time-inhomogeneous Markov state for the current curve, the affine line \(\{T \mapsto x + \sum_{T_n \le t} v_n (T - T_n)\}\), and for the conditional law of the future short-rate path, which is \(X_t\) plus independent future coordinates plus a deterministic function of time; the diffusion-factor count is zero. The short-rate jump at \(T_n\) is \(\Delta r_{T_n} = Z_n\), so for every \(t < T_n\) its risk-neutral event variance is \(v_n\), and the ideal call paying \((\Delta r_{T_n} - K)^+\) at \(T_n\), with price \(C_n(t, K) := E_Q[(B_t / B_{T_n})(\Delta r_{T_n} - K)^+ \mid \mathcal F_t]\), factors because the discount factor does not contain \(Z_n\):

\begin{equation*}
\frac{C_n(t, K)}{P(t, T_n)} = E[(Z_n - K)^+], \qquad v_n = 2\pi \left[ \frac{C_n(t, 0)}{P(t, T_n)} \right]^2. \tag{3.48}\label{eq:3.48}
\end{equation*}

As the parameters \(v_1, \dots, v_N\) vary, the vector of event variances ranges over all of \([0, \infty)^N\), and the same holds for the remaining meetings after any number of them have passed; in each fixed model it is constant on the state line, so the state image is the single point \((v_1, \dots, v_N)\). State dimension one and diffusion-factor count zero therefore give no restriction across meetings by themselves: a cross-meeting restriction has to fix the model's parameters and say which payoffs are observed, which is what the rest of this section does. The construction is known; the distinction between its fixed-model state image and the union over its parameters is not stated in the source.

\subsection{Gaussian meeting-variance cones}\label{sec:3.11}

Fix the Gaussian specialization of \cite{brace2024sofr} eqs. 35 to 39. Let \(W_1, \dots, W_d\) be independent Brownian motions, \(\gamma_{i,j}\) the ordinal meeting loadings, \(G_{i,j} := \sum_{h=1}^i \gamma_{h,j}\), \(\lambda_j \ge 0\) the decay parameters, and \(a_j(s)\) deterministic scales; put

\begin{equation*}
h_j(s, T) := a_j(s) e^{-\lambda_j (T - s)} G_{j(T) - j(s), j}, \qquad g_{n,j}(s) := a_j(s) e^{-\lambda_j (T_n - s)} \gamma_{n - j(s), j}. \tag{3.49}\label{eq:3.49}
\end{equation*}

The volatility is \(\sigma_j = h_j\), the drift is the differentiated HJM drift, and there is no curve jump at fixed maturity; the short-rate jump at \(T_n\) is the maturity discontinuity of a curve continuous in time. Its Brownian part is \(\sum_j \int_0^{T_n} g_{n,j}\,dW_j\), so its conditional variance at \(t < T_n\) is

\begin{equation*}
V_n(t) := \operatorname{Var}_Q(\Delta r_{T_n} \mid \mathcal F_t) = \sum_{j=1}^d \int_t^{T_n} a_j(s)^2 e^{-2\lambda_j (T_n - s)} \gamma_{n - j(s), j}^2\,ds. \tag{3.50}\label{eq:3.50}
\end{equation*}

\phantomsection\label{res:Proposition-13}
\textbf{Proposition 13.} Fix \(t\), the schedule, the loadings and the decay parameters, and retain the \(M\) meetings after \(t\). If \(a_j(s) = a_j\) and \(q_j = a_j^2\), then

\begin{equation*}
V(t) = A(t) q, \qquad q \in [0, \infty)^d, \qquad A_{n,j}(t) = \sum_{k = j(t)}^{n-1} \gamma_{n-k, j}^2\, H_j(T_n; \max(t, T_k), T_{k+1}), \tag{3.51}\label{eq:3.51}
\end{equation*}

where \(H_j(T; a, b) = b - a\) if \(\lambda_j = 0\) and \(H_j(T; a, b) = [e^{-2\lambda_j (T - b)} - e^{-2\lambda_j (T - a)}] / (2\lambda_j)\) otherwise. The attainable variance vectors form the cone \(A(t)[0, \infty)^d\), the linear equalities common to it are exactly \(w^\top V(t) = 0\) for \(w \in \ker A(t)^\top\), and its containing column space has dimension \(\operatorname{rank} A(t) \le d\); with three factors and \(M > 3\) meetings there are at least \(M - 3\) independent equalities. If instead \(a_j(s)^2 = q_{j,l}\) on a deterministic partition \([b_l, b_{l+1})\), the matrix has \(dK\) columns,

\begin{equation*}
A_{n,(j,l)}(t) = \int_{[t, T_n] \cap [b_l, b_{l+1}]} e^{-2\lambda_j (T_n - s)} \gamma_{n - j(s), j}^2\,ds, \tag{3.52}\label{eq:3.52}
\end{equation*}

and for one Brownian factor, \(t = 0\), \(b_k = T_k\) and \(\gamma_{1,1} \ne 0\) this \(N\)-by-\(N\) matrix is lower triangular with \(\det A = \prod_{n=1}^N \gamma_{1,1}^2 H_1(T_n; T_{n-1}, T_n) > 0\), so the cone has nonempty interior and no nonzero linear equality survives all meeting-dependent scale choices.

\emph{Proof.} Splitting the integral in (3.50) at the meeting dates and evaluating the exponential integral gives (3.51); the cone and equality statements are \((\operatorname{range} A)^\perp = \ker A^\top\). Splitting also at the scale partition gives (3.52). In the one-factor case an interval beginning at or after \(T_n\) cannot enter row \(n\), each diagonal entry is the displayed positive factor, and invertibility gives the interior and the trivial common left nullspace.\hfill$\square$

Loadings and decay are fixed in (3.51); recalibrating them changes \(A\), and Section \hyperref[sec:3.22]{3.22} determines what survives when the scales, and then the decays, vary. The derivation of (3.50) from the Brownian construction takes the maturity limits and the conditional-variance identity as given.

\subsection{The stochastic-volatility variance map}\label{sec:3.12}

Retain the stochastic-volatility specification of \cite{brace2024sofr} eqs. 35 to 41: Brownian motions \(W_j, U_j\) with

\begin{equation*}
d[W_i, W_j]_s = d[U_i, U_j]_s = \mathbf 1_{\{i = j\}}\,ds, \qquad d[W_i, U_j]_s = \mathbf 1_{\{i = j\}} \rho_j(s)\,ds, \tag{3.53}\label{eq:3.53}
\end{equation*}

and continuous nonnegative solutions, starting from one, of

\begin{equation*}
dv_j(s) = \theta_j(s)(1 - v_j(s))\,ds + \alpha_j(s) \sqrt{v_j(s)}\,dU_j(s), \tag{3.54}\label{eq:3.54}
\end{equation*}

with deterministic, locally bounded, piecewise constant coefficients, \(\theta_j, \alpha_j \ge 0\), \(|\rho_j| \le 1\), and \(E[\sup_{s \le H} v_j(s)^2] < \infty\) on finite horizons. Existence, nonnegativity and this moment bound are hypotheses; the stochastic calculus used below is cited in Section \hyperref[sec:4]{4} and enters the formal development as the fields of a hypothesis structure audited against those sources. Put

\begin{equation*}
\begin{aligned}
h_j(s, T) &:= a_j(s) e^{-\int_s^T \lambda_j(u)\,du} G_{j(T) - j(s), j}, & \sigma_j(s, T) &:= h_j(s, T) \sqrt{v_j(s)}, \\
g_{n,j}(s) &:= a_j(s) e^{-\int_s^{T_n} \lambda_j(u)\,du} \gamma_{n - j(s), j}, & b_{n,j}(s) &:= g_{n,j}(s) \int_s^{T_n} h_j(s, u)\,du, \\
R_{n,j}(u) &:= \int_u^{T_n} b_{n,j}(s) e^{-\int_u^s \theta_j(q)\,dq}\,ds, & F_{n,j}(u) &:= g_{n,j}(u)^2 + 2\rho_j g_{n,j} \alpha_j R_{n,j} + \alpha_j^2 R_{n,j}^2,
\end{aligned} \tag{3.55}\label{eq:3.55}
\end{equation*}

all unlabelled functions in the last expression being evaluated at \(u\).

\phantomsection\label{res:Proposition-14}
\textbf{Proposition 14 (conditional).} Under the preceding stochastic hypothesis, the short-rate jump and its conditional variance are

\begin{equation*}
\begin{gathered}
\Delta r_{T_n} = \sum_{j=1}^d \left[ \int_0^{T_n} b_{n,j}(s) v_j(s)\,ds + \int_0^{T_n} g_{n,j}(s) \sqrt{v_j(s)}\,dW_j(s) \right], \\
V_n(t) = \sum_{j=1}^d \int_t^{T_n} F_{n,j}(u) \left[ 1 + (v_j(t) - 1) e^{-\int_t^u \theta_j(q)\,dq} \right] du,
\end{gathered} \tag{3.56}\label{eq:3.56}
\end{equation*}

with the kernel nonnegative because \(F_{n,j} = (g_{n,j} + \rho_j \alpha_j R_{n,j})^2 + (1 - \rho_j^2)(\alpha_j R_{n,j})^2\). For fixed parameters, dates and observation time, with \(E_j(t, u) := e^{-\int_t^u \theta_j}\),

\begin{equation*}
A_{n,j}(t) := \int_t^{T_n} F_{n,j}(u) E_j(t, u)\,du, \qquad C_n(t) := \sum_j \int_t^{T_n} F_{n,j}(u) (1 - E_j(t, u))\,du, \qquad V(t) = C(t) + A(t) v(t), \tag{3.57}\label{eq:3.57}
\end{equation*}

with \(C_n(t) \ge 0\) and \(A_{n,j}(t) \ge 0\). Hence \(V(t) \in C(t) + A(t)[0, \infty)^d\), every \(w \in \ker A(t)^\top\) gives \(w^\top (V(t) - C(t)) = 0\), and the containing affine space has dimension at most \(d\).

\emph{Proof.} The maturity difference gives the jump. Since \(R' = \theta R - b\) and \(R(T_n) = 0\), the product rule gives

\begin{equation*}
\int_t^{T_n} b_{n,j}(s) v_j(s)\,ds = R_{n,j}(t) v_j(t) + \int_t^{T_n} \theta_j R_{n,j}\,ds + \int_t^{T_n} \alpha_j R_{n,j} \sqrt{v_j}\,dU_j, \tag{3.58}\label{eq:3.58}
\end{equation*}

so after subtracting the conditional mean the future random part is \(\sum_j \int_t^{T_n} \sqrt{v_j}\,(g_{n,j}\,dW_j + \alpha_j R_{n,j}\,dU_j)\); the conditional isometry and (3.53) give the kernel \(F\), and the integrating-factor solution of (3.54) gives \(E[v_j(u) \mid \mathcal F_t] = 1 + (v_j(t) - 1) E_j(t, u)\), proving (3.56). The rest is algebra.\hfill$\square$

The random HJM drift and its covariance with the Brownian term are in \(F\); replacing \(F\) by \(g^2\) is valid only when every volatility-of-volatility coefficient vanishes. The short-rate jump is measurable immediately before its meeting in this continuous filtration and is not an additional curve jump.

\subsection{Identification from ideal bond calls}\label{sec:3.13}

In the model of Section \hyperref[sec:3.10]{3.10} at observation time zero, for expiry \(S\), bond maturity \(U > S\) and strike \(K > 0\), define the ideal European bond call and the cumulative variance

\begin{equation*}
C(S, U, K) := E_Q[B_S^{-1} (P(S, U) - K)^+], \qquad q(S, U) := (U - S)^2 \sum_{T_n \le S} v_n. \tag{3.59}\label{eq:3.59}
\end{equation*}

The forward-measure calculation is that of \cite{fontana2022term} Lemma 4.11 and Proposition 4.12.

\phantomsection\label{res:Proposition-15}
\textbf{Proposition 15.} For \(q > 0\),

\begin{equation*}
C(S, U, K) = \Phi(d_1) - K \Phi(d_2), \qquad d_1 = \frac{-\log K + q/2}{\sqrt q}, \qquad d_2 = d_1 - \sqrt q, \tag{3.60}\label{eq:3.60}
\end{equation*}

and \(C(S, U, K) = (1 - K)^+\) for \(q = 0\); in particular \(C(S, U, 1) = 2\Phi(\sqrt q / 2) - 1\). At one expiry the complete call surface identifies exactly the cumulative variance \(\sum_{T_n \le S} v_n\): two parameter vectors give the same prices for all \(U > S\) and \(K > 0\) if and only if those sums agree, so for \(N = 3\) and \(\varepsilon > 0\) the vectors \((\varepsilon, 2\varepsilon, 3\varepsilon)\) and \((2\varepsilon, \varepsilon, 3\varepsilon)\) give the same full surface at \(S = T_3\) with different first and second event variances. Conversely, expiries \(T_n \le S_n < T_{n+1}\) with maturities \(U_n > S_n\) and the strike-one prices at each \(S_n\) identify every \(v_n\), including zero:

\begin{equation*}
z_0 := 0, \qquad z_n := \frac{4}{(U_n - S_n)^2} \left[ \Phi^{-1}\left( \frac{1 + C(S_n, U_n, 1)}{2} \right) \right]^2, \qquad v_n = z_n - z_{n-1}. \tag{3.61}\label{eq:3.61}
\end{equation*}

\emph{Proof.} By (3.47), \(E[B_S^{-1}] = E[B_S^{-1} P(S, U)] = 1\), so \(dQ^S = B_S^{-1} dQ\) is a probability measure on \(\mathcal F_S\) and the discounted payoff is integrable. Completing the square in each revealed coordinate preserves independence and shifts \(Z_n\) to mean \(-v_n (S - T_n)\); substitution in the bond formula cancels the meeting-date terms and gives \(\log P(S, U) \sim N(-q/2, q)\) under \(Q^S\). Integrating the lognormal positive part proves (3.60), and strict monotonicity of \(\Phi\) makes the strike-one price injective in \(q\). Equal sums give equal surfaces. At \(S_n\) exactly the first \(n\) meetings enter \(q\), so (3.61) recovers the cumulative sums and their differences.\hfill$\square$

\subsection{Exact support with zero variance mean reversion}\label{sec:3.14}

Specialize the setting of Section \hyperref[sec:3.12]{3.12} to \(\theta_j = 0\) and constant \(\alpha_j \ge 0\), so that the variance states solve

\begin{equation*}
dv_j(s) = \alpha_j \sqrt{v_j(s)}\,dU_j(s), \qquad v_j(0) = 1. \tag{3.62}\label{eq:3.62}
\end{equation*}

Fix \(0 < t < T_N\), retain the future meetings, take \(A(t)\) from (3.57) with the full kernel, and put \(J := \{j : \alpha_j > 0\}\), \(J' := \{j \in J : A_{\cdot j}(t) \ne 0\}\).

\phantomsection\label{res:Proposition-16}
\textbf{Proposition 16.} Assume adapted, nonnegative solutions of (3.62) with almost surely continuous paths on a filtered probability space satisfying the usual conditions, driven by Brownian motions with the covariations (3.53); the second-moment bound of Section \hyperref[sec:3.12]{3.12} is not needed. For nonnegative \(\lambda_j\),

\begin{equation*}
E_Q\left[ e^{-\sum_j \lambda_j v_j(t)} \right] = \exp\left( -\sum_j \frac{\lambda_j}{1 + \alpha_j^2 t \lambda_j / 2} \right). \tag{3.63}\label{eq:3.63}
\end{equation*}

If \(j \notin J\) then \(v_j(t) = 1\); if \(j \in J\) then, with \(c_j = \alpha_j^2 t / 2\), \(v_j(t)\) has the law of \(\sum_{k=1}^{L_j} E_{j,k}\) with \(L_j \sim \operatorname{Poisson}(1 / c_j)\) and \(E_{j,k} \sim \operatorname{Exponential}(\text{mean } c_j)\), independently across \(j\), an empty sum being zero. The support and affine hull of the meeting-variance vector are

\begin{equation*}
\operatorname{supp} V(t) = \sum_{j \notin J} A_{\cdot j}(t) + A_J(t)[0, \infty)^J, \qquad \operatorname{aff} \operatorname{supp} V(t) = \sum_{j \notin J} A_{\cdot j}(t) + \operatorname{range} A_J(t), \tag{3.64}\label{eq:3.64}
\end{equation*}

so the affine-hull dimension is \(\operatorname{rank} A_J(t) \le |J|\) and the affine equalities about the deterministic contribution are exactly \(w \in \ker A_J(t)^\top\); moreover

\begin{equation*}
\operatorname{Cov}_Q(V(t)) = \sum_{j \in J} \alpha_j^2 t\, A_{\cdot j}(t) A_{\cdot j}(t)^\top, \qquad Q\Bigl( V(t) = \sum_{j \notin J} A_{\cdot j}(t) \Bigr) = \exp\Bigl( -\sum_{j \in J'} \frac{2}{\alpha_j^2 t} \Bigr), \tag{3.65}\label{eq:3.65}
\end{equation*}

with \(\operatorname{rank} \operatorname{Cov}_Q(V(t)) = \operatorname{rank} A_J(t)\). For deterministic \(0 \le s < t\) and \(J_s := \{j \in J : v_j(s) > 0\}\), almost surely the conditional law of \(V(t)\) given \(\mathcal F_s\) has support \(\sum_{j \notin J} A_{\cdot j}(t) + A_{J_s}(t)[0, \infty)^{J_s}\), affine hull \(\sum_{j \notin J} A_{\cdot j}(t) + \operatorname{range} A_{J_s}(t)\), mean \(A(t) v(s)\), and

\begin{equation*}
\operatorname{Cov}_Q(V(t) \mid \mathcal F_s) = \sum_{j \in J} \alpha_j^2 (t - s) v_j(s)\, A_{\cdot j}(t) A_{\cdot j}(t)^\top, \qquad Q\Bigl( V(t) = \sum_{j \notin J} A_{\cdot j}(t) \Bigm| \mathcal F_s \Bigr) = \exp\Bigl( -\sum_{j \in J'} \frac{2 v_j(s)}{\alpha_j^2 (t - s)} \Bigr); \tag{3.66}\label{eq:3.66}
\end{equation*}

a coordinate that is zero at \(s\) stays zero and supplies no later support direction. For the printed constant parameters of \cite{brace2024sofr} Table 2, \(\theta = (0, 0, 0)\) and \(\alpha = (1.57, 0.82, 0)\), so \(\operatorname{supp} V(t) = A_{\cdot 3}(t) + \operatorname{cone}\{A_{\cdot 1}(t), A_{\cdot 2}(t)\}\), whose affine hull has dimension at most two; rank two is not asserted.

\emph{Proof.} For terminal time \(t\), set \(q_j(s) = \lambda_j / [1 + \alpha_j^2 \lambda_j (t - s) / 2]\); Itô's formula gives \(\exp(-\sum_j q_j(s) v_j(s))\) as its initial value plus stochastic integrals whose drift cancels, and stopping at the hitting times of compact intervals makes the stopped process a bounded martingale, so the exponential is a martingale on \([0, t]\), proving (3.63). The compound Poisson sum has the same Laplace transform, and multivariate Laplace uniqueness, \cite{li2019continuous} Theorem 1.1, identifies the law; its support is \([0, \infty)\) with an atom at zero and its variance is \(\alpha_j^2 t\). The continuous linear map in (3.57) gives (3.64), with \cite{fontana2025extended} Lemma A.3 the general continuous-image statement; nonnegativity of the columns makes the image cone closed. Independence gives (3.65), and repeating the bounded-martingale argument from \(v(s)\) gives the regular conditional distribution and (3.66).\hfill$\square$

The existence of the solutions of (3.62) is a hypothesis of the statement and is not constructed; the stochastic calculus in the proof is the cited standard one (Section \hyperref[sec:4]{4}), and the short-rate jump whose variance vector \(V(t)\) collects is the jump of the forward-rate construction of Section \hyperref[sec:3.12]{3.12} with the differentiated HJM drift. At \(t = 0\) and at \(s = t\) the laws are points; the support can have faces and an atom at its deterministic contribution. Proposition \hyperref[res:Proposition-26]{26} extends the exact-support calculation to nonnegative mean reversion and piecewise-constant volatility of volatility.

\subsection{Deterministic price intervals}\label{sec:3.15}

Use the model and the meeting-separated strike-one calls of Proposition \hyperref[res:Proposition-15]{15}, with \(h_n = U_n - S_n\), \(z_n = \sum_{i \le n} v_i\) and \(C_n(v) = 2\Phi(h_n \sqrt{z_n} / 2) - 1\). For deterministic intervals \(0 \le a_n \le b_n \le 1\) with \(a_n < 1\), define

\begin{equation*}
H_n(c) := \frac{4}{h_n^2} \left[ \Phi^{-1}\left( \frac{1 + c}{2} \right) \right]^2, \qquad l_n := H_n(a_n), \qquad u_n := H_n(b_n) \text{ if } b_n < 1, \quad u_n := +\infty \text{ if } b_n = 1, \tag{3.67}\label{eq:3.67}
\end{equation*}

where \(+\infty\) means no finite upper restriction, and set \(L_n = \max_{i \le n} l_i\), \(R_n = \min_{j \ge n} u_j\), \(L_0 = R_0 = 0\).

\phantomsection\label{res:Proposition-17}
\textbf{Proposition 17.} A finite nonnegative vector \(v\) has \(C_n(v) \in [a_n, b_n]\) for every \(n\) if and only if \(l_n \le z_n \le u_n\) for every \(n\), and such a vector exists if and only if \(L_n \le R_n\) for \(n = 1, \dots, N\). When feasible, the attainable values of \(v_n\) form the interval with endpoints

\begin{equation*}
\underline v_n = \max(0, L_n - R_{n-1}) \ (\text{or } 0 \text{ if } R_{n-1} = +\infty), \qquad \overline v_n = R_n - L_{n-1} \ (\text{or } +\infty \text{ if } R_n = +\infty), \tag{3.68}\label{eq:3.68}
\end{equation*}

every finite endpoint and intermediate value being attained by a full feasible vector; the coordinate intervals are projections, and choosing their values independently need not preserve the joint constraints. On a price range \([0, \beta_n]\) with \(\beta_n < 1\), \(H_n\) is Lipschitz with derivative \(H_n'(c) = 4y / [h_n^2 \varphi(y)]\), \(y = \Phi^{-1}((1 + c) / 2)\), \(\varphi\) the standard normal density, so two price vectors in these ranges differing by at most \(\varepsilon_n\) coordinatewise give

\begin{equation*}
|v_n - v_n'| \le k_n \varepsilon_n + k_{n-1} \varepsilon_{n-1}, \qquad k_n = \frac{4 y_n}{h_n^2 \varphi(y_n)}, \qquad k_0 = \varepsilon_0 = 0, \tag{3.69}\label{eq:3.69}
\end{equation*}

and there is no finite Lipschitz constant up to price one.

\emph{Proof.} Strict monotonicity of \(C_n\) makes each price interval equivalent to \(l_n \le z_n \le u_n\). Since the cumulative sums are nondecreasing, every feasible sequence satisfies \(L_n \le z_n \le R_n\); conversely \(z_n = L_n\) is finite, nondecreasing and feasible exactly when \(L_n \le R_n\). The differences \(z_n - z_{n-1}\) give (3.68): \(z_{n-1} = \min(L_n, R_{n-1})\) attains the lower endpoint, \(z_n = R_n\), or any large finite value when \(R_n = +\infty\), the upper one, and convexity of the constraints fills the intermediate values. Differentiating \(H_n\) and using that \(y / \varphi(y)\) increases for \(y \ge 0\), the mean value theorem gives (3.69); \(H_n\) and its derivative diverge as \(c \uparrow 1\).\hfill$\square$

\subsection{Identification from compounded-rate futures and calls}\label{sec:3.16}

In the model of Section \hyperref[sec:3.10]{3.10} at time zero, for \(0 \le a < b\), \(\delta = b - a\), define the ideal futures rate for the continuously compounded backward-looking rate, with the undiscounted conditional-expectation convention of \cite{skov2021dynamic} (13), (19) and (20), and the European cash-call price with expiry \(0 \le S \le b\) and strike \(K > 0\):

\begin{equation*}
\begin{gathered}
G_t(a, b) := E_Q\left[ \exp\left( \int_a^b r_u\,du \right) \Bigm| \mathcal F_t \right], \qquad F(t; a, b) := \frac{G_t(a, b) - 1}{\delta}, \\
C(S, a, b, K) := E_Q[B_S^{-1} (G_S(a, b) - K)^+].
\end{gathered} \tag{3.70}\label{eq:3.70}
\end{equation*}

Dividing \(C\) by \(\delta\) gives the price of a call on \(F(S; a, b)\) with strike \((K - 1) / \delta\). For each meeting set

\begin{equation*}
\begin{gathered}
w_i := (b - T_i)^+ - (a - T_i)^+, \qquad d_i := \tfrac12 \bigl( (b - T_i)^{+2} - (a - T_i)^{+2} \bigr), \qquad h_i := d_i + \tfrac12 w_i^2, \\
j_i := \mathbf 1\{T_i \le S\} w_i (S - T_i), \qquad k_i := \mathbf 1\{T_i \le S\} w_i^2,
\end{gathered} \tag{3.71}\label{eq:3.71}
\end{equation*}

and \(p := \sum_i h_i v_i\), \(z := \sum_i j_i v_i\), \(q := \sum_i k_i v_i\), \(m := e^{p - z}\). Write, for \(m > 0\), \(s \ge 0\) and \(K > 0\),

\begin{equation*}
\begin{gathered}
c_K(m, s) := m \Phi(D_1) - K \Phi(D_2), \qquad D_1 = \frac{\log(m / K) + s^2 / 2}{s}, \qquad D_2 = D_1 - s \quad (s > 0); \\
c_K(m, 0) := (m - K)^+.
\end{gathered} \tag{3.72}\label{eq:3.72}
\end{equation*}

The two-strike argument is the following lemma.

\phantomsection\label{res:Lemma-18}
\textbf{Lemma 18.} For any fixed \(K_* > 1\), the map \((m, s) \mapsto (c_1(m, s), c_{K_*}(m, s))\) is injective on \(\{m \ge 1, s \ge 0\}\).

\emph{Proof.} For \(s > 0\),

\begin{equation*}
\frac{\partial c_K}{\partial m} = \Phi(D_1) > 0, \qquad \frac{\partial c_K}{\partial s} = m \varphi(D_1) > 0, \qquad \left( \frac{\varphi}{\Phi} \right)'(x) < 0. \tag{3.73}\label{eq:3.73}
\end{equation*}

Along a level set \(c_1(m, s) = \text{const} > 0\), implicit differentiation with the last inequality shows that \(c_{K_*}\) is strictly increasing in \(s\), so the pair determines \((m, s)\) on \(\{c_1 > 0\}\). A zero strike-one price forces \(m = 1\) and \(s = 0\), since \(c_1(m, s) \ge (m - 1)^+\) with equality only at \(s = 0\); continuity at \(s = 0\) handles the remaining boundary.\hfill$\square$

\phantomsection\label{res:Proposition-19}
\textbf{Proposition 19.} (a) \(G(a, b)\) is an integrable \(Q\)-martingale,

\begin{equation*}
G_t(a, b) = \exp\Bigl( \sum_{T_i \le t} w_i Z_i + \sum_i d_i v_i + \tfrac12 \sum_{T_i > t} w_i^2 v_i \Bigr), \qquad 1 + \delta F(0; a, b) = e^p, \tag{3.74}\label{eq:3.74}
\end{equation*}

\(\log G_S \sim N(\log m - q/2, q)\) under \(dQ^S = B_S^{-1} dQ\), and \(C(S, a, b, K) = c_K(m, \sqrt q)\). The complete surface \(K > 0\) determines \((m, q)\), since \(m = \lim_{K \downarrow 0} C(S, a, b, K)\) and \(q = 4[\Phi^{-1}((1 + C(S, a, b, m) / m) / 2)]^2\); with the futures quote it determines \((p, z, q)\), and by Lemma \hyperref[res:Lemma-18]{18} the two calls at \(K = 1\) and any fixed \(K_* > 1\) do the same, including at \(q = 0\).

(b) Fix \(L\) triples \((S_\ell, a_\ell, b_\ell)\) and let \(M\) be the \(3L\)-by-\(N\) matrix whose rows for contract \(\ell\) are \(h^{(\ell)}, j^{(\ell)}, k^{(\ell)}\). For \(v, v' \in [0, \infty)^N\), equality of the \(L\) futures quotes and either the \(L\) complete surfaces or the \(2L\) fixed-strike calls is equivalent to \(Mv = Mv'\). The panel identifies every \(v\) if and only if \(\operatorname{rank} M = N\); if \(\operatorname{rank} M < N\), every strictly positive \(v\) has a distinct strictly positive \(v'\) with the same observations, and the set compatible with a model-generated panel is \(\{x \in [0, \infty)^N : Mx = Mv\}\).

(c) If one accrual interval lies after all meetings and \(T_n \le S_n < T_{n+1}\), with \(S_N = T_N\), then \(q_n = \delta^2 \sum_{i \le n} v_i\), so \(v_n = (q_n - q_{n-1}) / \delta^2\) with \(q_0 = 0\). Futures alone also identify: with \(a_n = T_n\), \(T_n < b_n < T_{n+1}\) for \(n < N\), \(b_N > T_N\), and \(\delta_n = b_n - T_n\),

\begin{equation*}
v_n = \frac{p_n - \delta_n \sum_{i < n} (b_n - T_i) v_i}{\delta_n^2}, \qquad p_n = \log\bigl( 1 + \delta_n F(0; T_n, b_n) \bigr). \tag{3.75}\label{eq:3.75}
\end{equation*}

(d) Full later-window panels need not identify: for \(N = 3\), \(T_i = i\) and \(\varepsilon > 0\), the vectors \(v = \varepsilon(2, 3, 2)\) and \(v' = \varepsilon(3, 1, 3)\) give the same futures quote for every \(3 \le a < b\) and the same call price for every \(3 \le S \le a < b\) and \(K > 0\), although their event variances differ.

\emph{Proof.} Integrating the short rate in (3.45) gives \(\int_a^b r_u\,du = \sum_i (w_i Z_i + d_i v_i)\); independence and the Gaussian exponential moment give (3.74), also during accrual, and the tower property the martingale assertion. The discount factor tilts every revealed coordinate to mean \(-v_i (S - T_i)\) without changing variances or independence, so \(\log G_S\) has the displayed \(Q^S\)-law with first moment \(m\), which gives integrability and the price; monotone convergence as \(K \downarrow 0\) recovers \(m\), and at \(K = m\) normal symmetry gives \(C / m = 2\Phi(\sqrt q / 2) - 1\). Since \(h_i - j_i \ge 0\), \(m \ge 1\), so Lemma \hyperref[res:Lemma-18]{18} applies. Each contract thus identifies \((p, z, q)\), three linear forms in \(v\), which is (b); full column rank gives injectivity, and under rank deficiency a kernel vector gives a small two-sided perturbation of every strictly positive \(v\). For (c), a post-meeting accrual interval has \(w_i = \delta\), so the rows of \(k\) at successive expiries are lower triangular; with \(a_n = T_n\) the futures coefficient vanishes for \(i > n\) and equals \(\delta_n (b_n - T_i)\) for \(i \le n\), giving (3.75). For (d), every triple's coefficients depend only on \(\sum_i v_i\) and \(\sum_i T_i v_i\), which are \(7\varepsilon\) and \(14\varepsilon\) for both vectors.\hfill$\square$

\subsection{Contracts settled after a date}\label{sec:3.17}

Replace the continuously compounded convention of Section \hyperref[sec:3.16]{3.16} by exact discrete bond-implied compounding and the European option by an American cash call. Fix \(A \le S \le a < b\), a partition \(a = u_0 < \dots < u_J = b\), \(d_j = u_{j+1} - u_j\), and define the bond-implied simple overnight rates, the exact compounded rate, its futures process and the cash call

\begin{equation*}
\begin{gathered}
L_j = \frac{P(u_j, u_{j+1})^{-1} - 1}{d_j}, \qquad R = \frac{\prod_{j=0}^{J-1} (1 + d_j L_j) - 1}{\delta}, \qquad F(t; a, b) = E_Q[R \mid \mathcal F_t], \\
U(A, S, a, b, K) = \sup_{\tau \in \mathcal T[A, S]} E_Q\bigl[ B_\tau^{-1} (F(\tau; a, b) - K)^+ \bigr],
\end{gathered} \tag{3.76}\label{eq:3.76}
\end{equation*}

with \(\delta = b - a\) and \(\mathcal T[A, S]\) the stopping times valued in \([A, S]\); this is the exact product convention of \cite{skov2021dynamic}, and exercise pays cash rather than entering a futures position.

The following result shows, for every contract settled after a fixed date and with any number of unrevealed meetings, that the revealed variances enter only through their sum and their date-weighted sum, and that the compatible set is a polytope with explicit sharp coordinate bounds. Corollaries \hyperref[res:Corollary-21]{21} and \hyperref[res:Corollary-22]{22} specialize it to American cash exercise after, and across, the last meeting.

In the model of Section \hyperref[sec:3.10]{3.10} fix a deterministic date \(A \ge 0\), let \(k = j(A)\) be the number of meetings at or before \(A\), and write

\begin{equation*}
V_k = \sum_{i \le k} v_i, \qquad H_k = \sum_{i \le k} T_i v_i, \qquad h_i = A - T_i \ (i \le k), \qquad y = r_A, \qquad Z_{>k} = (Z_{k+1}, \dots, Z_N). \tag{3.77}\label{eq:3.77}
\end{equation*}

\phantomsection\label{res:Proposition-20}
\textbf{Proposition 20.} (a) For \(t \ge A\) and \(U \ge t\),

\begin{equation*}
f(t, U) = y + V_k (U - A) + \sum_{k < i \le j(t)} [Z_i + v_i (U - T_i)], \qquad \frac{B_t}{B_A} = \exp\Bigl( \int_A^t r_s\,ds \Bigr), \tag{3.78}\label{eq:3.78}
\end{equation*}

and \(Z_i = r_{T_i} - r_{T_i-}\) for \(i > k\). Every forward rate, short rate, bond price, bank-account ratio and bond-implied fixing at or after \(A\) is therefore a fixed Borel function of the state \((y, Z_{>k})\), of \((V_k, v_{k+1}, \dots, v_N)\), and of the dates \(T_{k+1}, \dots, T_N\) and \(A\); the individual \(v_1, \dots, v_k\) and \(T_1, \dots, T_k\) enter only through \(V_k\). Conversely, for \(t \ge A\), \(\sigma(f(s, \cdot) : A \le s \le t) = \sigma(y, Z_i : k < i \le j(t))\).

(b) \(B_A^{-1} = \exp(-\sum_{i \le k} [h_i Z_i + v_i h_i^2 / 2])\) has \(E_Q[B_A^{-1}] = 1\). Under \(\tilde Q := B_A^{-1} \cdot Q\) the coordinates stay independent, \(Z_i \sim N(-v_i h_i, v_i)\) for \(i \le k\) and \(Z_i \sim N(0, v_i)\) for \(i > k\), so \(y \sim N(0, V_k)\) is independent of \(Z_{>k}\), and for every Borel \(\Phi \ge 0\)

\begin{equation*}
E_Q[B_A^{-1} \Phi(y, Z_{>k})] = E[\Phi(Y, Z')], \qquad Y \sim N(0, V_k), \quad Z'_i \sim N(0, v_i) \ (i > k), \quad \text{all independent}, \tag{3.79}\label{eq:3.79}
\end{equation*}

and likewise for real \(\Phi\) with \(E_Q[B_A^{-1} |\Phi|]\) finite. If \(V_k > 0\), the residuals \(\xi_i := Z_i + v_i h_i - (v_i / V_k) y\), \(i \le k\), are under \(\tilde Q\) jointly Gaussian and independent of \((y, Z_{>k})\), and \(\sigma(Z_1, \dots, Z_k) = \sigma(y, \xi_1, \dots, \xi_k)\); if \(V_k = 0\) then \(Z_1 = \dots = Z_k = 0\) and \(y = 0\) almost surely.

(c) Let \(S \ge A\) be deterministic and \(X = \Phi(y, Z_{k+1}, \dots, Z_{j(S)})\) a cash flow paid at \(S\), with \(\Phi\) Borel and allowed to depend on \((V_k, v_{k+1}, \dots, v_N)\) and the contract data. If \(E_Q[B_S^{-1} |X|] < \infty\), then

\begin{equation*}
E_Q[B_S^{-1} X] = E\Bigl[ \exp\Bigl( -\int_A^S \tilde r_s\,ds \Bigr) \Phi(Y, Z') \Bigr], \tag{3.80}\label{eq:3.80}
\end{equation*}

with \(\tilde r\) the short rate of (3.78) evaluated at \((Y, Z')\); the right side depends on the revealed variances only through \(V_k\). This covers every European payoff on bonds, bond-implied fixings and compounded rates over accrual windows \([a, b]\) with \(a \ge A\), windows crossing unrevealed meetings included, and every partition.

(d) Let \(A \le S\) and let \(\Phi : [A, S] \times \mathbb R \times \mathbb R^{N-k} \to \mathbb R\) be Borel with \(\Phi(t, y, z)\) depending on \(z\) only through the coordinates \(i \le j(t)\), so that \(X_t := \Phi(t, y, Z_{>k})\) is a function of the curve on \([A, t]\), and assume a Borel \(D \ge 0\) with \(|\exp(-\int_A^t r_s\,ds)\, \Phi(t, y, Z_{>k})| \le D(y, Z_{>k})\) for \(t \in [A, S]\) and \(E[D(Y, Z')] < \infty\). Then the time-zero American value

\begin{equation*}
U := \sup_{\tau \in \mathcal T[A, S]} E_Q[B_\tau^{-1} X_\tau] = \sup_{\rho \in \tilde{\mathcal T}[A, S]} E\Bigl[ \exp\Bigl( -\int_A^\rho \tilde r_s\,ds \Bigr) \Phi(\rho, Y, Z') \Bigr] \tag{3.81}\label{eq:3.81}
\end{equation*}

is finite, where \(\tilde{\mathcal T}[A, S]\) are the stopping times valued in \([A, S]\) of the filtration generated by \(Y\) and by the \(Z'_i\) revealed at \(T_i\). Hence \(U\) depends on \((v_1, \dots, v_k)\) only through \(V_k\). Every cash call (3.76), in particular those of Corollaries \hyperref[res:Corollary-21]{21} and \hyperref[res:Corollary-22]{22}, with any exercise window \([A, S]\) crossing any number of unrevealed meetings and any accrual window \([a, b]\) with \(a \ge A\), satisfies the domination and is included; no formula for the value is asserted.

(e) For an accrual window \(a < b\) with \(a \ge T_k\), a partition \(a = u_0 < \dots < u_J = b\) with no meeting strictly inside a fixing interval, and the exact compounding of (3.76), \(\prod_j (1 + d_j L_j) = \exp(\int_a^b r_s\,ds)\) and

\begin{equation*}
F(0; a, b) = \frac{1}{\delta} \Bigl[ \exp\Bigl( \delta (b V_k - H_k) + \delta \sum_{k < i,\ T_i \le a} v_i (b - T_i) + \sum_{a < T_i \le b} v_i (b - T_i)^2 \Bigr) - 1 \Bigr], \tag{3.82}\label{eq:3.82}
\end{equation*}

so the time-zero quote depends on \((v_1, \dots, v_k)\) only through \((V_k, H_k)\), also for windows that cross unrevealed meetings at partition points.

(f) Fix \(v_{k+1}, \dots, v_N\). Two vectors \((v_1, \dots, v_k)\) and \((v'_1, \dots, v'_k)\) with the same \(V_k\) and \(H_k\) give the same quote (3.82) for every window and partition in (e), the same price (3.80) for every contract in (c) and the same value (3.81) for every contract in (d). If the observation family contains two quotes (3.82) with \(a \ge T_k\) and distinct \(b\), then conversely equal observations force equal \((V_k, H_k)\), so the set of \((v_1, \dots, v_k)\) consistent with exact observation of the family is the polytope

\begin{equation*}
\Pi(V_k, H_k) = \Bigl\{ x \in [0, \infty)^k : \sum_i x_i = V_k,\ \sum_i T_i x_i = H_k \Bigr\}, \tag{3.83}\label{eq:3.83}
\end{equation*}

nonempty if and only if \(T_1 V_k \le H_k \le T_k V_k\), over which each coordinate attains its minimum and maximum at a point with at most two nonzero coordinates. For \(k = 3\), with \(V = V_3\) and \(H = H_3\), the sharp attained bounds are

\begin{equation*}
\begin{aligned}
s_* &= \min\Bigl\{ \frac{H - T_1 V}{T_2 - T_1}, \frac{T_3 V - H}{T_3 - T_2} \Bigr\} \ge 0, \qquad x_2 \in [0, s_*], \\
x_1 &\in \Bigl[ \frac{T_3 V - H - (T_3 - T_2) s_*}{T_3 - T_1}, \frac{T_3 V - H}{T_3 - T_1} \Bigr], \qquad
x_3 \in \Bigl[ \frac{H - T_1 V - (T_2 - T_1) s_*}{T_3 - T_1}, \frac{H - T_1 V}{T_3 - T_1} \Bigr].
\end{aligned} \tag{3.84}\label{eq:3.84}
\end{equation*}

For dates \((1, 2, 3)\) and every \(\varepsilon > 0\) the vectors \(\varepsilon(2, 3, 2)\) and \(\varepsilon(3, 1, 3)\), completed by any common \(v_4, \dots, v_N\), are two such points, with bounds \(x_1, x_3 \in [0, 7\varepsilon/2]\) and \(x_2 \in [0, 7\varepsilon]\).

\emph{Proof.} (a) For \(t \ge A\), \(j(t) \ge k\), \(y = \sum_{i \le k} [Z_i + v_i h_i]\) and \(\sum_{i \le k} v_i (U - T_i) = \sum_{i \le k} v_i h_i + V_k (U - A)\), which gives (3.78); the path of \(r\) on \([A, t]\) is right-continuous, piecewise affine and jumps by \(Z_i\) at \(T_i\), which is the jump identity. For the sigma-algebras, \(y = f(A, A)\), each \(Z_i\) with \(A < T_i \le t\) equals \(f(T_i, T_i) - f(\max(A, T_{i-1}), T_i)\), and each \(f(s, U)\) with \(A \le s \le t\) is a function of \((y, Z_i : k < i \le j(s))\) by (3.78).

(b) Only \(i \le k\) enter \(\int_0^A r_s\,ds = \sum_{i \le k} [h_i Z_i + v_i h_i^2 / 2]\). The density is a product of unit-mean factors, and completing the square tilts the earlier coordinates and leaves the later ones, which gives the laws, the independence and (3.79). The vector \((y, \xi)\) is a linear image of a Gaussian vector with \(\operatorname{Cov}_{\tilde Q}(\xi_i, y) = v_i - (v_i / V_k) V_k = 0\), so \(\xi\) is independent of \(y\) and, being a function of \(Z_{\le k}\), of \(Z_{>k}\); \(Z_i = \xi_i + (v_i / V_k) y - v_i h_i\) gives the sigma-algebra identity, and \(V_k = 0\) forces every earlier variance to vanish.

(c) \(B_S^{-1} = B_A^{-1} (B_S / B_A)^{-1}\), and \((B_S / B_A)^{-1} X\) is a Borel function of \((y, Z_{>k})\) by (3.78); apply (3.79) to its absolute value and to itself. Fixings are ratios of bond prices at times at or after \(A\), and the compounded rate is a finite product of them.

(d) For \(\tau \in \mathcal T[A, S]\), \(B_\tau^{-1} X_\tau = B_A^{-1} Y_\tau\) with \(Y_t := (B_t / B_A)^{-1} \Phi(t, y, Z_{>k})\) and \(|Y_\tau| \le D(y, Z_{>k})\), so \(U = \sup_\tau E_{\tilde Q}[Y_\tau]\) is finite by (3.79). For \(t \ge A\) the completed filtration is \(\mathcal F_t = \sigma(y, \xi, Z_i : k < i \le j(t)) \vee \mathcal N\), with \(y\), \(\xi\) and \(Z_{>k}\) independent under \(\tilde Q\), and \(Y\) is adapted to \(\mathcal G_t := \sigma(y, Z_i : k < i \le j(t)) \vee \mathcal N\); both filtrations are right-continuous because \(j\) is right-continuous with finitely many jumps. Realize \(\tilde Q\) on the product of the laws of \(y\), \(\xi\) and \(Z_{>k}\). Given \(\tau\), for each rational \(q \in [A, S]\) and for \(q = S\) there is \(C_q \in \sigma(y, \xi, Z_i : k < i \le j(q))\) with \(\tilde Q(\{\tau \le q\} \,\triangle\, C_q) = 0\); by Fubini, for \(\xi\) outside a null set every section \(\tau_\xi\) has \(\{\tau_\xi \le q\} \in \mathcal G_q\), and right-continuity extends this to every \(t\), so \(\tau_\xi\) is a \(\mathcal G\)-stopping time valued in \([A, S]\), a rule on the auxiliary space. By Fubini,

\begin{equation*}
E_{\tilde Q}[Y_\tau] = \int E[Y_{\tau_\xi}]\,(\text{law of } \xi)(d\xi) \le \sup_{\rho \in \tilde{\mathcal T}[A, S]} E\Bigl[ \exp\Bigl( -\int_A^\rho \tilde r_s\,ds \Bigr) \Phi(\rho, Y, Z') \Bigr], \tag{3.85}\label{eq:3.85}
\end{equation*}

and conversely every \(\rho \in \tilde{\mathcal T}[A, S]\) is an \(\mathcal F\)-stopping time on the model with \(E_{\tilde Q}[Y_\rho]\) equal to the auxiliary expectation by (3.79); when \(V_k = 0\) the residuals are dropped and \(\tau\) is almost surely its own section. The auxiliary problem is specified by \((V_k, v_{k+1}, \dots, v_N)\), the later dates, \(A\), \(S\) and \(\Phi\). For the cash calls (3.76), \(R\) is a Borel function of \((y, Z_{>k})\) by (a) and the telescoping in (e) on each fixing interval, \(F(t; a, b)\) has the version \(E_Q[R \mid \mathcal G_t]\) because \(\xi\) is independent of \((y, Z_{>k})\) under \(Q\) as well, and the bound \(\max_{0 \le u \le d} e^{-zu - qu^2/2} \le e^{d|z|}\) for the discount factors gives \(D = C \exp(c(|y| + \sum_{i > k} |Z_i|))\).

(e) On a fixing interval with no meeting inside, \(r_u = f(u_j, u)\), so \(P(u_j, u_{j+1})^{-1} = B_{u_{j+1}} / B_{u_j}\) and the product telescopes. With \(m = (a + b) / 2\), \(\int_a^b r_s\,ds = \sum_{T_i \le a} [\delta Z_i + v_i \delta (m - T_i)] + \sum_{a < T_i \le b} [(b - T_i) Z_i + v_i (b - T_i)^2 / 2]\), and the Gaussian exponential moment with \(m + \delta / 2 = b\) gives (3.82); since \(a \ge T_k\), the \(i \le k\) part of the first sum is \(\delta (b V_k - H_k)\).

(f) The first sentence collects (c), (d) and (e). Two quotes with distinct \(b, b'\) give \(b V_k - H_k\) and \(b' V_k - H_k\), hence \((V_k, H_k)\). The polytope is the nonnegative orthant cut by two affine equations; it is nonempty if and only if \((V_k, H_k)\) lies in the cone generated by the \((1, T_i)\), and a linear function attains its extrema at vertices, which have at most two nonzero coordinates because the constraint matrix has two rows. For \(k = 3\), eliminating \(x_1\) and \(x_3\) gives \(x_1 = (T_3 V - H - (T_3 - T_2) x_2) / (T_3 - T_1)\) and \(x_3 = (H - T_1 V - (T_2 - T_1) x_2) / (T_3 - T_1)\), both decreasing in \(x_2\), and nonnegativity of all three coordinates is exactly \(0 \le x_2 \le s_*\); both endpoints are attained. For dates \((1, 2, 3)\), \(V = 7\varepsilon\) and \(H = 14\varepsilon\) give \(s_* = 7\varepsilon\).\hfill$\square$

Specialized to \(k = N\) and \(k = N - 1\), (d) gives the dependence through \(V\) in Corollary \hyperref[res:Corollary-21]{21} and through \((V_-, w)\) in Corollary \hyperref[res:Corollary-22]{22}, whose value formulas follow. The randomization step in (d) is the classical fact that a stopping time of a filtration enlarged by an independent variable is a randomized stopping time in the sense of \cite{baxter1977compactness}; (3.85) proves the value inequality directly. The proposition says nothing about the unrevealed variances \(v_{k+1}, \dots, v_N\), jointly or individually; when they also vary, the equal-observation set is larger and is not described; and a fixing interval with a meeting strictly inside is excluded from (e) while its fixing remains covered by (c).

\subsection{Late American cash exercise after the last meeting}\label{sec:3.18}

Take \(T_N \le A \le S \le a < b\) in (3.76). Write \(V = \sum_{i=1}^N v_i\) and \(H = \sum_{i=1}^N T_i v_i\). After the last meeting, with \(X = \sum_i Z_i\), (3.45) gives \(r_t = X + tV - H\) and \(f(t, U) = X + UV - H\), so \(1 + d_j L_j = P(u_j, u_{j+1})^{-1} = B_{u_{j+1}} / B_{u_j}\), the product telescopes, and by the Gaussian exponential moment

\begin{equation*}
\prod_{j=0}^{J-1} (1 + d_j L_j) = \exp\Bigl( \int_a^b r_s\,ds \Bigr), \qquad F(0; a, b) = \frac{\exp\{\delta (bV - H)\} - 1}{\delta}. \tag{3.86}\label{eq:3.86}
\end{equation*}

For \(L = S - A\) put \(D_{V,L}(x) := \max_{0 \le u \le L} e^{-xu - Vu^2/2}\), which equals \(1\) for \(x \ge 0\), \(e^{x^2 / (2V)}\) for \(-VL < x < 0\) and \(e^{-xL - VL^2/2}\) for \(x \le -VL\) when \(V > 0\).

\phantomsection\label{res:Corollary-21}
\textbf{Corollary 21.} Since every \(Z_i\) is revealed by \(A\), \(R\) is \(\mathcal F_A\)-measurable and \(F(t; a, b) = R = [\exp\{\delta r_A + \delta V ((a + b)/2 - A)\} - 1] / \delta\) for \(t \in [A, b]\). If \(V > 0\), the time \(\tau_* = A + \min\{L, \max\{0, -r_A / V\}\}\) is optimal and

\begin{equation*}
U(A, S, a, b, K) = E_{Y \sim N(0, V)}\left[ \left( \frac{e^{\delta Y + \delta V ((a + b)/2 - A)} - 1}{\delta} - K \right)^+ D_{V,L}(Y) \right], \tag{3.87}\label{eq:3.87}
\end{equation*}

so every late American cash value depends on the meeting variances only through \(V\); if \(A < S\) it strictly exceeds the value under mandatory exercise at \(S\); if \(V = 0\) then \(F = 0\), \(B = 1\), \(U = (-K)^+\) and \(A\) is optimal. For \(N = 3\), \(T_i = i\) and \(\varepsilon > 0\), the vectors \(v = \varepsilon(2, 3, 2)\) and \(v' = \varepsilon(3, 1, 3)\) give every futures quote (3.86) with \(a \ge 3\) and every value (3.87) with \(3 \le A \le S \le a < b\), for every partition and strike, although their event variances differ.

\emph{Proof sketch.} The payoff is constant along each path over the exercise interval, while \(B_{A+u} / B_A = \exp\{r_A u + V u^2 / 2\}\) for \(0 \le u \le L\); minimizing the convex quadratic over \([0, L]\) gives \(\tau_*\) and \(D_{V,L}\), \(\tau_*\) is \(\mathcal F_A\)-measurable and hence a stopping time, and pathwise domination shows it attains the supremum. Under \(dQ^A = B_A^{-1} dQ\), \(r_A \sim N(0, V)\), which gives (3.87), and \(D_{V,L}(x) \le e^{L|x|}\) gives integrability. A Gaussian set of positive probability has an interior optimizer when \(A < S\), which is the strict premium; \(V = 0\) is direct. Both vectors have \(V = 7\varepsilon\) and \(H = 14\varepsilon\), and the dependence through \(V\) alone is Proposition \hyperref[res:Proposition-20]{20}(d) with \(k = N\). Appendix \hyperref[sec:A]{A} records the formal proof.\hfill$\square$

\subsection{American cash exercise across one unrevealed meeting}\label{sec:3.19}

Keep the contract (3.76) and move the exercise window across the last meeting: fix

\begin{equation*}
T_{N-1} \le A < T_N < S \le a < b, \tag{3.88}\label{eq:3.88}
\end{equation*}

so that exactly one meeting is unrevealed when exercise opens and falls strictly inside the window, while accrual still follows the last meeting; the partition, strike and dates are fixed before models are compared. Write \(T = T_N\), \(\ell = T - A\), \(L = S - T\), \(m = (a + b)/2\), \(V_- = \sum_{i < N} v_i\), \(w = v_N\), so \(V = V_- + w\), and

\begin{equation*}
f_-(x) = \frac{\exp\{\delta [x + V_-(m - A) + w(b - T)]\} - 1}{\delta}, \qquad f_+(y) = \frac{\exp\{\delta [y + V(m - T)]\} - 1}{\delta}. \tag{3.89}\label{eq:3.89}
\end{equation*}

For \(q \ge 0\) and \(d \ge 0\) let \(u_{q,d}(z) = \min\{d, \max\{0, -z/q\}\}\) if \(q > 0\), and \(u_{0,d}(z) = 0\) for \(z \ge 0\), \(u_{0,d}(z) = d\) for \(z < 0\); then

\begin{equation*}
D_{q,d}(z) := \exp\{-z u_{q,d}(z) - q u_{q,d}(z)^2 / 2\} = \max_{0 \le u \le d} e^{-zu - qu^2/2}, \tag{3.90}\label{eq:3.90}
\end{equation*}

which extends the discount maximizer of Section \hyperref[sec:3.18]{3.18} to \(q = 0\). Put

\begin{equation*}
p(x) = (f_-(x) - K)^+ D_{V_-, \ell}(x), \qquad c(x) = e^{-x\ell - V_- \ell^2 / 2}\, E_{Z \sim N(0, w)}\bigl[ (f_+(x + V_- \ell + Z) - K)^+ D_{V, L}(x + V_- \ell + Z) \bigr]. \tag{3.91}\label{eq:3.91}
\end{equation*}

\phantomsection\label{res:Corollary-22}
\textbf{Corollary 22.} With the dates (3.88), exact compounding and the initial quote are as in (3.86), and there is a simultaneous version of the futures process on the exercise window with

\begin{equation*}
F(t; a, b) = f_-(r_A) \text{ for } A \le t < T, \qquad F(t; a, b) = f_+(r_T) \text{ for } T \le t \le S, \tag{3.92}\label{eq:3.92}
\end{equation*}

so the underlying jumps at \(T\). The American cash value is finite and equals

\begin{equation*}
U(A, S, a, b, K) = E_{Y \sim N(0, V_-)}\bigl[ \max\{p(Y), c(Y)\} \bigr]; \tag{3.93}\label{eq:3.93}
\end{equation*}

for fixed contract data it depends on the variance vector only through \((V_-, w)\), while the futures quote also depends on \(H\). The rule that exercises at \(A + u_{V_-, \ell}(r_A)\) on \(E = \{p(r_A) > c(r_A)\}\) and otherwise waits for the revelation at \(T\) and exercises at \(T + u_{V, L}(r_T)\) is a stopping time valued in \([A, S]\) attaining the supremum; on \(E\) its time is strictly less than \(T\), ties wait, and no decision before \(T\) uses \(Z_N\). For \(N = 4\), \(T_i = i\) and \(\varepsilon > 0\), the strictly positive vectors

\begin{equation*}
v = \varepsilon(2, 3, 2, 1), \qquad v' = \varepsilon(3, 1, 3, 1) \tag{3.94}\label{eq:3.94}
\end{equation*}

give every futures quote with \(4 \le a < b\) and every value (3.93) with \(3 \le A < 4 < S \le a < b\), for every partition and strike, although their first three event variances differ; the crossed variance is the same in the pair.

\emph{Proof sketch.} Set \(x = r_A\). For \(0 \le u < \ell\), \(r_{A+u} = x + V_- u\); at \(T\), \(y := r_T = x + V_- \ell + Z_N\); for \(0 \le u \le L\), \(r_{T+u} = y + Vu\); and \(R = f_+(y)\). Since \(Z_N\) is independent of \(\mathcal F_A\), the Gaussian exponential moment gives

\begin{equation*}
E_Q[f_+(x + V_- \ell + Z_N) \mid \mathcal F_A] = f_-(x), \tag{3.95}\label{eq:3.95}
\end{equation*}

and since the filtration is constant on \([A, T)\) and after \(T\), (3.92) is a simultaneous adapted version of the futures process. Minimizing \(zu + qu^2/2\) over \([0, d]\) gives (3.90). The bank-account ratios are

\begin{equation*}
\frac{B_{A+u}}{B_A} = e^{xu + V_- u^2 / 2}, \ 0 \le u \le \ell, \qquad \frac{B_{T+u}}{B_T} = e^{yu + Vu^2/2}, \ 0 \le u \le L, \tag{3.96}\label{eq:3.96}
\end{equation*}

so, measured at \(A\), every payoff before \(T\) is at most \(p(x)\), and the largest discounted payoff after \(T\) has conditional expectation \(c(x)\). Since \(D_{V,L} \ge 1\), (3.95) and convexity of the positive part give

\begin{equation*}
c(x) \ge e^{-x\ell - V_- \ell^2 / 2} (f_-(x) - K)^+, \tag{3.97}\label{eq:3.97}
\end{equation*}

whose right side is \(p(x)\) when \(u_{V_-, \ell}(x) = \ell\); so \(p(x) > c(x)\) forces \(u_{V_-, \ell}(x) < \ell\), and early exercise, when strictly preferable, is attained at an allowed time before \(T\). For any \(\tau \in \mathcal T[A, S]\), \(\{\tau < T\} = \bigcup_{q \in \mathbb Q \cap [A, T)} \{\tau \le q\}\) belongs to \(\mathcal F_A\), because the filtration is constant on \([A, T)\); so conditioning on \(\mathcal F_A\) bounds every admissible payoff by \(\max\{p(x), c(x)\}\). The stated rule attains the bound, and it is a stopping time: \(E\) and its early exercise time are \(\mathcal F_A\)-measurable and its late exercise time is \(\mathcal F_T\)-measurable, so each event \(\{\tau \le t\}\) lies in \(\mathcal F_t\). The value is therefore \(E_Q[B_A^{-1} \max\{p(r_A), c(r_A)\}]\). With \(h_i = A - T_i \ge 0\),

\begin{equation*}
B_A^{-1} = \exp\Bigl\{ -\sum_{i < N} h_i Z_i - \tfrac12 \sum_{i < N} v_i h_i^2 \Bigr\}, \qquad E_Q[B_A^{-1}] = 1, \tag{3.98}\label{eq:3.98}
\end{equation*}

and under the measure with this density the earlier \(Z_i\) are independent \(N(-v_i h_i, v_i)\), so \(r_A \sim N(0, V_-)\), which is (3.93); the bounds \(1 \le D_{q,d}(z) \le e^{d|z|}\) give finiteness. For (3.94) both vectors have \(V_- = 7\varepsilon\), \(w = \varepsilon\), \(V = 8\varepsilon\) and \(H = 18\varepsilon\); the dependence through \((V_-, w)\) alone is Proposition \hyperref[res:Proposition-20]{20}(d) with \(k = N - 1\). Appendix \hyperref[sec:A]{A} records the formal proof, including the measurability and integrability steps.\hfill$\square$

Before \(T\) the holder knows every earlier shock and faces one revelation; the futures rate is constant until then while discounting evolves deterministically, so the choice at \(A\) is between exercising before \(T\) and waiting, the latter valued by a Gaussian expectation. Discounting to time zero centres the law of \(r_A\) and removes the separate dependence of the cash value on the dates of the revealed shocks, which is what (3.94) exploits while keeping \(V_-\), \(w\) and \(H\) fixed. The proposition says nothing about the crossed variance \(w\) on its own, gives no necessary and sufficient equality condition for arbitrary parameter pairs, and does not cover windows opening before \(T_{N-1}\), exercise across two unrevealed meetings or accrual spanning a meeting.

\subsection{Futures-style exercise}\label{sec:3.20}

The cash calls of Sections \hyperref[sec:3.17]{3.17} to 3.19 pay cash at the exercise date and are discounted. Under futures-style margining the premium is paid at expiry and variation margin is exchanged daily without remuneration, so the option price is the undiscounted conditional expectation of the payoff on the futures price, \cite{nastasi2018smile} §2.2 eq. 3, the same convention as the futures rate itself in \cite{skov2021dynamic} (20); \cite{itkin2024semi} footnote 2 leaves the discounting of cash-settled SOFR futures options open. That a futures-style option is never exercised early is known: \cite{lieu1990option} Properties B and C, pp. 330 to 331, prove it in the Black lognormal futures model with fixed rates, and \cite{chen1993pricing} inequalities (1) and (2), pp. 15 to 16, by Jensen's inequality under the undiscounted risk-adjusted expectation with stochastic rates. Part (b) below is the scheduled-jump instance of that result, proved directly on the piecewise-constant filtration, which is in neither source; the identification content of parts (c) to (e) is new. Nothing is asserted about which convention a listed contract follows.

In the model of Section \hyperref[sec:3.10]{3.10} fix a window \(a < b\), \(\delta = b - a\), \(m = (a + b)/2\), a partition \(a = u_0 < \dots < u_J = b\) with no meeting strictly inside a fixing interval, the exact compounding (3.76), and the coefficients \(w_i\), \(d_i\) of (3.71). Fix deterministic \(0 \le A \le S \le b\) and a real strike \(K\), and define the futures-style American value

\begin{equation*}
\tilde U(A, S, a, b, K) := \sup_{\tau \in \mathcal T[A, S]} E_Q\bigl[ (F(\tau; a, b) - K)^+ \bigr], \tag{3.99}\label{eq:3.99}
\end{equation*}

with no discount factor.

\phantomsection\label{res:Proposition-23}
\textbf{Proposition 23.} (a) The product telescopes, \(\int_a^b r_s\,ds = \sum_i (w_i Z_i + d_i v_i)\), and for \(0 \le t \le b\)

\begin{equation*}
F(t; a, b) = \frac{1}{\delta} \Bigl[ \exp\Bigl( \sum_{T_i \le t} w_i Z_i + \sum_i d_i v_i + \tfrac12 \sum_{T_i > t} w_i^2 v_i \Bigr) - 1 \Bigr] \tag{3.100}\label{eq:3.100}
\end{equation*}

is a version of \(E_Q[R \mid \mathcal F_t]\) that is constant on each \([T_n, T_{n+1}) \cap [0, b]\), measurable at \(T_{j(t)}\), square integrable and a martingale, with \(1 + \delta F(0; a, b) = E_Q[1 + \delta F(t; a, b)]\).

(b) For every \(\tau \in \mathcal T[A, S]\), \(E_Q[(F(\tau; a, b) - K)^+] \le E_Q[(F(S; a, b) - K)^+]\); exercise at \(S\) is optimal and

\begin{equation*}
\tilde U(A, S, a, b, K) = E_Q\bigl[ (F(S; a, b) - K)^+ \bigr], \tag{3.101}\label{eq:3.101}
\end{equation*}

for every \(A \le S\), including windows crossing any number of unrevealed meetings. This is the scheduled-jump instance of the known result of \cite{lieu1990option} Properties B and C and \cite{chen1993pricing} inequalities (1) and (2), whose mechanism, a martingale futures price and a convex payoff, is the same; it is proved here by the discrete optional sampling theorem over meeting indices on the scheduled filtration. By contrast, the cash value of Corollary \hyperref[res:Corollary-21]{21} strictly exceeds its mandatory-exercise value whenever \(V > 0\) and \(A < S\).

(c) With \(n = j(S)\), \(\sigma_S^2 = \sum_{i \le n} w_i^2 v_i\), \(\mu_S = \sum_i d_i v_i + \tfrac12 \sum_{i > n} w_i^2 v_i\), \(G = e^{\mu_S + \sigma_S^2 / 2} = 1 + \delta F(0; a, b)\) and \(K' = 1 + \delta K\), the value is \(c_{K'}(G, \sigma_S) / \delta\) in the notation (3.72):

\begin{equation*}
\begin{gathered}
\tilde U = \frac{1}{\delta} \bigl[ G \Phi(d_1) - K' \Phi(d_2) \bigr], \quad d_{1,2} = \frac{\log(G / K') \pm \sigma_S^2 / 2}{\sigma_S} \ (\sigma_S > 0,\ K' > 0); \\
\tilde U = \frac{G - K'}{\delta} \ (K' \le 0); \qquad \tilde U = \frac{(e^{\mu_S} - K')^+}{\delta} \ (\sigma_S = 0).
\end{gathered} \tag{3.102}\label{eq:3.102}
\end{equation*}

It is finite, continuous and nonincreasing in \(K\).

(d) Let \(k = j(A)\), suppose \(a \ge T_k\), and take \(V_k\), \(H_k\) from (3.77). Then \(w_i = \delta\) and \(d_i = \delta (m - T_i)\) for \(i \le k\), so

\begin{equation*}
\sigma_S^2 = \delta^2 V_k + \sum_{k < i \le n} w_i^2 v_i, \qquad \mu_S = \delta (m V_k - H_k) + \sum_{i > k} d_i v_i + \tfrac12 \sum_{i > n} w_i^2 v_i, \tag{3.103}\label{eq:3.103}
\end{equation*}

and every value (3.101), every futures-style European value and the quote \(F(0; a, b)\) depend on \((v_1, \dots, v_k)\) only through \((V_k, H_k)\), while variances with \(T_i > b\) do not enter. Conversely, fix \(v_{k+1}, \dots, v_N\). Two quotes with \(a \ge T_k\) and distinct \(b\) identify \((V_k, H_k)\). For a window with \(a \ge T_k\), let \(n_b\) be the largest \(n\) with \(T_n < b\) and take one expiry \(S\) in each interval \([T_n, T_{n+1}) \cap [A, b]\), \(n = k, \dots, n_b\), with \(S = b\) allowed; the surface \(\{\tilde U(A, S, a, b, K) : K \in \mathbb R\}\) at \(S\) identifies \(\sigma_S^2\), and successive differences give \(w_n^2 v_n\) for \(k < n \le n_b\), so every unrevealed variance with \(T_n < b\) is identified, while a variance at \(T_n \ge b\) enters neither the quote nor the surfaces of this window. Among vectors sharing the later block, the set consistent with exact observation of the quotes and surfaces is the polytope (3.83), as for the cash-settled contracts.

(e) For dates \((1, 2, 3, 4, 5)\), any \(3 \le A < 4\), any window with \(a \ge A\) whose partition contains the meetings inside it, and every allowed \(S\) and \(K\), the vectors \((2\varepsilon, 3\varepsilon, 2\varepsilon, v_4, v_5)\) and \((3\varepsilon, \varepsilon, 3\varepsilon, v_4, v_5)\), \(\varepsilon > 0\), \(v_4, v_5 \ge 0\), give identical quotes and values. The vectors \(\varepsilon(2, 3, 2, 3, 2)\) and \(\varepsilon(2, 3, 2, 1, 2)\), which differ only at the fourth meeting, are separated by every surface with expiry \(S \in [4, 5)\) and window \(a < 4 < b\): \(\sigma_S^2\) differs by \(2(b - 4)^2 \varepsilon\) and \(\mu_S\) by \((b - 4)^2 \varepsilon\), both larger for the first vector, the quotes differ, and the value differs at every strike.

\emph{Proof.} (a) Telescoping and the accrual integral are as in Proposition \hyperref[res:Proposition-20]{20}(e) with the coefficients of (3.71). For \(t \le b\) and \(n = j(t)\), \(\mathcal F_t = \sigma(Z_1, \dots, Z_n) \vee \mathcal N\), the \(Z_i\) with \(i > n\) are independent of it, and \(E_Q[e^{w_i Z_i}] = e^{w_i^2 v_i / 2}\), which gives (3.100); it is a function of \(n\) and \((Z_1, \dots, Z_n)\), hence constant between meetings and measurable at \(T_n\), square integrable by Gaussian exponential moments, and a martingale as a family of conditional expectations of the integrable \(R\); the mean identity is the tower property.

(b) The futures price is a martingale for the filtration generated at the meetings, so \((F - K)^+\) is a submartingale by Jensen's inequality and optional sampling at the bounded meeting index of any exercise time makes exercise at \(S\) optimal, as in \cite{lieu1990option} and \cite{chen1993pricing}; the comparison with the cash call is Corollary \hyperref[res:Corollary-21]{21}.

(c) By (3.100) at \(t = S\), \(F(S; a, b) = [e^{Y + \mu_S} - 1] / \delta\) with \(Y = \sum_{i \le n} w_i Z_i \sim N(0, \sigma_S^2)\), so \(\tilde U = E[(e^{Y + \mu_S} - K')^+] / \delta\), the Gaussian call integral (3.72) with forward \(G\); the three cases are its evaluation, and \(K \mapsto E[(W - K)^+]\) is continuous and nonincreasing for integrable \(W\).

(d) For \(i \le k\), \(T_i \le T_k \le a\), so (3.71) gives \(w_i = \delta\) and \(d_i = \delta(m - T_i)\); summing gives (3.103), and \(w_i = 0\) for \(T_i > b\). For the converse, the quote exponent is \(\delta(b V_k - H_k)\) plus terms fixed by the later block, so two distinct \(b\) give \((V_k, H_k)\). Next, \(\partial \tilde U / \partial K = -Q(Y + \mu_S > \log K')\), so the surface determines the law of \(Y + \mu_S\), a Gaussian, hence \(\sigma_S^2\), also when it is a point mass. Expiries \(S_n\) with \(j(S_n) = n\) give \(\sigma_{S_n}^2 - \sigma_{S_{n-1}}^2 = w_n^2 v_n\) with \(w_n = b - T_n > 0\) for \(T_n < b\), while \(T_n = b\) gives \(w_n = 0\). Equal observations therefore force equal \((V_k, H_k)\) and equal identified later coordinates; with the later block fixed, equal observations are equivalent to equal \((V_k, H_k)\), and the equal-observation set is the polytope by Proposition \hyperref[res:Proposition-20]{20}(f).

(e) The first pair has \(V_3 = 7\varepsilon\) and \(H_3 = 14\varepsilon\) with a common later block; apply (d) with \(k = 3\). The second pair shares \(V_3\), \(H_3\) and \(v_5\) and has \(v_4 = 3\varepsilon\) against \(\varepsilon\); for \(S \in [4, 5)\) and \(a < 4 < b\), \(w_4 = b - 4\), \(d_4 = (b - 4)^2 / 2\) and \(n = 4\), so by (3.103) \(\sigma_S^2\) differs by \(2(b - 4)^2 \varepsilon\) and \(\mu_S\) by \((b - 4)^2 \varepsilon\), both larger for the first vector, and \(G\) by the factor \(e^{2(b - 4)^2 \varepsilon}\). By (3.73) the value \(c_{K'}(G, \sigma_S) / \delta\) is strictly increasing in \(G\) and nondecreasing in \(\sigma_S\) for \(K' > 0\), and for \(K' \le 0\) it is \((G - K') / \delta\); so it is strictly larger for the first vector at every strike.\hfill$\square$

Under this convention the exercise payoff is a convex function of a martingale and no discounting rewards waiting, so the payoff process is a submartingale and the last allowed date is optimal, which is the argument of \cite{chen1993pricing} and the content of \cite{lieu1990option} Properties B and C; the scheduled model makes it discrete, since every stopping rule reduces to a choice of meeting index. The European value is a Black-type formula whose only variance input is \(\sigma_S^2\), which separates each unrevealed meeting inside the window as expiries pass it, while the revealed shocks enter \(\sigma_S^2\) and \(\mu_S\) only through their sum and date-weighted sum, exactly as the quote does; the polytope ambiguity survives the change of convention, and the cash convention's timing value, which depends on \(V\) alone by Corollary \hyperref[res:Corollary-21]{21}, is absent. \cite{skov2021dynamic} footnote 8 notes that the listed options are American and carry an early-exercise premium, which is the cash convention's timing value. Nothing is asserted about which settlement convention describes a listed SOFR futures option, and the hybrid with the premium paid at trade is not modelled; the identification in (d) uses expiries in every meeting interval of the window and partitions containing the meetings inside it; a variance at a meeting at or after \(b\) is invisible to that window's quote and surfaces, and which later windows identify it is not addressed here.

\subsection{The listed calendar}\label{sec:3.21}

The rank criterion of Proposition \hyperref[res:Proposition-19]{19} is now evaluated on the actual expiry calendar of the listed three-month SOFR futures options, with the ideal European cash payoffs of that proposition, exact prices, the zero initial curve and a stated quote precision; no market data enter. The contract rules are those of `cme2026sr3futures` (Chapter 460) and `cme2026sr3options` (Chapter 460A), the exercise-price rule also in `cme2020sr3filing` and `cme2020sr3optionsfiling`, and the meeting days are the decision days of \cite{fomc2026calendar}. The Reference Quarter of a future is the interval from the third Wednesday of the third month before its delivery month to the third Wednesday of the delivery month, settled on compounded daily SOFR over it; a quarterly option's underlying is the future whose Reference Quarter begins on the third Wednesday of the expiry month, a serial option's the future of the next quarterly month, a mid-curve option's the future twelve or more months later; every option other than a weekly mid-curve stops trading on the Friday before the third Wednesday of its expiry month; prices are quoted in index points, one hundredth of a point being one basis point, and the minimum fluctuation is one quarter or one half of a basis point depending on the expiry and the premium. Which expiries are listed is set by the exchange; the 2020 filings show four serial and the nearest quarterly expiries listed.

Observation is on 2 January 2026, time is measured in years on an actual-over-360 basis, the scheduled dates are the sixteen decision days of 2026 and 2027, and an option expiring on a Friday \(S\) sees a meeting \(T_i\) if and only if \(T_i \le S\). The decision days are 28 January, 18 March, 29 April, 17 June, 29 July, 16 September, 28 October and 9 December 2026, and 27 January, 17 March, 28 April, 9 June, 28 July, 15 September, 27 October and 8 December 2027. The futures and options in the window are as follows; every Reference Quarter spans exactly 91 days, \(\delta = 91/360\).

{\footnotesize
\begin{longtable}{lll}
\toprule
\textbf{Future} & \textbf{Reference Quarter} & \textbf{Option expiries (two serials and the quarterly)} \\
\midrule
\endhead
\bottomrule
\endlastfoot
Jun 2026 & 18 Mar to 17 Jun 2026 & 16 Jan, 13 Feb, 13 Mar 2026 \\
Sep 2026 & 17 Jun to 16 Sep 2026 & 10 Apr, 15 May, 12 Jun 2026 \\
Dec 2026 & 16 Sep to 16 Dec 2026 & 10 Jul, 14 Aug, 11 Sep 2026 \\
Mar 2027 & 16 Dec 2026 to 17 Mar 2027 & 16 Oct, 13 Nov, 11 Dec 2026 \\
Jun 2027 & 17 Mar to 16 Jun 2027 & 15 Jan, 12 Feb, 12 Mar 2027 \\
Sep 2027 & 16 Jun to 15 Sep 2027 & 16 Apr, 14 May, 11 Jun 2027 \\
Dec 2027 & 15 Sep to 15 Dec 2027 & 16 Jul, 13 Aug, 10 Sep 2027 \\
Mar 2028 & 15 Dec 2027 to 15 Mar 2028 & 15 Oct, 12 Nov, 10 Dec 2027 \\
\end{longtable}}

The assumptions are: (A1) every serial and quarterly expiry in the window is listed; (A2) no exchange holiday moves these Fridays; (A3) the at-the-money premium of each option is known to within \(\varepsilon \in \{0.25, 0.5\}\) basis points, one half to one minimum fluctuation; (A4) a reference meeting standard deviation \(s_0 \in \{10, 15, 25\}\) basis points for every meeting, a size against which the uncertainty is compared and not a market estimate. A meeting is called not practically identified when the propagated uncertainty of its variance is at least \(s_0^2\).

\phantomsection\label{res:Proposition-24}
\textbf{Proposition 24.} In the model of Section \hyperref[sec:3.10]{3.10} on this calendar, under (A1) to (A4), the following hold.

(a) Every listed option expires before its underlying accrues, \(S < a\), so for every meeting \(T_i \le S\) the coefficients of (3.71) are \(w_i = \delta\) and \(k_i = \delta^2\), and \(k_i = 0\) for \(T_i > S\); each complete surface at \((S, a, b)\) determines, by Proposition \hyperref[res:Proposition-19]{19}(a),

\begin{equation*}
q(S) = \delta^2 \sum_{T_i \le S} v_i, \tag{3.104}\label{eq:3.104}
\end{equation*}

the accumulated variance of the meetings revealed by \(S\). Meetings inside the accrual window enter no option surface.

(b) Merge the expiries in date order, \(S_1 < \dots < S_L\), with \(S_0\) the observation date. The surfaces identify every meeting variance with \(T_i \le S_L\), globally on the orthant, if and only if each gap \((S_{\ell-1}, S_\ell]\) contains at most one meeting; a gap with two or more meetings identifies only their sum, and meetings after \(S_L\) are identified by no option. On the calendar above every one of the sixteen meetings lies in its own gap, eleven within one future's strip and five only by a difference across futures (18 March, 17 June and 16 September 2026, 17 March and 15 September 2027), so all sixteen are identified, by

\begin{equation*}
v_i = \frac{q_\ell}{\delta_\ell^2} - \frac{q_{\ell-1}}{\delta_{\ell-1}^2}, \qquad T_i \text{ the meeting in } (S_{\ell-1}, S_\ell], \quad q_0 = 0, \tag{3.105}\label{eq:3.105}
\end{equation*}

with \(\delta_\ell\) the accrual length of the future underlying \(S_\ell\). Mid-curve options share the Fridays and revealed sets of the standard options and add proportional rows; weekly mid-curves subdivide gaps; both add redundancy only. With one row per meeting the coefficient matrix is \(K = \delta^2 L\), \(L\) the lower triangular matrix of ones, whose singular values are \(\delta^2 / (2\sin((2k-1)\pi/66))\), \(k = 1, \dots, 16\):

\begin{equation*}
\sigma_{\min}(K) = \frac{\delta^2}{2\cos(\pi/33)} = 0.0321, \qquad \sigma_{\max}(K) = \frac{\delta^2}{2\sin(\pi/66)} = 0.671, \qquad \kappa(K) = 20.9. \tag{3.106}\label{eq:3.106}
\end{equation*}

(c) At the at-the-money strike the cash price is \(c = m[2\Phi(\sqrt q / 2) - 1]\), so \(q = 4[\Phi^{-1}((1 + c/m)/2)]^2\) and

\begin{equation*}
\frac{dq}{dc} = \frac{2\sqrt q}{m\,\varphi(\sqrt q / 2)}, \tag{3.107}\label{eq:3.107}
\end{equation*}

nondecreasing in \(q\). A perturbation of \(\varepsilon\) in rate units of the at-the-money premium, whose cash price is \(\delta\) times the premium, perturbs \(q\) by at most \((dq/dc)\,\delta\varepsilon\) at the larger price, and by (3.105) the variance of the meeting in gap \(\ell\) by at most

\begin{equation*}
\Delta v_i \le 2\varepsilon \left[ \frac{\sigma_\Sigma(S_\ell)}{m_\ell\, \varphi(\delta_\ell \sigma_\Sigma(S_\ell)/2)} + \frac{\sigma_\Sigma(S_{\ell-1})}{m_{\ell-1}\, \varphi(\delta_{\ell-1} \sigma_\Sigma(S_{\ell-1})/2)} \right], \qquad \sigma_\Sigma(S)^2 = \sum_{T_i \le S} v_i. \tag{3.108}\label{eq:3.108}
\end{equation*}

For accumulated standard deviations below one percentage point both \(\varphi(\delta\sigma_\Sigma/2)/\varphi(0)\) and \(m\) are within \(2 \times 10^{-5}\) of one, so

\begin{equation*}
\Delta v_i \approx \frac{2}{\varphi(0)}\, \varepsilon\, [\sigma_\Sigma(S_\ell) + \sigma_\Sigma(S_{\ell-1})] = 5.01\, \varepsilon\, [\sigma_\Sigma(S_\ell) + \sigma_\Sigma(S_{\ell-1})]: \tag{3.109}\label{eq:3.109}
\end{equation*}

the amplification grows with the accumulated standard deviation, not with the meeting's own variance. Under (A4), \(\sigma_\Sigma(S) = s_0 \sqrt{n(S)}\) with \(n(S)\) the number of meetings revealed by \(S\), and the \(n\)-th meeting has

\begin{equation*}
\frac{\Delta v_n}{s_0^2} \approx 5.01\, \frac{\varepsilon}{s_0}\, \bigl( \sqrt n + \sqrt{n - 1} \bigr), \tag{3.110}\label{eq:3.110}
\end{equation*}

nondecreasing in \(n\). Evaluating (3.108) exactly on the calendar: at \(\varepsilon = 0.5\) and \(s_0 = 10\) basis points, meetings 1 to 4 (28 January to 17 June 2026) are identified to within their size and meetings 5 to 16 (29 July 2026 onward) are not practically identified; at \(\varepsilon = 0.5\) and \(s_0 = 15\), meetings 1 to 9 are identified and 10 to 16 (17 March 2027 onward) are not; at \(\varepsilon = 0.5\) and \(s_0 = 25\), and at \(\varepsilon = 0.25\) for all three sizes, all sixteen are identified, the closest case being the sixteenth meeting at \(\varepsilon = 0.25\), \(s_0 = 10\) with ratio \(0.98\). In \(\ell^2\), \(\|\Delta v\|_2 \le \|\Delta q\|_2 / \sigma_{\min}(K)\).

(d) The futures quote of each future gives \(p = \sum_i h_i v_i\) with \(h_i = \delta(b - T_i)\) for \(T_i \le a\) and \(h_i = (b - T_i)^2\) for \(a < T_i < b\). The quotes are redundant for the meetings identified in (b), and for a meeting inside an accrual window after the last option expiry the quote would be the only observation, with a small coefficient: for the 8 December 2027 meeting inside the December 2027 Reference Quarter, which the 10 December 2027 expiry does identify, \(h_i = (7/360)^2 = 3.78 \times 10^{-4}\) against \(h = \delta(b - T) = 0.482\) for the 28 January 2026 meeting in the same row, a ratio of about \(1270\); a half-basis-point perturbation of that futures rate moves \(p\) by \(1.26 \times 10^{-5}\) and hence \(v_i\) by \(0.033\), a standard deviation of about 18 percentage points.

(e) If instead only the four nearest serials are listed on the observation date, the serial expiries available are those of January, February, April and May 2026 and consecutive listed expiries are quarterly thereafter; the gap from 12 June to 11 September 2026 then contains the meetings of 17 June and 29 July, and the gap from 11 September to 11 December 2026 those of 16 September, 28 October and 9 December, so by (b) only their sums are identified from mid-2026 on, unless mid-curve or weekly expiries fill the gaps.

\emph{Proof.} (a) For a quarterly option expiring in month \(M\) the Reference Quarter begins on the third Wednesday of \(M\) and the expiry is the Friday before it; for a serial option the underlying's Reference Quarter begins on the third Wednesday of the next quarterly month, later still; for mid-curves twelve or more months later. Hence \(S < a\), and (3.71) gives \(w_i = (b - T_i)^+ - (a - T_i)^+ = \delta\) for \(T_i \le a\), the indicator selecting the revealed meetings; Proposition \hyperref[res:Proposition-19]{19}(a) recovers \((m, q)\) from the surface.

(b) By (3.104) the map \(v \mapsto (q_1, \dots, q_L)\) is linear with rows \(\delta_\ell^2 \mathbf 1\{T_i \le S_\ell\}\), prefix indicators, so it has full column rank if and only if the prefix sets grow by at most one element per step and reach every meeting; when a gap holds two meetings the vector with \(+1\) and \(-1\) on them is in the kernel and Proposition \hyperref[res:Proposition-19]{19}(b) gives strictly positive indistinguishable pairs. The gap assignment is read off the calendar; the cross-future cases are those whose bounding expiries belong to different futures, for which (3.105) uses the two accrual lengths, here equal. With one row per meeting, \(L^{-1}\) is the bidiagonal difference matrix, whose Gram matrix is tridiagonal with \(2\) on the diagonal except \(1\) in the last entry and \(-1\) off the diagonal, with eigenvalues \(4\sin^2((2k-1)\pi/(2(2n+1)))\), \(k = 1, \dots, n\), the eigenvectors having entries \(\sin(j\theta_k)\) with \(\theta_k = (2k-1)\pi/(2n+1)\); the singular values of \(L\) are their reciprocal square roots, and \(n = 16\) gives (3.106).

(c) Proposition \hyperref[res:Proposition-19]{19}(a) at \(K = m\) gives \(c = m[2\Phi(\sqrt q/2) - 1]\) by normal symmetry, and differentiating the inverse as in Proposition \hyperref[res:Proposition-17]{17} gives (3.107), nondecreasing in \(q\) because \(y/\varphi(y)\) increases; the mean value theorem turns a premium perturbation into (3.108) with \(dc = \delta\varepsilon\) and the triangle inequality on (3.105). The approximation (3.109) uses \(\varphi(y)/\varphi(0) = e^{-y^2/2}\) with \(y = \delta\sigma_\Sigma/2 \le 0.0013\) for \(\sigma_\Sigma \le 0.01\), and \(m = \exp(\delta(b - S)\sum_{T_i \le S} v_i) \le 1 + 1.1 \times 10^{-5}\) for sixteen meetings of 25 basis points, \(b - S\) being at most about \(0.44\) years for a serial expiry. Substituting \(\sigma_\Sigma(S) = s_0\sqrt n\) gives (3.110); the exact expression (3.108) dominates it by a factor at most \(1 + 2 \times 10^{-6}\) at these sizes, and its evaluation at every meeting, tick and size gives the thresholds. The \(\ell^2\) bound is the definition of the smallest singular value.

(d) The coefficients are (3.71) with \(d_i = (b - T_i)^2/2\) and \(w_i = b - T_i\) inside the window, so \(h_i = (b - T_i)^2\); the ratio is exactly \(1274\), and \(\Delta p = \Delta \log(1 + \delta F) \approx \delta\,\Delta F\).

(e) The reduced list of expiries is a sub-list of the calendar in date order whose gaps up to 12 June 2026 hold at most one meeting each, and the two later gaps hold two and three meetings; (b) applies.\hfill$\square$

One calendar fact carries the structure: every listed option stops trading on the Friday before the third Wednesday on which its underlying's Reference Quarter starts, so an option never sees a meeting inside its own accrual window, and its surface reports the accumulated variance of the meetings it has seen. Differences between consecutive expiries peel off one meeting at a time as long as no two meetings fall between consecutive Fridays, which on this calendar never happens, the expiries falling in mid-month and the meetings late in the month. Rank is the easy part: a tick-size error in the at-the-money premium becomes an error in the accumulated variance proportional to the accumulated standard deviation, and a single meeting's variance is the difference of two such numbers, so the sixteenth meeting is recovered as the difference of two variances of order sixteen times its own, each known to a fixed fraction of a tick. On the full calendar the listed surfaces therefore identify every meeting variance before the last expiry exactly in principle; beyond the listed serials identification collapses to sums; and beyond roughly the first year it is empirical in practice. The result assumes (A1) to (A4); real premiums carry bid-ask spreads and stale quotes beyond one tick; the model has no diffusion between meetings, and a diffusive variance would enter (3.108) in the same way and worsen the amplification; the listed options are American and exercise into a futures position, and Corollaries \hyperref[res:Corollary-21]{21} and \hyperref[res:Corollary-22]{22} and Proposition \hyperref[res:Proposition-23]{23} say that American exercise and either settlement convention add no identification among revealed variances, which is cited here and not used; nothing is asserted about meetings after the last listed expiry, unknown initial curves, calibration or noisy prices beyond the stated perturbation; and the precision statement is a worst-case single-quote propagation, which using all strikes or overlapping futures can reduce in its random part but not in a common tick-level bias.

\subsection{Restrictions surviving scale calibration}\label{sec:3.22}

Proposition \hyperref[res:Proposition-13]{13} fixes the scales and asks which linear equalities the cone leaves; here the scales vary freely and the question is which restrictions survive. Use the model of Section \hyperref[sec:3.11]{3.11} with constant decays \(\lambda_j \ge 0\), fix an observation time \(t < T_N\) with \(k_0 = j(t)\), retain the rows \(n = k_0 + 1, \dots, N\), and write \(\Delta_n = T_n - T_{n-1}\). Admissible scales are nonnegative Borel functions bounded on bounded intervals. Let \(A(t)\) be the matrix (3.52) with the partition at the meeting dates, its columns indexed by \((j, k)\), \(k = k_0, \dots, N - 1\),

\begin{equation*}
A_{n,(j,k)}(t) = \gamma_{n-k,j}^2\, H_j(T_n; \max(t, T_k), T_{k+1}) \ \text{ for } n > k, \qquad A_{n,(j,k)}(t) = 0 \ \text{ for } n \le k, \tag{3.111}\label{eq:3.111}
\end{equation*}

and put \(K(t) := \{A(t) q : q \ge 0\}\).

\phantomsection\label{res:Proposition-25}
\textbf{Proposition 25.} (a) The set of vectors \(V(t)\) of (3.50) attained as the scales range over all admissible functions is exactly the closed convex cone \(K(t)\). Refining the partition or allowing measurable scales adds nothing: on each meeting interval the kernel \(s \mapsto (e^{-2\lambda_j (T_n - s)} \gamma_{n - j(s), j}^2)_n\) has a fixed direction, so every sub-interval contributes a multiple of the same column.

(b) A linear inequality \(w^\top V(t) \ge 0\) holds for every admissible scale choice if and only if \(w^\top A(t) \ge 0\) componentwise, and \(V(t)\) is attainable if and only if it satisfies all of them; the surviving equalities are those with \(A(t)^\top w = 0\), as in Proposition \hyperref[res:Proposition-13]{13}. For one factor with \(\gamma_{1,1} \ne 0\), \(A(t)\) is lower triangular with positive diagonal and the surviving inequalities are exactly the \(N - k_0\) conditions \(q = A(t)^{-1} V(t) \ge 0\), the recursion

\begin{equation*}
q_{k_0+1} = \frac{V_{k_0+1}}{A_{k_0+1,k_0}} \ge 0, \qquad q_{k+1} = \frac{V_{k+1} - \sum_{l < k} A_{k+1,l}\, q_{l+1}}{A_{k+1,k}} \ge 0, \quad k = k_0 + 1, \dots, N - 1, \tag{3.112}\label{eq:3.112}
\end{equation*}

whose rows are the rows of \(A(t)^{-1}\); they are irredundant, and every surviving inequality is a nonnegative combination of them.

(c) If \(\gamma_{i,j} = \gamma_j\) for all \(i\), each factor loading every future meeting equally, with \(\gamma_j \ne 0\), then for each factor \(A_{n,(j,k)} = e^{-2\lambda_j \Delta_n} A_{n-1,(j,k)}\) for \(k < n - 1\), and with \(\lambda_{\max} = \max_j \lambda_j\)

\begin{equation*}
V_n(t) \ge e^{-2\lambda_{\max} \Delta_n}\, V_{n-1}(t), \quad n = k_0 + 2, \dots, N, \qquad V_{k_0+1}(t) \ge 0, \tag{3.113}\label{eq:3.113}
\end{equation*}

hold for every admissible scale choice and characterize \(K(t)\): every vector satisfying (3.113) is attained, already by the scales of one factor with \(\lambda_j = \lambda_{\max}\). As the scales are recalibrated freely, the event variance of a meeting can fall below that of the preceding meeting at most by the factor \(e^{-2\lambda_{\max} \Delta_n}\); with zero decay, \(V(t)\) is nondecreasing in the meeting index.

(d) If the loadings may vary over all real arrays, nothing survives beyond \(V(t) \ge 0\): the loading \(\gamma_{1,j} \ne 0\), \(\gamma_{i,j} = 0\) for \(i \ge 2\) makes \(A(t)\) diagonal with positive diagonal, so every nonnegative vector is attained. If the decays may vary without bound, the closure of the union of the cones is the nonnegative orthant; if they are bounded by \(\Lambda\), then (3.113) with \(\Lambda\) in place of \(\lambda_{\max}\) survives every joint recalibration of scales and decays under constant ordinal loadings and characterizes the union.

\emph{Proof.} (a) Fix \(j\) and a meeting interval \([\max(t, T_k), T_{k+1})\). On it \(j(s) = k\), so the row-\(n\) entry of the kernel is \(e^{2\lambda_j s}\) times the fixed vector \((e^{-2\lambda_j T_n} \gamma_{n-k,j}^2 \mathbf 1\{n > k\})_n\); for any admissible \(a_j\), (3.50) restricted to the interval is \(\int a_j(s)^2 e^{2\lambda_j s}\,ds\) times that vector, a nonnegative multiple of the column (3.111), with the multiplier ranging over \([0, \infty)\), the column itself being the case \(a_j \equiv 1\) and any multiple attained by a constant scale on the interval. Summing over intervals and factors gives \(V(t) = A(t) q\) with \(q \ge 0\), every such \(q\) attained by piecewise-constant scales; the cone is closed because it is finitely generated.

(b) For a finitely generated cone \(\{Aq : q \ge 0\}\), \(w^\top V \ge 0\) on the cone if and only if \(w^\top A \ge 0\), and \(V\) lies in the cone if and only if \(w^\top V \ge 0\) for all such \(w\), by the separating hyperplane theorem applied to the closed convex cone. Equalities are the inequalities holding with both signs. For \(d = 1\) the matrix is square and lower triangular with positive diagonal by Proposition \hyperref[res:Proposition-13]{13}, so \(q = A^{-1} V\) is unique and \(V\) is in the cone if and only if \(q \ge 0\), which is (3.112); the rows of \(A^{-1}\) are linearly independent, so none is implied by the others, and a surviving \(w\) satisfies \(w^\top = (w^\top A) A^{-1}\) with \(w^\top A \ge 0\).

(c) With constant loadings, \(A_{n,(j,k)} = \gamma_j^2 H_j(T_n; \max(t, T_k), T_{k+1})\), and \(H_j(T_n; a, b) = e^{-2\lambda_j (T_n - T_{n-1})} H_j(T_{n-1}; a, b)\) whenever \(b \le T_{n-1}\), directly from the definition of \(H_j\); this gives the row relation for \(k \le n - 2\) and, for each factor's contribution, \(V^{(j)}_n = e^{-2\lambda_j \Delta_n} V^{(j)}_{n-1} + A_{n,(j,n-1)} q_{j,n-1}\) with a nonnegative diagonal term. Summing over \(j\), \(V_n \ge \sum_j e^{-2\lambda_j \Delta_n} V^{(j)}_{n-1} \ge e^{-2\lambda_{\max} \Delta_n} V_{n-1}\), which is (3.113). Conversely, given \(V\) satisfying (3.113), take a factor \(j^*\) with \(\lambda_{j^*} = \lambda_{\max}\) and set the other scales to zero; the one-factor matrix is lower triangular and (3.112) reads \(q_{j^*, n-1} = [V_n - e^{-2\lambda_{\max} \Delta_n} V_{n-1}] / A_{n,(j^*,n-1)} \ge 0\), with \(q_{j^*, k_0} = V_{k_0+1} / A_{k_0+1,(j^*,k_0)} \ge 0\). Irredundancy: the inequalities (3.113) are the rows of the one-factor \(A^{-1}\) up to positive scaling, by (b).

(d) With \(\gamma_{1,j} \ne 0\) and \(\gamma_{i,j} = 0\) for \(i \ge 2\), (3.111) vanishes unless \(n = k + 1\), so \(A(t)\) is diagonal with positive entries and \(q = A^{-1} V \ge 0\) for every \(V \ge 0\). For unbounded decays the constants \(e^{-2\lambda \Delta_n}\) tend to zero, so a nonnegative \(V\) is the limit of \(V + \eta \mathbf 1\), which satisfies (3.113) for a large enough decay and so lies in that cone by (c); with \(\lambda_j \le \Lambda\) each cone satisfies (3.113) with \(\Lambda\), and the cone with decay \(\Lambda\) attains every such vector by (c).\hfill$\square$

Remark 12 removed the restrictions one might infer from state dimension alone, and Proposition \hyperref[res:Proposition-13]{13} showed that with time-dependent scales no linear equality survives. What remains are inequalities, and they are exact: within an interval the decay and loading kernels move every meeting's coefficient in the same proportion, so the variance vector is a nonnegative combination of finitely many fixed directions, one per meeting interval and factor, and the attainable set is a polyhedral cone whatever the scales do; its facets are the restrictions that no scale calibration can violate. When each factor loads all future meetings equally, the cone is generated by the shifted copies of one direction and its facets are the consecutive-meeting ratio bounds: a vector of event variances that falls faster than the largest decay allows between two meetings cannot be produced by this family at any scale calibration. The ordinal loadings of \cite{brace2024sofr} are not the constant pattern of (c); for them the surviving restrictions are the triangular ones of (3.112), which depend on the loadings, and no closed form is given beyond the recursion. Under the stochastic volatility of Section \hyperref[sec:3.12]{3.12} the fixed-model vector is random and nothing is claimed here about which restrictions on its support survive scale variation. A violation of (3.113) by estimated variances would show only that no scale calibration of this family reproduces them. Farkas' lemma and the closedness of finitely generated cones are used as standard facts.

\subsection{Exact support with positive variance mean reversion}\label{sec:3.23}

Return to the stochastic-volatility construction of Section \hyperref[sec:3.12]{3.12}. Suppose each \(\theta_j\) is constant and nonnegative, and each \(\alpha_j(s)\) is nonnegative and piecewise constant with finitely many pieces on bounded intervals. In addition to the solution and supremum-square hypotheses of Section \hyperref[sec:3.12]{3.12}, assume that the noise integrands \(\alpha_j\sqrt{v_j}\) and the forward-rate loading integrands of (3.55) are predictable and have almost surely finite integrated square on every finite horizon. The SDE and forward curve use their Itô integrals. These domain conditions are imposed on the given almost surely continuous state process; existence of such a process is an assumption.

The one-piece square-root transition law with positive mean reversion is stated in \cite{malham2008square} Proposition 1, and the time-dependent parameters used below are those of \cite{brace2024sofr} Table 2. Neither source gives the support of the meeting-variance vector after composition over pieces.

\phantomsection\label{res:Proposition-26}
\textbf{Proposition 26 (conditional).} Under the preceding hypotheses, let \(0 \le s < t < T_N\), retain \(M\) future meetings, and take \(C(t)\) and \(A(t)\) from (3.57). On a single piece of length \(h\), with starting value \(x \ge 0\), constant \(\alpha > 0\) and \(\theta > 0\), the conditional Laplace transform at the piece's end is

\begin{equation*}
E_Q[e^{-\lambda v_j(s+h)}\mid v_j(s)=x]
= \left[1+\frac{\alpha^2\lambda}{2\theta}(1-e^{-\theta h})\right]^{-2\theta/\alpha^2}
\exp\!\left(-\frac{\lambda e^{-\theta h}x}{1+\frac{\alpha^2\lambda}{2\theta}(1-e^{-\theta h})}\right),\qquad \lambda\ge0. \tag{3.114}\label{eq:3.114}
\end{equation*}

For \(\theta=0\) and \(\alpha>0\) this is \(\exp[-\lambda x/(1+\alpha^2h\lambda/2)]\). For \(\alpha=0\), the coordinate follows \(v_j(s+h)=1+(x-1)e^{-\theta h}\). Composing these one-piece laws over the breakpoints gives the joint conditional law of \(v(t)\) given \(\mathcal F_s\): its coordinates are independent, and their joint Laplace transform is the product of the composed coordinate transforms.

Let \(J=\{j:\alpha_j>0\text{ on some piece of }(0,t]\}\), and let \(\ell_j\) be the lower endpoint of the support of \(v_j(t)\). If the last piece before \(t\) has \(\alpha_j>0\), then \(\ell_j=0\). If an earlier piece is stochastic but the last piece is deterministic, \(\ell_j\) is the flow of zero under \(u\mapsto1+(u-1)e^{-\theta_j h}\) across the final deterministic pieces. If every piece is deterministic, \(v_j(t)=\ell_j=1\). The support, affine hull and covariance of the actual vector of conditional short-rate jump variances are

\begin{equation*}
\begin{aligned}
V(t)&=C(t)+A(t)v(t),\\
\operatorname{supp}V(t)&=C(t)+A(t)\ell+A_J(t)[0,\infty)^J,\\
\operatorname{aff}\operatorname{supp}V(t)&=C(t)+A(t)\ell+\operatorname{range}A_J(t),\\
\operatorname{Cov}_Q(V(t))&=\sum_{j\in J}\operatorname{Var}_Q(v_j(t))A_{\cdot j}(t)A_{\cdot j}(t)^\top .
\end{aligned} \tag{3.115}\label{eq:3.115}
\end{equation*}

This support is a closed translated cone. Every variance in the sum is positive; hence the affine-hull dimension and covariance rank both equal \(\operatorname{rank}A_J(t)\), and the almost-sure affine equalities are exactly those given by \(w\in\ker A_J(t)^\top\). The mass at the vertex \(C(t)+A(t)\ell\) is the product of the lower-endpoint atom masses for active coordinates with nonzero columns. A factor with \(\theta_j>0\) has no such atom. A factor with \(\theta_j=0\) is absorbed when its starting value is zero; on a single stochastic piece of length \(h\) starting at \(x>0\), its atom at zero is \(\exp[-2x/(\alpha_j^2h)]\).

For the conditional law given \(\mathcal F_s\), replace \(\ell\) by the lower endpoints obtained by starting at \(v(s)\) and using the pieces on \((s,t]\), and replace \(J\) by

\begin{equation*}
J_s=\{j:\alpha_j>0\text{ on some piece of }(s,t],\quad \theta_j>0\text{ or }v_j(s)>0\}. \tag{3.116}\label{eq:3.116}
\end{equation*}

Then (3.115) holds for the regular conditional support, affine hull and covariance with \(A_{J_s}\) and conditional coordinate variances. In particular,

\begin{equation*}
E_Q[V(t)\mid\mathcal F_s]=C(t)+A(t)\bigl(1+(v_j(s)-1)e^{-\theta_j(t-s)}\bigr)_j,\qquad
\operatorname{Cov}_Q(V(t)\mid\mathcal F_s)=\sum_{j\in J_s}\operatorname{Var}_Q(v_j(t)\mid\mathcal F_s)A_{\cdot j}(t)A_{\cdot j}(t)^\top . \tag{3.117}\label{eq:3.117}
\end{equation*}

A positive-speed coordinate that is zero at \(s\) becomes active on any later stochastic piece; a zero-speed coordinate at zero stays absorbed.

For the time-dependent column of \cite{brace2024sofr} Table 2, the printed speeds are \((0.1,0,11.0)\) and every printed volatility-of-volatility value on its four pieces is strictly positive. At any observation time whose last piece carries one of those positive values, for any completion of the missing calibration inputs satisfying the hypotheses,

\begin{equation*}
\operatorname{supp}V(t)=C(t)+A(t)[0,\infty)^3,\qquad
\operatorname{aff}\operatorname{supp}V(t)=C(t)+\operatorname{range}A(t),\qquad
\operatorname{rank}\operatorname{Cov}_Q(V(t))=\operatorname{rank}A(t)\le3. \tag{3.118}\label{eq:3.118}
\end{equation*}

Thus the two-active-factor bound of Proposition \hyperref[res:Proposition-16]{16} for the printed constant column does not extend to the time-dependent column. For \(M>3\), at least \(M-3\) affine equalities remain. The vertex has no atom if either positive-speed factor has a nonzero column. Equation (3.118) does not assert \(\operatorname{rank}A(t)=3\); the loadings, dates and time conventions needed to determine that rank were not published.

\emph{Proof.} On one coefficient piece, use the exponential whose terminal value is \(e^{-\lambda v_j(s+h)}\) and whose two backward coefficients are the coefficient of \(x\) and the constant term in the exponent of (3.114). Differentiation in actual time gives for the first coefficient its mean-reversion multiple plus half its squared value times \(\alpha_j^2\), and for the second minus \(\theta_j\) times the first. Itô's formula then cancels the drift against these derivatives and the drift and quadratic variation of (3.54). Localize where the nonnegative state and coefficients are bounded; the stopped stochastic integral is a true martingale. The exponential lies in \((0,1]\), so bounded convergence removes the localization. Its conditional expectation gives (3.114), and the same argument across each finite common coefficient partition gives the product transform. The integrand assumptions make the integral identities valid for the original state; an everywhere-continuous adapted version gives predictable square-root integrands, and indistinguishability transfers the SDE, the forward curve and its short-rate jumps between the versions.

For positive speed and volatility of volatility, (3.114) is the Laplace transform of a scaled noncentral chi-square law, equivalently a Poisson mixture of laws with strictly positive density on the positive half-line. It has a positive density on \((0,\infty)\), support \([0,\infty)\) and no atom at zero even when \(x=0\). At zero speed it reduces to the absorbed law of Proposition \hyperref[res:Proposition-16]{16}; a deterministic piece applies the increasing affine flow to its input support. Composing these facts proves the coordinate supports, lower endpoints and atoms, and Laplace uniqueness gives the conditional product law.

Proposition \hyperref[res:Proposition-14]{14} gives the affine identity for the actual conditional short-rate jump variances. The image of the product of coordinate half-lines is (3.115); nonnegative columns make this finitely generated image closed, and the continuous-image support identity then makes it exact. Conditional independence makes the covariance a sum of positive weighted column outer products, so its kernel is \(\ker A_J(t)^\top\) and its rank equals \(\operatorname{rank}A_J(t)\). The same argument applied to the regular conditional product law proves (3.116) and (3.117); differentiating the one-piece transform at zero gives the displayed conditional mean, and iteration gives it across pieces. Finally all three printed factors have a positive last piece, so all three support half-lines contribute to (3.118), while positive mean reversion in factors one and three removes their lower-endpoint atoms.\hfill$\square$

The result is conditional on the stipulated nonnegative solutions, their supremum-square moment bound and the admissible Itô integrands. It does not construct the solutions, prove a calibrated rank, identify variances from listed prices, or establish restrictions surviving variation of the stochastic-volatility parameters.

\subsection{A bounded curve with a variance state and a scheduled jump}\label{sec:3.24}

Fix \(\lambda>0\) and one scheduled date \(\tau>0\). Let \(W_1,W_2\) be independent standard Brownian motions under \(Q\), with the completed right-continuous joint filtration, and let the smaller filtration be the completed natural filtration of \(W_2\) with the same null sets. Assume the cited Itô calculus is available on this basis in both filtrations: a supplied one-driver structure \(R\) for \(W_2\) in the smaller filtration and a supplied two-driver structure \(S\) for \((W_1,W_2)\) in the joint filtration, with the same measure and actual \(W_2\) path, unit diagonal and zero off-diagonal covariation densities, and the integral, Itô formula, Lipschitz SDE, continuous-adapted predictability and left-endpoint integral-approximation properties recorded in AX-03 to AX-06, AX-09 and AX-11. Assume also the bounded-quadratic-variation exponential-martingale property AX-10 in the joint calculus. These are cited stochastic-calculus inputs from \cite{bauerschmidt2020stochastic}; constructing both calculus structures on this Brownian basis is not part of the result.

\phantomsection\label{res:Proposition-27}
\textbf{Proposition 27 (conditional).} Under these inputs there is, up to indistinguishability, a unique continuous solution with zero initial coefficients, stopped at \(\tau\), of

\begin{equation*}
\begin{gathered}
v_t=2+\tanh Z_2(t),\qquad
Z_1(t)=\frac{t\wedge\tau}{\lambda}+W_1(t\wedge\tau),\\
Z_2(t)=\int_0^{t\wedge\tau}\frac{v_s-2}{2\lambda}\,ds+\int_0^{t\wedge\tau}\sqrt{v_s}\,dW_2(s),\qquad
Z_4(t)=-\int_0^{t\wedge\tau}\frac{v_s}{2\lambda}\,ds.
\end{gathered} \tag{3.119}\label{eq:3.119}
\end{equation*}

Define, with the right value at \(T=\tau\),

\begin{equation*}
f(t,T)=\begin{cases}0,&T<\tau,\\
Z_1(t)e^{-\lambda(T-\tau)}+Z_2(t)e^{-2\lambda(T-\tau)}+Z_4(t)e^{-4\lambda(T-\tau)},&T\ge\tau.
\end{cases} \tag{3.120}\label{eq:3.120}
\end{equation*}

Each curve is bounded in maturity by \(|Z_1(t)|+|Z_2(t)|+|Z_4(t)|\), is exponential-polynomial of degree zero to the right of the scheduled breakpoint, and has fixed exponents \(\lambda,2\lambda,4\lambda\). Boundedness is per realized curve in the sense of \cite{filipovic2000exponential} Definition 3.1, not uniform over all coefficient vectors or sample points. The fixed exponents and the two doubling pairs satisfy, rather than contradict, that source's exponent restrictions. Its fixed-maturity HJM drift equals the sum of volatility times integrated volatility. With \(r_t=f(t,t)\), \(B_t=\exp(\int_0^t r_s\,ds)\) and \(P(t,T)=\exp(-\int_t^T f(t,u)\,du)\), every \(P(\cdot,T)/B\) is a positive true \(Q\)-martingale on \([0,T]\). The curve is continuous in time at fixed maturity, while the short rate has an almost surely nonzero jump at \(\tau\):

\begin{equation*}
\Delta r_\tau=Z_1(\tau)+Z_2(\tau)+Z_4(\tau)
=W_1(\tau)+\int_0^\tau\sqrt{v_s}\,dW_2(s)\ne0\quad Q\text{-a.s.} \tag{3.121}\label{eq:3.121}
\end{equation*}

Before \(\tau\), \(1<v_t<3\), \(d[Z_2]_t=v_t\,dt\), and \(d[v]_t=(1-\tanh^2 Z_2(t))^2v_t\,dt>0\). The three forward rates at maturities \(\tau,\tau+\log(2)/\lambda,\tau+2\log(2)/\lambda\) recover \(Z_1(t),Z_2(t),Z_4(t)\) for \(t\le\tau\), hence also \(v_t\). Thus the stochastic variance state enters the curve itself, including the coefficient of a shape generated by the first driver's drift. This bounded family gives a counterexample to the exclusion conjecture for curve-entering variance with scheduled breakpoints.

\emph{Proof.} The scalar drift \(\tanh z/(2\lambda)\) and diffusion \(\sqrt{2+\tanh z}\) are bounded and globally Lipschitz: their derivatives have absolute values at most \(1/(2\lambda)\) and \(1/2\). The cited Lipschitz theorem in the smaller filtration supplies \(Z_2\), and the other coordinates follow from (3.119). Continuous-adapted predictability and boundedness put the integrands in the stated Itô domains. Left-endpoint approximation identifies the \(W_2\) integrals in the two filtrations; uniqueness of their limits at each rational time and path continuity give one almost-sure identity at all times. Stopping at \(\tau\) and uniqueness in the joint filtration give the asserted stopped system.

For \(t<\tau\le T\), set \(x=T-\tau\). The two volatility shapes are \(\sigma_1=e^{-\lambda x}\) and \(\sigma_2=\sqrt{v_t}e^{-2\lambda x}\); their maturity integrals are \((1-e^{-\lambda x})/\lambda\) and \(\sqrt{v_t}(1-e^{-2\lambda x})/(2\lambda)\). Equation (3.119) gives

\begin{equation*}
\alpha(t,T)=\frac{e^{-\lambda x}-e^{-2\lambda x}}{\lambda}
+\frac{v_t(e^{-2\lambda x}-e^{-4\lambda x})}{2\lambda}
=\sum_{j=1}^2\sigma_j(t,T)\int_t^T\sigma_j(t,u)\,du. \tag{3.122}\label{eq:3.122}
\end{equation*}

For \(T<\tau\) or \(t\ge\tau\), the fixed-maturity drift and volatility vanish. There is no fixed-maturity curve jump. For \(T\ge\tau\), put \(A_k(T)=(1-e^{-k\lambda(T-\tau)})/(k\lambda)\), \(k=1,2,4\). Before \(\tau\), \(P(t,T)=\exp(-\sum_k Z_k(t)A_k(T))\) and \(B_t=1\). Itô's formula and

\begin{equation*}
\frac{A_1}{\lambda}+\frac{(v_t-2)A_2}{2\lambda}-\frac{v_t A_4}{2\lambda}
=\frac{A_1^2+v_t A_2^2}{2} \tag{3.123}\label{eq:3.123}
\end{equation*}

give \(d(P/B)/(P/B)=-A_1\,dW_1-\sqrt{v_t}A_2\,dW_2\). After \(\tau\) the coefficients are frozen, \(d\log P(t,T)/dt=r_t\), and the discounted bond is constant; the values match at \(\tau\). Its stochastic logarithm has quadratic variation at most \(7\tau/(4\lambda^2)\), since \(v<3\). AX-10 makes the resulting exponential a true martingale. If \(T<\tau\), the discounted bond is one.

Itô's formula gives the two bracket densities in the statement and the inverse \(Z_2=\tfrac12\log[(v-1)/(3-v)]\). The evaluation matrix at the three maturities has rows \((1,1,1)\), \((1/2,1/4,1/16)\), \((1/4,1/16,1/256)\), with determinant \(-21/1024\), so the curve recovers \(v\). The \(e^{-2\lambda x}\) coefficient \(Z_2\) has martingale part \(\sqrt v\,dW_2\), whereas \(Z_4\) has finite variation and closes the HJM drift without further exponents. Finally the drift of \(Z_1+Z_2+Z_4\) cancels. Its remaining \(W_2\) integral is measurable for the completed \(W_2\)-path sigma-algebra, by the smaller-filtration construction and integral comparison. The independent \(W_1(\tau)\) has a non-atomic Gaussian law, proving (3.121).\hfill$\square$

The proposition is conditional on the supplied two-filtration calculus. Their existence on the Brownian basis is not constructed here or established by the degenerate Upstream instances. It covers one scheduled date with a maturity breakpoint and short-rate jump, but no fixed-maturity curve jump, no post-meeting diffusion and no classification of other exponent sets or variance dynamics.

\subsection{Separable volatility with meeting loadings}\label{sec:3.25}

Fix a horizon \(H>T_N\); all times and maturities in this section lie in \([0,H]\), and \(I_N\) is used only up to \(H\). Let \(W^1,\dots,W^d\) be independent standard Brownian motions on a filtered probability space satisfying the usual conditions. Assume the cited Itô calculus of this basis: for predictable integrands \(h\) with \(\int_0^H h(s)^2\,ds<\infty\) almost surely, the integral \(\int_0^t h\,dW^j\) has a continuous adapted version, is linear in \(h\), and commutes with stopping at deterministic times, as recorded in \cite{bauerschmidt2020stochastic} §2.2 and §3.1 to §3.3; the driver covariations are the identity. Constructing this calculus is not part of the result. Let each scale \(\chi_j\) be predictable with \(\int_0^H\chi_j(s)^2\,ds<\infty\) almost surely; no equation for \(\chi_j\) is assumed. Let \(\varphi_j:[0,H]\to\mathbb R\) be continuous deterministic maturity shapes, which may vanish or change sign, and \(a_{j,k,m}\) real constants, \(1\le j\le d\), \(0\le k,m\le N\). Put

\begin{equation*}
g_{j,k}(T):=a_{j,k,m}\varphi_j(T)\ (T\in I_m),\qquad G_{j,k}(T):=\int_0^T g_{j,k}(u)\,du,\qquad \Phi_{j,m}(T):=\int_{T_m}^T\varphi_j(u)\,du. \tag{3.124}\label{eq:3.124}
\end{equation*}

Here \(\Phi_{j,m}\) is a deterministic primitive, not the accumulated variance \(\Phi\) of \cite{brace2024sofr} eq. 28, and the functions \(G_{j,k}(\cdot)\) are not the constants \(G_{i,j}\) of Section \hyperref[sec:3.11]{3.11}. The constants let the maturity loading change when time crosses a meeting. The ordinal loadings of Section \hyperref[sec:3.11]{3.11} are the case \(a_{j,k,m}=G_{m-k,j}\) for \(m\ge k\) and \(a_{j,k,m}=0\) otherwise, which includes zero volatility on the current maturity interval. For \(t\in I_k\) and \(t\le T\) set

\begin{equation*}
\sigma_j(t,T):=\chi_j(t)\,g_{j,k}(T),\qquad \alpha(t,T):=\sum_{j=1}^d\chi_j(t)^2\,g_{j,k}(T)\bigl[G_{j,k}(T)-G_{j,k}(t)\bigr], \tag{3.125}\label{eq:3.125}
\end{equation*}

take every curve jump \(\xi_n\) to be zero, take a bounded deterministic initial curve \(f(0,\cdot)\), continuous within each \(I_m\) with finite one-sided limits, and define \(f\) by (2.2). Values at the finitely many meeting times do not affect any integral. Put

\begin{equation*}
A_{j,k}(t):=\int_0^t\mathbf 1_{I_k}(s)\chi_j(s)^2\,ds,\qquad M_{j,k}(t):=\int_0^t\mathbf 1_{I_k}(s)\chi_j(s)\,dW^j_s-\int_0^t\mathbf 1_{I_k}(s)\chi_j(s)^2G_{j,k}(s)\,ds, \tag{3.126}\label{eq:3.126}
\end{equation*}

and, for each maturity interval \(I_m\),

\begin{equation*}
C_{j,m}(t):=\sum_{k}a_{j,k,m}\bigl[M_{j,k}(t)+G_{j,k}(T_m)A_{j,k}(t)\bigr],\qquad V_{j,m}(t):=\sum_k a_{j,k,m}^2A_{j,k}(t). \tag{3.127}\label{eq:3.127}
\end{equation*}

\phantomsection\label{res:Proposition-28}
\textbf{Proposition 28 (conditional).} Under the stated calculus inputs:

(a) The coefficients (3.125) satisfy Assumption 2.1 for every \(t\le T\), not only almost everywhere, and Assumption 2.2 holds trivially because there are no curve jumps. There is a version of the curve, continuous in \(t\), such that simultaneously for all \(t\le T\le H\)

\begin{equation*}
f(t,T)=f(0,T)+\sum_{j,k}g_{j,k}(T)\bigl[M_{j,k}(t)+G_{j,k}(T)A_{j,k}(t)\bigr], \tag{3.128}\label{eq:3.128}
\end{equation*}

and for each fixed \(T\) it agrees with the process (2.2). The current curve is therefore determined by the initial curve and at most \(2d(N+1)\) real coefficients, which do not depend on \(T\). Each \(A_{j,k}\) is adapted, continuous, nonnegative, nondecreasing and of finite variation; each \(M_{j,k}\) is an adapted continuous semimartingale. Both vanish up to \(T_k\) and are constant after \(T_{k+1}\) when \(k<N\): crossing a meeting retains the earlier coefficients and starts the next integrals at zero.

(b) For \(T\in I_m\) and \(t\le T\),

\begin{equation*}
f(t,T)=f(0,T)+\sum_{j=1}^d\varphi_j(T)\bigl[C_{j,m}(t)+\Phi_{j,m}(T)\,V_{j,m}(t)\bigr]. \tag{3.129}\label{eq:3.129}
\end{equation*}

The two shapes per factor are \(\varphi_j\) and \(\varphi_j\Phi_{j,m}\), restarted at each maturity meeting. The coefficient \(V_{j,m}\) is realized accumulated variance, adapted, continuous, nonnegative and nondecreasing, and has no martingale part; it is not the conditional variance of a future short-rate jump. For \(t\in I_k\), in integral form,

\begin{equation*}
dC_{j,m}=a_{j,k,m}\chi_j\,dW^j+a_{j,k,m}\chi_j^2\bigl[G_{j,k}(T_m)-G_{j,k}(t)\bigr]dt,\qquad dV_{j,m}=a_{j,k,m}^2\chi_j^2\,dt. \tag{3.130}\label{eq:3.130}
\end{equation*}

(c) If \(\varphi_j=1\) for every \(j\) and \(f(0,\cdot)\) is piecewise affine, (3.129) is a step-plus-ramp curve. On \(I_m\) its level at \(T_m\) is \(f(0,T_m)+\sum_jC_{j,m}(t)\), and its slope is the initial slope plus \(\sum_jV_{j,m}(t)\). The additional slope is adapted and of finite variation with derivative \(\sum_ja_{j,k,m}^2\chi_j^2\) on \(I_k\); it need not be deterministic. This is the local slope on \(I_m\), not the coefficient of one ramp in a sum over overlapping maturity intervals.

The separable device \(\sigma(t,T)=\chi(t)\varphi(T)\) with a stochastic scale, and its exponential meeting instance with ordinal indicators, are \cite{brace2024sofr} eqs. 20 to 22 and 36 to 43. Finite-dimensional realizations of calendar-fixed shapes, with running time as a state, are known: \cite{bjork2001existence} §3.3 with Theorem 3.4 and Proposition 6.2, and, for the Gaussian case with scheduled jumps, \cite{fontana2022term} Example 3.10. Proposition \hyperref[res:Proposition-28]{28} extends that known computation to loadings that restart at the meetings, stochastic scales and continuous shapes, and is not new in substance. No exponential restriction on the shapes is needed for (a) to (c).

\emph{Proof.} Continuity of \(\varphi_j\) on \([0,H]\) and finitely many loadings make every \(g_{j,k}\) bounded, and \(G_{j,k}\) bounded and absolutely continuous with derivative \(g_{j,k}\) almost everywhere. Masking \(\chi_j\) by a deterministic time interval and multiplying by bounded deterministic functions preserves predictability and almost sure finiteness of \(\int_0^H\chi_j^2\,ds\), so every stochastic integrand below lies in the domain of the cited integral, and the drift integrals exist pathwise. Each masked stochastic integral is the difference of the integral stopped at \(T_{k+1}\) and at \(T_k\), hence zero up to \(T_k\) and constant after \(T_{k+1}\); its finite-variation companion in (3.126) has the same property. There are finitely many such integrals, so their continuous versions can be chosen on one event of probability one.

For \(t\in I_k\) put \(D_j(T):=G_{j,k}(T)-G_{j,k}(t)\). Absolute continuity gives

\begin{equation*}
\int_t^T g_{j,k}(u)\,D_j(u)\,du=\tfrac12D_j(T)^2, \tag{3.131}\label{eq:3.131}
\end{equation*}

even where \(g_{j,k}\) jumps at a maturity meeting or changes sign, since \(G_{j,k}\) is continuous and its square satisfies the ordinary chain rule. Multiplying by \(\chi_j(t)^2\) and summing over \(j\) gives \(\int_t^T\alpha(t,u)\,du=\tfrac12\sum_j(\int_t^T\sigma_j(t,u)\,du)^2\), which is (2.3) for every \(t\le T\); independent drivers give the sum of squares without cross terms.

Split each time integral in (2.2) into the finitely many intervals \(I_k\). On the \(k\)-th piece \(g_{j,k}(T)\) and \(G_{j,k}(T)\) do not depend on the integration variable, so

\begin{equation*}
\int_0^t\mathbf 1_{I_k}(s)\,\chi_j(s)^2g_{j,k}(T)\bigl[G_{j,k}(T)-G_{j,k}(s)\bigr]ds=g_{j,k}(T)G_{j,k}(T)A_{j,k}(t)-g_{j,k}(T)\int_0^t\mathbf 1_{I_k}(s)\chi_j(s)^2G_{j,k}(s)\,ds, \tag{3.132}\label{eq:3.132}
\end{equation*}

and the diffusion term is \(g_{j,k}(T)\) times the stochastic integral in (3.126), by linearity of the integral in the integrand. Summing gives (3.128) with finite sums only; no stochastic Fubini theorem is used. The right side of (3.128) is one version for all maturities; for each fixed \(T\) it agrees with the integral process (2.2), so no uncountable family of almost sure events is intersected.

For \(T\in I_m\) the primitive splits as

\begin{equation*}
G_{j,k}(T)=G_{j,k}(T_m)+a_{j,k,m}\Phi_{j,m}(T). \tag{3.133}\label{eq:3.133}
\end{equation*}

Substituting (3.133) and \(g_{j,k}(T)=a_{j,k,m}\varphi_j(T)\) into (3.128) and collecting the two shapes gives (3.127) and (3.129); the constants \(G_{j,k}(T_m)\) carry the contributions of all earlier maturity intervals and are retained in \(C_{j,m}\). Differentiating (3.126) in integral form gives (3.130). The coefficients of \(V_{j,m}\) are squares and \(A_{j,k}\) is an integral of a square, which gives the monotonicity and finite-variation assertions. Finally \(\varphi_j=1\) gives \(\Phi_{j,m}(T)=T-T_m\), and (3.129) with (3.130) gives (c).\hfill$\square$

The proposition is conditional on the cited calculus and on the given scales; it constructs neither, and no stochastic-volatility equation for \(\chi_j\) is solved. It is a finite representation of the current curve, not a finite-dimensional Markov realization: finitely many current coefficients do not determine the conditional law of the future scales, and a model for \(\chi_j\) together with a proof that an augmented state closes is required; Proposition \hyperref[res:Proposition-29]{29} supplies both for a Lipschitz scale state, with a time-inhomogeneous realization. The statement that every deterministic shape gives a finite-dimensional realization is therefore established only in this weaker sense. The deterministic time dependence in (3.130) is not removed, so the proposition says nothing about time homogeneity and neither proves nor refutes Conjecture 2. The bound \(2d(N+1)\) is not a minimal dimension, and the coefficients are not claimed to be recoverable from the curve. Not covered: nonzero fixed-maturity curve jumps and restart laws; the law of the short-rate jump that crossing a maturity breakpoint produces; true bond martingales and bond prices; correlated drivers, and several independently scaled shapes within one driver on one time interval, which would add cross terms.

\subsection{A Markov realization with a Lipschitz scale state}\label{sec:3.26}

Keep the setting and notation of Section \hyperref[sec:3.25]{3.25}. Assume in addition that the drivers are \(m\ge d\) independent standard Brownian motions \(B^1,\dots,B^m\) whose increments after each time \(s\) are independent of \(\mathcal F_s\), with \(W^j=B^j\) for \(j\le d\). Assume that the cited calculus supplies, besides the integral of Section \hyperref[sec:3.25]{3.25}, strong solutions of equations with Lipschitz coefficients (\cite{bauerschmidt2020stochastic} §5.2), predictability of real adapted processes all of whose paths are continuous (\cite{bauerschmidt2020stochastic} p. 14), and the Markov property of solutions of equations whose coefficients are Borel in time and Lipschitz in the state (\cite{vanhandel2007stochastic} Theorem 5.2.2). Constructing this calculus is not part of the result. Fix \(p\ge1\), \(z_0\in\mathbb R^p\), and coefficients \(\beta:\mathbb R_{\ge0}\times\mathbb R^p\to\mathbb R^p\) and \(\Sigma:\mathbb R_{\ge0}\times\mathbb R^p\to\mathbb R^{p\times m}\) that are Borel in \((t,z)\), Lipschitz in \(z\) with one constant for all \(t\), and bounded on every set \([0,T']\times\{|z|\le R\}\). The scale state solves

\begin{equation*}
Z_t=z_0+\int_0^t\beta(u,Z_u)\,du+\sum_{l=1}^m\int_0^t\Sigma_{\cdot l}(u,Z_u)\,dB^l_u, \tag{3.134}\label{eq:3.134}
\end{equation*}

and may use the curve drivers, so leverage is allowed. Let \(\psi_j:\mathbb R_{\ge0}\times\mathbb R^p\to\mathbb R\) be jointly continuous, bounded by \(K_\psi\) and Lipschitz in \(z\) uniformly in \(t\), and take the scales of Section \hyperref[sec:3.25]{3.25} to be \(\chi_j(t)=\psi_j(t,Z_t)\). With the coefficients (3.126), put \(n=p+2d(N+1)\) and

\begin{equation*}
X_t=\bigl(Z_t,(M_{j,k}(t))_{j,k},(A_{j,k}(t))_{j,k}\bigr)\in\mathbb R^n. \tag{3.135}\label{eq:3.135}
\end{equation*}

\phantomsection\label{res:Proposition-29}
\textbf{Proposition 29 (conditional).} Under these inputs:

(a) Equation (3.134) has a solution every path of which is continuous, and every solution is indistinguishable from one; with such a \(Z\), each \(\chi_j\) is predictable and bounded, so Proposition \hyperref[res:Proposition-28]{28} applies.

(b) Write \(e_H=\mathbf 1_{[0,H]}\). The frozen state \(X_{t\wedge H}\) solves, from \(x_0=(z_0,0,0)\), the equation whose \(Z\)-block has the coefficients \(e_H\beta\) and \(e_H\Sigma\), and whose other blocks have, at \(x=(z,\mu,a)\), the drifts and the diffusion against \(B^j\)

\begin{equation*}
\begin{gathered}
\beta_{M_{j,k}}(t,x)=-\mathbf 1_{I_k}(t)e_H(t)\psi_j(t,z)^2G_{j,k}(t\wedge H),\qquad \beta_{A_{j,k}}(t,x)=\mathbf 1_{I_k}(t)e_H(t)\psi_j(t,z)^2, \\
\Sigma_{M_{j,k},j}(t,x)=\mathbf 1_{I_k}(t)e_H(t)\psi_j(t,z),
\end{gathered} \tag{3.136}\label{eq:3.136}
\end{equation*}

with every other diffusion entry zero. These coefficients are Borel in \((t,x)\), Lipschitz in \(x\) uniformly in \(t\), bounded on bounded sets, and do not depend on \((\mu,a)\); they switch at the meetings.

(c) The frozen state is Markov for the filtration: for deterministic \(s\le t_1,\dots,t_r\) and bounded Borel \(F\), \(E[F(X_{t_1\wedge H},\dots,X_{t_r\wedge H})\mid\mathcal F_s]\) is almost surely a Borel function of \(X_{s\wedge H}\).

(d) The affine map, which does not depend on \(t\),

\begin{equation*}
\Lambda(x)(T)=f(0,T)+\sum_{j,k}g_{j,k}(T)\bigl[\mu_{j,k}+G_{j,k}(T)a_{j,k}\bigr] \tag{3.137}\label{eq:3.137}
\end{equation*}

satisfies \(f(t,T)=\Lambda(X_t)(T)\) simultaneously for all \(0\le t\le T\le H\). Hence, for finitely many pairs \(s\le t_i\le T_i\le H\) and a bounded Borel function \(F\) of the values \(f(t_i,T_i)\), \(E[F\mid\mathcal F_s]\) is almost surely a Borel function of \(X_s\). The state dimension, running time excluded, is at most \(p+2d(N+1)\), for every choice of continuous shapes.

(e) If \(a_{j,k,k}=0\) for all \(j,k\), as for the ordinal loadings of Section \hyperref[sec:3.11]{3.11}, and \(\beta\), \(\Sigma\) and \(\psi_j\) do not depend on \(t\), then on each \(I_k\cap[0,H]\) the coefficients of (b) are functions of \(x\) alone: their only time dependence is the switch at each meeting and the freeze at \(H\), for every continuous \(\varphi_j\). If instead \(\psi_j\) is not identically zero and \(a_{j,k,k}\varphi_j\) is not identically zero on \(I_k\cap[0,H]\), the drift of \(M_{j,k}\) depends on \(t\) within \(I_k\).

In calendar maturity the map \(\Lambda\) is fixed and the time dependence sits in the coefficients (3.136), through the meeting indicators and \(G_{j,k}(t)\); in time to maturity \(x=T-t\) the curve is \(\Lambda(X_t)(t+x)\), which depends on \(t\) explicitly. Realizations of time-varying systems with running time as a state are \cite{bjork2001existence} §3.3 (Theorem 3.4) with Proposition 6.2, \cite{fontana2022term} Example 3.10 is the Gaussian case with scheduled jumps, and \cite{brace2024sofr} §3.5 and §3.6 is the exponential instance with a Heston scale. Proposition \hyperref[res:Proposition-29]{29} extends these known calendar-maturity realizations to meeting restarts, arbitrary continuous shapes and a Lipschitz scale state, and is not new in substance.

\emph{Proof.} (a) A solution with almost surely continuous paths exists by the cited strong-existence result. Its paths are continuous and start at \(z_0\) off a null set, which lies in \(\mathcal F_0\) by the usual conditions; setting the process equal to \(z_0\) on that set gives an adapted process with every path continuous, indistinguishable from the original. Each coordinate of it, and running time, are then predictable by the cited predictability result, so \(\Sigma(s,Z_s)\), \(\beta(s,Z_s)\) and \(\psi_j(s,Z_s)\) are predictable, and the box bounds make them square integrable in time on every path. The replaced and original diffusion integrands differ by an admissible integrand that vanishes off a null set, so its integral vanishes by the isometry and the two stochastic integrals agree almost surely at each time; both sides of (3.134) are continuous in \(t\), so agreement at rational times gives agreement at all times. Finally \(|\chi_j|\le K_\psi\).

(b) Each block of (3.136) is a product of Borel functions of \(t\) with a continuous function of \((t,z)\). From \(|\psi_j(t,z)^2-\psi_j(t,z')^2|\le2K_\psi L_\psi|z-z'|\), with \(L_\psi\) the Lipschitz constant of \(\psi_j\), and the bound \(\bar G\) on every \(|G_{j,k}|\) over \([0,H]\), the blocks are Lipschitz in \(z\) with constants \(2K_\psi L_\psi\bar G\), \(2K_\psi L_\psi\) and \(L_\psi\), uniformly in \(t\), and bounded by \(K_\psi^2(1+\bar G)+K_\psi\). The integrands are admissible by (a). On \([0,H]\) the definitions (3.126) with \(\chi_j=\psi_j(\cdot,Z)\) are, coordinate by coordinate, the equations of the blocks (3.136); the drivers other than \(B^j\) enter with the zero integrand. The \(Z\)-block is (3.134) at \(t\wedge H\), by the stopped-integral identity at the deterministic time \(H\). After \(H\) the coefficients vanish, which is the freeze; the identities at each time extend to all times by continuity, as in (a).

(c) The cited Markov property applies to the equation of (b), since its coefficients are Borel in time and Lipschitz in the state and its initial value is deterministic; it gives the case \(r=1\), and a Borel function of \(X_{s\wedge H}\) by the Doob--Dynkin lemma. For products \(F_1(x_1)\cdots F_r(x_r)\) with the times ordered, the tower property at the next-to-last time reduces \(r\) to \(r-1\), with \(F_{r-1}\) replaced by \(F_{r-1}\) times the bounded Borel function from the case \(r=1\). Products of indicators of Borel sets form a \(\pi\)-system generating the Borel sets of \((\mathbb R^n)^r\), and the functional monotone class theorem extends the property to every bounded Borel \(F\).

(d) Formula (3.128) is \(f(t,T)=\Lambda(X_t)(T)\) with the coordinates of \(X_t\) substituted, and \(X=X_{\cdot\wedge H}\) on \([0,H]\); each \(\Lambda(\cdot)(T_i)\) is affine, so (c) applies. The dimension is that of (3.135).

(e) For \(t\in I_k\cap[0,H]\), (3.133) gives \(G_{j,k}(t)=G_{j,k}(T_k)+a_{j,k,k}\Phi_{j,k}(t)\), which is constant when \(a_{j,k,k}=0\); the indicators are then constant on that set, and every block of (3.136) is a function of \(x\) alone. Conversely, choose \(z\) with \(\psi_j(z)\neq0\); \(\Phi_{j,k}\) is constant on the nondegenerate interval \(I_k\cap[0,H]\) only if \(\varphi_j\) vanishes there, by continuity.\hfill$\square$

The proposition is conditional on the cited calculus, the Brownian drivers and the three cited results above; it constructs none of them. It is a finite-dimensional, time-inhomogeneous Markov realization of the current curve, with running time as a state. It is not the realization in time to maturity with a schedule entering only through the distances to the meetings that Conjecture 2 concerns, and it proves no nonexistence of such a realization. Part (e) concerns the coordinates (3.135) in calendar maturity: it answers only the weaker reading in which the state equations between meetings have no explicit time dependence, and says nothing about other realizations. Not covered: transition functions indexed by the initial state, their composition law, the Feller or strong Markov property; a scale built from a solution whose paths are continuous only almost surely, unless its masked scales are separately known to be admissible integrands; square-root or unbounded scales, such as the Heston scale of \cite{brace2024sofr} eqs. 29 and 30, whose coefficients are not globally Lipschitz; minimality of the dimension and recovery of the state from the curve, since frozen coordinates remain in the state; true bond martingales; and correlated curve drivers.

\subsection{Meeting loadings obstruct realizations in time to maturity}\label{sec:3.27}

\textbf{Realizations in time to maturity.} Anchor the curve at the current time, with time to maturity \(x=T-t\). On the maturities between consecutive remaining meetings, the volatility of factor \(j\) is an ordinal loading, indexed by the number of meetings in \((t,t+x]\), times a fixed continuous function \(\lambda_j(x)\) times a function of a scale state that solves an autonomous Lipschitz equation. A realization is a finite-dimensional state with the same equation coefficients on every scheduled interval, a fixed restart map applied at each meeting, and a fixed map from the state and the distances to the remaining meetings to the curve; these are the same for every schedule, which enters only through those distances, and the state dimension counts every coordinate, running time included. The curve map is required to be locally Lipschitz in the state for each fixed distance vector and maturity, as the smoothness in \cite{bjork2001existence} Definition 3.1 implies. \cite{brace2024sofr} eqs. 36 to 38 with constant mean reversion is the case \(\lambda_j(x)=e^{-\kappa_jx}\).

This section treats one factor, a constant scale and no fixed-maturity curve jumps. Fix real ordinal loadings \(a_0,a_1,a_2,\dots\), independent of the schedule, and a continuous \(\lambda:[0,\infty)\to\mathbb R\). For a schedule \(0=T_0<T_1<\dots<T_N\) and \(t\ge0\), let \(i(t,T)\) be the number of meetings in \((t,T]\) and \(D(t)\) the vector of distances from \(t\) to the remaining meetings. The volatility is

\begin{equation*}
\sigma(t,T)=a_{i(t,T)}\,\lambda(T-t),\qquad 0\le t\le T. \tag{3.138}\label{eq:3.138}
\end{equation*}

A realization with state dimension \(q\) consists of a state \(Y_t\) in a set \(Z\subseteq\mathbb R^q\) and a map \(G\), the same for every schedule, with

\begin{equation*}
f(t,t+x)=G(Y_t,D(t))(x)\quad Q\text{-a.s., for every }t\ge0\text{ and }x\ge0, \tag{3.139}\label{eq:3.139}
\end{equation*}

with the null set allowed to depend on \((t,x)\), and such that \(y\mapsto G(y,D)(x)\) is locally Lipschitz on \(Z\) for each fixed \(D\) and \(x\), as the smoothness of the curve map in \cite{bjork2001existence} Definition 3.1 implies. The identity is required in this version form because the curve at each maturity involves its own stochastic integral, which the cited calculus fixes only up to a null set of its own; one null set for all \(x\) does not follow from it. Nothing about the state dynamics or the restart map is used. For \(R,n\ge0\) the Hankel section \(H_{R,n}\) is the \((R+1)\times(n+1)\) matrix with entries \((H_{R,n})_{i,l}=a_{i+l}\).

\phantomsection\label{res:Proposition-30}
\textbf{Proposition 30 (conditional).} Let \(W\) be a standard Brownian motion on a filtered probability space with the usual conditions, with the cited Itô calculus of Section \hyperref[sec:3.25]{3.25}, and let \(\lambda\equiv1\).

(a) Fix \(R,n\ge0\), the schedule \(T_m=m\) for \(m=1,\dots,n+R\), and \(t=n+\tfrac12\). The curve (2.2) with volatility (3.138), the drift that satisfies Assumption 2.1, and any initial curve satisfies, at the maturities \(T^{(i)}=n+i+\tfrac34\), \(i=0,\dots,R\),

\begin{equation*}
\bigl(f(t,T^{(i)})\bigr)_{i=0}^R=c+H_{R,n}\,\Delta W\quad Q\text{-a.s.}, \tag{3.140}\label{eq:3.140}
\end{equation*}

with \(c\) deterministic, \(\Delta W_0=W_t-W_{T_n}\) and \(\Delta W_l=W_{T_{n-l+1}}-W_{T_{n-l}}\) for \(1\le l\le n\). If a realization with state dimension \(q\) reproduces this curve at \(t\), then \(\operatorname{rank}H_{R,n}\le q\).

(b) If a realization with state dimension \(q\) serves every schedule, then \(\operatorname{rank}H_{R,n}\le q\) for all \(R,n\), and there are \(L\le q\) and real \(c_0,\dots,c_{L-1}\) with

\begin{equation*}
a_{i+L}=\sum_{l=0}^{L-1}c_l\,a_{i+l}\qquad\text{for every }i\ge0. \tag{3.141}\label{eq:3.141}
\end{equation*}

(c) For the cumulative ordinal loadings of \cite{brace2024sofr} eq. 36 with meeting weights \(\gamma_h=1/h\), that is \(a_0=0\) and \(a_i=\sum_{h=1}^i1/h\), no recurrence (3.141) holds. Hence for \(\lambda\equiv1\), which is quasi-exponential and even exponential (\(\kappa=0\)), no realization serves every schedule, for any \(q\).

The obstruction comes from the loadings, not from the shape: without meetings only \(a_0\) enters and the Hankel condition is vacuous, and the answer is the quasi-exponential condition on \(\lambda\) of \cite{bjork2001existence} Proposition 5.1 and Corollary 5.1. None of the realization results cited in Section \hyperref[sec:4]{4} has scheduled maturity breakpoints. Finite Hankel rank and linear recurrences are equivalent by classical linear algebra; the direction needed is proved below. The condition (3.141) is the discrete analogue, for the meeting loadings, of the quasi-exponential condition on the shape.

\emph{Proof.} (a) At each fixed maturity the volatility \(\sigma(\cdot,T)\) is deterministic and piecewise constant. For \(s\in I_k\) and \(T\in I_m\), \(m\ge k\), the number of meetings in \((s,T]\) is \(m-k\). The maturity \(T^{(i)}\) lies strictly inside \(I_{n+i}\), and \(t\in I_n\), so the driver integral up to \(t\) is \(\sum_{k=0}^na_{n+i-k}\) times the increment of \(W\) over \(I_k\cap[0,t]\), by the elementary values and linearity of the integral; changing the integrand at the finitely many meeting times changes the integral by one of zero variance, hence not at all. Putting \(l=n-k\) gives (3.140); the initial curve and the drift integral enter only through \(c\). The pieces are disjoint intervals of positive length, so \(\Delta W\) is a vector of independent centered Gaussian variables with positive variances, and its law is absolutely continuous on \(\mathbb R^{n+1}\). Let \(r=\operatorname{rank}H_{R,n}\) and choose \(r\) independent rows, with coordinate projection \(P\). Then \(PH_{R,n}\) maps onto \(\mathbb R^r\), so the law of \(V=P(c+H_{R,n}\Delta W)\) is absolutely continuous on \(\mathbb R^r\). By (3.139), \(V=\Phi(Y_t)\) almost surely, with \(\Phi(y)=(G(y,D(t))(T^{(i)}-t))_{i}\) over the chosen rows and \(D(t)\) deterministic; \(\Phi\) is locally Lipschitz on \(Z\). If \(r>q\), cover \(Z\) by countably many compact cubes on each of which \(\Phi\) is Lipschitz; splitting a cube into \(K^q\) subcubes of side \(1/K\) covers its image by \(K^q\) balls of total volume of order \(K^{q-r}\to0\). So \(\Phi(Z)\) lies in a Lebesgue-null set, which contradicts the absolute continuity of the law of \(V\). Hence \(r\le q\).

(b) Apply (a) to every \((R,n)\); one \(q\) serves all, because the realization is the same for every schedule. Fix \(n=q\) and let \(K_R\) be the set of unit vectors \(c\in\mathbb R^{q+1}\) with \(\sum_{l=0}^qc_la_{i+l}=0\) for all \(i\le R\). It is nonempty because the \(q+1\) columns of \(H_{R,q}\) are dependent, compact, and decreasing in \(R\), so the \(K_R\) share a unit vector \(c\). Let \(L\) be the largest index with \(c_L\neq0\). If \(L=0\), then \(a\equiv0\) and (3.141) holds with \(L=0\); otherwise dividing by \(-c_L\) gives (3.141).

(c) A recurrence (3.141) for \(a_i\) passes to \(\gamma_{i+1}=a_{i+1}-a_i\): \(1/(i+L+1)=\sum_{l<L}c_l/(i+l+1)\) for every \(i\ge0\). Both sides are rational functions of \(i\) that agree at infinitely many points, hence identical; but the left side has a pole at \(-L-1\) and the right side has poles only in \(\{-1,\dots,-L\}\). The case \(L=0\) would force \(a\equiv0\). By (b), no \(q\) works.\hfill$\square$

The same argument excludes every ordinal loading whose meeting weights satisfy no linear recurrence. The proposition is conditional on the cited calculus. It treats only \(\lambda\equiv1\), one factor and a constant scale; Proposition \hyperref[res:Proposition-33]{33} extends necessity of (3.141) to every nonzero quasi-exponential shape with a constant scale; for other shapes, several factors and stochastic scales it is not shown. Sufficiency, whether a quasi-exponential shape with loadings satisfying (3.141) always admits a realization, is not shown here. The obstruction needs the regularity of \(G\): with a merely measurable curve map and restart map dimension arguments fail. It uses (3.139) at one time and so does not cover local realizations up to a stopping time that may precede \(t\). It needs one realization to serve schedules with arbitrarily many meetings: on a single schedule with \(N\) meetings the Hankel sections that arise have rank at most \(N+1\). Whether that requirement suits the SOFR use, where only a few years of meetings are listed, is a question of scope. Nonzero fixed-maturity curve jumps and restart laws are not treated.

\subsection{Quasi-exponential shapes with recurrent loadings are realized}\label{sec:3.28}

Keep the realization notion and the notation \(i(t,T)\), \(D(t)\) of Section \hyperref[sec:3.27]{3.27}, and write \(i(x,D)\) for the number of distances in \(D\) that are at most \(x\), so that \(i(t,t+x)=i(x,D(t))\). Assume

\begin{itemize}
  \item (H1) \(\lambda(x)=c\,e^{Ax}b\) for an \(r\times r\) matrix \(A\), a row vector \(c\) and a column vector \(b\), that is, \(\lambda\) is quasi-exponential (\cite{bjork2001existence} Remark 5.1);
  \item (H2) \(a_i=u\,M^iv\) for all \(i\ge0\), for a \(p\times p\) matrix \(M\), a row vector \(u\) and a column vector \(v\).
\end{itemize}

Condition (H2) is equivalent to a recurrence (3.141), through the companion matrix of its coefficients. The matrices \(A\) and \(M\) of this section are unrelated to the coefficients \(A_{j,k}\), \(M_{j,k}\) of Section \hyperref[sec:3.25]{3.25}. Write \(m=pr\), \(\otimes\) for the Kronecker product, \(\hat A=I_p\otimes A\), \(\hat M=M\otimes I_r\), \(\hat b=v\otimes b\) and \(\hat c=u^\top\otimes c^\top\) in \(\mathbb R^m\), and for \(x\ge0\) and a distance vector \(D\) the row vectors

\begin{equation*}
k(x,D)=\bigl(u\,M^{i(x,D)}\bigr)\otimes\bigl(c\,e^{Ax}\bigr),\qquad K(x,D)=\int_0^xk(y,D)\,dy. \tag{3.142}\label{eq:3.142}
\end{equation*}

The basis is that of Section \hyperref[sec:3.26]{3.26}, with the curve driven by \(W=B^1\) on a finite horizon \(H\) and the cited calculus of Sections \hyperref[sec:3.25]{3.25} and \hyperref[sec:3.26]{3.26} together with Itô's product rule (\cite{bauerschmidt2020stochastic} §3.5). The scale state \(Z\in\mathbb R^{p'}\) solves an autonomous equation \(dZ=\mu_Z(Z)\,dt+\Sigma_Z(Z)\,dB\) from a deterministic \(z_0\), with \(\mu_Z\), \(\Sigma_Z\) globally Lipschitz and bounded on bounded sets, and is taken with every path continuous, as Proposition \hyperref[res:Proposition-29]{29}(a) provides. It may use \(B^1\), so any correlation between the scale and the curve driver is allowed. Let \(\psi:\mathbb R^{p'}\to\mathbb R\) be bounded by \(K_\psi\) and Lipschitz, the initial curve zero, and

\begin{equation*}
\sigma(t,T)=\psi(Z_t)\,a_{i(t,T)}\,\lambda(T-t),\qquad 0\le t\le T\le H. \tag{3.143}\label{eq:3.143}
\end{equation*}

\phantomsection\label{res:Proposition-31}
\textbf{Proposition 31 (conditional).} Under (H1), (H2) and these inputs:

(a) The drift \(\alpha(t,T)=\sigma(t,T)\int_t^T\sigma(t,u)\,du\) satisfies Assumption 2.1 on every path, and the curve (2.2) with zero initial curve and no curve jumps is well defined for \(0\le t\le T\le H\).

(b) Let \(Y=(Z,\Xi,\Pi,\Theta)\in\mathbb R^{p'}\times\mathbb R^m\times\mathbb R^{m\times m}\times\mathbb R^m\), with \((\Xi,\Pi,\Theta)_0=0\), solve between meetings

\begin{equation*}
\begin{gathered}
d\Xi=\hat A\,\Xi\,dt+\psi(Z)\,\hat b\,dW,\qquad d\Pi=\bigl(\hat A\Pi+\Pi\hat A^\top+\psi(Z)^2\hat b\hat b^\top\bigr)dt,\qquad d\Theta=\bigl(\hat A\Theta+\Pi\hat c\bigr)dt,
\end{gathered} \tag{3.144}\label{eq:3.144}
\end{equation*}

together with the equation of \(Z\), and at each meeting apply the restart map

\begin{equation*}
\mathcal R(Z,\Xi,\Pi,\Theta)=\bigl(Z,\hat M\Xi,\hat M\Pi\hat M^\top,\hat M\Theta\bigr). \tag{3.145}\label{eq:3.145}
\end{equation*}

Then, for every \(t\) and \(x\) with \(t+x\le H\), almost surely, with a null set that may depend on \(x\) and one random state for all \(x\),

\begin{equation*}
f(t,t+x)=G(Y_t,D(t))(x),\qquad G\bigl((Z,\Xi,\Pi,\Theta),D\bigr)(x)=k(x,D)\bigl[\Xi+\Pi K(x,D)^\top+\Theta\bigr], \tag{3.146}\label{eq:3.146}
\end{equation*}

with \(\Xi\) on its version continuous between meetings. The coefficients of (3.144) are globally Lipschitz in \(Y\), the same on every interval and free of \(t\), \(D\) and the schedule; \(\mathcal R\) is a fixed linear map; \(G\) is linear in the state, does not depend on \(Z\), and depends on the schedule only through \(D\). Since the horizon \(H\) is arbitrary, this gives (3.139) for every \(t\) and \(x\), so this is a realization in the sense of Section \hyperref[sec:3.27]{3.27}, serving every schedule, with state dimension \(p'+m+m(m+1)/2+m\), counting \(\Pi\) as symmetric. The constant scale \(\psi\equiv1\) is the case in which \(Z\) may be dropped.

(c) Consequently, for one factor and every scale of this class, a quasi-exponential shape with linearly recurrent ordinal loadings has a realization serving every schedule, with no condition on the correlation between the scale and the curve driver.

Without meetings, \cite{bjork2004finite} Proposition 5.2 gives sufficiency of quasi-exponential shapes for \(\sigma=\sum_j\varphi_{ij}(r,y)\lambda_j(x)\), and its Proposition 4.16 inserts a scale independent of the driving noise when the curve map does not depend on it, under smoothness and genericity assumptions; the constant-scale case without meetings is \cite{bjork2001existence} Proposition 5.1. Proposition \hyperref[res:Proposition-31]{31} adds the meeting breakpoints and restarts, with the recurrent loadings as the new ingredient; its argument is pathwise in the drift and uses only predictability of the scale in the diffusion, so no independence is needed.

\emph{Proof.} For \(s\le t\) let \(n(s,t)\) be the number of meetings in \((s,t]\), so that \(i(s,T)=n(s,t)+i(t,T)\), and put \(w(s,t)=(M^{n(s,t)}v)\otimes(e^{A(t-s)}b)\in\mathbb R^m\). By (H1), (H2) and the mixed-product rule, for \(s\le t\le T=t+x\),

\begin{equation*}
\sigma(s,T)=\psi(Z_s)\,k(x,D(t))\,w(s,t). \tag{3.147}\label{eq:3.147}
\end{equation*}

(a) The scale \(\psi(Z_s)\) is adapted with continuous paths, hence predictable, and bounded; the coordinates of \(w(\cdot,t)\) are bounded deterministic Borel functions, so every integrand is admissible. For fixed \(s\), \(\sigma(s,\cdot)=\psi(Z_s)\tilde\sigma(s,\cdot)\) with \(\tilde\sigma\) bounded and piecewise continuous, so \(S(s,T)=\int_s^T\sigma(s,u)\,du\) is absolutely continuous and \(\int_s^T\alpha(s,u)\,du=\tfrac12S(s,T)^2\), which is (2.3) for one driver.

(b), diffusion. By (3.147) and linearity, \(\int_0^t\sigma(s,t+x)\,dW_s=k(x,D(t))\,\Xi_t\) with \(\Xi_t=\int_0^t\psi(Z_s)w(s,t)\,dW_s\), one random vector for all \(x\). With \(\tau\le t\) the last meeting, \(w(s,t)=e^{\hat A(t-\tau)}w(s,\tau)\) for \(s\le\tau\) and \(w(s,t)=e^{\hat A(t-s)}\hat b\) for \(\tau<s\le t\), so \(\Xi_t=e^{\hat A(t-\tau)}\Xi_\tau+e^{\hat At}\int_\tau^te^{-\hat As}\psi(Z_s)\hat b\,dW_s\). The right side is continuous in \(t\), and the product rule gives the first equation of (3.144) between meetings. At a meeting \(T_j\), \(w(s,T_j)=\hat Mw(s,T_j-)\) for \(s<T_j\); the two integrands differ only at \(s=T_j\), so their integrals agree almost surely, which is the restart (3.145) for \(\Xi\).

(b), drift. Fix a path. Split \(S(s,T)=\psi(Z_s)[\tilde J(s,t)+K(x,D(t))w(s,t)]\) with \(\tilde J(s,t)=\int_s^t\tilde\sigma(s,u)\,du\). Then

\begin{equation*}
\int_0^t\alpha(s,T)\,ds=k(x,D(t))\bigl[\Pi_tK(x,D(t))^\top+\Theta_t\bigr],\qquad \Pi_t=\int_0^t\psi(Z_s)^2w\,w^\top ds,\qquad \Theta_t=\int_0^t\psi(Z_s)^2w\,\tilde J\,ds, \tag{3.148}\label{eq:3.148}
\end{equation*}

with \(w=w(s,t)\) and \(\tilde J=\tilde J(s,t)\). Between meetings \(\partial_tw=\hat Aw\), \(w(t,t)=\hat b\), \(\tilde J(t,t)=0\) and \(\partial_t\tilde J(s,t)=\hat c^\top w(s,t)\), because \(k(0,D)=\hat c^\top\); since \(s\mapsto\psi(Z_s)^2\) is continuous, Leibniz's rule gives the equations of \(\Pi\) and \(\Theta\) in (3.144). At a meeting \(\tilde J(s,\cdot)\) is continuous and \(w\) jumps by \(\hat M\), which gives (3.145). With the zero initial state this proves (3.146) at each \(t\); (3.146) defines the version of the curve, which for each fixed \(T\) agrees almost surely with (2.2).

(b), conditions. The drift of \(\Xi\) is linear and its diffusion \(\psi(Z)\hat b\) is bounded and Lipschitz in \(Z\); the source of \(\Pi\) is Lipschitz with constant \(2K_\psi L_\psi|\hat b|^2\), with \(L_\psi\) the Lipschitz constant of \(\psi\); the drift of \(\Theta\) is linear. None involves \(t\), \(D\) or the schedule. \(\Pi\) is symmetric, which gives the dimension. (c) restates (b).\hfill$\square$

The proposition is conditional on the cited calculus and on the scale state of Section \hyperref[sec:3.26]{3.26}; it constructs neither. It gives sufficiency only: necessity of the recurrence and of quasi-exponential shapes under a stochastic scale is not shown, and Propositions \hyperref[res:Proposition-30]{30} and \hyperref[res:Proposition-33]{33} cover only a constant scale. Not covered: several factors, several scales per factor, and scales that move the shape, which admit no realization even without meetings (\cite{bjork2004finite} Propositions 4.11 and 4.12); scales outside the stated class, such as unbounded or non-Lipschitz \(\psi\) and square-root or CIR scale states, so the Heston-type scale of \cite{brace2024sofr} eqs. 29 and 30 is excluded; nonzero initial curves, nonzero fixed-maturity curve jumps, and minimality of the dimension; true bond martingales.

\subsection{Recurrent loadings are necessary for every quasi-exponential shape}\label{sec:3.29}

Keep the realization notion and the notation of Section \hyperref[sec:3.27]{3.27}, with one factor, the constant scale \(\psi\equiv1\) and no fixed-maturity curve jumps, and let

\begin{equation*}
\sigma(t,T)=a_{i(t,T)}\,\lambda(T-t),\qquad \lambda(x)=c\,e^{Ax}b, \tag{3.149}\label{eq:3.149}
\end{equation*}

with \(\lambda\) quasi-exponential and not identically zero. The basis is that of Proposition \hyperref[res:Proposition-30]{30}, and the initial curve is any deterministic curve.

\phantomsection\label{res:Lemma-32}
\textbf{Lemma 32.} Let \(F:[0,\tau]\to\mathbb R^N\) be bounded, with finitely many pieces on each of which it is continuous. Then \(\int_0^\tau F(s)\,dW_s\), taken coordinatewise, is centered Gaussian with covariance \(\int_0^\tau F(s)F(s)^\top ds\).

\emph{Proof.} Left-endpoint step functions \(F_k\) on refining partitions converge to \(F\) in \(L^2([0,\tau])\). The integral of a step function is a linear combination of Brownian increments by the elementary values and linearity of the integral, hence centered Gaussian with covariance \(\int F_kF_k^\top ds\). The isometry gives \(E|\int(F-F_k)\,dW|^2=\int|F-F_k|^2ds\to0\); along an almost surely convergent subsequence the characteristic functions converge by bounded convergence, and their limit \(\exp(-u^\top Cu/2)\), \(C=\int FF^\top ds\), identifies the law.\hfill$\square$

Fix \(n\ge1\), \(R\ge0\), the schedule \(T_m=m\) for \(m=1,\dots,n+R\), \(t=n+\tfrac12\) and \(T^{(i)}=t+\tfrac14+i\) for \(i=0,\dots,R\). After the reduction to a controllable pair \((A,b)\) in the proof, put \(h_k=a_k\,c\,e^{kA}\in\mathbb R^{1\times r}\) and let \(\mathcal H_{R,n}\) be the \((R+1)\times nr\) block Hankel matrix with \(1\times r\) blocks

\begin{equation*}
(\mathcal H_{R,n})_{i,l}=h_{i+l},\qquad 0\le i\le R,\ 1\le l\le n. \tag{3.150}\label{eq:3.150}
\end{equation*}

\phantomsection\label{res:Proposition-33}
\textbf{Proposition 33 (conditional).} Under these assumptions:

(a) If a realization with state dimension \(q\) reproduces the curve at \(t\) on this schedule, then \(\operatorname{rank}\mathcal H_{R,n}\le q\).

(b) If a realization with state dimension \(q\) serves every schedule, then the ordinal loadings satisfy a linear recurrence (3.141), of order at most \(q+1\).

(c) Consequently, for one factor, the constant scale and a nonzero quasi-exponential \(\lambda\), a realization serving every schedule exists, from the zero initial curve, if and only if the ordinal loadings satisfy a linear recurrence. Necessity holds for every deterministic initial curve; sufficiency is Proposition \hyperref[res:Proposition-31]{31} with \(\psi\equiv1\). In particular, for \(\lambda\equiv1\), Propositions \hyperref[res:Proposition-30]{30} and \hyperref[res:Proposition-31]{31} together characterize the ordinal loadings that admit a realization as the linearly recurrent ones. This settles the one-factor constant-scale case of Conjecture 2.

Proposition \hyperref[res:Proposition-30]{30} is the case \(\lambda\equiv1\), where the pieces carry Brownian increments. Without meetings the answer is the quasi-exponential condition alone (\cite{bjork2001existence} Proposition 5.1). The reduction to a controllable pair is the standard restriction to the reachable subspace.

\emph{Proof.} \emph{Reduction.} The span of \(b,Ab,A^2b,\dots\) is \(A\)-invariant and contains \(b\), and \(\lambda\) is unchanged when \(A\) and \(c\) are restricted to it; so \((A,b)\) may be taken controllable, and then \(c\neq0\). If \(y^\top e^{Au}b=0\) on an interval, differentiation gives \(y^\top A^kb=0\) for all \(k\), so \(y=0\); hence

\begin{equation*}
\Gamma_\delta=\int_0^\delta e^{Au}bb^\top e^{A^\top u}\,du \tag{3.151}\label{eq:3.151}
\end{equation*}

is positive definite for every \(\delta>0\).

(a) As in Proposition \hyperref[res:Proposition-30]{30}, split \([0,t]\) into piece 0, \([n,t]\), and pieces \(l\ge1\), \([n-l,n-l+1]\), with right ends \(\tau_0=t\) and \(\tau_l=n-l+1\); on piece \(l\), \(i(s,T^{(i)})=i+l\), and changing the integrand at the finitely many piece ends does not change the integral. Put \(\zeta_l=\int_{\text{piece }l}e^{A(\tau_l-s)}b\,dW_s\in\mathbb R^r\). The constant matrix \(e^{A(T^{(i)}-\tau_l)}\) comes out of the integral, and \(T^{(i)}-\tau_l=i+l-\tfrac14\) for \(l\ge1\), so

\begin{equation*}
f(t,T^{(i)})=\bar f_i+a_i\,c\,e^{A/4}e^{iA}\zeta_0+\sum_{l=1}^nh_{i+l}\,e^{-A/4}\zeta_l \tag{3.152}\label{eq:3.152}
\end{equation*}

with \(\bar f_i\) deterministic. By Lemma \hyperref[res:Lemma-32]{32}, \((\zeta_0,\dots,\zeta_n)\) is centered Gaussian with covariance \(\operatorname{diag}(\Gamma_{1/2},\Gamma_1,\dots,\Gamma_1)\), positive definite by (3.151), so its law is absolutely continuous. The coefficient matrix of \((\zeta_1,\dots,\zeta_n)\) in (3.152) is \(\mathcal H_{R,n}(I_n\otimes e^{-A/4})\), of rank \(\operatorname{rank}\mathcal H_{R,n}\), and adding the columns of \(\zeta_0\) cannot lower the rank \(\rho\) of the whole linear map. So the vector of curve values has an absolutely continuous law on a \(\rho\)-dimensional affine subspace, and the covering argument of Proposition \hyperref[res:Proposition-30]{30}(a) gives \(\operatorname{rank}\mathcal H_{R,n}\le\rho\le q\).

(b) By (a), \(\operatorname{rank}\mathcal H_{R,n}\le q\) for all \(R\) and \(n\). For each \(n\) the first \(n\) blocks of the \(q+1\) rows \((h_{i+1},h_{i+2},\dots)\), \(i=0,\dots,q\), are rows of \(\mathcal H_{q,n}\) and hence dependent; the annihilating unit vectors form nonempty compact sets decreasing in \(n\), so a unit \(d\in\mathbb R^{q+1}\) satisfies \(\sum_id_ih_{i+l}=0\) for every \(l\ge1\). Right multiplication by \(e^{-lA}\) gives \(\sum_id_ia_{i+l}\,c\,e^{iA}=0\). For \(i_0\) with \(d_{i_0}\neq0\), \(c\,e^{i_0A}\neq0\), so some coordinate gives scalars \(e_i=d_i(c\,e^{iA})_s\), not all zero, with \(\sum_ie_ia_{i+l}=0\) for all \(l\ge1\). This is a recurrence of order at most \(q\) for \((a_{k+1})_{k\ge0}\), or \(a_l=0\) for all \(l\ge1\); prepending \(a_0\) with a zero coefficient gives (3.141) of order at most \(q+1\).

(c) Sufficiency is Proposition \hyperref[res:Proposition-31]{31} with \(\psi\equiv1\), necessity is (b).\hfill$\square$

The proposition is conditional on the cited calculus. It treats one factor and the constant scale; necessity under a stochastic scale and with several factors remains open: with a stochastic scale independent of \(W\) the law is a Gaussian mixture, and with leverage it is not Gaussian, and neither case is treated. Several factors, necessity of quasi-exponential shapes when meetings are present, realizations without the Lipschitz regularity of the curve map, and local realizations are not covered. The order bound \(q+1\) is not claimed to be sharp.

\subsection{The splice with a correlated step factor}\label{sec:3.30}

Fix a finite horizon \(H>T_N\), with \(I_N\) used up to \(H\). Let \(B^1,B^2\) be independent standard Brownian motions with the cited Itô calculus of Section \hyperref[sec:3.25]{3.25}, and for a fixed \(\rho\in[-1,1]\) put \(W^S=B^1\) and \(W^B=\rho B^1+\sqrt{1-\rho^2}\,B^2\). The step factor has the deterministic volatility \(\sigma^S(t,T)=s_m(t)\) for \(T\in I_m\), with bounded Borel \(s_m\), and the block factor \(\sigma^B(t,T)=b\,e^{-a(T-t)}\) with \(a\neq0\) and \(b\neq0\). The curve is (2.2) with these two factors, no curve jumps, and the drift of Assumption 2.1 written with the independent drivers:

\begin{equation*}
\alpha=\sigma^SS^S+\sigma^BS^B+\rho\bigl(\sigma^SS^B+\sigma^BS^S\bigr),\qquad S^X(u,T)=\int_u^T\sigma^X(u,v)\,dv. \tag{3.153}\label{eq:3.153}
\end{equation*}

Write \(\Delta_m=s_m-s_{m-1}\), \(1\le m\le N\), for the jump of the step volatility at the maturity \(T_m\), and call the step factor genuine if some \(\Delta_m\) is nonzero on a subset of \([0,T_m)\) of positive measure. A block is a finite-dimensional space \(\mathcal E\) of real-analytic functions on \(\mathbb R\), invariant under the shifts \(g\mapsto g(\cdot+h)\), \(h\ge0\), and containing \(e^{-ax}\) and \(e^{-2ax}\); a forward shift is injective on \(\mathcal E\) by analyticity, so \(\mathcal E\) is invariant under every shift. The family at time \(t\) is \(\mathcal M_t(\mathcal E)=\{p+g(\cdot-t):p\in S^+,\ g\in\mathcal E\}\) restricted to \([t,H]\), with \(S^+\) the functions affine on each \(I_m\cap[0,H]\); it is the restriction of one space \(\mathcal M(\mathcal E)\) of functions of calendar maturity. The curve is consistent with \(\mathcal M(\mathcal E)\) if \(f(0,\cdot)\in\mathcal M(\mathcal E)\) and for every \(t\in[0,H]\) some random element \(g_t\) of \(\mathcal M_t(\mathcal E)\) satisfies \(g_t(T)=f(t,T)\) almost surely for every \(T\in[t,H]\). This is the invariance of \cite{bjork1999interest} Definition 2.2 from the initial time, read in the version form of (3.139).

\phantomsection\label{res:Proposition-34}
\textbf{Proposition 34 (conditional).} Under the cited calculus:

(a) For \(t\le T\) with \(T\in I_m\), writing \(S^S(u,T)=s_m(u)T+\kappa_m(u)\) on \(I_m\),

\begin{equation*}
\int_0^t\alpha(u,T)\,du=p_t(T)+g^B_t(T)+e^{-aT}\bigl(\eta_m(t)+\theta_m(t)T\bigr), \tag{3.154}\label{eq:3.154}
\end{equation*}

with \(p_t\in S^+\), \(g^B_t\in\operatorname{span}\{e^{-aT},e^{-2aT}\}\), \(\theta_m(t)=\rho b\int_0^te^{au}s_m(u)\,du\) and \(\eta_m(t)=\rho b\int_0^te^{au}(\kappa_m(u)-s_m(u)/a)\,du\). For \(t<T_m\),

\begin{equation*}
\theta_m(t)-\theta_{m-1}(t)=\rho b\int_0^te^{au}\Delta_m(u)\,du,\qquad \eta_m(t)-\eta_{m-1}(t)=-\rho b\int_0^te^{au}\bigl(T_m+\tfrac1a\bigr)\Delta_m(u)\,du. \tag{3.155}\label{eq:3.155}
\end{equation*}

For every \(t\) and \(T\), almost surely, the random part of \(f(t,T)\) equals \(\sum_m\mathbf 1_{I_m}(T)\int_0^ts_m\,dW^S+b\,e^{-aT}\int_0^te^{au}\,dW^B\), one version for all \(T\), which lies in \(S^++\operatorname{span}\{e^{-aT}\}\).

(b) If the curve is consistent with \(\mathcal M(\mathcal E)\) for some block \(\mathcal E\), then \(\rho\Delta_m=0\) almost everywhere on \([0,T_m)\) for every \(m\). A genuine step factor therefore splices onto a fixed block only if \(\rho=0\).

(c) Let \(\mathcal E_0=\operatorname{span}\{e^{-ax},xe^{-ax},e^{-2ax}\}\), a block, and \(f(0,\cdot)\in\mathcal M(\mathcal E_0)\). (i) If \(\rho\Delta_m=0\) almost everywhere on \([0,T_m)\) for every \(m\), the curve is consistent with \(\mathcal M(\mathcal E_0)\). (ii) For every \(\rho\), the curve is consistent with the enlarged family \(\mathcal M(\mathcal E_0)+\operatorname{span}\{\mathbf 1_{I_m}(T)e^{-aT},\,\mathbf 1_{I_m}(T)\,Te^{-aT}:0\le m\le N\}\).

(d) If the step factor is genuine, \(\rho\neq0\) and \(\rho\Delta_m\) is not almost everywhere zero on \([0,T_m)\), then for every block \(\mathcal E\) there is a nonempty open set of times \(t<T_m\) at which \(f(t,\cdot)\) has no version in \(\mathcal M_t(\mathcal E)\), and almost surely no element of \(\mathcal M_t(\mathcal E)\) agrees with \(f(t,\cdot)\) at every rational \(T\). At such \(t\) the coefficient \(\theta\) of \(Te^{-aT}\) in (3.154) jumps across \(T_m\).

Case (c)(i) with \(\rho=0\) is the independence splice of \cite{gellert2021short} §3, and the indicator-times-exponential terms of (c)(ii) are those of \cite{brace2024sofr} eq. 36. Proposition \hyperref[res:Proposition-34]{34} shows that, with deterministic volatilities, one step factor and one exponential block factor, these are the only two ways: a step factor with maturity steps and nonzero correlation forces the indicator terms into any fixed block. A correlation that vanishes wherever the step factor has maturity steps is the independence splice locally in time.

\emph{Proof.} (a) With \(S^B(u,T)=b(1-e^{-a(T-u)})/a\), the cross part of (3.153) is \(\rho b[\sigma^S(u,T)/a+e^{-a(T-u)}(S^S(u,T)-\sigma^S(u,T)/a)]\). On \(I_m\), \(\sigma^S(u,T)=s_m(u)\) and \(S^S(u,\cdot)\) is continuous and piecewise affine, so \(S^S(u,T)=s_m(u)T+\kappa_m(u)\) with \(\kappa_m\) bounded. Integrating over \(u\in[0,t]\), the first term is constant in \(T\) on \(I_m\), hence in \(S^+\), and the second gives \(e^{-aT}(\eta_m(t)+\theta_m(t)T)\). The term \(\sigma^SS^S\) is affine in \(T\) on \(I_m\), and \(\sigma^BS^B=(b^2/a)(e^{-a(T-u)}-e^{-2a(T-u)})\). This is (3.154). Continuity of \(S^S(u,\cdot)\) at \(T_m\) gives \(\kappa_m-\kappa_{m-1}=-\Delta_mT_m\), hence (3.155). The integrands of the random part are bounded deterministic Borel functions, and the integral splits by linearity; the result involves finitely many random variables, so it is one version for all \(T\).

(b) Fix \(t\) and a version \(g_t\in\mathcal M_t(\mathcal E)\). On one event of probability one, \(g_t\) agrees at every rational \(T\) with the version from (a) plus the deterministic parts; both are affine plus real-analytic on each \(I_j\cap[t,H]\), hence they agree there. The version of the random part, \(p_t\), \(g^B_t\) and \(f(0,\cdot)\) lie in \(\mathcal M_t(\mathcal E)\), so \(X_t(T)=e^{-aT}(\eta_m(t)+\theta_m(t)T)\) on \(I_m\cap[t,H]\) lies in \(\mathcal M_t(\mathcal E)\). For \(t<T_m\), write \(X_t=p+g(\cdot-t)\); the real-analytic \(e(T)=g(T-t)\) differs from \(e^{-aT}(\eta_m+\theta_mT)\) by an affine function on \(I_m\), hence on \(\mathbb R\), and similarly on \(I_{m-1}\). Subtracting, \(e^{-aT}((\eta_m-\eta_{m-1})+(\theta_m-\theta_{m-1})T)\) is affine on \(I_{m-1}\), and since \(1,T,e^{-aT},Te^{-aT}\) are linearly independent on every interval, both differences vanish. By (3.155) and Lebesgue's differentiation theorem, \(\rho b\,e^{au}\Delta_m(u)=0\) for almost every \(u<T_m\).

(c) Under (i) the differences (3.155) vanish for \(t<T_m\), and for \(t\ge T_m\) the interval \(I_{m-1}\) does not meet \((t,H]\); so on \([t,H]\), \(X_t\) is \(e^{-aT}(\eta+\theta T)\) with common coefficients, which lies in \(\operatorname{span}\{e^{-a(T-t)},(T-t)e^{-a(T-t)}\}\subseteq\mathcal E_0(\cdot-t)\). With the version from (a) this is a version in \(\mathcal M_t(\mathcal E_0)\). Under (ii) the terms \(\mathbf 1_{I_m}(T)e^{-aT}(\eta_m(t)+\theta_m(t)T)\) lie in the added span.

(d) The function \(\varphi(t)=\int_0^te^{au}\Delta_m(u)\,du\) is continuous and, by Lebesgue's differentiation theorem, not identically zero on \([0,T_m)\), so \(\{t<T_m:\varphi(t)\neq0\}\) is open and nonempty. There \(\theta_m\neq\theta_{m-1}\), and the argument of (b), applied on any event of positive outer probability on which an element of \(\mathcal M_t(\mathcal E)\) agrees with \(f(t,\cdot)\) at every rational \(T\), gives a contradiction.\hfill$\square$

The proposition is conditional on the cited calculus. It is a classification only in its setting: deterministic volatilities, one step factor, one block factor with a single exponent, a correlation constant in time, and families of the form \(S^+\) plus a block, possibly enlarged by indicator terms. Not covered: stochastic volatilities, under which the ramp slope and the cross-term coefficients are adapted (Proposition \hyperref[res:Proposition-28]{28}(c); Proposition \hyperref[res:Proposition-41]{41} treats predictable bounded scales); blocks with several exponents or polynomial factors, several block or step factors, and time-varying correlation; blocks with constant or linear terms, where the front end could absorb part of the block's equation; fixed-maturity curve jumps; consistency from every initial time; and the exponent restrictions of \cite{sharef2004conditions} inside the block.

\subsection{The splice with a quasi-exponential block factor}\label{sec:3.31}

Keep the setting and notation of Section \hyperref[sec:3.30]{3.30}, and replace the block factor by

\begin{equation*}
\sigma^B(t,T)=\lambda(T-t),\qquad \lambda(x)=c\,e^{Ax}b, \tag{3.156}\label{eq:3.156}
\end{equation*}

with an invertible \(r\times r\) matrix \(A\), so that no exponent is zero, \((A,b)\) controllable and \(c\neq0\). Put \(\Lambda(y)=\int_0^y\lambda=c\,A^{-1}(e^{Ay}-I)\,b\). Section \hyperref[sec:3.30]{3.30} is the case \(r=1\), \(A=-a\). A block is as there, now required to contain \(\lambda\) and \(\lambda\Lambda\), and \(\mathcal E_1\) is the smallest block containing \(\lambda\), \(x\lambda(x)\) and \(\lambda\Lambda\); it is finite-dimensional because each of these functions has finitely many independent translates. Consistency is that of Section \hyperref[sec:3.30]{3.30}. For \(1\le m\le N\) put

\begin{equation*}
v_m(t)=\int_0^t\Delta_m(u)\,e^{-Au}b\,du\in\mathbb R^r. \tag{3.157}\label{eq:3.157}
\end{equation*}

\phantomsection\label{res:Proposition-35}
\textbf{Proposition 35 (conditional).} Under the cited calculus:

(a) The time integral of the drift is as in (3.154), with \(e^{-aT}(\eta_m+\theta_mT)\) replaced by a function \(X_{t,m}\) on \(I_m\) that extends to an entire function of \(T\). For \(t<T_m\) the jump of the extensions across \(T_m\) is

\begin{equation*}
X_{t,m}(T)-X_{t,m-1}(T)=\rho\,c\,e^{AT}\bigl(A^{-1}+(T-T_m)I\bigr)v_m(t)-\rho\,c\,A^{-1}b\int_0^t\Delta_m(u)\,du. \tag{3.158}\label{eq:3.158}
\end{equation*}

The random part has, for every \(t\) and \(T\) almost surely, the version \(\sum_m\mathbf 1_{I_m}(T)\int_0^ts_m\,dW^S+c\,e^{AT}\int_0^te^{-Au}b\,dW^B_u\), one for all \(T\).

(b) If the curve is consistent with \(\mathcal M(\mathcal E)\) for some block \(\mathcal E\), then \(\rho\Delta_m=0\) almost everywhere on \([0,T_m)\) for every \(m\).

(c) Let \(f(0,\cdot)\in\mathcal M(\mathcal E_1)\). (i) If \(\rho\Delta_m=0\) almost everywhere on \([0,T_m)\) for every \(m\), the curve is consistent with \(\mathcal M(\mathcal E_1)\). (ii) For every \(\rho\), the curve is consistent with \(\mathcal M(\mathcal E_1)+\operatorname{span}\{\mathbf 1_{I_m}(T)g(T):g\in\operatorname{span}\{c\,e^{AT}y,\,Tc\,e^{AT}y:y\in\mathbb R^r\},\ 0\le m\le N\}\).

(d) If the step factor is genuine and \(\rho\neq0\), the curve is consistent with \(\mathcal M(\mathcal E)\) for no block \(\mathcal E\).

Controllability and \(c\neq0\) are no loss of generality: restricting \(A\) and \(c\) to the span of \(b,Ab,A^2b,\dots\) leaves \(\lambda\) unchanged and \(A\) invertible, and then \(\lambda\not\equiv0\) means \(c\neq0\). Proposition \hyperref[res:Proposition-35]{35} extends Proposition \hyperref[res:Proposition-34]{34} from one exponent to every quasi-exponential block factor without an exponent zero: a correlated step factor with maturity steps still forces indicator-times-quasi-exponential terms into the family, and the independence splice of \cite{gellert2021short} §3 remains the only alternative. Quasi-exponential functions are those of \cite{bjork2001existence} Remark 5.1.

\emph{Proof.} (a) For \(T\in I_m\), \(S^S(u,T)=s_m(u)T+\kappa_m(u)\) and \(S^B(u,T)=\Lambda(T-u)\), so the cross part of the drift is \(\rho[s_m(u)\Lambda(T-u)+\lambda(T-u)(s_m(u)T+\kappa_m(u))]\), entire in \(T\) and bounded on compacts uniformly in \(u\); its integral \(X_{t,m}\) is entire. Subtracting the expressions for \(m\) and \(m-1\), with \(\kappa_m-\kappa_{m-1}=-\Delta_mT_m\), gives the integrand \(\rho\Delta_m(u)[\Lambda(T-u)+\lambda(T-u)(T-T_m)]\), and \(e^{A(T-u)}=e^{AT}e^{-Au}\) gives (3.158). The term \(\sigma^SS^S\) lies in \(S^+\). The term \(\sigma^BS^B=\lambda\Lambda(T-u)\) integrates into \(\mathcal E(\cdot-t)=\mathcal E\): a forward shift is bijective on a finite-dimensional space of analytic functions, and the coordinates of \(g(\cdot+h)\) in a basis are continuous in \(h\). The random part is as in Proposition \hyperref[res:Proposition-34]{34}(a).

(b) As in Proposition \hyperref[res:Proposition-34]{34}(b), consistency puts the deterministic cross term in \(\mathcal M_t(\mathcal E)\) for every \(t\), and analytic continuation from \(I_{m-1}\) and \(I_m\) makes (3.158) affine in \(T\) for \(t<T_m\). Its constant part is affine, and \(c\,e^{AT}(A^{-1}+(T-T_m)I)v_m(t)\) is an exponential polynomial whose exponents are eigenvalues of \(A\), all nonzero; since \(x^ke^{\mu x}\) for distinct nonzero \(\mu\), together with \(1\) and \(x\), are linearly independent, it vanishes. So \(v_m(t)\in\mathcal K=\{v:c\,e^{AT}(TI+N_m)v=0\text{ for all }T\}\) with \(N_m=A^{-1}-T_mI\), for every \(t<T_m\). The function \(v_m\) is absolutely continuous with derivative \(\Delta_m(t)e^{-At}b\) almost everywhere, and a \(\mathcal K\)-valued absolutely continuous function has its derivative in \(\mathcal K\) almost everywhere; so \(e^{-At}b\in\mathcal K\) for almost every \(t\) with \(\Delta_m(t)\neq0\). If that set had positive measure, the projection of \(e^{-At}b\) onto \(\mathcal K^\perp\), real-analytic in \(t\), would vanish on a set with an accumulation point, hence identically; differentiating at \(0\) puts every \(A^kb\) in \(\mathcal K\), so \(\mathcal K=\mathbb R^r\) by controllability. Then \(c\,e^{AT}(TI+N_m)=0\) for all \(T\); at \(T=0\), \(cA^{-1}=T_mc\), and the derivative at \(0\) gives \(c(AN_m+I)=0\), that is \(2c=T_mcA\). Multiplying the first identity by \(A\) gives \(c=T_mcA=2c\), so \(c=0\), a contradiction. Hence \(\rho\Delta_m=0\) almost everywhere on \([0,T_m)\).

(c) Under (i) the jump (3.158) vanishes for \(t<T_m\), and for \(t\ge T_m\) the interval \(I_{m-1}\) does not meet \((t,H]\); so on \([t,H]\) the cross term is one function, the integral of \(\rho[s(u)\Lambda(T-u)+\lambda(T-u)(s(u)T+\kappa(u))]\) with coefficients common to the intervals that meet \([t,H]\). Here \(\Lambda(T-u)\) is a constant plus an element of the translates of \(\lambda\), because \(A^{-1}\) is a polynomial in \(A\) by the Cayley--Hamilton theorem, and \(\lambda(T-u)T=\lambda(T-u)(T-t)+t\lambda(T-u)\). Both lie in \(S^++\mathcal E_1(\cdot-t)\), because a block containing \(\lambda\) contains \(c\,e^{Ax}y\) for every \(y\), by differentiation and controllability, and \(\mathcal E_1\) contains \(x\lambda(x)\). With the version from (a), \(f(t,\cdot)\) has a version in \(\mathcal M_t(\mathcal E_1)\). Under (ii), on each \(I_m\) the cross term is a constant plus \(c\,e^{AT}y_m(t)+Tc\,e^{AT}z_m(t)\), which lies in the added span.

(d) is (b) for a genuine step factor with \(\rho\neq0\).\hfill$\square$

The proposition is conditional on the cited calculus, with the scope of Proposition \hyperref[res:Proposition-34]{34} otherwise. Unlike Proposition \hyperref[res:Proposition-34]{34}(d), it does not locate an open set of times at which no version exists. Not covered: block factors with an exponent zero, whose polynomial part lets the front end absorb affine parts of the jump; several block or step factors; stochastic volatilities, treated for predictable bounded scales in Proposition \hyperref[res:Proposition-41]{41}; curve jumps; and consistency from every initial time.

\subsection{The block's restrictions survive the uncorrelated splice}\label{sec:3.32}

Fix \(\beta>0\), integers \(n_1,n_2\ge0\), and the biexponential-polynomial family of \cite{sharef2004conditions} §2.2,

\begin{equation*}
F(x,z)=\Bigl(\sum_{\mu=0}^{n_1}z_{1,\mu}x^\mu\Bigr)e^{-\beta x}+\Bigl(\sum_{\mu=0}^{n_2}z_{2,\mu}x^\mu\Bigr)e^{-2\beta x}, \tag{3.159}\label{eq:3.159}
\end{equation*}

with parameters \(z\in\mathbb R^{n_1+n_2+2}\) indexed by \((i,\mu)\). For a real vector \(b\) and a symmetric nonnegative definite matrix \(a\), both indexed like \(z\), the consistency equation of \cite{sharef2004conditions} §2.1, quoting \cite{filipovic1999note} Proposition 3.2, is

\begin{equation*}
\partial_xF(x,z)=\sum_ib^i\partial_{z_i}F(x,z)+\sum_{i,j}a^{i,j}\Bigl[\tfrac12\partial^2_{z_iz_j}F(x,z)-\partial_{z_i}F(x,z)\int_0^x\partial_{z_j}F(\eta,z)\,d\eta\Bigr]. \tag{3.160}\label{eq:3.160}
\end{equation*}

\phantomsection\label{res:Lemma-36}
\textbf{Lemma 36.} If \((z,a,b)\) satisfies (3.160) for every \(x\ge0\), then \(a_{(i,\mu),(j,\nu)}=0\) unless \(i=j=1\) and \(\mu,\nu\le\min(n_1,\lfloor n_2/2\rfloor)\). So a process consistent with the family has at most \(\min(n_1,\lfloor n_2/2\rfloor)+1\) components with nonzero diffusion, all among the \(Z^{1,\mu}\).

This is the first proposition of \cite{sharef2004conditions} §4.2; the proof below is ours.

\emph{Proof.} \(F\) is linear in \(z\), so the second derivatives vanish, and the two sides of (3.160) are exponential polynomials with exponents \(-\beta,-2\beta,-3\beta,-4\beta\) and no exponent zero apart from constants multiplying exponentials. Since \(\int_0^x\eta^\nu e^{-c\eta}\,d\eta=\nu!/c^{\nu+1}-e^{-cx}\bigl(x^\nu/c+\text{lower powers}\bigr)\), the term \(-a_{(i,\mu),(j,\nu)}\partial_{z_{i,\mu}}F\int_0^x\partial_{z_{j,\nu}}F\) contributes \(a_{(i,\mu),(j,\nu)}x^{\mu+\nu}e^{-(i+j)\beta x}/(j\beta)\) plus lower powers. Only the pairs \(i=j=2\) produce \(e^{-4\beta x}\). If some \(a_{(2,k),(2,k)}\neq0\), take \(k\) largest; nonnegative definiteness gives \(a_{(2,\mu),(2,\nu)}=0\) when \(\mu\) or \(\nu\) exceeds \(k\), so the coefficient of \(x^{2k}e^{-4\beta x}\) is \(a_{(2,k),(2,k)}/(2\beta)\neq0\), which contradicts (3.160). Hence every diagonal entry with \(i=2\) vanishes, and with it every entry in its row and column. The coefficient of \(e^{-2\beta x}\) then receives, from \(a\), the terms \(a_{(1,\mu),(1,\nu)}x^{\mu+\nu}/\beta\) plus lower powers, while \(\partial_xF\) and the \(b\)-terms contribute degree at most \(n_2\). If some \(a_{(1,k),(1,k)}\neq0\) with \(2k>n_2\), take \(k\) largest; the coefficient of \(x^{2k}e^{-2\beta x}\) is \(a_{(1,k),(1,k)}/\beta\neq0\), a contradiction. So \(a_{(1,\mu),(1,\mu)}=0\) for \(\mu>\lfloor n_2/2\rfloor\), and nonnegative definiteness removes the rest.\hfill$\square$

Now splice the family onto a step-plus-ramp front end. Keep the dates and horizon of Section \hyperref[sec:3.25]{3.25}, let \(W=(W^S,W^B)\) be a standard Brownian motion in \(\mathbb R^{d_S+d_B}\) with the cited calculus of Section \hyperref[sec:3.25]{3.25} and Itô's formula (\cite{bauerschmidt2020stochastic} §3.5), and let

\begin{equation*}
f(t,T)=\sum_m\mathbf 1_{I_m}(T)\bigl[L_m(t)+C_m(t)(T-T_m)\bigr]+F(T-t,Z_t),\qquad 0\le t\le T\le H, \tag{3.161}\label{eq:3.161}
\end{equation*}

where each \(L_m\) and \(C_m\) is a continuous Itô process driven by \(W^S\) alone, and \(dZ=b\,dt+\sigma_Z\,dW^B\) with \(\sigma_Z\) predictable with locally square-integrable paths and \(b\) progressively measurable with locally integrable paths; put \(a=\sigma_Z\sigma_Z^\top\).

\phantomsection\label{res:Proposition-37}
\textbf{Proposition 37 (conditional).} Suppose the drift and volatility that Itô's formula gives for (3.161) satisfy Assumption 2.1.

(a) Then for almost every \(t\), almost surely, \((Z_t,a_t,b_t)\) satisfies (3.160) for every \(x\ge0\), not only for \(x\le H-t\), and the front end satisfies Assumption 2.1 for its own drivers.

(b) Hence Lemma \hyperref[res:Lemma-36]{36} applies to \(Z\): at most \(\min(n_1,\lfloor n_2/2\rfloor)+1\) components of \(Z\) have nonzero diffusion. Conversely, if \((Z,a,b)\) solves (3.160) for every \(x\ge0\) and the front end satisfies Assumption 2.1 for its own drivers, as a front end of Theorem \hyperref[res:Theorem-11]{11} driven by \(W^S\) does, then the curve (3.161) satisfies Assumption 2.1. So the uncorrelated splice imposes on the block exactly the restrictions of the block alone.

(c) The front end absorbs no part of (3.160), because the functions \(x^ke^{-i\beta x}\), \(i=1,\dots,4\), are linearly independent of the polynomials on every interval, and the splice has no cross terms, because the two groups of drivers are independent.

\cite{gellert2021short} §3 splices by independence without discussing the block's restrictions. So, in the uncorrelated splice, the exponent restrictions of \cite{sharef2004conditions} survive unchanged inside the block.

\emph{Proof.} For \(T\in I_m\) the front-end part is an Itô process with volatility \(\sigma^S\) in the \(W^S\) coordinates and drift \(D^S(t,\cdot)\) affine on \(I_m\). Itô's formula for \((t,z)\mapsto F(T-t,z)\), which is smooth, gives the drift \(D^B=-\partial_xF+\sum_ib^i\partial_{z_i}F+\tfrac12\sum_{i,j}a^{i,j}\partial^2_{z_iz_j}F\) and the volatility \(\sigma^B=\sum_i\partial_{z_i}F\,\sigma_Z^{i,\cdot}\) in the \(W^B\) coordinates, at \((T-t,Z_t)\). Because the driver groups are disjoint, \(|\int\sigma|^2=|\int\sigma^S|^2+|\int\sigma^B|^2\), and differentiating (2.3) in \(T\) gives \(D^S+D^B=\sigma^S\cdot\int_t^T\sigma^S+\sigma^B\cdot\int_t^T\sigma^B\) for almost every \(T\). With \(x=T-t\), \(\sigma^B\cdot\int\sigma^B=\sum_{i,j}a^{i,j}\partial_{z_i}F(x,Z)\int_0^x\partial_{z_j}F(\eta,Z)\,d\eta\), so the residual \(R(t,x)=D^B-\sigma^B\cdot\int\sigma^B\) of (3.160) equals \(-(D^S-\sigma^S\cdot\int\sigma^S)(t,t+x)\). The right side is a polynomial in \(x\) on each \(I_m\cap[t,H]\), while \(R(t,\cdot)\) is an exponential polynomial with exponents \(-\beta,\dots,-4\beta\) whose coefficients are polynomials; by linear independence both vanish on each such interval, hence on \([0,H-t]\), and by analyticity \(R(t,x)=0\) for all \(x\ge0\). This is (a) and (c); the vanishing of the right side is the front end's own Assumption 2.1. Lemma \hyperref[res:Lemma-36]{36} gives the bound in (b). For the converse, (3.160) and the front end's condition give the differentiated identity, and integrating in \(T\) gives (2.3).\hfill$\square$

The proposition is conditional on the cited calculus. It assumes Assumption 2.1 for the drift of the representation (3.161); identifying it with the drift of another Itô decomposition of \(f(\cdot,T)\) uses uniqueness of Itô decompositions, which is standard but not among the cited inputs. Not covered: correlated splices, where cross terms appear (Propositions \hyperref[res:Proposition-34]{34} and \hyperref[res:Proposition-35]{35} treat deterministic volatilities); blocks with terms of exponent zero, such as the constant of the Svensson and Nelson--Siegel families, where the front end absorbs the part of (3.160) that is affine in \(T\) (Proposition \hyperref[res:Proposition-38]{38}); the extended Svensson families of \cite{sharef2004conditions} §5 as such; the front end's own consistency when its volatility is stochastic; and existence of a parameter process with nontrivial factors. Sufficient and exact conditions for \(\min(n_1,\lfloor n_2/2\rfloor)+1\) nondegenerate factors of the block alone are Proposition \hyperref[res:Proposition-39]{39}.

\subsection{Blocks with constant or linear terms}\label{sec:3.33}

Keep the uncorrelated splice (3.161) of Section \hyperref[sec:3.32]{3.32} with two changes. The block is any family linear in its parameters, \(F(x,z)=\sum_{i=1}^Nz_i\varphi_i(x)\), with fixed real-analytic \(\varphi_i\) linearly independent on every interval. The front end is of the type of Theorem \hyperref[res:Theorem-11]{11}: each \(C_m\) has finite variation, so the front-end volatility is constant on each \(I_m\), while the drifts of \(L_m\) and \(C_m\) are free. The block residual is

\begin{equation*}
R(t,x)=\sum_ib^i\varphi_i(x)-\sum_{i,j}a^{ij}\varphi_i(x)\int_0^x\varphi_j(\eta)\,d\eta-\partial_xF(x,Z_t), \tag{3.162}\label{eq:3.162}
\end{equation*}

the difference of the two sides of (3.160), since \(\partial^2_zF=0\).

\phantomsection\label{res:Proposition-38}
\textbf{Proposition 38 (conditional).} Suppose the drift and volatility that Itô's formula gives for (3.161) satisfy Assumption 2.1.

(a) Then for almost every \(t\), almost surely, \(R(t,\cdot)\) is affine: \(R(t,x)=c_0(t)+c_1(t)x\) for all \(x\ge0\). Conversely, if \(R(t,\cdot)\) is affine and the front-end drifts are chosen so that \(D^S(t,T)-\sigma^S(t,T)\cdot\int_t^T\sigma^S(t,u)\,du=-R(t,T-t)\) on each \(I_m\), the splice satisfies Assumption 2.1. If the block contains no nonzero affine function this gives \(R=0\), as in Proposition \hyperref[res:Proposition-37]{37}.

(b) For the Nelson--Siegel family with fixed exponent \(\beta>0\),

\begin{equation*}
F(x,z)=z_1+(z_2+z_3x)e^{-\beta x}, \tag{3.163}\label{eq:3.163}
\end{equation*}

\(R(t,\cdot)\) is affine if and only if \(a\) is supported on the single entry \(a^{11}\), \(b^2=z_3-\beta z_2\) and \(b^3=-\beta z_3\), with \(a^{11}\ge0\) and \(b^1\) arbitrary. Then \(R=b^1-a^{11}x\).

(c) Consequently, in the uncorrelated splice only the level \(z_1\) of a Nelson--Siegel block may carry a factor. A move of \(z_1\) is the same curve move as adding one common constant to every \(L_m\), a parallel shift that the front-end family already contains, driven by \(W^B\); the front-end drifts absorb its convexity, \(C_m\) acquiring the ramp \(a^{11}\), so that the front end is then not consistent on its own. The block's own parameters \(z_2\) and \(z_3\) carry no factor, as when the block stands alone.

Alone, a consistent Itô process for the Nelson--Siegel family is deterministic, with a constant level and \(z_2\), \(z_3\) following the drifts in (b) (\cite{filipovic1999note} Proposition 3.2 and Theorem 4.1). Proposition \hyperref[res:Proposition-38]{38} shows that the scheduled-step front end changes this only through the overlap of the two families, the constants.

\emph{Proof.} (a) As in the proof of Proposition \hyperref[res:Proposition-37]{37}, Assumption 2.1 gives \(R(t,T-t)=-(D^S-\sigma^S\cdot\int\sigma^S)(t,T)\) for almost every \(T\in[t,H]\). The front-end volatility is constant on each \(I_m\) and \(D^S\) is affine there, so the right side is affine in \(T\), hence in \(x=T-t\), on each \(I_m\cap[t,H]\), almost everywhere and then everywhere by continuity. \(R(t,\cdot)\) is real-analytic, because the primitive of a real-analytic function is real-analytic, so it is affine on \(\mathbb R\). Conversely, if \(R=c_0+c_1x\), take the drift of \(L_m\) to be \(-(c_0+c_1(T_m-t))\) plus the value of \(\sigma^S\cdot\int\sigma^S\) at \(T_m\), and the drift of \(C_m\) to be \(-c_1\) plus the slope of that term on \(I_m\); then the differentiated drift condition holds and integrating in \(T\) gives (2.3).

(b) Here \(\varphi_1=1\), \(\varphi_2=e^{-\beta x}\), \(\varphi_3=xe^{-\beta x}\), and \(\int_0^x\eta^ke^{-\beta\eta}\,d\eta=k!/\beta^{k+1}-P_k(x)e^{-\beta x}\) for polynomials \(P_k\) with \(x^k=\beta P_k-P_k'\). Only the pairs \(i,j\in\{2,3\}\) produce \(e^{-2\beta x}\), with coefficient \(\sum a^{ij}x^{i-2}P_{j-2}=\beta S-S'/2\), where \(S=\sum a^{ij}P_{i-2}P_{j-2}\) and \(a\) is symmetric. If \(R\) is affine this coefficient vanishes; the polynomial solutions of \(\beta S-S'/2=0\) are zero, so \(S=0\). The \(\{2,3\}\) block of \(a\) is nonnegative definite, so \(S\) is a sum of squares of elements of the span of \(P_0\) and \(P_1\), which are independent; hence that block vanishes, and nonnegative definiteness gives \(a^{12}=a^{13}=0\). With only \(a^{11}\) left, the \(e^{-\beta x}\) part of \(R\) is \((b^2-z_3+\beta z_2)+(b^3+\beta z_3)x\), which must vanish, and the exponent-zero part is \(b^1-a^{11}x\). The converse is the same computation.

(c) \(a\) supported on \(a^{11}\) means \(\sigma_Z\) has no loading on \(z_2\) or \(z_3\), and moving \(z_1\) by \(c\) moves \(F(\cdot,z)\) by the constant \(c\), as does adding \(c\) to every \(L_m\); the ramp is the \(c_1=-a^{11}\) of (a)'s converse.\hfill$\square$

The proposition is conditional on the cited calculus, with Assumption 2.1 for the drift of the representation (3.161) as in Proposition \hyperref[res:Proposition-37]{37}. Not covered: a varying Nelson--Siegel exponent; the Svensson family, whose second exponential needs its own computation (Proposition \hyperref[res:Proposition-40]{40} for fixed exponents); blocks whose polynomial part has degree two or more, where only part of the exponent-zero equation is absorbed (Proposition \hyperref[res:Proposition-45]{45}); front ends with stochastic ramp volatility, which absorb polynomials of degree up to three; correlated splices; and existence of a process with the stated coefficients.

\subsection{Correlated factors for the biexponential-polynomial family}\label{sec:3.34}

Keep the family (3.159) and the consistency equation (3.160) of Section \hyperref[sec:3.32]{3.32}, and put \(n=\lfloor n_2/2\rfloor\). Assume the conditions and conventions of \cite{sharef2004conditions} §4.2: (C1) \(n_1\ge n\); (C2) \(z_{1,i}\neq0\) for \(0\le i\le n_1\); (C3) for \(0\le i\le2n\), \(z_{2,i}\neq0\) or \(z_{2,i+1}\neq0\); and \(z_{2,n_2}\neq0\). Let \(K_0=\{k\le2n:z_{2,k}=0\}\); as in the source's reduced system, a coefficient \(z_{2,k}\) with \(k\in K_0\) is structurally absent, not a parameter, and its drift \(b_{2,k}\) is zero. By (C3), \(K_0\) contains no two consecutive integers. Put \(r_k=(k+1)z_{2,k+1}\) for \(k\in K_0\). By Lemma \hyperref[res:Lemma-36]{36}, \(a\) is supported on the block \(A=(a_{(1,\mu),(1,\nu)})_{0\le\mu,\nu\le n}\), and \(n+1\) nondegenerate factors means \(A\) positive definite. Let \(P_\nu\) be the polynomials with \(P_\nu(x)e^{-\beta x}=\int_x^\infty\eta^\nu e^{-\beta\eta}\,d\eta\), and for a polynomial \(S=\sum_{j\le2n}s_jx^j\) put

\begin{equation*}
t_k(S)=\beta s_k-\tfrac12(k+1)s_{k+1},\qquad S_A=\sum_{\mu,\nu\le n}a_{\mu\nu}P_\mu P_\nu. \tag{3.164}\label{eq:3.164}
\end{equation*}

\phantomsection\label{res:Proposition-39}
\textbf{Proposition 39.} Under these conditions:

(a) \((z,a,b)\) satisfies (3.160) for every \(x\ge0\) if and only if the drifts \(b\) are given by explicit formulas in \((z,A)\) and

\begin{equation*}
t_k(S_A)=r_k\qquad\text{for every }k\in K_0. \tag{3.165}\label{eq:3.165}
\end{equation*}

(b) A positive definite \(A\) satisfying (a) exists if and only if there is a polynomial \(S\) of degree at most \(2n\) with \(S>0\) on \(\mathbb R\), \(s_{2n}>0\) and \(t_k(S)=r_k\) for every \(k\in K_0\). Equivalently, there is no \(\lambda\neq0\) in \(\mathbb R^{K_0}\) such that the coefficient sequence \(\ell\) of \(\sum_k\lambda_kt_k\) has a positive semidefinite Hankel matrix \((\ell_{i+j})_{0\le i,j\le n}\) and \(\sum_k\lambda_kr_k\le0\).

(c) If \(n_2=2n\) or \(z_{2,2n}\neq0\), a positive definite \(A\) exists for every \(z\) satisfying the conditions: (3.160) has a solution with \(n+1\) nondegenerate, in general correlated, factors at every such parameter point.

(d) If \(n_2=2n+1\) and \(z_{2,2n}=0\), then \(z_{2,2n+1}>0\) is necessary, even for a nonnegative definite \(A\) with \(a_{nn}>0\). For \(n=0\) it is sufficient. For \(n=1\) the exact condition is \(z_{2,3}>0\), together with \(3z_{2,3}+4\beta^2z_{2,1}>0\) when also \(z_{2,0}=0\).

(e) At \(\beta=1\), \(n_1=0\), \(n_2=1\), \(z_1=(1)\), \(z_2=(0,-1)\), which satisfies the conditions, no nonnegative definite \(a\) and no \(b\) satisfy (3.160); and no path whose coordinate \(Z^{2,1}\) is right-continuous at a time where it is negative satisfies (3.160) at almost every later time near it.

Part (e) shows that the central proposition of \cite{sharef2004conditions} §4.2, which asserts \(\min(n_1,\lfloor n_2/2\rfloor)+1\) nontrivial factors under these conditions, is false as the source's reduced system states it, with \(z_{2,k}\) structurally absent and \(b_{2,k}=0\) for \(k\in K_0\); if every \(b_{2,k}\) were free, the proposition would hold trivially. Part (c) shows that its conclusion does hold, with correlated rather than diagonal factors, whenever \(2n\notin K_0\). The source's proof treats only diagonal \(A\), and its reduced system carries the wrong sign on the \(e^{-2\beta x}\) terms in \(a\). The identity behind (3.165) and the reduction to positive polynomials are ours; that nonnegative polynomials in one variable are sums of squares, and the Hankel description of the dual cone, are classical.

\emph{Proof.} (a) Since \(\int_0^x\eta^\nu e^{-\beta\eta}\,d\eta=\nu!/\beta^{\nu+1}-P_\nu(x)e^{-\beta x}\), the \(a\)-term of (3.160) is \(-\sum a_{\mu\nu}(\nu!/\beta^{\nu+1})x^\mu e^{-\beta x}+[\sum a_{\mu\nu}x^\mu P_\nu]e^{-2\beta x}\). The coefficient of \(x^ke^{-\beta x}\) fixes \(b_{1,k}\), since every \(z_{1,k}\) is a parameter by (C2); that of \(x^ke^{-2\beta x}\) fixes \(b_{2,k}\) for \(k\notin K_0\), and for \(k\in K_0\) it reads \(r_k=[x^k]\sum a_{\mu\nu}x^\mu P_\nu\); nothing arises for \(k>2n\). Writing \(A=\sum_i\varepsilon_iv_iv_i^\top\) with \(\varepsilon_i=\pm1\), \(q_i=\sum_\mu v_{i\mu}x^\mu\) and \(Q_i=\sum_\nu v_{i\nu}P_\nu\), differentiation gives \(q_i=\beta Q_i-Q_i'\), so \(\sum a_{\mu\nu}x^\mu P_\nu=\sum_i\varepsilon_i(\beta Q_i^2-Q_iQ_i')=\beta S_A-S_A'/2\), which is (3.165).

(b) Since \(P_\nu=x^\nu/\beta+\)lower terms, \(A\mapsto S_A\) is the Gram map \(M\mapsto m^\top Mm\), \(m=(1,x,\dots,x^n)\), composed with an invertible congruence. The Gram map sends the positive semidefinite cone onto the nonnegative polynomials of degree at most \(2n\), because these are sums of squares, and, being linear and surjective, sends the positive definite cone onto the interior \(U=\{S>0\text{ on }\mathbb R,\ s_{2n}>0\}\). This is the first form. For the second, separate the open convex cone \(U\) from the affine space \(\{t_k(S)=r_k,\ k\in K_0\}\): a separating functional is \(\sum_k\lambda_kt_k\), nonnegative on the nonnegative polynomials, which means that its coefficient sequence has a positive semidefinite Hankel matrix, and \(\lambda\cdot r\le0\); the \(t_k\) are independent because no two indices in \(K_0\) are consecutive.

(c) It suffices to find \(S_h\in U\) with \(t_k(S_h)=0\) on \(K_0\): then \(S_p+NS_h\in U\) for a particular solution \(S_p\) and large \(N\). By (b) with \(r=0\), this fails only if some \(\lambda\neq0\) gives a positive semidefinite Hankel matrix. Here \(\ell_j=\beta\lambda_j[j\in K_0]-\tfrac j2\lambda_{j-1}[j-1\in K_0]\), and a zero diagonal entry kills its row, so either \(\ell=\ell_0\delta_0\) or \(\ell_{2n}>0\). The first forces \(0\in K_0\), \(1\notin K_0\), \(\ell_1=-\lambda_0/2=0\), hence \(\ell_0=0\). In the second, \(2n\notin K_0\) forces \(2n-1\in K_0\), and descending induction with the \(2\times2\) principal minors gives \(\ell_{2j}>0\) and \(2j-1\in K_0\) for \(j=n,\dots,1\); at \(j=1\), \(\ell_0>0\) needs \(0\in K_0\) next to \(1\in K_0\), a contradiction. For \(n=0\), \(K_0=\emptyset\) and \(S_h=1\).

(d) \(\lambda=\delta_{2n}\) gives \(\ell=\beta\delta_{2n}\), whose Hankel matrix is positive semidefinite, so (b) requires \(r_{2n}=(2n+1)z_{2,2n+1}>0\); for nonnegative definite \(A\) with \(a_{nn}>0\) the leading coefficient \(a_{nn}/\beta^2\) of \(S_A\) gives the same. For \(n=0\), \(U=\{s_0>0\}\) and the condition is \(\beta s_0=r_0\). For \(n=1\), \(U=\{s_2>0,\ 4s_0s_2>s_1^2\}\); with \(K_0=\{2\}\) the condition is \(r_2>0\), and with \(K_0=\{0,2\}\) the quadratic condition on \(s_1\) has a solution if and only if \(s_2/\beta+4r_0>0\), which is the stated inequality.

(e) Here \(n=0\), \(K_0=\{0\}\), and (3.165) reads \(a_{00}/\beta=z_{2,1}=-1\), impossible for \(a_{00}\ge0\). A right-continuous \(Z^{2,1}\) that is negative at a time stays negative on an interval of positive length after it, where (3.160) would have to hold almost everywhere.\hfill$\square$

The proposition is pointwise algebra at one parameter point, as is the source's; existence of an Itô process with these coefficients along its path is not shown. Not covered: the explicit form of the condition in (b) for \(n\ge2\) when \(2n\in K_0\), which is a semidefinite feasibility condition; degenerate nonnegative definite \(A\) with positive diagonal but rank below \(n+1\); structurally absent \(z_{1,k}\), and coefficients that vanish at isolated times rather than structurally. With Proposition \hyperref[res:Proposition-37]{37} the conclusions transfer to the uncorrelated splice, but they are stated here for the block alone.

\subsection{The Svensson block with fixed exponents}\label{sec:3.35}

Keep the setting of Section \hyperref[sec:3.33]{3.33}, with the Svensson family with fixed exponents \(0<\beta_1\neq\beta_2\),

\begin{equation*}
F(x,z)=z_1+(z_2+z_3x)e^{-\beta_1x}+z_4\,x\,e^{-\beta_2x}, \tag{3.166}\label{eq:3.166}
\end{equation*}

and basis \(\varphi_1=1\), \(\varphi_2=e^{-\beta_1x}\), \(\varphi_3=xe^{-\beta_1x}\), \(\varphi_4=xe^{-\beta_2x}\). Alone means \(R(t,\cdot)=0\) in (3.162); in the splice means \(R(t,\cdot)\) affine, which by Proposition \hyperref[res:Proposition-38]{38}(a) is consistency of the splice. The statements hold at one parameter point \((z,a,b)\) with \(a\) nonnegative definite.

\phantomsection\label{res:Proposition-40}
\textbf{Proposition 40.} (a) If \(\beta_2\neq2\beta_1\) and \(z_4\neq0\), no \(a\) and no \(b\) make \(R\) affine. No path whose coordinate \(Z^4\) is right-continuous at a time where it is nonzero makes \(R\) affine at almost every later time near it, alone or in the splice.

(b) Let \(\beta_2=2\beta_1=:2\beta\). Then \(R\) is affine if and only if every entry of \(a\) involving \(z_3\) or \(z_4\) vanishes, \(a_{22}=\beta z_4\) (so \(z_4\ge0\)), \(b_2=z_3+z_4-\beta z_2-a_{12}/\beta\), \(b_3=a_{12}-\beta z_3\), \(b_4=-2\beta z_4\), and \(a_{11}a_{22}\ge a_{12}^2\), with \(b_1\) arbitrary; then \(R=b_1-a_{12}/\beta-a_{11}x\). Alone, in addition \(a_{11}=a_{12}=b_1=0\).

(c) Consequently the splice adds only the level \(z_1\) as a factor, a parallel shift already in the front-end family, correlated with \(z_2\) up to the Cauchy--Schwarz bound \(a_{12}^2\le a_{11}\beta z_4\); for \(z_4>0\) a correlated splice point with \(a_{12}\neq0\) exists. The block's own nontrivial factor is \(z_2\), with variance \(\beta z_4\), as alone.

The Svensson family alone is treated in \cite{filipovic2000exponential} §9.2: by its Theorems 8.1 and 8.2 only \(\beta_1=2\beta_2\) or \(\beta_2=2\beta_1\) can carry a nontrivial process, and in the second case its equations (30) and (31) are the case alone of (b). Part (a) makes explicit that the first case also needs \(z_4=0\). Parts (b) and (c) in the splice follow the affine reduction of Proposition \hyperref[res:Proposition-38]{38}.

\emph{Proof.} Write \(R\) as a sum of polynomials times \(e^{-\kappa x}\), \(\kappa\in\{0,\beta_1,\beta_2,2\beta_1,\beta_1+\beta_2,2\beta_2\}\), using \(\int_0^x\eta^ke^{-\gamma\eta}\,d\eta=k!/\gamma^{k+1}-P^{(\gamma)}_k(x)e^{-\gamma x}\) with \(P^{(\gamma)}_0=1/\gamma\) and \(P^{(\gamma)}_1=x/\gamma+1/\gamma^2\). Polynomials times exponentials with distinct real exponents are linearly independent on every interval, so every coefficient with \(\kappa\neq0\) vanishes, in the splice as alone.

The only pair with exponent \(2\beta_2\) is \(\varphi_4\) with \(\int\varphi_4\), contributing \(a_{44}(x^2/\beta_2+x/\beta_2^2)e^{-2\beta_2x}\), unless \(2\beta_2=\beta_1\). In the first case \(a_{44}=0\). If \(2\beta_2=\beta_1\), the group \(e^{-\beta_1x}\) also receives \(-a_{13}x^2e^{-\beta_1x}\), and the group \(e^{-2\beta_1x}\) has \(x^2\)-coefficient \(a_{33}/\beta_1\) and no other source; so \(a_{33}=0\), hence \(a_{13}=0\) and \(a_{44}=0\). Either way row 4 of \(a\) vanishes by nonnegative definiteness.

(a) With row 4 zero, the terms with exponent \(\beta_2\) are \(-z_4(1-\beta_2x)e^{-\beta_2x}\) and \(b_4xe^{-\beta_2x}\), together with terms from \(\varphi_2,\varphi_3\) only if \(2\beta_1=\beta_2\). Otherwise the constant coefficient \(-z_4\) must vanish. A right-continuous \(Z^4\) nonzero at a time stays nonzero on an interval of positive length after it, where the condition must hold almost everywhere.

(b) Row 4 vanishes since \(4\beta\neq\beta\). The \(x^2e^{-2\beta x}\) coefficient is \(a_{33}/\beta-a_{14}\), so \(a_{33}=0\) and row 3 vanishes. The remaining \(e^{-2\beta x}\) terms are \(-z_4(1-2\beta x)\), \(b_4x\) and \(a_{22}/\beta\), giving \(a_{22}=\beta z_4\) and \(b_4=-2\beta z_4\). The \(e^{-\beta x}\) terms are \(-(z_3-\beta z_2-\beta z_3x)\), \(b_2+b_3x\) and \(a_{12}/\beta-a_{12}x-a_{22}/\beta\), giving \(b_2\) and \(b_3\). The exponents \(3\beta\) and \(4\beta\) involve rows 3 and 4 only, and the exponent-zero part is \(b_1-a_{11}x-a_{12}/\beta\). Conversely these values make every part with \(\kappa\neq0\) vanish. Alone, the exponent-zero part vanishes too, so \(a_{11}=0\), hence \(a_{12}=0\) and \(b_1=0\).

(c) The bound is nonnegative definiteness of the \(\{1,2\}\) block, and \(a=vv^\top\) with \(v=(1,\sqrt{\beta z_4},0,0)\) gives \(a_{12}\neq0\); the level is a parallel shift as in Proposition \hyperref[res:Proposition-38]{38}(c).\hfill$\square$

The proposition is pointwise algebra, with the pathwise part (a); existence of a process with these coefficients is not shown. Not covered: varying exponents; correlated splices; a front end with stochastic ramp volatility, which absorbs polynomials of degree up to three; the case \(z_4=0\) is Proposition \hyperref[res:Proposition-38]{38}(b).

\subsection{The correlated splice with stochastic scales}\label{sec:3.36}

Keep the setting of Section \hyperref[sec:3.31]{3.31}, including the quasi-exponential \(\lambda=c\,e^{Ax}b\) with \(A\) invertible, \((A,b)\) controllable and \(c\neq0\), the blocks and consistency, and replace the deterministic volatilities by predictable bounded processes:

\begin{equation*}
\sigma^S(u,T)=s_m(u)\ (T\in I_m),\qquad \sigma^B(u,T)=\psi_u\,\lambda(T-u),\qquad dW^S=dB^1,\qquad dW^B=\rho_u\,dB^1+\sqrt{1-\rho_u^2}\,dB^2, \tag{3.167}\label{eq:3.167}
\end{equation*}

with \(s_0,\dots,s_N\), \(\psi\) and \(\rho\) predictable and bounded, \(|\rho|\le1\), and \(\Delta_m=s_m-s_{m-1}\). The drift is (3.153) with \(\rho\) replaced by \(\rho_u\), path by path; it satisfies Assumption 2.1 with respect to the independent drivers. The initial curve lies in \(\mathcal M(\mathcal E)\).

\phantomsection\label{res:Proposition-41}
\textbf{Proposition 41 (conditional).} Under the cited calculus:

(a) For \(t<T_m\), on every path, the jump across \(T_m\) of the analytic extensions of the cross term is

\begin{equation*}
c\,e^{AT}\bigl(A^{-1}+(T-T_m)I\bigr)v_m(t)-c\,A^{-1}b\int_0^t\rho_u\psi_u\Delta_m(u)\,du,\qquad v_m(t)=\int_0^t\rho_u\psi_u\Delta_m(u)\,e^{-Au}b\,du. \tag{3.168}\label{eq:3.168}
\end{equation*}

For every \(t\) and \(T\), almost surely, the random part of \(f(t,T)\) equals \(\sum_m\mathbf 1_{I_m}(T)\int_0^ts_m\,dW^S+c\,e^{AT}\int_0^t\psi_ue^{-Au}b\,dW^B_u\), one version for all \(T\).

(b) If the curve is consistent with \(\mathcal M(\mathcal E)\) for some block \(\mathcal E\), then almost surely \(\rho_u\psi_u\Delta_m(u)=0\) for almost every \(u<T_m\), for every \(m\).

(c) (i) If \(f(0,\cdot)\in\mathcal M(\mathcal E_1)\) and almost surely \(\rho_u\psi_u\Delta_m(u)=0\) for almost every \(u<T_m\), for every \(m\), the curve is consistent with \(\mathcal M(\mathcal E_1)\). (ii) For all predictable bounded scales, the curve is consistent with the indicator-enlarged family of Proposition \hyperref[res:Proposition-35]{35}(c)(ii).

(d) If, with positive probability, \(\rho_u\psi_u\Delta_m(u)\neq0\) on a set of \(u\) of positive measure for some \(m\), the curve is consistent with \(\mathcal M(\mathcal E)\) for no block \(\mathcal E\).

So with predictable bounded scales on the step factor, the block and the correlation, the independence splice and the indicator-times-quasi-exponential terms remain the only two ways; a correlation that is switched off wherever the step factor has maturity steps is the independence splice locally in time. \cite{brace2024sofr} eqs. 35 to 41 construct the indicator-times-exponential model with Heston scales and do not claim that it is forced; \cite{bjork2004finite} has neither breakpoints nor a splice.

\emph{Proof.} The integrands are predictable and bounded, hence admissible, and the drift is bounded, so its time integral is a pathwise Lebesgue integral. For a fixed path and \(u\), (3.167) is the volatility of Section \hyperref[sec:3.31]{3.31} with the constants replaced by the numbers \(s_m(u)\), \(\psi_ub\) and \(\rho_u\), so Assumption 2.1 holds as there, and the cross part of the drift is \(\rho_u\psi_u\) times the integrand of Proposition \hyperref[res:Proposition-35]{35}(a); integrating in \(u\) gives (3.168) on every path. The random part follows by linearity, with finitely many random variables. For (b), fix \(t<T_m\); on one event of probability one the version in \(\mathcal M_t(\mathcal E)\) agrees with the version from (a) plus the drift parts, the parts other than the cross term lie in \(\mathcal M_t(\mathcal E)\), and the argument of Proposition \hyperref[res:Proposition-35]{35}(b), which is pathwise, gives \(v_m(t,\omega)\in\mathcal K\). Intersecting over rational \(t<T_m\) and using continuity of \(v_m(\cdot,\omega)\) and closedness of \(\mathcal K\), \(v_m(t,\omega)\in\mathcal K\) for every \(t<T_m\); then the derivative \(\rho_u\psi_u\Delta_m(u)e^{-Au}b\) lies in \(\mathcal K\) for almost every \(u\), and controllability, as in Proposition \hyperref[res:Proposition-35]{35}(b), forces \(\rho_u\psi_u\Delta_m(u)=0\) almost everywhere. Part (c) is Proposition \hyperref[res:Proposition-35]{35}(c) path by path; the version is measurable because predictable scales are jointly measurable in \((u,\omega)\), and in (i) it is set to zero on a measurable null set. For (d), on the event named there some rational \(t<T_m\) has \(v_m(t)\notin\mathcal K\) with positive probability, and (b)'s argument excludes a version at that \(t\).\hfill$\square$

The proposition is conditional on the cited calculus, with controllability, \(c\neq0\) and \(A\) invertible as standing hypotheses as in Proposition \hyperref[res:Proposition-35]{35}. Not covered: a block volatility that depends on the block's own state, where the jump carries a random vector in place of \(\psi_ub\) (Proposition \hyperref[res:Proposition-43]{43}); several step or block factors; unbounded or square-root scales, such as the Heston scale of \cite{brace2024sofr}; a diffusive front-end ramp, which the front end's own Assumption 2.1 excludes (Theorem \hyperref[res:Theorem-11]{11}, Proposition \hyperref[res:Proposition-28]{28}(c)); and existence of processes beyond the integrals themselves.

\subsection{The correlated splice with a state-dependent block volatility}\label{sec:3.37}

Keep the splice setting of Section \hyperref[sec:3.33]{3.33}, with independent drivers \(B=(B^1,\dots,B^k)\) and the cited calculus with Itô's formula. The front end is of the type of Theorem \hyperref[res:Theorem-11]{11}, \(L_m(t)+C_m(t)(T-T_m)\) on \(I_m\) with \(C_m\) of finite variation and free drifts, and \(L_m\) has the noise \(s_m(t)H^S(t)\cdot dB\) with \(s_m\) predictable and bounded, \(\Delta_m=s_m-s_{m-1}\), and a predictable bounded row vector \(H^S\). The block is a quasi-exponential family in minimal form,

\begin{equation*}
F(x,z)=c\,e^{Ax}z,\qquad z\in\mathbb R^r, \tag{3.169}\label{eq:3.169}
\end{equation*}

with \(A\) invertible and \((A,c)\) observable, that is, \(c\,e^{Ax}y\equiv0\) only for \(y=0\); every finite-dimensional shift-invariant family of exponential polynomials without an exponent-zero part has this form. The curve is

\begin{equation*}
f(t,T)=\sum_m\mathbf 1_{I_m}(T)\bigl[L_m(t)+C_m(t)(T-T_m)\bigr]+c\,e^{A(T-t)}Z_t, \tag{3.170}\label{eq:3.170}
\end{equation*}

with \(dZ=b\,dt+H^Z\,dB\), \(H^Z\) a predictable bounded \(r\times k\) matrix and \(b\) locally integrable, so the block volatility depends on the block's own state through \(H^Z\). Put

\begin{equation*}
w(u)=H^Z(u)H^S(u)^\top\in\mathbb R^r, \tag{3.171}\label{eq:3.171}
\end{equation*}

the covariation density of the block's noise with the step noise. All correlation enters through the weights: the independence splice has \(H^S\) and the rows of \(H^Z\) on disjoint drivers, and a correlation \(\rho_u\) with a single block direction \(\sigma_Z(u)\) gives \(w=\rho_u\sigma_Z(u)\). Assumption 2.1 is assumed for the drift and volatility that Itô's formula gives for (3.170).

\phantomsection\label{res:Lemma-42}
\textbf{Lemma 42.} For \(v\in\mathbb R^r\) and real \(T_m\), \(c\,e^{AT}(TI+A^{-1}-T_mI)v\equiv0\) only if \(v=0\).

\emph{Proof.} With \(k(T)=c\,e^{AT}A^{-1}v\), the function is \((T-T_m)k'(T)+k(T)=\frac d{dT}[(T-T_m)k(T)]\). If it vanishes, \((T-T_m)k(T)\) is constant and zero at \(T_m\), so \(k\equiv0\), and observability gives \(A^{-1}v=0\).\hfill$\square$

\phantomsection\label{res:Proposition-43}
\textbf{Proposition 43 (conditional).} (a) Almost surely \(\Delta_m(u)w(u)=0\) for almost every \(u<T_m\), for every \(m\): the block's noise is orthogonal to the step noise wherever the step volatility has a maturity step.

(b) If \(\Delta_mw=0\) almost everywhere, almost surely, the integrated cross term has no jump across the meetings, and the curve satisfies Assumption 2.1 if and only if the cross term together with the block residual is affine in the time to maturity and the front-end drifts absorb that affine part, as in Proposition \hyperref[res:Proposition-38]{38}(a).

(c) If \(\Delta_mw\neq0\) on a set of positive \(dt\otimes dQ\) measure, no choice of the drifts makes (3.170) satisfy Assumption 2.1; on each maturity interval the jump to be absorbed has the form \(c\,e^{AT}y+Tc\,e^{AT}z\), which a family can only absorb by indicator-times-exponential terms as in \cite{brace2024sofr} eq. 36.

Proposition \hyperref[res:Proposition-35]{35} is the case \(w(u)=\rho\psi b\) with \((A,b)\) controllable, and Proposition \hyperref[res:Proposition-41]{41} its extension to random scales; here observability of \((A,c)\) replaces controllability and makes the subspace of Proposition \hyperref[res:Proposition-35]{35}(b) trivial. \cite{brace2024sofr} eqs. 35 to 41 make each factor's volatility state-dependent through its Heston variance but keep the factors independent; \cite{bjork2004finite} Propositions 4.11, 4.12 and 5.1 concern scale states that move the shapes, which are fixed here.

\emph{Proof.} The block volatility is \(c\,e^{A(T-u)}H^Z(u)\) and the front-end volatility \(s_m(u)H^S(u)\), both in driver form, so the differentiated drift condition is \(D^S+D^B=\sigma\cdot\int\sigma\) with \(\sigma=\sigma^S+\sigma^B\). With \(S^S(u,T)=s_m(u)T+\kappa_m(u)\) on \(I_m\), the cross term is \(s_m(u)cA^{-1}(e^{A(T-u)}-I)w(u)+c\,e^{A(T-u)}w(u)(s_m(u)T+\kappa_m(u))\), the integrand of Proposition \hyperref[res:Proposition-35]{35} with \(\rho b\) replaced by \(w(u)\). Integrating over \(u\in[0,t]\) path by path, the front-end parts are affine in \(T\) on each \(I_m\cap[t,H]\), the block parts are real-analytic and the same on every interval, and the cross term \(X_{t,m}\) is entire and depends on \(m\). So \(X_{t,m}-X_{t,m-1}=c\,e^{AT}(A^{-1}+(T-T_m)I)v_m(t)-cA^{-1}\int_0^t\Delta_mw\,du\), with \(v_m(t)=\int_0^t\Delta_m(u)e^{-Au}w(u)\,du\), must be affine; its exponential part has only nonzero exponents and vanishes, and Lemma \hyperref[res:Lemma-42]{42} gives \(v_m(t)=0\). This holds for rational \(t<T_m\) on one event, hence for all \(t<T_m\) by continuity, and differentiation gives (a). Under (a)'s condition the jump vanishes, and (b) is the affine reduction of Proposition \hyperref[res:Proposition-38]{38}(a) with the cross term added to the block residual. For (c), on an event of positive probability some rational \(t\) gives a jump that is not affine.\hfill$\square$

The proposition is conditional on the cited calculus. Not covered: blocks with an exponent-zero part, such as the Nelson--Siegel or Svensson level, whose covariation with the step noise enters the affine part; several step factors (Proposition \hyperref[res:Proposition-44]{44}); unbounded volatilities and the Heston scale; and existence of the processes \(Z\) and \(L_m\) with given coefficients.

\subsection{The correlated splice with several step and block factors}\label{sec:3.38}

Keep the setting of Section \hyperref[sec:3.37]{3.37} in driver form, now with \(n_S\) step factors: on \(I_m\) the front-end noise is \(\sum_{\ell=1}^{n_S}s_{\ell,m}(t)H^S_\ell(t)\cdot dB\), with \(s_{\ell,m}\) predictable and bounded and predictable bounded row vectors \(H^S_\ell\). Put \(\Delta_{\ell,m}=s_{\ell,m}-s_{\ell,m-1}\) and define the aggregate maturity-step integrand at \(T_m\),

\begin{equation*}
G_m(u)=\sum_{\ell=1}^{n_S}\Delta_{\ell,m}(u)H^S_\ell(u)\in\mathbb R^{1\times k}. \tag{3.172}\label{eq:3.172}
\end{equation*}

The block is a sum of quasi-exponential factors,

\begin{equation*}
F(x,z)=\sum_{j=1}^Jc_je^{A_jx}z_j=\mathcal C\,e^{\mathcal Ax}z,\qquad \mathcal A=\operatorname{diag}(A_1,\dots,A_J)\text{ invertible},\ \mathcal C=(c_1,\dots,c_J), \tag{3.173}\label{eq:3.173}
\end{equation*}

with \(dZ=b\,dt+H^Z\,dB\) as in Section \hyperref[sec:3.37]{3.37} and \(H^Z_j\) the rows of \(H^Z\) belonging to \(z_j\). Assumption 2.1 is assumed as in Proposition \hyperref[res:Proposition-43]{43}.

\phantomsection\label{res:Proposition-44}
\textbf{Proposition 44 (conditional).} (a) If \((\mathcal A,\mathcal C)\) is observable, then almost surely

\begin{equation*}
H^Z(u)\,G_m(u)^\top=0\qquad\text{for almost every }u<T_m,\text{ every }m: \tag{3.174}\label{eq:3.174}
\end{equation*}

the block's noise is orthogonal to the front end's aggregate maturity-step noise. Conversely, under (3.174) the integrated cross term has no jump across the meetings, so (3.174) is exactly the condition for a jump-free cross term; the curve then satisfies Assumption 2.1 if and only if that cross term together with the block residual is affine in the time to maturity and the front-end drifts absorb that affine part, as in Proposition \hyperref[res:Proposition-38]{38}(a) applied to the residual plus the cross term. The cross term need not be affine, so (3.174) is necessary for consistency but does not characterize it; Proposition \hyperref[res:Proposition-52]{52}(b)(iv) and (e) make the converse explicit and show that, for diagonalizable \(\mathcal A\), consistency also requires \(H^\zeta V_{j(u)}^\top=0\).

(b) If each \((A_j,c_j)\) is observable and the spectra of the \(A_j\) are pairwise disjoint, then \((\mathcal A,\mathcal C)\) is observable, and (3.174) holds block by block: \(H^Z_j(u)G_m(u)^\top=0\) for each \(j\).

(c) In general, a minimal realization \(F(x,z)=\mathcal C'e^{\mathcal A'x}Pz\) with \((\mathcal A',\mathcal C')\) observable, \(\mathcal A'\) invertible and \(P\mathcal A=\mathcal A'P\) exists, and (3.174) holds for the reduced noise \(PH^Z\): blocks with shared exponents act as one effective factor.

(d) If (3.174), or its reduced form in (c), fails on a set of positive \(dt\otimes dQ\) measure, no choice of the drifts makes the curve satisfy Assumption 2.1, and indicator-times-exponential terms are forced as in Proposition \hyperref[res:Proposition-43]{43}(c).

We count (3.174) as the independence splice in vector form; this is a classification convention, not a result. Condition (3.174) can hold even when single step factors are correlated with the block, as long as the block is uncorrelated with the aggregate step, and the conclusion of Sections \hyperref[sec:3.36]{3.36} to 3.38 that there are only two ways depends on this convention. A reader who counts such cancellation separately obtains a third construction, in which correlations of single step factors with the block cancel in the aggregate, and Proposition \hyperref[res:Proposition-44]{44} characterizes by (3.174) exactly when the cross term has no jumps; consistency then needs, in addition, the affine condition on the block residual plus the cross term. With one step factor it is Proposition \hyperref[res:Proposition-43]{43}(a). \cite{gellert2021short} §3.1.3 decomposes correlated deterministic factors into independent ones; (3.174) is the condition that survives that decomposition at the scheduled maturities. Observability of block-diagonal pairs with disjoint spectra and minimal realizations are standard linear systems theory; the arguments are written out below.

\emph{Proof.} (a) The front-end volatility on \(I_m\) is \(\sum_\ell s_{\ell,m}H^S_\ell\) and the block volatility \(\mathcal C\,e^{\mathcal A(T-u)}H^Z\), so the cross term of \(\sigma\cdot\int\sigma\) is the sum over \(\ell\) of the cross terms of Proposition \hyperref[res:Proposition-43]{43} with \(w_\ell=H^Z(H^S_\ell)^\top\). The jump across \(T_m\) of the integrated cross term is therefore \(\mathcal C\,e^{\mathcal AT}(\mathcal A^{-1}+(T-T_m)I)v_m(t)-\mathcal C\mathcal A^{-1}\int_0^tW_m\,du\) with \(W_m=\sum_\ell\Delta_{\ell,m}w_\ell=H^ZG_m^\top\) and \(v_m(t)=\int_0^te^{-\mathcal Au}W_m(u)\,du\), and the argument of Proposition \hyperref[res:Proposition-43]{43}(a), with Lemma \hyperref[res:Lemma-42]{42} for \((\mathcal A,\mathcal C)\), gives \(W_m=0\) almost everywhere, almost surely. The converse is Proposition \hyperref[res:Proposition-43]{43}(b).

(b) The set \(\{y:\mathcal C\,e^{\mathcal Ax}y\equiv0\}\) is invariant under every polynomial in \(\mathcal A\). Let \(p_j\) be the characteristic polynomial of \(A_j\) and \(q_j=\prod_{i\neq j}p_i\); disjoint spectra make \(p_j\) and \(q_j\) coprime. By the Cayley--Hamilton theorem \(q_j(\mathcal A)\) annihilates every block but the \(j\)-th, and a Bézout identity \(\alpha p_j+\gamma q_j=1\) shows that \(q_j(\mathcal A)\) is invertible on the \(j\)-th block. So for \(y\) in the set, \((\gamma q_j)(\mathcal A)y\) lies in it and has \(y_j\) as its only nonzero block, whence \(c_je^{A_jx}y_j\equiv0\) and \(y_j=0\) by observability of \((A_j,c_j)\). The components of \(H^ZG_m^\top\) along \(z_j\) are \(H^Z_jG_m^\top\).

(c) The unobservable subspace \(N\) is \(\mathcal A\)-invariant, and \(\mathcal AN=N\) since \(\mathcal A\) is invertible, so \(\mathcal A\) induces an invertible \(\mathcal A'\) on \(\mathbb R^r/N\); with \(P\) the quotient map, \(P\mathcal A=\mathcal A'P\) passes through the exponential series and gives \(\mathcal C\,e^{\mathcal Ax}z=\mathcal C'e^{\mathcal A'x}Pz\). The curve is unchanged with \((\mathcal A',\mathcal C',PZ)\), and (a) applies.

(d) As in Proposition \hyperref[res:Proposition-43]{43}(c), with \(W_m\) in place of \(\Delta_mw\).\hfill$\square$

The proposition is conditional on the cited calculus and is proved on one path. Not covered: blocks with an exponent-zero part; unbounded scales and the Heston scale of \cite{brace2024sofr}; front ends with a diffusive ramp; and existence of the processes.

\subsection{Polynomial parts in the uncorrelated splice}\label{sec:3.39}

Keep the uncorrelated splice of Section \hyperref[sec:3.33]{3.33} with a front end of the type of Theorem \hyperref[res:Theorem-11]{11}, which absorbs exactly the part of the block residual that is affine in \(x\) (Proposition \hyperref[res:Proposition-38]{38}(a)). The block is an exponential-polynomial family in the sense of \cite{filipovic2000exponential} eq. 5 with a full polynomial part,

\begin{equation*}
F(x,z)=p_0(x)+\sum_{i=1}^Kp_i(x)\,e^{-\kappa_ix},\qquad p_0(x)=\sum_{\mu=0}^dz_{0,\mu}x^\mu, \tag{3.175}\label{eq:3.175}
\end{equation*}

where every coefficient \(z_{0,0},\dots,z_{0,d}\) is a parameter, the \(p_i\) are polynomials with parameter coefficients, and the exponents \(\kappa_i\) are fixed or are themselves parameters, positive at the point considered. \(Z\) is driven by \(W^B\) with \(a=\sigma_Z\sigma_Z^\top\), and \(R\) is the residual of the consistency equation; by \cite{filipovic2000exponential} Remark 8.4 it extends from \(x\le H-t\) to all \(x\ge0\) by analyticity. Alone means \(R=0\), in the splice means \(R\) affine. The statements hold at one parameter point.

\phantomsection\label{res:Proposition-45}
\textbf{Proposition 45.} Let \(m=\lfloor(d-1)/2\rfloor\).

(a) Let \(d\ge1\). Then \((a,b)\) makes \(R\) affine if and only if \((a,b')\) makes \(R=0\), where \(b'\) agrees with \(b\) except in \(b_{0,0}\) and \(b_{0,1}\). So the admissible covariances, and all drifts other than those two, are the same alone and in the splice: the front end changes none of the block's restrictions.

(b) For \(d\ge1\), alone and in the splice, \(a_{(0,\mu),(0,\nu)}=0\) unless \(\mu,\nu\le m\): at most \(\lfloor(d+1)/2\rfloor\) polynomial coefficients have nonzero diffusion.

(c) If the block is polynomial only (\(K=0\)), every nonnegative definite \(a\) supported on \(\{(0,\mu),(0,\nu):\mu,\nu\le m\}\) is admissible with suitable \(b\), alone and in the splice, so exactly \(\lfloor(d+1)/2\rfloor\) possibly correlated polynomial factors can occur: \(z_{0,0}\) for \(d=1,2\), \(z_{0,0}\) and \(z_{0,1}\) for \(d=3,4\), and so on.

(d) For \(d=0\), a pure level, \(R=0\) forces \(a_{(0,0),(0,0)}=0\), while in the splice the level is free, as in Proposition \hyperref[res:Proposition-38]{38}. This is the only case in which the front end changes the restrictions of a full polynomial part.

The case \(d=1\) of (c) is the nontrivial diffusion for \(z_0+z_1x\) noted in \cite{filipovic2000exponential} Remark 8.3. \cite{bjork1999interest} §7 shows that the full exponential-polynomial family with varying exponents is inconsistent under deterministic volatility and treats fixed exponents, but gives no count of polynomial factors; we found the count \(\lfloor(d+1)/2\rfloor\) stated in neither source nor in \cite{filipovic2001consistency}.

\emph{Proof.} Sort \(R\) by exponential factor. For an exponential parameter with exponent \(\kappa>0\), its derivatives carry \(e^{-\kappa x}\), and \(\int_0^x\) of such a derivative is a constant \(c_j\) minus \(e^{-\kappa x}\) times a polynomial; so terms involving only exponential parameters have positive exponents. The exponent-zero part of \(R\) comes from the polynomial part:

\begin{equation*}
\sum_{k\le d}\Bigl[b_{0,k}-(k+1)z_{0,k+1}-\sum_{j\text{ exp}}a_{(0,k),j}c_j\Bigr]x^k-\sum_{\mu,\nu\le d}a_{(0,\mu),(0,\nu)}\frac{x^{\mu+\nu+1}}{\nu+1}, \tag{3.176}\label{eq:3.176}
\end{equation*}

with \(z_{0,d+1}=0\), since second derivatives in \(z_0\) vanish and \(x^\mu\int_0^x\eta^\nu\,d\eta=x^{\mu+\nu+1}/(\nu+1)\).

(a) The parts with positive exponents are the same alone and in the splice, because the front end absorbs only polynomials and \(x^ke^{-\kappa x}\), \(\kappa>0\), are independent of the polynomials on every interval. In (3.176) the coefficient of \(x^k\), \(k\le d\), contains \(b_{0,k}\) with coefficient one, while for \(k\ge d+1\ge2\) it involves no drift. The front end absorbs exactly the coefficients of \(x^0\) and \(x^1\), which alone are set to zero by \(b_{0,0}\) and \(b_{0,1}\).

(b) Let \(\mu^*\) be the largest \(\mu\) with \(a_{(0,\mu),(0,\mu)}>0\); nonnegative definiteness kills every row beyond it, so the coefficient of \(x^{2\mu^*+1}\) in (3.176) is \(-a_{(0,\mu^*),(0,\mu^*)}/(\mu^*+1)\neq0\). If \(2\mu^*+1>d\) this coefficient has no drift and degree above one, a contradiction; so \(\mu^*\le m\).

(c) With \(K=0\) and \(a\) supported on \(\mu,\nu\le m\), every term of the double sum has degree at most \(2m+1\le d\), so each coefficient of (3.176) contains a drift that can be chosen to make it zero, or is left free for \(k\le1\) in the splice.

(d) For \(d=0\), (3.176) is \(b_{0,0}-\sum_ja_{(0,0),j}c_j-a_{(0,0),(0,0)}x\), whose coefficient of \(x\) has no drift; alone it vanishes, and in the splice it is absorbed.\hfill$\square$

The proposition is conditional on the cited calculus through the reduction of Proposition \hyperref[res:Proposition-38]{38}(a); the statements themselves are algebra at a parameter point. Not covered: polynomial parts that are not full, where a structurally absent coefficient of degree at most one removes a drift and the front end lifts a constraint; front ends with a diffusive ramp, which absorb polynomials of degree up to three; correlated splices with a polynomial part; and existence of \(Z\).

\subsection{Varying exponents in the uncorrelated splice}\label{sec:3.40}

Keep the uncorrelated splice (3.161) of Section \hyperref[sec:3.32]{3.32} with a general front end, \(L_m\) and \(C_m\) Itô processes driven by \(W^S\), which absorbs on each \(I_m\) exactly a polynomial in \(T\) of degree at most three. The block is the bounded exponential-polynomial family of \cite{filipovic2000exponential} Definition 3.1,

\begin{equation*}
F(x,z)=\sum_{i=1}^Kp_i(x,z)\,e^{-z_{i,n_i+1}x},\qquad p_i(x,z)=\sum_{\mu\le n_i}z_{i,\mu}x^\mu, \tag{3.177}\label{eq:3.177}
\end{equation*}

in which every exponent \(z_{i,n_i+1}\) is a parameter that may vary, and \(Z\), driven by \(W^B\), stays in the set of parameters with \(\sup_x|F(x,z)|<\infty\). Assumption 2.1 is assumed for the drift of the representation, as in Proposition \hyperref[res:Proposition-37]{37}. Hypothesis (P) says that \(dt\otimes dQ\)-almost every \((t,\omega)\) has every exponent \(Z^{i,n_i+1}_t\) positive.

\phantomsection\label{res:Proposition-46}
\textbf{Proposition 46 (conditional).} (a) At every point with all exponents positive at which the splice satisfies Assumption 2.1, \(Z\) satisfies the consistency equation of \cite{filipovic2000exponential} Proposition 2.2 for the block alone: the front end absorbs no part of it.

(b) Under (P), and conditional on the cited conclusions of \cite{filipovic2000exponential} Theorem 3.2 and Corollaries 3.4 and 3.5 for a process satisfying that equation almost everywhere, those conclusions hold for \(Z\) in the splice. In particular the diffusion and drift of the exponent \(Z^{i,n_i+1}\) vanish where \(p_i(Z)\neq0\), respectively where \(p_i(Z)\) and the polynomial of the paired exponent do not vanish, so the exponents are constant there, and under the conditions of Corollary 3.5 the block is quasi-deterministic.

So the exponent restrictions of \cite{filipovic2000exponential} for the block alone hold unchanged in the uncorrelated splice. Proposition \hyperref[res:Proposition-37]{37} is the case of fixed exponents. Hypothesis (P) also requires positivity where \(p_i(Z)=0\), where the source lets the exponent be an arbitrary Itô process (its Lemma 5.1), so (b) covers a narrower class of processes than Theorem 3.2 does.

\emph{Proof.} (a) As in the proof of Proposition \hyperref[res:Proposition-37]{37}, Assumption 2.1 gives \(R(t,x)=-(D^S-\sigma^S\cdot\int\sigma^S)(t,t+x)\), a polynomial of degree at most three in \(x\) on each \(I_m\cap[t,H]\), where \(R\) is the residual of the consistency equation, with the parameter derivatives of \(F\) including those in the exponents. At a point with exponents \(\beta_i>0\), every term of \(R\) carries a factor \(e^{-\beta_ix}\) or \(e^{-(\beta_i+\beta_j)x}\), since \(\int_0^x\eta^\mu e^{-\beta\eta}\,d\eta=\mu!/\beta^{\mu+1}-e^{-\beta x}\)(polynomial). So \(R(t,\cdot)\) is a finite sum of \(x^ke^{-\kappa x}\) with \(\kappa>0\), independent of the polynomials on every interval; it vanishes on \([0,H-t]\), hence for all \(x\ge0\) by analyticity. (b) By (a) and (P) the consistency equation holds \(dt\otimes dQ\)-almost everywhere, which is the hypothesis under which the cited results are stated; their proofs use the equation at fixed points, the support condition and lemmas about the Itô process \(Z\) alone, none of which involves the front end.\hfill$\square$

The proposition is conditional on the cited calculus and, for (b), on the cited results of \cite{filipovic2000exponential}. Not covered: points where an exponent vanishes, such as the Nelson--Siegel and Svensson level, where the front end can absorb part of an exponent-zero group (Propositions \hyperref[res:Proposition-38]{38}, \hyperref[res:Proposition-40]{40} and \hyperref[res:Proposition-45]{45} treat fixed exponents); correlated splices with varying exponents; and existence of \(Z\).

\subsection{Exponent-zero block parts in the correlated splice}\label{sec:3.41}

Keep the driver-form setting of Section \hyperref[sec:3.37]{3.37}, with one step factor, and let the block have a full polynomial part and an observable quasi-exponential part,

\begin{equation*}
F(x,z)=\sum_{\mu=0}^dz_{0,\mu}x^\mu+\mathcal C\,e^{\mathcal Ax}\zeta, \tag{3.178}\label{eq:3.178}
\end{equation*}

with \(\mathcal A\) invertible and \((\mathcal A,\mathcal C)\) observable. \(Z=(z_0,\zeta)\) is an Itô process in driver form with predictable bounded integrand rows \(H^{0,\mu}\) for \(z_{0,\mu}\) and \(H^\zeta\) for \(\zeta\), and the covariation densities with the step integrand are

\begin{equation*}
w_\mu(u)=H^{0,\mu}(u)H^S(u)^\top\in\mathbb R,\qquad w_\zeta(u)=H^\zeta(u)H^S(u)^\top\in\mathbb R^r. \tag{3.179}\label{eq:3.179}
\end{equation*}

Assumption 2.1 is assumed for the drift of the representation, as in Proposition \hyperref[res:Proposition-43]{43}.

\phantomsection\label{res:Proposition-47}
\textbf{Proposition 47 (conditional).} (a) Almost surely, for every \(m\) and almost every \(u<T_m\), \(\Delta_m(u)w_\mu(u)=0\) for \(\mu=1,\dots,d\) and \(\Delta_m(u)w_\zeta(u)=0\). No condition is imposed on \(w_0\), the covariation of the level.

(b) Under these conditions the jumps of the cross term across the meetings are affine in \(T\) and come only from the level, and the front end absorbs them. The curve satisfies Assumption 2.1 if and only if the block residual plus the jump-free cross terms is affine in the time to maturity, as in Propositions \hyperref[res:Proposition-38]{38}(a) and 45; the bound and the admissible covariances of Proposition \hyperref[res:Proposition-45]{45} carry over, with only the drifts changed.

(c) Every block direction outside the constants, that is every polynomial coefficient of degree at least one and every exponential part, must be uncorrelated with the step noise wherever the step volatility jumps, or no choice of drifts satisfies Assumption 2.1.

The level is a constant in calendar maturity, a direction the front-end family already contains (Proposition \hyperref[res:Proposition-38]{38}(c)), so its correlation with the step noise is internal to the front end. We count a correlated level as part of the front end rather than as a third construction; this is a classification convention, not a result. A reader who counts it separately obtains a third construction, which Proposition \hyperref[res:Proposition-47]{47} characterizes exactly: only the level may correlate with the step noise where the step volatility jumps. With this convention the answer is that of Sections \hyperref[sec:3.30]{3.30} to 3.38, the overlap of the two families being exempt. \cite{brace2024sofr} has no level component correlated with meeting steps, and \cite{gellert2021short} §3 splices by independence.

\emph{Proof.} The block volatility is \(\sum_\mu(T-u)^\mu H^{0,\mu}(u)+\mathcal C\,e^{\mathcal A(T-u)}H^\zeta(u)\), so the cross part of \(\sigma\cdot\int\sigma\) splits by component: for \(z_{0,\mu}\) it is \(w_\mu(u)[s_m(u)(T-u)^{\mu+1}/(\mu+1)+(T-u)^\mu(s_m(u)T+\kappa_m(u))]\), and the exponential part gives the cross term of Proposition \hyperref[res:Proposition-43]{43} with \(w_\zeta\). As in Proposition \hyperref[res:Proposition-43]{43}(a), the jump across \(T_m\) of the integrated cross term must be affine in \(T\). With \(\kappa_m-\kappa_{m-1}=-\Delta_mT_m\) it is \(J=\sum_\mu J_\mu+J_\zeta\), where \(J_\mu(T)=\int_0^tw_\mu\Delta_m[(T-u)^{\mu+1}/(\mu+1)+(T-u)^\mu(T-T_m)]\,du\) is a polynomial of degree at most \(\mu+1\) with leading coefficient \(\frac{\mu+2}{\mu+1}\int_0^tw_\mu\Delta_m\,du\), and \(J_\zeta\) is the jump of Proposition \hyperref[res:Proposition-43]{43}. A quasi-exponential with nonzero exponents that equals a polynomial on an interval vanishes, so the exponential part of \(J_\zeta\) vanishes and Lemma \hyperref[res:Lemma-42]{42} gives \(\Delta_mw_\zeta=0\). For \(d\ge1\) the coefficient of \(T^{d+1}\) comes from \(J_d\) alone and must vanish, so \(\int_0^tw_d\Delta_m=0\) for rational and then all \(t<T_m\), and \(w_d\Delta_m=0\) almost everywhere; descending through \(\mu=d-1,\dots,1\), whose leading degree is at least two, gives (a). \(J_0(T)=\int_0^tw_0\Delta_m(2T-u-T_m)\,du\) is affine for every \(w_0\). For (b), \(J=J_0\) is absorbed by the front end, and the remaining cross terms form one analytic function joining the block residual. For the bound, \(w_\mu\neq0\) needs \(a_{(0,\mu),(0,\mu)}\neq0\) by Cauchy--Schwarz, so the extra cross terms have degree at most \(\mu^*+1\), below the top degree \(2\mu^*+1\) of Proposition \hyperref[res:Proposition-45]{45}(b) when \(\mu^*\ge1\). Part (c) is (a) with (b).\hfill$\square$

The proposition is conditional on the cited calculus and is proved on one path. Not covered: polynomial parts that are not full; several step factors, treated in Proposition \hyperref[res:Proposition-52]{52}; front ends with a diffusive ramp; varying exponents in the correlated splice; and existence of the processes.

\subsection{The compensating shape of a scheduled jump}\label{sec:3.42}

Proposition \hyperref[res:Proposition-9]{9} treats shapes that are affine on each interval. This section removes that restriction on the interval after one scheduled date. Let \((\Omega,\mathcal F,P)\) be a probability space, \(\mathcal G\subseteq\mathcal F\) a sub-sigma-algebra and \(L\in(0,\infty]\). Let \(X\) be a real random variable, the level jump, and \(h:[0,L)\times\Omega\to\mathbb R\) a \(\mathcal B([0,L))\otimes\mathcal G\)-measurable function, the compensating shape, with \(\int_0^\tau|h(u,\omega)|\,du<\infty\) for every \(\tau<L\) and every \(\omega\). The curve jump at maturity \(T_n+u\) and its primitive are

\begin{equation*}
\xi(u)=X+h(u),\qquad H(\tau)=\int_0^\tau h(u)\,du. \tag{3.180}\label{eq:3.180}
\end{equation*}

Let \(\kappa\) be a regular conditional distribution of \(X\) given \(\mathcal G\), and put

\begin{equation*}
M(\omega,\tau)=\int e^{-\tau x}\,\kappa(\omega,dx)\in(0,\infty],\qquad K(\omega,\tau)=\log M(\omega,\tau); \tag{3.181}\label{eq:3.181}
\end{equation*}

where \(M(\omega,\tau)<\infty\), let \(\tilde E_\tau[X]\) and \(\tilde V_\tau[X]\) be the mean and variance of the tilted law \(e^{-\tau x}\kappa(\omega,dx)/M(\omega,\tau)\). Hypothesis (H) is Assumption 2.2 for this jump, with \(\mathcal G=\mathcal F_{T_n-}\) and \(L\) the distance to the next scheduled date:

\begin{equation*}
E\Bigl[\exp\Bigl(-\int_0^\tau\xi(u)\,du\Bigr)\Bigm|\mathcal G\Bigr]=1\quad\text{a.s., for every }\tau\in[0,L). \tag{3.182}\label{eq:3.182}
\end{equation*}

\phantomsection\label{res:Proposition-48}
\textbf{Proposition 48.} (a) (H) holds if and only if there is one \(\mathcal G\)-measurable null set off which \(M(\omega,\tau)<\infty\) and \(H(\omega,\tau)=K(\omega,\tau)\) for every \(\tau\in[0,L)\).

(b) Under (H), off that null set, \(K(\omega,\cdot)\) is infinitely differentiable on \((0,L)\) with \(K'(\tau)=-\tilde E_\tau[X]\) and \(K''(\tau)=\tilde V_\tau[X]\ge0\), and \(K''\) vanishes somewhere if and only if \(\kappa(\omega,\cdot)\) is a Dirac mass. The shape is fixed by the law of the level jump:

\begin{equation*}
h(\omega,\tau)=-\tilde E_\tau[X](\omega)\quad\text{for almost every }\tau\in(0,L)\text{ and every right-continuity point }\tau, \tag{3.183}\label{eq:3.183}
\end{equation*}

so \(h(\omega,\cdot)\) agrees almost everywhere with a nondecreasing function, strictly increasing unless the law is a Dirac mass; if \(h(\omega,\cdot)\) is right-continuous at \(0\), then \(X\) is conditionally integrable and \(h(\omega,0)=-\int x\,\kappa(\omega,dx)\). Conversely, (H) holds if \(M(\omega,\cdot)<\infty\) on \([0,L)\) and (3.183) holds almost everywhere, for almost every \(\omega\). Two level jumps satisfying (H) with the same \(\mathcal G\) and \(h\) have the same conditional law, and every conditional law with finite Laplace transform on \([0,L)\) has a measurable locally integrable shape satisfying (H).

(c) (i) An affine shape \(h(u)=c+yu\) with \(\mathcal G\)-measurable \(c,y\) satisfies (H) if and only if, almost surely, \(y\ge0\) and \(\kappa(\omega,\cdot)=N(-c,y)\); with \(c=0\) this is the first interval of Proposition \hyperref[res:Proposition-9]{9}, and with \(y=0\) the one-date form of Lemma \hyperref[res:Lemma-5]{5}. (ii) For trivial \(\mathcal G\), \(L=\infty\) and \(X=\Delta\neq0\) with probability \(p\in(0,1)\) and \(0\) otherwise, the continuous shape is

\begin{equation*}
h(\tau)=-\Delta q(\tau),\qquad q(\tau)=\frac{pe^{-\tau\Delta}}{1-p+pe^{-\tau\Delta}}, \tag{3.184}\label{eq:3.184}
\end{equation*}

with \(h'=\Delta^2q(1-q)>0\); it is strictly concave on \([0,\infty)\) if and only if \(\Delta(1-2p)\ge0\), and otherwise convex before and concave after \(\tau^*=\log(p/(1-p))/\Delta\), and never affine. The Gaussian level jump with the same mean and variance has the ramp \(h(0)+h'(0)\tau\), the tangent of \(h\) at \(0\); for \(p=\tfrac12\) the gap \(g=h(0)+h'(0)\tau-h(\tau)\) satisfies

\begin{equation*}
0\le g(\tau)\le\frac{|\Delta|^4\tau^3}{48},\qquad 0\le\frac1\tau\int_0^\tau g\le\frac{|\Delta|^4\tau^3}{192}. \tag{3.185}\label{eq:3.185}
\end{equation*}

(d) Let \(t<T_n\), \(\mathcal G=\mathcal F_{T_n-}\supseteq\mathcal F_t\), \(r_t\) \(\mathcal F_t\)-measurable, \(J\) \(\mathcal F_{T_n}\)-measurable, and \(f(t,\cdot)\) a \(\mathcal B\otimes\mathcal F_t\)-measurable locally integrable curve with \(f(t,T)=r_t\) for \(T\in[t,T_n)\). In the model in which the curve is frozen off \(T_n\) and jumps there to \(r_t+J\) on \([T_n,T_n+L)\), (H) holds if and only if, almost surely, \(E[e^{-\tau J}\mid\mathcal F_{T_n-}]<\infty\) for \(\tau\in[0,L)\) and

\begin{equation*}
f(t,T)=r_t+\frac{E_{T_n-}[Je^{-\tau J}]}{E_{T_n-}[e^{-\tau J}]},\qquad\tau=T-T_n,\ \text{for almost every }\tau\in(0,L), \tag{3.186}\label{eq:3.186}
\end{equation*}

with the expectations computed from the regular conditional distribution of \(J\) given \(\mathcal F_{T_n-}\). Under (H), the conditional laws of \(J\) given \(\mathcal F_{T_n-}\) and given \(\mathcal F_t\) agree, so (3.186) is (5.1) as printed; if \(f(t,\cdot)\) is right-continuous it holds for every \(\tau\in[0,L)\), at \(\tau=0\) as \(f(t,T_n)=r_t+E_t[J]\) with \(E_t|J|<\infty\); and \(P(t,T)=e^{-r_t(T-t)}E_t[e^{-(T-T_n)J}]\) for \(T\in[T_n,T_n+L)\). For a general initial curve, \(f(t,T)\) is the tilted conditional expectation of the post-meeting short rate \(r_T\), and (5.1) holds for every \(\tau\) if and only if \(r_{T_n-}=r_t\) and the short rate is constant after \(T_n\).

(e) On a later interval, if \(E[e^{-S}\mid\mathcal G]=1\) for the exponent \(S\) accumulated before it, the shape there is the derivative of the conditional cumulant generating function of its level jump under the measure \(e^{-S}\cdot P\), which is the mechanism of Proposition \hyperref[res:Proposition-9]{9} on later intervals.

The Gaussian case of (c)(i) is the compensator of \cite{fontana2022term} Example 3.10, and the Gaussian case of the bond formula in (d) is \cite{kim2014jumps} Proposition 1; the characterization with one null set for general shapes, the uniqueness of the law and the exact form and range of (5.1) are ours. The identity of a forward rate with a forward-measure expectation of the short rate is textbook, and (5.1) is its specialization. As a numerical check, not part of the statement, the gap in (3.185) at \(\Delta=25\) basis points is of order \(10^{-11}\) basis points at six weeks and \(10^{-8}\) at one year, against about \(12.5\) basis points for \(h\) itself, so a two-point meeting law and its Gaussian ramp are indistinguishable at pricing precision, as Section \hyperref[sec:3.9]{3.9} found for the ramps.

\emph{Proof.} For fixed \(\tau\), \(\exp(-\int_0^\tau\xi)=e^{-\tau X}e^{-H(\tau)}\) with \(e^{-H(\tau)}\) positive and \(\mathcal G\)-measurable, so (3.182) at \(\tau\) is \(M(\cdot,\tau)=e^{H(\cdot,\tau)}\) almost surely. (a) Take the union of the exceptional sets over rational \(\tau\). Off it \(M(\omega,\cdot)\) is finite at the rationals, hence by convexity finite and continuous on \((0,L)\), and \(H(\omega,\cdot)\) is continuous, so the identity extends to every \(\tau\). (b) Domination of \(|x|^je^{-sx}\) by exponentials at nearby arguments gives \(M\in C^\infty((0,L))\) with \(M^{(j)}(s)=\int(-x)^je^{-sx}\kappa(dx)\), hence \(K'=-\tilde E_\tau[X]\) and \(K''=\tilde V_\tau[X]\); the tilted law and \(\kappa\) are equivalent, so \(K''\) vanishes only for a Dirac mass. Since \(H=K\), the Lebesgue differentiation theorem gives (3.183) almost everywhere, and right-continuity gives it at the point; at \(0\), Fatou and dominated convergence give integrability and \(h(0)=-\int x\,\kappa\). The converse integrates \(K'\). For uniqueness, tilt both laws at some \(\tau_0\in(0,L)\): the tilted Laplace transforms are finite and equal on an open interval containing \(0\), so the laws agree (\cite{curtiss1942note}). For existence, \(K'\) is jointly measurable and locally integrable. (c)(i) \(K(\tau)=c\tau+\tfrac12y\tau^2\) forces \(y=K''\ge0\), and the tilted law has the Laplace transform of a Gaussian, so by uniqueness \(\kappa=N(-c,y)\); conversely \(N(-c,y)\) has \(K=H\). (ii) Differentiate \(K=\log(1-p+pe^{-\tau\Delta})\); the sign of \(h''=-\Delta^3q(1-q)(1-2q)\) follows the monotone \(q\), and for \(p=\tfrac12\), \(g=(|\Delta|/2)(z-\tanh z)\) with \(z=|\Delta|\tau/2\) and \(0\le z-\tanh z\le z^3/3\). (d) The jump is \(\xi(T)=r_t+J-f(t,T)\), which is (3.180) with \(X=J\) and \(h(u)=r_t-f(t,T_n+u)\), so (a) and (b) give (3.186). Under (H), \(E[e^{-\tau J}\mid\mathcal F_t]=e^{H(\tau)}\) by the tower property, and uniqueness identifies the two conditional laws; right-continuity and the statement at \(0\) are (b); the bond price is \(\exp(-r_t\tau+H(\tau))\) times the pre-meeting discount. For a general initial curve the tilt by \(\exp(-\int_{T_n}^Tr_s\,ds)\) is \(e^{-\tau X}\) times a positive \(\mathcal G\)-measurable factor, which gives the tilted expectation, and comparing it with (5.1) at every \(\tau\) gives the stated equivalence. (e) The measure \(e^{-S}\cdot P\) agrees with \(P\) on \(\mathcal G\), and the abstract Bayes formula reduces the condition to (a) under it.\hfill$\square$

The proposition treats one scheduled date. Its model in (d) has zero drift and volatility off \(T_n\); the assembly into a model satisfying Assumptions 2.1 and 2.2 on all of \([0,\infty)\), the induction over several meetings for general shapes, and a characterization of which functions are cumulant generating functions are not covered. The implication in (d) fails if the laws are taken given \(\mathcal F_t\) while information about \(J\) arrives on \((t,T_n)\).

\subsection{Meeting variances with a diffusion between meetings}\label{sec:3.43}

The identification results of Sections \hyperref[sec:3.13]{3.13} to 3.21 use a model in which the curve moves only at meetings. This section adds a parallel Gaussian diffusion between meetings and an initial curve. Keep the jumps \(Z_1,\dots,Z_N\sim N(0,v_n)\) of Section \hyperref[sec:3.10]{3.10}, \(N\ge3\), and let \(W\) be an independent Brownian motion, with the usual augmentation of the filtration they generate at their revelation times. Fix a partition \(0=c_0<c_1<\dots<c_P\), variances \(u_p\ge0\) with \(\sigma(s)=\sqrt{u_p}\) on \([c_{p-1},c_p)\) and \(\sigma=0\) after \(c_P\), and a bounded piecewise continuous initial curve. With \(Y_t=\int_0^t\sigma\,dW\), the curve is

\begin{equation*}
f(t,T)=f(0,T)+\sum_{T_n\le t}\bigl[Z_n+v_n(T-T_n)\bigr]+Y_t+\int_0^t\sigma(s)^2(T-s)\,ds, \tag{3.187}\label{eq:3.187}
\end{equation*}

and the unknown parameter is \(\theta=(v,u)\in[0,\infty)^{N+P}\). The contracts are the ideal futures and European cash calls of Section \hyperref[sec:3.16]{3.16}. With \(w(s)=(b-s)^+-(a-s)^+\), \(d(s)=((b-s)^{+2}-(a-s)^{+2})/2\), \(h(s)=d(s)+w(s)^2/2\), \(A=\int_a^bf(0,u)\,du\) and the coefficients of (3.71), put

\begin{equation*}
p=\sum_ih_iv_i+\int_0^b\sigma^2h,\qquad z=\sum_ij_iv_i+\int_0^S\sigma(s)^2w(s)(S-s)\,ds,\qquad q=\sum_ik_iv_i+\int_0^S\sigma^2w^2,\qquad m=e^{A+p-z}. \tag{3.188}\label{eq:3.188}
\end{equation*}

For a panel of complete surfaces at expiries \(S_1<\dots<S_L\) before their accrual windows, let \(G_\ell=(S_{\ell-1},S_\ell]\), \(S_0=0\), be the gaps, \(E\) the set of meeting-free gaps, \(\lambda_p(J)\) the length of \(J\cap[c_{p-1},c_p)\), and \(Q(S)=\sum_{T_i\le S}v_i+\int_0^S\sigma^2\).

\phantomsection\label{res:Proposition-49}
\textbf{Proposition 49 (conditional).} (a) The curve (3.187) satisfies the regularity conditions and Assumptions 2.1 and 2.2, and every discounted bond is a true martingale, the product of \(P(0,T)\), a Gaussian stochastic exponential and the discounted bond of the model of Section \hyperref[sec:3.10]{3.10}.

(b) \(G(a,b)\) is an integrable martingale with \(G_0=e^{A+p}\), the call price is \(P(0,S)\) times the Black formula (3.72) with forward \(m\) and variance \(q\), and the complete surface at one expiry determines exactly \(P(0,S)\), \(m\) and \(q\). For \(S<a\),

\begin{equation*}
q(S)=\delta^2Q(S),\qquad z(S)=\delta\int_0^SQ(x)\,dx, \tag{3.189}\label{eq:3.189}
\end{equation*}

so the initial curve does not enter \(q\) or \(z\); the futures quote adds \(z=\log G_0-\log m\), and \(p\) only when the bond curve is known.

(c) The panel \((Q(S_\ell))_\ell\) equals \(\mathcal A\theta\) with \(\mathcal A_{\ell,i}=\mathbf 1\{T_i\le S_\ell\}\) and \(\mathcal A_{\ell,N+p}=\lambda_p((0,S_\ell])\). It identifies \(\theta\) on \([0,\infty)^{N+P}\) if and only if \(\operatorname{rank}\mathcal A=N+P\), and otherwise every strictly positive \(\theta\) has a distinct strictly positive \(\theta'\) with the same panel. Moreover

\begin{equation*}
\operatorname{rank}\mathcal A=N+P\iff T_N\le S_L,\ \text{no gap contains two meetings, and}\ \operatorname{rank}\bigl(\lambda_p(G_\ell)\bigr)_{\ell\in E,\,p\le P}=P, \tag{3.190}\label{eq:3.190}
\end{equation*}

so identification needs at least \(P\) meeting-free gaps; the diffusion is read off the meeting-free gaps and each meeting variance off its own gap.

(d) On the listed calendar of Proposition \hyperref[res:Proposition-24]{24}, under its assumptions (A1) to (A4), exactly eight of the 24 gaps are meeting-free. With a constant diffusion variance (\(P=1\)) the sixteen meeting variances and \(u_1\) are identified. No identifying partition has more than eight cells, and eight are attained, for example with cell boundaries at the meetings of 28 January, 18 March, 17 June and 16 September 2026 and 27 January, 17 March and 28 July 2027. A partition by calendar quarter (rank 23 of 24) and a partition at every meeting (rank 24 of 33) do not identify, and explicit pairs of parameter vectors with equal panels exhibit this; with the futures quotes the stacked matrix of \((Q,z)\) has full column rank for both, and even with one cell per gap.

(e) With \(P=1\) on that calendar, whose first gap is meeting-free, an error \(\varepsilon\) in the at-the-money premium \(C/\delta\) perturbs \(Q(S)\) to first order by

\begin{equation*}
\eta(S)=\frac{2\varepsilon\,\sigma_\Sigma(S)}{P(0,S)\,m\,\varphi(\delta\sigma_\Sigma(S)/2)},\qquad\sigma_\Sigma(S)^2=Q(S), \tag{3.191}\label{eq:3.191}
\end{equation*}

and the estimator \(\hat u=Q(S_1)/\lambda_1\), \(\hat v_i=Q(S_\ell)-Q(S_{\ell-1})-\lambda_\ell\hat u\) for \(T_i\in G_\ell\), with \(\lambda_\ell=|G_\ell|\), has the sharp first-order error bound

\begin{equation*}
|\Delta\hat v_i|\le\eta_\ell+\eta_{\ell-1}+\frac{\lambda_\ell}{\lambda_1}\eta_1,\qquad\eta_\ell=\eta(S_\ell),\ \eta_0=0. \tag{3.192}\label{eq:3.192}
\end{equation*}

The model is that of Remark 12 with a diffusion, the contracts those of Proposition \hyperref[res:Proposition-19]{19} and the calendar that of Proposition \hyperref[res:Proposition-24]{24}; the form \(q=\delta^2(\sum v+\int\sigma^2)\) is the standard Gaussian computation (\cite{kim2014jumps} Proposition 1 has its bond-price analogue). \cite{wright2020event} separates event from background variance in one-day options by learning the background where there is no event and subtracting it where there is one, the statistical counterpart of (3.190), without an identification statement for a fixed calendar of expiries. The failing partitions in (d) do not fail because each meeting trades against the diffusion over its own gap: a cell whose diffusion enters no meeting-free gap trades against the meetings of every gap it overlaps. The cell \([18\text{ Mar},29\text{ Apr }2026)\) overlaps the gap of the 18 March meeting by 23 days and that of 29 April by 19, so the meeting-date pair moves both meetings and \(u\) on that cell, with every other variance fixed; the quarterly pair moves the last three meetings of 2027 and \(u\) on the last quarter. The futures-assisted identification of (d) is formal only: the meeting-date pair changes every \(z(S_\ell)\) by at most about \(1.9\times10^{-9}\), while an error of half a basis point in a futures rate moves \(\log G_0\) by about \(1.3\times10^{-5}\). As a numerical check, not part of the statement, the number of meetings practically identified in the sense of Proposition \hyperref[res:Proposition-24]{24}, (3.192) below \(s_0^2\), always the first ones in date order, is, with zero initial curve:

{\scriptsize
\begin{longtable}{lllllll}
\toprule
\textbf{\(\sigma\) (bp per year)} & \textbf{\(s_0=10\), \(\varepsilon=0.25\)} & \textbf{10, 0.5} & \textbf{15, 0.25} & \textbf{15, 0.5} & \textbf{25, 0.25} & \textbf{25, 0.5} \\
\midrule
\endhead
\bottomrule
\endlastfoot
none & 16 & 4 & 16 & 9 & 16 & 16 \\
50 & 2 & 0 & 12 & 2 & 16 & 13 \\
100 & 0 & 0 & 3 & 0 & 16 & 5 \\
\end{longtable}}

Two effects combine: the accumulated standard deviation includes the diffusion, which dominates the meeting variances at two years, and the diffusion variance is learned only from meeting-free gaps, best from the first two-week gap, with its error entering each meeting weighted by that meeting's gap length. In every case of the table, (3.192) coincides with the smallest worst-case bound over all linear estimators exact on the model, computed by linear programming, which is also a check. With a nonzero initial curve, \(\eta\) is multiplied by \(1/(P(0,S)m)\), a few percent away from one at current rate levels.

\emph{Proof.} (a) The drift \(\sigma(s)^2(T-s)\) is the differentiated HJM drift, so (2.3) holds at every \((\omega,t)\), and the jump condition is the computation of Remark 12, since \(Z_n\) is independent of \(\mathcal F_{T_n-}\). Itô's product rule for the deterministic integrand \(T-s\) and Fubini split \(\int_0^tr+\int_t^Tf(t,\cdot)\) into the initial-curve part, the jump part of Remark 12 and \(\int_0^t(T-s)\sigma\,dW+\tfrac12\int_0^t\sigma^2(T-s)^2\); the Gaussian factor is a true martingale, and the two factors have independent increments. (b) Stochastic Fubini gives \(\int_a^bY=\int_0^b\sigma w\,dW\), and conditioning on \(\mathcal F_t\) leaves independent Gaussians, so \(\log G_t\) is explicit and \(G_0=e^{A+p}\). Under the discount measure \(\log G_S\) is Gaussian with variance \(q\) and mean shifted by the covariance \(-z\) with \(-\int_0^Sr\), which gives the Black formula; the surface returns \(P(0,S)\) and \(P(0,S)m\) as limits at zero strike and \(q\) at the strike \(m\), as in Proposition \hyperref[res:Proposition-19]{19}. For \(S<a\), \(w=\delta\) on \([0,S]\), which is (3.189). (c) Differencing the rows gives, on a gap with no meeting, \(\sum_pu_p\lambda_p(G_\ell)\), and on a gap with one meeting \(v_i\) plus that sum; a meeting after \(S_L\) has a zero column, two meetings in one gap give a kernel vector, and a kernel vector of the meeting-free block extends to one of \(\mathcal A\), while conversely the rank condition on the meeting-free block and the meeting rows force a kernel vector to vanish. (d) is exact computation on the dates of Proposition \hyperref[res:Proposition-24]{24}. (e) The at-the-money price is \(P(0,S)m(2\Phi(\sqrt q/2)-1)\), and differentiating in \(q=\delta^2Q\) gives (3.191); the error of \(\hat v_i\) is a combination of the errors in \(Q(S_\ell)\), \(Q(S_{\ell-1})\) and \(Q(S_1)\) with coefficients \(1\), \(-1\) and \(-\lambda_\ell/\lambda_1\), the last two merging for the first meeting, which lies in \(G_2\), and its worst case, attained with aligned signs, is (3.192).\hfill$\square$

Parts (a) and (b) are conditional on the Wiener-integral facts the proof uses: with deterministic integrands, Wiener integrals are linear, jointly Gaussian and independent of the jumps, \(Y\) is continuous and satisfies the stochastic Fubini identity, and the filtration has the stated independence properties; no Brownian motion is constructed. Parts (c) to (e) are unconditional; the table, the linear-programming comparison and the size of the \(z\)-difference are numerical checks. The proposition does not treat an exponentially decaying diffusion factor, under which the diffusion's contribution to \(q\) is reweighted and (3.189) changes. Part (e) is first order in \(\varepsilon\), and practical identification is the convention of Proposition \hyperref[res:Proposition-24]{24}, not a statistical statement; nonlinear estimators are not considered. The listed contracts' settlement, exercise and daily compounding are outside the model, as in Propositions \hyperref[res:Proposition-19]{19} and \hyperref[res:Proposition-24]{24}.

\subsection{A maturity-dependent diffusion between meetings}\label{sec:3.44}

The diffusion of Section \hyperref[sec:3.43]{3.43} moves every forward rate by the same amount. Keep that section's setting, notation and contracts, with one change: the diffusion volatility is \(\sigma(s)\phi(s,T)\) for \(T\ge s\), where \(\sigma\) is as before and the shape \(\phi\) is a known bounded Borel function on \(\{0\le s\le T\}\). With \(\Phi(s,T)=\int_s^T\phi(s,x)\,dx\) and the differentiated drift, the curve is

\begin{equation*}
f(t,T)=f(0,T)+\sum_{T_n\le t}\bigl[Z_n+v_n(T-T_n)\bigr]+\int_0^t\sigma(s)\phi(s,T)\,dW_s+\int_0^t\sigma(s)^2\phi(s,T)\Phi(s,T)\,ds, \tag{3.193}\label{eq:3.193}
\end{equation*}

and \(\phi\equiv1\) is (3.187). The meeting shocks stay parallel. Two shapes are treated in detail: (S1) exponential decay, \(\phi(s,T)=e^{-\kappa(T-s)}\) with \(\kappa>0\); and (S2) no loading before the next meeting, \(\phi(s,T)=\mathbf 1\{T\ge T_{j(s)+1}\}\) with \(T_{N+1}=\infty\), the loading of the first meeting-indexed component of \cite{brace2024sofr} §3.1, eqs. (4) to (7). For a window \([a,b]\) and \(s<b\) put \(\tilde w(s)=\int_{a\vee s}^b\phi(s,x)\,dx\), \(\tilde d(s)=(\Phi(s,b)^2-\Phi(s,a\vee s)^2)/2\), \(\tilde h=\tilde d+\tilde w^2/2\), the window weight \(\beta(s)=\tilde w(s)/\delta\) for \(s<a\), and, with the coefficients of (3.71),

\begin{equation*}
\tilde p=\sum_ih_iv_i+\int_0^b\sigma^2\tilde h,\qquad\tilde z=\sum_ij_iv_i+\int_0^S\sigma(s)^2\Phi(s,S)\tilde w(s)\,ds,\qquad q=\sum_ik_iv_i+\int_0^S\sigma^2\tilde w^2,\qquad m=e^{A+\tilde p-\tilde z}. \tag{3.194}\label{eq:3.194}
\end{equation*}

For a panel of expiries \(S_\ell\) with their own windows \([a_\ell,b_\ell]\), \(S_\ell<a_\ell\), and weights \(\beta_\ell\), put \(D_{\ell,p}=\int_{(0,S_\ell]\cap[c_{p-1},c_p)}\beta_\ell^2\), \(D_{0,p}=0\) and \(\tilde\Lambda_{\ell,p}=D_{\ell,p}-D_{\ell-1,p}\).

\phantomsection\label{res:Proposition-50}
\textbf{Proposition 50 (conditional).} (a) For every bounded Borel shape, the curve (3.193) satisfies the regularity conditions and Assumptions 2.1 and 2.2, the process \(s\mapsto\int_0^s\sigma(v)\phi(v,s)\,dW_v\) has a jointly measurable version, and with it every discounted bond is the true martingale
\(P(0,T)\exp\bigl(-\int_0^t\sigma\Phi(\cdot,T)\,dW-\tfrac12\int_0^t\sigma^2\Phi(\cdot,T)^2\bigr)\) times the discounted bond of Section \hyperref[sec:3.10]{3.10}.

(b) Proposition \hyperref[res:Proposition-49]{49}(b) holds with \(\tilde p\), \(\tilde z\), \(q\) and \(m\) of (3.194), and for \(S<a\)

\begin{equation*}
Q(S):=\frac{q(S)}{\delta^2}=\sum_{T_i\le S}v_i+\int_0^S\sigma^2\beta^2,\qquad\tilde z(S)=\delta\Bigl[\sum_{T_i\le S}(S-T_i)v_i+\int_0^S\sigma(s)^2\Phi(s,S)\beta(s)\,ds\Bigr]. \tag{3.195}\label{eq:3.195}
\end{equation*}

(c) The panel \((Q(S_\ell))_\ell\) is \(\mathcal A\theta\), with the meeting columns of Proposition \hyperref[res:Proposition-49]{49}(c) and the diffusion columns \(D_{\cdot,p}\). For any real diffusion columns,

\begin{equation*}
\operatorname{rank}\mathcal A=N+P\iff T_N\le S_L,\ \text{no gap contains two meetings, and}\ \operatorname{rank}\bigl(\tilde\Lambda_{\ell,p}\bigr)_{\ell\in E,\,p\le P}=P, \tag{3.196}\label{eq:3.196}
\end{equation*}

and every other conclusion of Proposition \hyperref[res:Proposition-49]{49}(c) holds with \(\tilde\Lambda\) in place of \(\lambda_p(G_\ell)\); in particular \(P\le|E|\) for every shape.

(d) On the listed calendar of Proposition \hyperref[res:Proposition-24]{24}, of the eight meeting-free gaps exactly two, \((11\text{ Dec }2026,15\text{ Jan }2027]\) and \((11\text{ Jun},16\text{ Jul }2027]\), lie between expiries on different futures.

(d1) With \(P=1\), every shape with \(D_{1,1}>0\) identifies the sixteen meeting variances and \(u_1\); this holds for (S1) at every \(\kappa>0\) and for (S2), where \(D_{1,1}=14/360\).

(d2) For every shape, the eight-cell partition \(\mathcal P_8\) of Proposition \hyperref[res:Proposition-49]{49}(d) gives a lower-triangular \(\tilde\Lambda_E\), whose six rows between expiries on the same future have only the diagonal entry \(\int_{G_\ell}\beta_\ell^2\); \(\mathcal P_8\) identifies if and only if the eight diagonal entries are nonzero.

(d3) Under (S2), \(\beta_\ell=1\) except on \((9\text{ Dec},11\text{ Dec }2026]\) and \((9\text{ Jun},11\text{ Jun }2027]\) for the expiries ending them, where \(\beta=49/91\), and on \((8\text{ Dec},10\text{ Dec }2027]\) for the last expiry, where \(\beta=0\). The matrix \(\tilde\Lambda_E(\mathcal P_8)\) is diagonal with entries \(14,28,28,28,35+w,28,35+w,28\), divided by 360, \(w=2(1-(49/91)^2)\), so \(\mathcal P_8\) identifies and no identifying partition has more than eight cells. The partition by calendar quarter does not identify, as the pair \(v_i=100\), \(u_p=2500\) against the same except \(v=86,72,74\) for the meetings of 15 September, 27 October and 8 December 2027 and \(u=2860\) on the last cell shows; merging the last two quarters identifies; and the partition at every meeting does not identify, with the pair of Proposition \hyperref[res:Proposition-49]{49}(d).

(d4) Under (S1), with \(\mathcal P_8\), the six same-window diagonal entries are positive and the two others both equal \((g(\kappa)/2\kappa)\,F(\kappa/360)\), where \(g(\kappa)=((1-e^{-\kappa\delta})/\kappa\delta)^2\) and, with exponents in days,

\begin{equation*}
F(x)=e^{-122x}-e^{-364x}-e^{-10x}+e^{-182x}. \tag{3.197}\label{eq:3.197}
\end{equation*}

\(F\) has exactly one positive zero \(x^*\), and \(\mathcal P_8\) identifies if and only if \(\kappa\neq\kappa^*:=360x^*\in(0.828,0.864)\).

(d5) With \(P=1\) and any shape of (d1), take instead the twelve expiries listed on 2 January 2026: the four nearest serials, 16 January, 13 February, 10 April and 15 May 2026, and the quarterlies of Proposition \hyperref[res:Proposition-24]{24} through 10 December 2027. Their gaps contain \(0,1,0,1,1,0,2,3,1,3,1,3\) meetings in date order. For strictly positive \(\theta\), a linear functional of \(\theta\) is determined by the panel if and only if it lies in the span of \(u_1\) and the gap sums \(\sum_{T_i\in G_\ell}v_i\), which has dimension 10, so \(\operatorname{rank}\mathcal A=10\) of 17. The panel identifies \(u_1\) and, individually, the meetings of 28 January, 18 March and 29 April 2026 and of 27 January and 28 July 2027; it identifies the other eleven only through the four sums over 17 June and 29 July 2026, over 16 September, 28 October and 9 December 2026, over 17 March, 28 April and 9 June 2027, and over 15 September, 27 October and 8 December 2027, and no other combination of them.

The diffusion now loads on each window with the weight \(\beta\), and a meeting-free gap measures

\begin{equation*}
\tilde\Lambda_{\ell,p}=\int_{[c_{p-1},c_p)\cap(0,S_{\ell-1}]}\bigl(\beta_\ell^2-\beta_{\ell-1}^2\bigr)+\int_{[c_{p-1},c_p)\cap G_\ell}\beta_\ell^2, \tag{3.198}\label{eq:3.198}
\end{equation*}

a combination of the diffusion rates that includes cells wholly before the gap when the two expiries have different windows. Under (S1) a later window has the smaller weight, \(\beta_\ell(s)^2=g(\kappa)e^{-2\kappa(a_\ell-s)}\), and that term is negative; under (S2) the cell \((9\text{ Dec},11\text{ Dec }2026)\) has the positive entry \((2/360)(1-(49/91)^2)\) in the row of the gap ending 15 January 2027. This is why (d4) turns on the two window changes, whose day counts \((61,182,5,91)\) coincide and give the same \(F\). Part (d5) depends only on which gap holds each meeting. The 2027 meeting dates are tentative, and (d5) is unchanged if any of them moves within its gap; every decision day is a Wednesday, the next business day a Thursday and every expiry a Friday, so using effective dates in place of decision days changes no gap assignment, here or for the 24 expiries of (d1) to (d4). Proposition \hyperref[res:Proposition-24]{24}(e) is the special case of the two gaps from 12 June to 11 December 2026. Parts (d1) and (d4) are exact statements about rank, not about practical identification: at \(\kappa=50\), \(D_{1,1}\approx10^{-9}\) days, and above \(\kappa^*\approx0.8527\) the two window-change entries are negative. The identity \(\tilde z/\delta=\int_0^SQ\) of (3.189) survives exactly when \(\int_0^S\sigma(s)^2[\Phi(s,S)\beta(s)-(S-s)\beta(s)^2]\,ds=0\), which \(\phi\equiv1\) satisfies. Proposition \hyperref[res:Proposition-49]{49} is the case \(\phi\equiv1\), and the proof of (3.190) never uses the form of the diffusion columns; (S2) is the loading of \cite{brace2024sofr} §3.1, where no identification from option panels is made; \cite{wright2020event} is the statistical counterpart, as in Section \hyperref[sec:3.43]{3.43}.

As a numerical check, not part of the statement, redo Proposition \hyperref[res:Proposition-49]{49}(e) with \(P=1\), \(\hat u=Q(S_1)/D_{1,1}\) and \(\hat v_i=Q(S_\ell)-Q(S_{\ell-1})-\tilde\Lambda_{\ell,1}\hat u\). Under (S2) the counts of practically identified meetings are those of the table in Section \hyperref[sec:3.43]{3.43}, since \(\beta=1\) almost everywhere on the calendar; under (S1) they are never smaller and often larger, because the diffusion loads on each window with \(\beta<1\), and less as the window moves out:

{\scriptsize\setlength{\tabcolsep}{4pt}
\begin{longtable}{lllllll}
\toprule
\textbf{(S1), \(\sigma\) (bp per year)} & \textbf{\(s_0=10\), \(\varepsilon=0.25\)} & \textbf{10, 0.5} & \textbf{15, 0.25} & \textbf{15, 0.5} & \textbf{25, 0.25} & \textbf{25, 0.5} \\
\midrule
\endhead
\bottomrule
\endlastfoot
\(\kappa=0.1\), 50 & 2 & 0 & 14 & 2 & 16 & 15 \\
\(\kappa=0.1\), 100 & 0 & 0 & 4 & 0 & 16 & 6 \\
\(\kappa=0.5\), 50 & 4 & 0 & 16 & 4 & 16 & 16 \\
\(\kappa=0.5\), 100 & 0 & 0 & 6 & 1 & 16 & 9 \\
\(\kappa=1\), 50 & 4 & 1 & 16 & 4 & 16 & 16 \\
\(\kappa=1\), 100 & 2 & 0 & 10 & 2 & 16 & 12 \\
\end{longtable}}

Each estimator's first-order bound agrees with the linear-programming optimum over linear estimators exact on the model. In basis points of standard deviation, an error \(\Delta v\) in a meeting variance of reference size \(s_0^2\) is, to first order, \(\Delta v/(2s_0)\), and a meeting is practically identified exactly when this is below \(s_0/2\). For the first and the sixteenth meetings it is 11.9 and 40.7 basis points under (S2) with \(\sigma=100\), \(s_0=10\) and \(\varepsilon=0.5\), and 1.3 and 5.6 under (S1) with \(\kappa=1\), \(\sigma=50\), \(s_0=25\) and \(\varepsilon=0.25\). The condition number of the \(24\times17\) matrix \(\mathcal A\) with \(P=1\), time in years, is 94 for the parallel diffusion, 109, 103 and 102 under (S1) at \(\kappa=0.1\), 0.5 and 1, and 93 under (S2).

\emph{Proof.} (a) For fixed \(t\), \(x\mapsto\Phi(t,x)\) is absolutely continuous with derivative \(\phi(t,x)\) almost everywhere, so \(\int_t^T\sigma^2\phi\Phi=\sigma^2\Phi(t,T)^2/2\), which is (2.3); the jump condition is that of Proposition \hyperref[res:Proposition-49]{49}(a). Stochastic Fubini for the bounded measurable integrand \(\sigma(v)\phi(v,s)\mathbf 1\{v\le s\}\) gives the jointly measurable version and turns the diffusion part of \(\int_0^tr+\int_t^Tf(t,\cdot)\) into \(\int_0^t\sigma\Phi(\cdot,T)\,dW\), and Fubini with \(\partial_x\Phi^2/2=\phi\Phi\) turns the drift part into \(\tfrac12\int_0^t\sigma^2\Phi(\cdot,T)^2\). (b) Stochastic Fubini gives \(\int_a^b\int_0^x\sigma(s)\phi(s,x)\,dW_s\,dx=\int_0^b\sigma\tilde w\,dW\), and the rest is the argument of Proposition \hyperref[res:Proposition-49]{49}(b), with variance \(\int_0^S\sigma^2\tilde w^2\) and, by (a) at \(t=T=S\), covariance \(-\tilde z\); for \(S<a\), \(\tilde w=\delta\beta\) on \([0,S]\). (c) The proof of (3.190) uses only that the meeting columns are prefix indicators and that differencing maps row \(\ell\) to \((\sum_{T_i\in G_\ell}e_i^\top,\tilde\Lambda_{\ell,\cdot})\); (3.198) is the definition of \(D\) split at \(S_{\ell-1}\). (d1) The first gap is meeting-free and \(D_{0,1}=0\). (d2) Row \(\ell\) involves only \((0,S_\ell]\), the \(k\)-th meeting-free gap lies in cell \(k\) of \(\mathcal P_8\), and cell \(k+1\) starts after it; in a same-window row the first term of (3.198) vanishes. (d3) \(\beta<1\) only where the first meeting after \(S_\ell\) falls inside the window or does not exist, which on the calendar happens at the three stretches listed, each inside the cell of its row's gap; the pairs are kernel vectors. (d4) The window-change entry is \(g[\int_c^{S_\ell}e^{-2\kappa(a_\ell-s)}ds-\int_c^{S_{\ell-1}}e^{-2\kappa(a_{\ell-1}-s)}ds]\) with \(c\) the start of the cell. With \(G(x)=e^{10x}F(x)\) and \(H(x)=e^{112x}G'(x)=-112-172e^{-60x}+354e^{-242x}\), \(H(0)=70\), \(H\to-112\) and \(H'\) changes sign once, so \(H\) has one zero, \(G\) rises and then falls to \(-1\), and \(F\) has one positive zero; \(F(0.0023)>0>F(0.0024)\) gives the bracket. (d5) After differencing, the first row is \((0,D_{1,1})\) with \(D_{1,1}>0\), so \(u_1\) lies in the row space, and subtracting multiples of that row from the others leaves the sums over the nine gaps with meetings; these ten vectors have disjoint supports, so the rank is 10 and \(v_i\) lies in the span exactly when \(T_i\) is alone in its gap. A functional outside the row space changes along a kernel vector, and a small step along it keeps a strictly positive \(\theta\) positive.\hfill$\square$

Parts (a) and (b) are conditional on the Wiener-integral facts of Proposition \hyperref[res:Proposition-49]{49} and on stochastic Fubini for bounded jointly measurable deterministic integrands with a jointly measurable version; no Brownian motion is constructed. Parts (c) and (d) are unconditional; the value of \(\kappa^*\), the table, the standard-deviation errors, the linear-programming comparison and the condition numbers are numerical checks. Part (d4) is proved for \(\mathcal P_8\) only. In (S2) the model has no meeting after 8 December 2027; adding the January 2028 meeting would change only the last coefficient of the quarterly pair and no rank. If the jump at \(T_i\) loaded on window \(\ell\) with a known factor \(\omega_{\ell,i}\), the meeting columns would no longer be prefix indicators, (3.196) would not apply, and a surface would identify \(\omega_{\ell,i}^2v_i\); that case is not treated.

\subsection{Approximation by recurrent loadings on a finite horizon}\label{sec:3.45}

Propositions \hyperref[res:Proposition-30]{30} and \hyperref[res:Proposition-33]{33} show that loadings satisfying no linear recurrence admit no realization serving every schedule. On a finite horizon a model with recurrent loadings can stand in for such a model, each keeping its own drift. Keep the notation \(i(s,T)\) of Section \hyperref[sec:3.27]{3.27}, one Brownian factor \(W\) with the cited calculus of Section \hyperref[sec:3.25]{3.25}, a horizon \(H>0\), a finite schedule \(0<T_1<\dots<T_N<H\), and a Borel shape \(\lambda\) on \([0,H]\), not identically zero, with \(\Lambda=\sup_{[0,H]}|\lambda|<\infty\); a constant scale is absorbed into \(\lambda\). Take original loadings \(a_0,\dots,a_N\) and approximating loadings \(\tilde a_0,\dots,\tilde a_N\), the only ones that occur on \([0,H]\), and put

\begin{equation*}
\varepsilon=\max_{i\le N}|a_i-\tilde a_i|,\qquad\bar a=\max_{i\le N}\max(|a_i|,|\tilde a_i|). \tag{3.199}\label{eq:3.199}
\end{equation*}

Each model has the volatility (3.138) and the differentiated drift generated from its own volatility, \(\sigma(s,T)=a_{i(s,T)}\lambda(T-s)\) and \(\alpha(s,T)=\sigma(s,T)\int_s^T\sigma(s,u)\,du\), with \(\tilde\sigma\), \(\tilde\alpha\) likewise from \(\tilde a\). With a common bounded Borel initial curve \(f_0\), \(F_0=\sup|f_0|\), a common driver and no curve jumps, the curves are \(f(t,T)=f_0(T)+\int_0^t\alpha(s,T)\,ds+\int_0^t\sigma(s,T)\,dW_s\) and \(\tilde f\) likewise, for \(0\le t\le T\le H\).

\phantomsection\label{res:Proposition-51}
\textbf{Proposition 51 (conditional).} (a) For \(0\le s\le T\le H\), \(|\sigma-\tilde\sigma|\le\varepsilon\Lambda\) and \(|\alpha-\tilde\alpha|\le2\varepsilon\bar a\Lambda^2(T-s)\). For \(0\le t\le T\le H\), with the deterministic \(m(t,T)=\int_0^t(\alpha-\tilde\alpha)(s,T)\,ds\),

\begin{equation*}
E|f(t,T)-\tilde f(t,T)|^2=m(t,T)^2+\int_0^t(\sigma-\tilde\sigma)(s,T)^2\,ds\le\varepsilon^2\bigl[\Lambda^2t+4\bar a^2\Lambda^4(tT-t^2/2)^2\bigr], \tag{3.200}\label{eq:3.200}
\end{equation*}

\begin{equation*}
\sup_{0\le t\le T\le H}E|f(t,T)-\tilde f(t,T)|^2\le\varepsilon^2\bigl(\Lambda^2H+\bar a^2\Lambda^4H^4\bigr). \tag{3.201}\label{eq:3.201}
\end{equation*}

For \(a\equiv\bar a\), \(\tilde a\equiv\bar a-\varepsilon\) and \(\lambda\equiv\Lambda\), \(E|f-\tilde f|^2=\varepsilon^2[\Lambda^2t+(2\bar a-\varepsilon)^2\Lambda^4(tT-t^2/2)^2]\), so (3.200) is attained as \(\varepsilon/\bar a\to0\).

(b) Assume (H1). Each curve \(f(t,\cdot)\) has a version measurable in \((\omega,u)\) and bounded on \([t,H]\), which defines \(P(t,T)=\exp(-\int_t^Tf(t,u)\,du)\) and \(\tilde P\), and for \(0\le t\le T\le H\)

\begin{equation*}
E|\log P-\log\tilde P|^2\le\varepsilon^2\bigl(\Lambda^2H^3+\tfrac14\bar a^2\Lambda^4H^6\bigr),\qquad E|P-\tilde P|^2\le2\sqrt3\,e^{2HF_0+5\bar a^2\Lambda^2H^3}\,\varepsilon^2\bigl(\Lambda^2H^3+\tfrac14\bar a^2\Lambda^4H^6\bigr). \tag{3.202}\label{eq:3.202}
\end{equation*}

(c) Assume (H1), with \(A\) of size \(r\), and (H2) for the approximating loadings, \(\tilde a_i=uM^iv\) with \(M\) of size \(p\), and take \(f_0=0\). For every schedule in \((0,H)\), the approximating model with its own drift \(\tilde\alpha\) is realized by Proposition \hyperref[res:Proposition-31]{31} with \(\psi\equiv1\), with the state \((\Xi,\Pi,\Theta)\) of dimension \(m+m(m+1)/2+m\), \(m=pr\), the same for every schedule and every \(H\), of which only \(\Xi\) is random; and its realized curve and bond price satisfy (3.201) and (3.202) with \(F_0=0\) against the original model.

The constants depend only on \(\Lambda\), \(\bar a\), \(H\) and, for prices, \(F_0\): not on the number or spacing of the meetings, on the initial curve, which cancels, or on \(\lambda\) beyond \(\Lambda\). The schedule enters only through the indices compared, since \(\varepsilon\) and \(\bar a\) are maxima over \(0,\dots,N\). So once bounds on \(\sup_i|a_i|\) and \(\sup_i|\tilde a_i|\) are fixed, the bounds are free of the number and spacing of meetings. For \(\tilde a_i=uM^iv\) with \((u,M,v)\) minimal, the sequence is bounded if and only if the spectral radius of \(M\) is below one, or equal to one with every eigenvalue of modulus one semisimple (standard linear systems theory); otherwise \(\bar a\), and with it the factor \(e^{5\bar a^2\Lambda^2H^3}\), grows with \(N\), and the comparison holds on each schedule with constants that are not schedule-free. The price bound needs \(E\,e^{-4X}<\infty\) for \(X=-\log P\) in the zero-curve model; \(X\) is Gaussian with \(|EX|\le\frac12\bar a^2\Lambda^2H^3\) and \(\operatorname{Var}X\le\bar a^2\Lambda^2H^3\), so \(E\,e^{-4X}\le e^{10\bar a^2\Lambda^2H^3}\), automatically because the volatility is deterministic and bounded; at \(\bar a\Lambda\approx0.01\) per square-root year and \(H=3\) years, \(5\bar a^2\Lambda^2H^3\approx0.014\). Part (c) is a comparison conditional on a supplied recurrent approximation. Every \(a_0,\dots,a_N\) is matched exactly by a recurrence of order \(N+1\), whose dimension grows with \(N\); that is interpolation, not compression, and nothing here says which loading sequences have accurate recurrent approximations of small order. Only the supremum of pointwise mean-square errors is bounded, not an expected supremum over the curve.

Parts (a) and (b) are the standard stability argument for Gaussian HJM models with deterministic volatility, through the isometry and the locally Lipschitz drift map \(\sigma\mapsto\sigma\int\sigma\), and the standard passage from Gaussian log-prices to prices through exponential moments. What is specific is the schedule-free constant in the ordinal-loading parametrization and the pairing with Proposition \hyperref[res:Proposition-31]{31}, which realizes the approximating model with its own drift in a dimension independent of the schedule. The nearest precedent, \cite{benth2018approximation} Theorem 4.1, approximates forward-curve dynamics in commodity markets by arbitrage-free models with finite-dimensional state through a Riesz basis, with rates of convergence; its forward prices are martingales, so there is no drift map to control, and it has no schedule. As a numerical check, not part of the statement, on six schedules with 2 to 200 meetings, evenly spaced and random, the largest ratio of the mean-square error to (3.201) is 0.115.

\emph{Proof.} (a) \(\sigma-\tilde\sigma=(a_i-\tilde a_i)\lambda\), and \(\alpha-\tilde\alpha=(\sigma-\tilde\sigma)S+\tilde\sigma(S-\tilde S)\) with \(S=\int_s^T\sigma(s,u)\,du\), \(|S|\le\bar a\Lambda(T-s)\) and \(|S-\tilde S|\le\varepsilon\Lambda(T-s)\). Since \(f_0\) cancels, \(f-\tilde f=m+\int_0^t(\sigma-\tilde\sigma)\,dW\), and the isometry and the mean-zero property give the equality in (3.200); the bounds follow, with \(tT-t^2/2\le H^2/2\). In the example, \(\sigma-\tilde\sigma\equiv\varepsilon\Lambda\) and \(\alpha-\tilde\alpha=(2\bar a-\varepsilon)\varepsilon\Lambda^2(T-s)\). (b) Fix \(t\) and split \([0,t]\) at the meetings into pieces with right ends \(\tau_l\); on each the number \(n_l\) of meetings in \((s,t]\) is constant and \(i(s,u)=n_l+i(t,u)\) for \(u\ge t\), so by (H1) \(\int_0^t\sigma(s,u)\,dW_s=\sum_la_{n_l+i(t,u)}\,c\,e^{A(u-\tau_l)}\zeta_l\) with \(\zeta_l=\int_{\text{piece }l}e^{A(\tau_l-s)}b\,dW_s\). This defines the version, the same \(\zeta_l\) serve \(\tilde f\), and by Lemma \hyperref[res:Lemma-32]{32} they are jointly Gaussian. So \(X=\int_t^T(f(t,u)-f_0(u))\,du\) and \(\tilde X\) are jointly Gaussian, with \(EX=\frac12\int_0^t(S(s,T)^2-S(s,t)^2)\,ds\) and \(\operatorname{Var}X=\int_0^tU^2\), \(U(s)=\int_t^T\sigma(s,u)\,du\), and \(D=X-\tilde X\) has \(|ED|\le\frac12\varepsilon\bar a\Lambda^2H^3\) and \(\operatorname{Var}D\le\varepsilon^2\Lambda^2H^3\), which is the first bound. For the second, \(|e^{-X}-e^{-\tilde X}|\le|D|\max(e^{-X},e^{-\tilde X})\), and Cauchy--Schwarz with \(ED^4\le3(ED^2)^2\) and \(E\,e^{-4X}=e^{-4EX+8\operatorname{Var}X}\) gives the constant. (c) Proposition \hyperref[res:Proposition-31]{31} with \(\psi\equiv1\) and loadings \(\tilde a\) has the volatility \(\tilde\sigma\) and the drift \(\tilde\alpha\). Its realized curve is jointly measurable and agrees with \(\tilde f(t,t+x)\) almost surely for each \(x\), hence, by Fubini, for almost every \(x\) almost surely, so the realized bond price is \(\tilde P\) almost surely; (ii) is (a) and (b).\hfill$\square$

Part (a) is conditional on the cited calculus of Section \hyperref[sec:3.25]{3.25}, parts (b) and (c) in addition on \(W\) being a Brownian motion, which gives the Gaussian law of Lemma \hyperref[res:Lemma-32]{32}, and part (c) on the inputs of Proposition \hyperref[res:Proposition-31]{31}. For a bounded Borel \(\lambda\) that is not quasi-exponential, (3.202) would follow from a jointly measurable version with stochastic Fubini, which is not shown. A realization from a nonzero initial curve is not given, since Proposition \hyperref[res:Proposition-31]{31} starts from the zero curve; (a) and (b) hold for any common bounded initial curve. Stochastic scales, several factors, fixed-maturity curve jumps, the choice of recurrence, approximation rates, minimal state dimension and option-payoff errors are not treated.

\subsection{The unified correlated splice}\label{sec:3.46}

Proposition \hyperref[res:Proposition-44]{44} treats several step factors and blocks without an exponent zero, and Proposition \hyperref[res:Proposition-47]{47} a full polynomial part with one step factor. This section combines them and states the result without a counting convention. Keep the driver-form setting of Section \hyperref[sec:3.38]{3.38}: independent drivers \(B\) with the cited calculus, a front end of the type of Theorem \hyperref[res:Theorem-11]{11} whose level on \(I_m\) has the noise \(V_m\cdot dB\), \(V_m=\sum_\ell s_{\ell,m}H^S_\ell\), and the aggregate maturity-step integrand \(G_m\) of (3.172). The block is (3.178), with \(\mathcal A\) invertible and \((\mathcal A,\mathcal C)\) observable, and \(Z=(z_0,\zeta)\) in driver form with predictable bounded rows \(H^{0,\mu}\) and \(H^\zeta\); Assumption 2.1 is assumed for the drift and volatility of the representation (3.161), as in Propositions \hyperref[res:Proposition-43]{43}, \hyperref[res:Proposition-44]{44} and \hyperref[res:Proposition-47]{47}. Write \(\varphi_k\) for the block's coordinate functions (\(x^\mu\) for \(z_{0,\mu}\), the components of \(\mathcal C\,e^{\mathcal Ax}\) for \(\zeta\)), \(\Phi_k=\int_0^x\varphi_k\), \(H^k\) for the corresponding rows, \(a=(H^k\cdot H^l)\) and \(R\) for the residual (3.162). Let

\begin{equation*}
\bar\omega_k(u)=H^k(u)\,V_{j(u)}(u)^\top,\qquad H^{0,0}_\sharp=H^{0,0}+V_{j(u)},\quad H^k_\sharp=H^k\ \text{otherwise}, \tag{3.203}\label{eq:3.203}
\end{equation*}

the covariation of each block coordinate with the current front-end noise and the effective level, and let \(a_\sharp=(H^k_\sharp\cdot H^l_\sharp)\) and \(R_\sharp\) be \(R\) with \(a_\sharp\) in place of \(a\).

\phantomsection\label{res:Proposition-52}
\textbf{Proposition 52 (conditional).} (a) If the curve satisfies Assumption 2.1, then almost surely, for every \(m\) and almost every \(u<T_m\),

\begin{equation*}
H^{0,\mu}(u)\,G_m(u)^\top=0\quad(\mu=1,\dots,d),\qquad H^\zeta(u)\,G_m(u)^\top=0; \tag{3.204}\label{eq:3.204}
\end{equation*}

no condition is imposed on the level's covariation \(H^{0,0}G_m^\top\).

(b) Assume (3.204). (i) The jump across \(T_m\) of the cross term between front end and block comes from the level only and is the affine function \(T\mapsto(H^{0,0}G_m^\top)(u)\,(2T-u-T_m)\). (ii) For every \(m\ge j(u)\) the non-level part of the cross term on \(I_m\) is one function of \(x=T-u\),

\begin{equation*}
K(u,x)=\sum_{\mu=1}^d\bar\omega_\mu(u)\,\frac{\mu+2}{\mu+1}\,x^{\mu+1}+\mathcal C\bigl[x\,e^{\mathcal Ax}+\mathcal A^{-1}(e^{\mathcal Ax}-I)\bigr]\bar\omega_\zeta(u), \tag{3.205}\label{eq:3.205}
\end{equation*}

and the level part is affine in \(T\) on each \(I_m\). (iii) For every \(u\) and \(x\),

\begin{equation*}
R(u,x)-K(u,x)=R_\sharp(u,x)+(\text{an affine function of }x): \tag{3.206}\label{eq:3.206}
\end{equation*}

on the current interval the front end's noise moves the curve by a constant, exactly as the level does. (iv) The curve satisfies Assumption 2.1 if and only if (3.204) holds and, for \(dt\otimes dQ\)-almost every \((u,\omega)\), \(x\mapsto R_\sharp(u,x)\) is affine on \([0,\infty)\); the front-end drifts then absorb that affine part, the level's cross terms including the jumps of (i), and \(\sigma^S\cdot\int\sigma^S\), as in Proposition \hyperref[res:Proposition-38]{38}(a).

(c) If \(\mathcal A\) is diagonalizable over \(\mathbb C\), then \(R_\sharp\) affine forces, for \(dt\otimes dQ\)-almost every \((u,\omega)\),

\begin{equation*}
H^\zeta H^{0,\mu\top}=0\quad(\mu=1,\dots,d),\qquad H^\zeta\bigl(H^{0,0}+V_{j(u)}\bigr)^\top=0, \tag{3.207}\label{eq:3.207}
\end{equation*}

and then \(R_\sharp\) affine is the exponent-zero condition (3.176) with \(a_\sharp\) in place of \(a\) and without polynomial--exponential constants, together with the nonzero-exponent condition of the block alone, which involves only \(a_{\zeta\zeta}\), \(b_\zeta\) and \(\zeta\); pairs of eigenvalues with \(\lambda_i+\lambda_j=0\) contribute constants to the exponent-zero part, which are affine.

(d) Condition (3.204) alone is not sufficient: for \(d=1\), no exponential part, a block uncorrelated with the front end and \(a_{(0,1),(0,1)}>0\), \(R_\sharp\) contains \(-\frac12a_{(0,1),(0,1)}x^3\), which no drift reaches.

(e) With \(d=0\) and \(H^{0,0}=0\), which changes no curve, (3.204) is (3.174), (b)(iv) is the converse of Proposition \hyperref[res:Proposition-44]{44}(a) with \(K=\mathcal C[x\,e^{\mathcal Ax}+\mathcal A^{-1}(e^{\mathcal Ax}-I)]H^\zeta V_{j(u)}^\top\), and for diagonalizable \(\mathcal A\) consistency also requires \(H^\zeta V_{j(u)}^\top=0\): the block noise must be orthogonal to the current front-end noise, not only to the aggregate step. With one step factor, (a) and (b) are Proposition \hyperref[res:Proposition-47]{47}(a) and (b) with \(K\) and the effective level explicit. In the uncorrelated splice \(K=0\) and \(R_\sharp=R\) up to an affine function. Proposition \hyperref[res:Proposition-45]{45} holds with \(a_\sharp\): the bound (b) is unchanged, since the entries with \(\mu,\nu\ge1\) are the same in \(a_\sharp\); for a polynomial-only block every bounded predictable choice of \(H^{0,\mu}\), zero for \(\mu>\lfloor(d-1)/2\rfloor\) and satisfying (3.204), is admissible with suitable drifts, whatever its covariation with the current front-end noise; and for \(d\ge1\) the splice is consistent if and only if (3.204) holds and \((a_\sharp,b')\) makes the block-alone residual vanish, \(b'\) differing from \(b\) only in \(b_{0,0}\) and \(b_{0,1}\).

The step factors enter (3.204) only through \(G_m\), so single step factors may covary with any block direction as long as their maturity steps cancel in the aggregate; two step factors on the same noise with \(\Delta_{1,m}=-\Delta_{2,m}\) give \(G_m=0\) and no restriction. The current front-end noise is constrained only through the effective level. If \((\mathcal A,\mathcal C)\) is not observable, the minimal realization of Proposition \hyperref[res:Proposition-44]{44}(c) leaves the curve unchanged, and (a) to (c) hold with \(PH^\zeta\) in place of \(H^\zeta\). Relative to the uncorrelated splice, the drift \(b_{0,\mu+1}\) shifts by \(\frac{\mu+2}{\mu+1}\bar\omega_\mu\) for \(1\le\mu\le d-1\), and \(b_{0,0}\) and \(b_\zeta\) absorb the constant and the \(\mathcal C\mathcal A^{-1}e^{\mathcal Ax}\bar\omega_\zeta\) part of \(K\); the term \(\mathcal C\,x\,e^{\mathcal Ax}\bar\omega_\zeta\) cancels, for diagonalizable \(\mathcal A\), only through (3.207). For a non-diagonalizable \(\mathcal A\), such as the Nelson--Siegel term \(x\,e^{-\beta x}\), a drift reaches \(x^ke^{\lambda x}\), (3.207) is not forced by this argument, and the condition stays "\(R_\sharp\) affine". The aggregate integrand and observability are those of Proposition \hyperref[res:Proposition-44]{44}, the descending-degree argument and the level exemption those of Proposition \hyperref[res:Proposition-47]{47}, the affine absorption that of Proposition \hyperref[res:Proposition-38]{38} and the polynomial bound that of Proposition \hyperref[res:Proposition-45]{45}. New are the effective-level identity (3.206), which turns the correlated splice's residual condition into the block-alone condition for \(a_\sharp\) and makes the carry-over of Proposition \hyperref[res:Proposition-45]{45} exact, and (c), which sharpens Proposition \hyperref[res:Proposition-44]{44}'s converse. \cite{gellert2021short} §3.1.3 decomposes correlated deterministic factors into independent ones, \cite{brace2024sofr} has no level correlated with meeting steps, and \cite{filipovic2000exponential} and \cite{bjork1999interest} §7 treat exponential-polynomial blocks alone; the drift of a two-factor Gaussian model with a correlated level and an exponential factor, which produces \(x\,e^{-\kappa x}\), is standard in form, and its reading as an effective level is what (3.206) adds.

\emph{Proof.} Itô's formula for the representation and Assumption 2.1 at rational maturities, extended by continuity, give for \(dt\otimes dQ\)-almost every \((u,\omega)\) that \(R(u,T-u)-c_m(u,T)\) is affine in \(T\) on each \(I_m\cap[u,H]\), where \(c_m\) is the cross part of \(\sigma\cdot\int\sigma\) on \(I_m\); everything below is at such a point. With \(\Sigma^S(u,T)=\sum_\ell(s_{\ell,m}T+\kappa_{\ell,m})H^S_\ell\), \(\kappa_{\ell,j(u)}=-s_{\ell,j(u)}u\) and \(\kappa_{\ell,m}-\kappa_{\ell,m-1}=-\Delta_{\ell,m}T_m\), the jump is \(c_m-c_{m-1}=\sum_k(H^kG_m^\top)[\varphi_k(x)(T-T_m)+\Phi_k(x)]\). (a) Its part with nonzero exponents, \(\mathcal C\,e^{\mathcal Ax}[(x+u-T_m)I+\mathcal A^{-1}]H^\zeta G_m^\top\), must vanish, and Lemma \hyperref[res:Lemma-42]{42} gives \(H^\zeta G_m^\top=0\); in the polynomial part the \(\mu\)-th term has degree \(\mu+1\) and leading coefficient \(\frac{\mu+2}{\mu+1}H^{0,\mu}G_m^\top\), and descending from \(\mu=d\) to \(\mu=1\) gives (3.204), while the level term \((H^{0,0}G_m^\top)(2x+u-T_m)\) is affine. (b) (i) is this computation under (3.204). For (ii), under (3.204) the sums \(\sum_\ell s_{\ell,m'}H^kH^{S\top}_\ell\) equal \(\bar\omega_k\) and \(\sum_\ell\kappa_{\ell,m'}H^kH^{S\top}_\ell=-u\bar\omega_k\) for every later interval and non-level \(k\), so the non-level cross part is \(\sum_k\bar\omega_k[x\varphi_k+\Phi_k]\), which is (3.205). For (iii), \(a_\sharp-a\) vanishes outside the level row and column, where it is \(\bar\omega_k\) and \(2\bar\omega_0+|V_{j(u)}|^2\), so \(R-R_\sharp=\sum_{k\neq0}\bar\omega_k(\Phi_k+x\varphi_k)+(2\bar\omega_0+|V_{j(u)}|^2)x\). For (iv), \(R-K\) is affine on the current interval, which has positive length, and is real-analytic in \(x\), hence affine on \([0,\infty)\); conversely the drifts are chosen interval by interval as in Proposition \hyperref[res:Proposition-38]{38}(a). (c) In a complex eigenbasis with eigenvalues \(\lambda_i\neq0\) and \(c_i\neq0\), a term \(x^ke^{\lambda_ix}\) with \(k\ge1\) comes only from the pair \((\mu,\zeta_i)\), with coefficient \(-c_i[a_{\sharp,(k,i)}/\lambda_i+a_{\sharp,(k-1,i)}/k]\); descending from \(k=d+1\) gives (3.207). (d) The coefficient of \(x^3\) in (3.176) is \(-a_{(0,1),(0,1)}/2\), and no drift has degree 3. (e) is substitution, with the top-degree argument of Proposition \hyperref[res:Proposition-45]{45}(b), which uses only entries unchanged in \(a_\sharp\).\hfill$\square$

The proposition is conditional on the cited calculus through Itô's formula for the representation, and on Assumption 2.1 for that representation; the statements are algebra at a point, and no process \(Z\), \(L_m\) or \(C_m\) is constructed. Not covered: varying exponents, polynomial parts that are not full, front ends with a diffusive ramp, a non-diagonalizable exponential part beyond "\(R_\sharp\) affine", and unbounded coefficients, which Proposition \hyperref[res:Proposition-53]{53} treats.

\subsection{The unified splice beyond bounded coefficients}\label{sec:3.47}

Proposition \hyperref[res:Proposition-52]{52} assumes bounded predictable integrands. Keep its setting on the horizon \(H\) and replace boundedness by the class (L):

\begin{itemize}
  \item (L1) each front-end noise \(V_m\) and each block row \(H^k\) is predictable with \(\int_0^H|\cdot|^2\,du<\infty\) almost surely;
  \item (L2) the block drifts \(b_k\) and the drifts of \(L_m\) and \(C_m\) are progressively measurable with \(\int_0^H|\cdot|\,du<\infty\) almost surely, \(C_m\) having no noise.
\end{itemize}

Condition (L1) is placed on the sums \(V_m=\sum_\ell s_{\ell,m}H^S_\ell\), not on their factors: local square integrability is not closed under products, since \(s(u)=h(u)=u^{-1/3}\) are square integrable on \((0,1]\) and \(s\,h\) is not. One locally bounded factor times a square-integrable one suffices; for example, a Heston-type scale \(\sqrt{v_u}\) with \(v\) continuous, adapted and nonnegative, times a bounded row. The aggregate \(G_m\) enters only through the pointwise products \(H^kG_m^\top\) and needs no integrability of its own.

\phantomsection\label{res:Proposition-53}
\textbf{Proposition 53 (conditional).} (a) Under (L), if the curve satisfies Assumption 2.1, then (3.204) holds, \(Q\)-almost surely, for every \(m\) and almost every \(u<T_m\); nothing is imposed on the level.

(b) Parts (b) to (e) of Proposition \hyperref[res:Proposition-52]{52} hold verbatim under (L), with the quantifier "for \(dt\otimes dQ\)-almost every \((u,\omega)\)".

(c) Let \(Z\) be of class (L), satisfy (3.204), and have \(R_\sharp(u,\cdot)\) affine for \(dt\otimes dQ\)-almost every \((u,\omega)\). Then the front-end drifts of Proposition \hyperref[res:Proposition-52]{52}(b)(iv) are locally integrable, and the processes \(L_m\) and \(C_m\) defined from arbitrary initial values by these drifts and the noises \(V_m\) give a curve that satisfies the regularity conditions of Section \hyperref[sec:2]{2} and Assumption 2.1.

No stopping argument is used, and the quantifiers are those of Proposition \hyperref[res:Proposition-52]{52}. After Itô's formula, the proof of Proposition \hyperref[res:Proposition-52]{52} is algebra at a fixed \((u,\omega)\) and never uses a bound on the coefficients, and Itô's formula and the stochastic integral are already cited for locally square-integrable integrands and locally integrable drifts (\cite{bauerschmidt2020stochastic} §3.3 and §3.5). Localizing times \(\tau_n=\inf\{t:\int_0^t(\sum_m|V_m|^2+\sum_k|H^k|^2+\sum_k|b_k|+\text{front-end drifts})\,du\ge n\}\wedge H\) would reduce (L) to square integrability in expectation, but they give no pointwise bound, which needs locally bounded coefficients; local boundedness implies local square integrability, not conversely (\cite{bauerschmidt2020stochastic} p. 45). Stopping the parameters is also not stopping the curve: after \(\tau\) the block \(F(T-t,Z_\tau)\) still moves through the drift \(-\partial_xF\), and a stopped consistent model must continue with \(z(t)=S_{t-\tau}Z_\tau\), where \(F(x+h,z)=F(x,S_hz)\) with \(S_h\) the binomial shift on \(z_0\) and \(e^{\mathcal Ah}\) on \(\zeta\); under that continuation the curve is frozen. Localization of stochastic integrals is textbook; what the proposition adds is that no bound is needed after Itô's formula, that the drifts of the converse stay integrable, and where the integrability hypothesis belongs. It extends Proposition \hyperref[res:Proposition-52]{52} to Heston-type scales, conditional on their existence.

\emph{Proof.} For \(T\in I_m\), \(\sigma(u,T)=V_m(u)+\sum_k\varphi_k(T-u)H^k(u)\) is square integrable in \(u\), a finite sum with bounded continuous multipliers, and \(\alpha(\cdot,T)\) is locally integrable: it collects the drifts of \(L_m\) and \(C_m\), the terms \(b_k\varphi_k(T-u)\) and \(-\partial_xF(T-u,Z_u)\), which is continuous in \(u\) because \(Z\) has continuous paths. So Itô's formula applies to the representation as in Proposition \hyperref[res:Proposition-52]{52}. By (L1) the coefficients are finite for almost every \(u\), almost surely, and the rest of the proof of Proposition \hyperref[res:Proposition-52]{52} uses only their values at a fixed \((u,\omega)\); Fubini restores the stated quantifiers. (c) On each \(I_m\) the chosen drifts are the value at \(T_m\) and the slope of \(-(R-c_m)+\sigma^S\cdot\int\sigma^S\), finite combinations, with coefficients continuous in \((u,T_m)\), of the \(b_k\), the products \(H^k\cdot H^l\), \(\bar\omega_k\), \(H^{0,0}\cdot V_{m'}\) and \(V_m\cdot V_{m'}\), and the components of \(Z_u\); \(|x\cdot y|\le(|x|^2+|y|^2)/2\) makes them locally integrable. The regularity of Section \hyperref[sec:2]{2} follows from \(\sup_{T\le H}|\sigma(u,T)|^2\le C(\sum_m|V_m(u)|^2+\sum_k|H^k(u)|^2)\) and the same bookkeeping for \(\sup_{T\le H}|\alpha(u,T)|\).\hfill$\square$

The proposition is conditional on the cited calculus and on the existence of \(Z\) and of the front-end scales with the stated integrability; in particular neither existence nor nonnegativity of the square-root variance processes of \cite{brace2024sofr} eqs. 29 and 30 is shown, and the remark on \(\sqrt v\) assumes a continuous \(v\). The true-martingale property of the discounted bonds is not addressed, since local integrability gives no exponential moments. Coefficients outside (L) have no stochastic integral in the cited calculus, and no finite-dimensional realization for general stochastic scales is given.

\section{Related work}\label{sec:4}

\textbf{Drift conditions with stochastic discontinuities.} \cite{fontana2018term} Theorem 3.7 gives necessary and sufficient conditions for a risk-neutral measure in a multi-curve HJM model whose forward rates carry a pure-jump term with jumps only at fixed dates and whose bond prices integrate against a measure that may have atoms at those dates; its conditions (iii) and (iv) are the ones specific to that setting. \cite{fontana2022term} Theorem 3.5 is the single-curve version with an overnight numéraire and separate sets of roll-over dates and expected jump dates; condition (ii) is the drift condition, condition (iv) the jump condition, and Remark 3.9 with Example 3.10 is the no-atoms case that Section \hyperref[sec:2]{2} adopts. \cite{kellerressel2018affine} Proposition 6.11 eq. 74 and Corollary 6.12 give the drift condition in affine-semimartingale form with a deterministic càdlàg clock. \cite{gehmlich2016dynamic} introduced the atoms-at-predictable-times device, in credit. \cite{fontana2022term} §3 notes that a risk-neutral measure ensures no asymptotic free lunch with vanishing risk for the large market of bonds of all maturities, citing \cite{cuchiero2014perspective}, whose Theorem 3.3 makes that condition equivalent to the existence of an equivalent separating measure and whose Section 6 shows that in a large market it need not give a \(\sigma\)-martingale measure. The conditions of Section \hyperref[sec:2]{2} are therefore sufficient for absence of arbitrage in that sense; nothing here shows them necessary.

\textbf{Jumps at fixed dates in short-rate and affine models.} \cite{kellerressel2018affine} §6 (Examples 6.15 and 6.16) has Gaussian jumps at fixed dates in Vasiček-type models; \cite{fontana2022term} §4.3 and §4.4 treat Hull--White with jumps at scheduled dates (Lemma 4.3, Propositions 4.5 and 4.6), and its Lemma 4.11 and Proposition 4.12 derive the Gaussian forward-measure tilt and a forward-looking caplet price without studying which meeting variances a chosen set of strikes and expiries identifies or what meeting-variance bounds follow when those prices are known only within intervals; \cite{fontana2025extended} a CIR process with jumps at deterministic dates; \cite{calvia2025short} short-rate models with fixed-time random-size jumps by a PDE approach; \cite{silva2024discretely} affine jump-diffusions with Skellam-distributed scheduled jumps and closed-form characteristic functions; \cite{kim2014jumps} a Gaussian affine model with jumps of the whole state at deterministic times, its Proposition 1 a backward recursion across the jump dates; \cite{piazzesi2005bond} a target rate that is a pure jump process with intensity elevated on meeting days, so jumps at random rather than deterministic times; \cite{das2002surprise} §4.2 likewise lets the jump probability of a Poisson--Gaussian model of the Fed funds rate depend on FOMC meeting days (Table 7); and \cite{andersen2024spike} models spikes and hikes in the overnight rate for SOFR caplets. In the scalar Markov construction of \cite{fontana2022term} Example 3.10, eq. 3.14, the Gaussian variance at each scheduled date is a free deterministic model parameter; after those parameters are fixed, the vector of meeting variances is independent of the current scalar state. None poses the consistency question for a curve family.

\textbf{Consistency and finite-dimensional realizations.} \cite{bjork1999interest} Theorem 4.1 characterizes invariance of a curve manifold under Wiener-driven HJM, Propositions 5.1 to 5.3 show Nelson--Siegel inconsistent with Ho--Lee and Hull--White, and Theorem 6.1 extends the criterion to marked-point-process jumps with an intensity-form compensator. \cite{filipovic1999note} Proposition 3.2 and Theorem 4.1 show that no nontrivial Itô process is consistent with Nelson--Siegel; spliced onto a scheduled-step front end with independent drivers, a Nelson--Siegel block with fixed exponent can carry only a level factor (Proposition \hyperref[res:Proposition-38]{38}). \cite{filipovic2000exponential} Theorem 3.2 constrains the exponents of a consistent bounded exponential-polynomial family to be constant, and Corollary 3.5 and Theorem 8.2 require an exponent-doubling pair for a nontrivial process; these restrict exponents only and permit polynomial factors, such as \(\lambda(x)=xe^{-\kappa x}\), and exponent sets such as \(\lambda,2\lambda,4\lambda\). \cite{bjork2001existence} gives Lie-algebra conditions for finite-dimensional realizations with a fixed state-to-curve map and autonomous state equations (its Definition 3.1). For a deterministic volatility \(\lambda(x)\) in time to maturity its Proposition 5.1 and Corollary 5.1 give a realization if and only if \(\lambda\) is quasi-exponential, and Proposition 6.1 the same for \(\gamma(r)\lambda(x)\); Proposition 6.2, with running time as a state, realizes every calendar-fixed shape that is \(C^1\) and nowhere zero, with no scheduled dates. \cite{tappe2010alternative} Corollary 8.1 characterizes affine realizations by the quasi-exponential condition, and Corollaries 7.1 and 7.3 treat volatility of constant direction; in these results the scale is a constant or a functional of the curve. \cite{bjork2004finite} lets the volatility depend on a separate Markov diffusion \(y\): for \(\sigma=\varphi(r,y)\lambda(x)\) with scalar \(y\), its Proposition 5.3 gives a realization if and only if \(\lambda\) is quasi-exponential when the curve and scale noises are not perfectly correlated, and Propositions 5.1 and 5.2 give the necessary form under orthogonal noise and sufficiency of quasi-exponential shapes. None of these covers a scheduled maturity breakpoint. \cite{sharef2004conditions} bounds the number of nontrivial factors for a biexponential-polynomial family by \(\min(n_1,\lfloor n_2/2\rfloor)+1\) (its first proposition of §4.2, proved here as Lemma \hyperref[res:Lemma-36]{36}), states a sufficient condition for that many factors, and proposes extensions of the Svensson family. Its central proposition of §4.2, that sufficient condition, is not established by its proof, whose reduced system carries the wrong sign on the \(e^{-2\beta x}\) terms in the diffusion matrix and which treats only diagonal diffusion matrices; as the reduced system states it the proposition is false (Proposition \hyperref[res:Proposition-39]{39}(e)), and with correlated factors its conclusion holds whenever \(z_{2,2n}\neq0\) or \(n_2=2n\), Proposition \hyperref[res:Proposition-39]{39}(b) giving the exact condition. \cite{kellerressel2024term} Theorems 2.5 and 2.7 and Corollary 2.6 classify the curve shapes that the Svensson family can keep under consistent dynamics with fixed exponents, complex shapes disappearing after a deterministic horizon; it counts no factors, does not address the sufficient condition of \cite{sharef2004conditions} §4.2, and has no scheduled breakpoints. \cite{tappe2012existence} gives affine realizations for Lévy-driven HJM, \cite{filipovic2012invariant} gives necessary and sufficient conditions for invariance of finite-dimensional submanifolds with boundary under HJM-type equations driven by a Wiener process and a compensated Poisson random measure, whose jumps are not at scheduled dates, \cite{fontana2024geometry} treats multi-curve HJM, \cite{filipovic2010term} and \cite{filipovic2010jump} give existence and positivity for HJM-type SPDEs with Wiener and Poisson noise, and \cite{filipovic2001consistency} is the monograph of the framework. \cite{benth2018approximation} Theorem 4.1 approximates forward-curve dynamics in commodity markets by arbitrage-free models with finite-dimensional state, with rates of convergence; its forward prices are martingales, so no drift map enters, and there are no scheduled dates (compare Proposition \hyperref[res:Proposition-51]{51}). Every jump in this literature is stochastically continuous.

\textbf{Separable and stochastic volatility.} \cite{brace2024sofr} §3.5 takes a one-factor volatility \(\sigma(t, T) = \chi(t)\varphi(T)\) with \(\chi\) stochastic (eq. 20), exponential \(\varphi\) and the accumulated-variance state (eq. 28), obtaining with a Heston variance (eq. 29) the three-state system of eq. 30 and the curve (Appendix A.1); §3.6 repeats this per meeting factor with the indicators of eq. 36, and eqs. 35 to 41 make the drift and the diffusion part of a short-rate meeting jump depend on the Heston variance, with within-factor leverage. Its drift is taken from the HJM result and its bond prices are obtained by integrating the curve, so, given solutions of its variance equations, Assumption 2.1 holds by construction and Assumption 2.2 trivially for every parameter value, and the only cross-meeting consequences are the fixed-parameter ones of Propositions \hyperref[res:Proposition-14]{14}, \hyperref[res:Proposition-16]{16} and \hyperref[res:Proposition-26]{26}. The paper does not prove that its square-root variance equations have solutions on a given basis, nor that its discounted bonds are true rather than local martingales. \cite{li2019continuous} gives Laplace-transform uniqueness and transition laws for continuous-state branching processes, and \cite{fontana2025extended} Lemma A.3 the support of a continuous image; they supply tools, not SOFR variance restrictions. \cite{brace2024sofr} footnote 28 attributes the Markov property under separable volatility to \cite{carverhill1994when}, whose Section \hyperref[sec:1]{1} and Proposition 2.1(iii) assume deterministic bond-price volatilities and stationarity in time to maturity. \cite{fontana2022term} Example 3.10 shows that the deterministic-volatility Cheyette class, with separable calendar volatility, stays finite-dimensional and time-inhomogeneous Markov under Gaussian scheduled jumps when the jump shape matches the diffusion shape; it has no maturity breakpoints. Proposition \hyperref[res:Proposition-28]{28} extends these calendar-shape results to meeting restarts, stochastic scales and continuous shapes and is not new in substance; it is a coefficient representation, not a realization.

\textbf{Stochastic-calculus inputs.} The cited calculus of Section \hyperref[sec:5.1]{5.1} comes from \cite{bauerschmidt2020stochastic}: predictability of adapted left-continuous processes (p. 14), the Itô integral with its isometry and local-martingale property (§3.1 to §3.3), the multivariate Itô formula (§3.5), strong existence and pathwise uniqueness under Lipschitz coefficients (§5.2), convergence of left-endpoint sums (p. 48) and an exponential-martingale criterion under a deterministic quadratic-variation bound (pp. 60--61); the Markov property of solutions with coefficients Borel in time is \cite{vanhandel2007stochastic} Theorem 5.2.2. A textbook account of this calculus is \cite{revuz1999continuous}. None of these treats the square-root equation (3.62): \cite{altay2013lecture} Corollary 2.19 proves its pathwise uniqueness but not existence, \cite{malham2008square} Proposition 1 quotes the transition law with positive mean reversion from \cite{cox1985theory}, and \cite{feller1951singular} Lemmas 8 and 9 solve its forward equation, zero mean reversion included, without constructing a solution on a given filtered space, so the solutions of (3.62) and (3.54) remain hypotheses. The Laplace-transform uniqueness of Section \hyperref[sec:3.14]{3.14} is \cite{li2019continuous} Theorem 1.1, and \cite{maghsoodi1996solution} gives closed-form solutions of the square-root equation with time-dependent coefficients, on which nothing here depends.

\textbf{SOFR and meeting-date models.} \cite{gellert2021short} builds a short-rate model with step volatilities at meeting dates (eq. 26), derives the risk-neutral solution with its ramp drift (eq. 28), states informally in §3.1.5 and §5 that piecewise-constant short-rate paths constrain the initial curve, and splices the step block onto a Vasiček block by independence (§3). The remark after its eq. 35 distinguishes diffusion-factor count from full curve-state dimension: factor values frozen at successive meetings remain in the state. \cite{brace2024sofr} continues it with the ramp drift eqs. 18 and 19, the stochastic-volatility finite-dimensional model of §3.5 and §3.6, principal-component dimension reduction (§4.1.1), joint calibration to futures and options (§4.2), and the drift term structure of Appendix C. Its eqs. 35 to 41 permit time-dependent calibration functions, but it does not characterize the meeting-variance vectors attainable as those scales vary or the restrictions that survive every such recalibration. For the deterministic-volatility family of Proposition \hyperref[res:Proposition-13]{13}, Proposition \hyperref[res:Proposition-25]{25} characterizes the attainable cone at fixed decay and loadings and its surviving inequalities; no such result is established here for stochastic volatility. Neither \cite{gellert2021short} nor \cite{brace2024sofr} asks whether its way of attaching the long-end block is forced. With deterministic or predictable bounded volatilities, one or several step and block factors, and fixed exponential-polynomial blocks, Propositions \hyperref[res:Proposition-34]{34}, \hyperref[res:Proposition-35]{35}, \hyperref[res:Proposition-41]{41}, \hyperref[res:Proposition-43]{43} and \hyperref[res:Proposition-44]{44} show that the independence splice of the first and the indicator-times-exponential terms of the second (eq. 36) are the only two constructions, provided a block uncorrelated with the aggregate maturity-step noise is counted as the independence splice; in the uncorrelated splice the front end frees only the level direction that the two families share (Propositions \hyperref[res:Proposition-38]{38}, \hyperref[res:Proposition-40]{40} and \hyperref[res:Proposition-45]{45}). \cite{heitfield2019inferring} fits a static piecewise-flat curve between meeting dates to futures, and \cite{gurkaynak2007market} §5, eqs. (5) to (8), unwinds monthly-average federal funds futures meeting by meeting into the expected rate after each of the next two meetings; (3.105) and Proposition \hyperref[res:Proposition-49]{49}(c) unwind option variances across expiry gaps in the same triangular way, for second moments and within a stated model. \cite{lyashenko2019looking} extends the LIBOR market model to backward-looking in-arrears rates, with the compounded rate of (3.70) (its eq. 5), a volatility that decays to zero over the accrual period (eqs. 13 and 14) and the futures convexity as an integrated drift (§6.1, eq. 23); its rates are diffusions, without scheduled jumps. \cite{backwell2021expected} is a short-rate model with expected jumps at scheduled dates whose sizes have a discrete law controlled by a state (§2.3) and unexpected jumps, its curve computed by pricing rather than assumed as a family. \cite{skov2021dynamic} fits arbitrage-free Nelson--Siegel models to SOFR futures and finds the convexity correction significant only beyond two years. \cite{pennanen2025statistical} is a statistical model under the physical measure.

\textbf{Bounds from finitely many prices.} \cite{bertsimas2002relation} bounds option prices and moments of one underlying from finitely many call prices, and \cite{beiglbock2011model} gives model-independent bounds from complete call surfaces of one asset at several dates; the panels of Sections \hyperref[sec:3.15]{3.15} to 3.19 have several accrual underlyings, discounting and expiry-dependent forward measures, and a variance functional as target, so neither supplies their bounds.

\textbf{Empirical side of the restrictions.} \cite{wright2020event} and \cite{zhong2026spanning} estimate event variance by specified empirical procedures, \cite{kim2014jumps}, \cite{skov2021dynamic} and \cite{collindufresne2002bonds} price jump risk, futures convexity and unspanned volatility, \cite{itkin2024semi}, \cite{lieu1990option}, \cite{chen1993pricing}, \cite{nastasi2018smile} and \cite{leung2014accounting} price options on futures and American options around scheduled events, and \cite{brace2024sofr} Table 2 does not publish the inputs needed to reproduce a calibrated meeting-variance matrix; none of them proves or refutes that equal specified SOFR option prices give equal event variances.

\section{Outlook}\label{sec:5}

This section states what is not done. Section \hyperref[sec:5.1]{5.1} lists what the results rest on; Sections \hyperref[sec:5.2]{5.2} to 5.5 take each group of results in turn, with what is known from the literature, what remains open, and the conjectures; Section \hyperref[sec:5.6]{5.6} is the empirical programme, which is not started. Assumptions 2.1 and 2.2 are in force throughout.

\subsection{What the results rest on}\label{sec:5.1}

\emph{Standing assumptions.} The drift condition and the jump martingale condition of Section \hyperref[sec:2]{2} are taken from \cite{fontana2022term} Theorem 3.5 as hypotheses. No result proves that they characterize absence of arbitrage in the present setting beyond what that theorem states.

\emph{Cited stochastic calculus.} Propositions \hyperref[res:Proposition-14]{14}, \hyperref[res:Proposition-16]{16}, \hyperref[res:Proposition-26]{26} to 31, 33 to 35, 37, 38, 41, 43, 44, 46, 47 and 51 to 53, and Lemma \hyperref[res:Lemma-32]{32}, use stochastic integration, Itô's formula, strong solutions of Lipschitz equations, predictability of continuous adapted processes, the Markov property of solutions, left-endpoint approximation of the Itô integral or an exponential-martingale criterion. These are cited from \cite{bauerschmidt2020stochastic} and \cite{vanhandel2007stochastic}, and Proposition \hyperref[res:Proposition-46]{46}(b) also cites \cite{filipovic2000exponential} Theorem 3.2 and Corollaries 3.4 and 3.5. They are not proved here. In the formalization each enters as a field of a hypothesis structure, checked against its source. Every such structure has only a degenerate instance, a model of every field on a trivial base, or, where the fields assume Brownian drivers, a vacuous one, accompanied by separate checks that the premises and conclusions are inhabited. Neither kind of instance constructs a nontrivial model of the fields, so these results are conditional on the cited inputs, in the text and in Lean alike. Only the standing-hypothesis structure of Section \hyperref[sec:2]{2} has a nonzero instance.

\emph{Solutions taken as given.} Propositions \hyperref[res:Proposition-14]{14}, \hyperref[res:Proposition-16]{16} and \hyperref[res:Proposition-26]{26} assume nonnegative solutions of their variance equations, and Proposition \hyperref[res:Proposition-26]{26} a moment bound; the square-root equations are not solved here. Where a parameter process \(Z\) of a block appears, as in Propositions \hyperref[res:Proposition-37]{37} to 40 and 45 to 47, its existence with the stated coefficients is not shown; those propositions are statements about any process that has them, or algebra at one parameter point.

\emph{Unconditional results.} Theorem \hyperref[res:Theorem-1]{1}, Lemmas \hyperref[res:Lemma-3]{3} to 7, Propositions \hyperref[res:Proposition-8]{8} to 10, Theorem \hyperref[res:Theorem-11]{11} and Remark 12 use only the standing assumptions. Propositions \hyperref[res:Proposition-15]{15}, \hyperref[res:Proposition-17]{17}, \hyperref[res:Proposition-19]{19}, \hyperref[res:Proposition-20]{20} and \hyperref[res:Proposition-23]{23} to 25, Corollaries \hyperref[res:Corollary-21]{21} and \hyperref[res:Corollary-22]{22} and Lemma \hyperref[res:Lemma-18]{18} are about explicit Gaussian models and ideal payoffs, and Proposition \hyperref[res:Proposition-48]{48} uses only conditional expectations and the uniqueness of laws from Laplace transforms, and Lemmas \hyperref[res:Lemma-36]{36} and \hyperref[res:Lemma-42]{42} and Propositions \hyperref[res:Proposition-39]{39}, \hyperref[res:Proposition-40]{40} and \hyperref[res:Proposition-45]{45} are algebra at a parameter point with, in two cases, a statement along right-continuous paths. Proposition \hyperref[res:Proposition-13]{13} is proved in the text from Wiener integrals of deterministic integrands; its formalization assumes the resulting variance-integral representation (Appendix \hyperref[sec:A]{A}). Propositions \hyperref[res:Proposition-49]{49}(a), 49(b), 50(a) and 50(b) rest on the same Wiener-integral facts together with stochastic Fubini for deterministic integrands; their formalizations take them as one explicit hypothesis each, whose only instance has no diffusion, and the remaining parts are unconditional.

\subsection{The step-plus-ramp family}\label{sec:5.2}

The step-plus-ramp family (2.5) is a consistent front-end curve for carry and roll-down.

Theorem \hyperref[res:Theorem-11]{11} answers the deterministic-volatility classification: consistency is equivalent to the differentiated HJM drift and the tilted-Gaussian jump conditions, and Proposition \hyperref[res:Proposition-10]{10} gives a nonzero multi-date instance. Condition (ii) of the theorem reduces to "Gaussian jump laws with matching covariances", eq. (3.43), when the level jumps are jointly Gaussian; that form is a natural guess, and Proposition \hyperref[res:Proposition-9]{9}(d) shows that its necessity direction fails, since the jump martingale condition does not force joint Gaussianity. Explicit instances of the drift condition are \cite{gellert2021short} eq. 28, a member of \(S^+\) with deterministic pairwise coefficients \(C_{qi}(t) = -\xi_q \xi_i \sum_j \lambda_{qj} \lambda_{ij}\,(t \wedge x_q \wedge x_i)\) for constant step volatilities and without curve jumps, and \cite{brace2024sofr} eq. 18 with time-dependent parameters; the jump instances in \cite{fontana2022term} Example 3.10 eq. 3.14, \cite{kellerressel2018affine} Examples 6.15 and 6.16 and \cite{kim2014jumps} Proposition 1 are all jointly Gaussian. The representation of the forward rate before a scheduled date as the exponentially tilted mean of the jump,

\begin{equation*}
f(t, T) = r_t + \frac{E_t[J_n e^{-\tau J_n}]}{E_t[e^{-\tau J_n}]}, \qquad \tau = T - T_n, \tag{5.1}\label{eq:5.1}
\end{equation*}

is Proposition \hyperref[res:Proposition-48]{48}(d): it holds exactly in the model in which the jump at \(T_n\) is the only randomness up to the next date, with the conditional law of the jump given \(\mathcal F_{T_n-}\), which (H) makes equal to the law given \(\mathcal F_t\). It is implicit in \cite{fontana2022term} Propositions 4.2 and 4.5 and in \cite{kim2014jumps} Proposition 1 for Gaussian jumps. By Proposition \hyperref[res:Proposition-48]{48}(b) the compensating shape is the derivative of the conditional cumulant generating function of the level jump, so ramps linear in \(T\) force a Gaussian level jump (Proposition \hyperref[res:Proposition-48]{48}(c)), with the tilt on later intervals as the correction (Proposition \hyperref[res:Proposition-9]{9}). Jumps in multiples of a tick, as in \cite{silva2024discretely}, are excluded by (ii); the difference between a two-point jump law and the Gaussian ramp with the same mean and variance is at most \(|\Delta|^4\tau^3/48\) at time to maturity \(\tau\) past the meeting, for a symmetric jump of size \(\Delta\) (Proposition \hyperref[res:Proposition-48]{48}(c), (3.185)), of order \(10^{-8}\) basis points at one year for \(\Delta\) of 25 basis points, and convexity effects of this kind are below quote precision (\cite{brace2024sofr} Appendix C; \cite{skov2021dynamic}), so the distinction is a matter of classification, not of pricing. With multiplicative stochastic step volatility \(\sigma(t, T) = \sum_n s_n \sqrt{v(t)}\, \mathbf 1\{T \ge T_n\}\), \(s_n\) deterministic, the drift is again a ramp with the realized coefficient \(C_{nm}(t) = \int_0^t \langle s_n, s_m \rangle\, v(s)\,ds\), so \(S^+\) stays closed with \(C\) adapted and of finite variation in place of deterministic. For constant \(s_n\) and a predictable \(v\) with \(\int_0^H v(s)\,ds<\infty\) almost surely, the local form of this is Proposition \hyperref[res:Proposition-28]{28}(c) with \(\varphi_j=1\), \(\chi_j=\sqrt v\) and \(a_{j,k,m}=\sum_{n\le m}s_{n,j}\), conditional on the cited Itô calculus: the slope on \(I_m\) is the initial slope plus \(|\sum_{n\le m}s_n|^2\int_0^t v(s)\,ds\). Time-dependent step volatilities and the pairwise coefficients \(C_{nm}\) themselves, of which only the sums attached to a date are identifiable, are not covered there.

\subsection{Cross-meeting restrictions on event variances}\label{sec:5.3}

The programme behind this paper sought, besides the consistent front-end curve, cross-meeting restrictions on option-implied event variances usable as a tradeable signal: deviations from a restriction forced by consistency would be event-volatility calendar spreads and flies. The results of Sections \hyperref[sec:3.10]{3.10} to 3.23 say that consistency theory does not supply such a signal. Remark 12 rules out a restriction inferred from state dimension alone: scalar-state members of the jointly Gaussian class fill the nonnegative orthant of meeting variances as their parameters vary, so any restriction has to come from a fixed model's parameters and from the payoffs observed. Proposition \hyperref[res:Proposition-13]{13} gives the fixed-parameter Gaussian cone when loadings, decay and constant scales are fixed, and shows why meeting-dependent scale freedom removes every common linear equality; Proposition \hyperref[res:Proposition-14]{14}, conditional on its stochastic-solution hypotheses, gives the affine variance map of the stochastic-volatility specification, and Proposition \hyperref[res:Proposition-16]{16} makes its support and affine hull exact at zero variance mean reversion, with affine-hull dimension at most two for the printed constant parameters. Proposition \hyperref[res:Proposition-26]{26} extends the exact support to nonnegative mean reversion and piecewise-constant volatility of volatility: all three factors in the printed time-dependent column are active, while the affine-hull dimension remains the rank of its incompletely specified loading matrix. These are restrictions of a fixed model, not of consistency, and they change under recalibration. Proposition \hyperref[res:Proposition-25]{25} makes exact what survives recalibration of the scales: the attainable set is a polyhedral cone whose facets are the surviving inequalities; with constant ordinal loadings they are the consecutive-meeting ratio bounds set by the largest decay, they survive joint recalibration of scales and bounded decays, and free loadings or unbounded decays leave only nonnegativity.

For deterministic volatility, ramp coefficients equal the risk-neutral conditional covariances of the meeting jumps plus inter-meeting diffusion, so curve convexity and risk-neutral event variance are the same quantity; this is implicit in \cite{gellert2021short} eqs. 28 and 32, \cite{fontana2022term} Proposition 4.6 and \cite{kim2014jumps} Proposition 1, and is not stated as an identity in any source. Under multiplicative stochastic volatility it is not an identity: the local slope is the initial convexity plus realized accumulated variance, a finite-variation process (Proposition \hyperref[res:Proposition-28]{28}(c), for loadings constant between meetings and conditional on the cited Itô calculus), while expected remaining variance has a martingale part, so the difference is a quantity of its own (the martingale part of expected remaining variance is not proven).

At a fixed observation time, Lemma \hyperref[res:Lemma-7]{7} and \(-\log P(t,T)=\int_t^T f(t,u)\,du\) show that three exact positive bond prices at distinct maturities inside one scheduled interval recover its current affine slope. This does not identify a future short-rate event variance from a listed panel.

Propositions \hyperref[res:Proposition-15]{15} and \hyperref[res:Proposition-17]{17} show in the scalar Gaussian model that a full ideal bond-call surface at one expiry identifies only cumulative variance, that expiries separating successive meetings identify every event variance, and that deterministic price intervals give an exact feasibility test and sharp bounds. Proposition \hyperref[res:Proposition-19]{19} adds ideal compounded futures and European cash calls: equality of a finite panel is a matrix equation in the event variances, full column rank is necessary and sufficient for identification, and two fixed strikes per contract suffice (Lemma \hyperref[res:Lemma-18]{18}); its positive pair shows that every later-window quote and surface can leave individual variances unidentified. Corollaries \hyperref[res:Corollary-21]{21} and \hyperref[res:Corollary-22]{22} show that exact discrete compounding and American cash exercise, after the last meeting or across exactly one unrevealed meeting, do not remove such a pair, and Proposition \hyperref[res:Proposition-20]{20} shows why: every contract settled after a date, with American exercise across any number of unrevealed meetings, sees the revealed variances only through their sum, initial futures quotes add the date-weighted sum, and two quotes leave exactly a polytope of earlier variances with sharp coordinate bounds. Proposition \hyperref[res:Proposition-23]{23} shows that futures-style margining removes the early-exercise value, as is known for futures-style options in general, and lets the surfaces identify each unrevealed variance inside the window, but leaves the same polytope of revealed variances. Proposition \hyperref[res:Proposition-24]{24} closes the listed-calendar question for the ideal payoffs: on the calendar of every serial and quarterly expiry the surfaces identify every meeting variance of 2026 and 2027, the limit being precision rather than rank, since a tick-size premium error propagates in proportion to the accumulated standard deviation; with only the nearest serials listed, identification collapses to sums from mid-2026 on. Proposition \hyperref[res:Proposition-49]{49} adds a piecewise-constant diffusion between meetings and an arbitrary initial curve: the surfaces do not depend on the initial curve, the diffusion is identified from the expiry gaps without a meeting, and on the listed calendar a constant diffusion is identified with every meeting variance, but at 50 to 100 basis points a year its error leaves only the first meetings practically identified at small reference sizes. Proposition \hyperref[res:Proposition-50]{50} keeps these answers for a known maturity shape: exactly for a loading that starts at the next meeting, and for exponential decay at every rate but one, with somewhat better precision. Whether a crossed variance is itself unidentified, what happens with two unrevealed meetings or accrual across a meeting, and whether identification survives listed settlement and exercise rules, an unknown diffusion shape, meeting shocks that are not parallel and noisy prices, are open; a transfer to listed instruments needs the contract family with its accrual state, exercise payoff, discounting and allowed curve and parameter variation specified before the question is posed.

\textbf{Conjecture 1.} For a fixed consistent model and a specified collection of listed SOFR futures and options with their settlement, exercise and discounting rules, equal prices among admissible states and allowed parameter choices imply equal short-rate event-variance vectors.

\cite{gellert2021short} (remark after eq. 35) notes that full curve-state dimension need not equal diffusion-factor count. \cite{brace2024sofr} reduces the diffusion factors by principal components and calibrates jointly to futures and options without reporting a cross-meeting restriction. \cite{wright2020event} and \cite{zhong2026spanning} estimate event variance under specified empirical procedures but do not prove the listed-contract implication. An infeasible restriction would reject the model; it would not by itself be an arbitrage certificate.

\emph{The listed-market gap.} Every identification result of Sections \hyperref[sec:3.13]{3.13} to 3.21, 3.24, 3.43 and 3.44 concerns ideal payoffs priced exactly in a stated model. Propositions \hyperref[res:Proposition-24]{24}, \hyperref[res:Proposition-49]{49} and \hyperref[res:Proposition-50]{50} take the expiry calendar and tick size of the listed contracts as inputs, but not their settlement, margining, exercise into futures or accrual conventions, nor noisy quotes; only Propositions \hyperref[res:Proposition-49]{49} and \hyperref[res:Proposition-50]{50} allow an unknown initial curve. None of these results is a statement about listed SOFR instruments; the transfer is Conjecture 1, and it is open.

\subsection{Stochastic volatility and finite-dimensional closure}\label{sec:5.4}

With separable factor volatilities \(\sigma_j(t,T)=\chi_j(t)\varphi_j(T)\), arbitrary continuous deterministic maturity shapes and meeting-dependent loadings, Proposition \hyperref[res:Proposition-28]{28} shows, conditional on the cited Itô calculus, that the drift condition holds and the current curve is the initial curve plus at most \(2d(N+1)\) coefficients on the shape pairs \((\varphi_j,\varphi_j\Phi_{j,m})\) restarted at each meeting, the drift-generated coefficient being realized accumulated variance of finite variation. In \cite{brace2024sofr} eq. 30 the state \(\Phi\) has no Brownian term, and its Appendix A.1 shows that its variance state enters neither the curve nor the bond price. For exponential \(\varphi\), the drift-generated shape gives the exponent pair \((\lambda,2\lambda)\) that \cite{filipovic2000exponential} Theorem 8.2 requires for a nontrivial bounded family. Proposition \hyperref[res:Proposition-28]{28} is a representation of the current curve, not a Markov realization. Proposition \hyperref[res:Proposition-29]{29} adds a Lipschitz scale state and obtains a finite-dimensional, time-inhomogeneous Markov realization with running time as a state, an extension of known calendar-maturity realizations; it is not the realization in time to maturity of Conjecture 2. Time homogeneity is untouched, and the separable form \(\chi_j(t)\varphi_j(T)\) cannot carry a classification of shapes. With running time as a state every calendar-fixed shape that is \(C^1\) and nowhere zero has a realization in the time-varying sense of \cite{bjork2001existence} Proposition 6.2 (Section \hyperref[sec:4]{4}). A single nowhere-zero shape whose translates agree up to scalar normalization, \(\varphi(a+x)/\varphi(a)\) independent of \(a\), is exponential by an elementary calculation, so imposing translation invariance on a shape anchored at the current time decides the answer by the setting, not by a theorem.

The question that remains uses the realizations in time to maturity defined at the start of Section \hyperref[sec:3.27]{3.27}, which serve every schedule. Without meetings and with a constant scale the answer is known: a realization exists exactly for quasi-exponential \(\lambda_j\) (\cite{bjork2001existence} Proposition 5.1 and Corollary 5.1, \cite{tappe2010alternative} Corollary 8.1). So exponential shapes are not the only admissible ones, already without meetings: \(\lambda(x)=xe^{-\kappa x}\) is quasi-exponential and not exponential (\cite{bjork2001existence} Corollary 5.1).

The loadings matter as well. Proposition \hyperref[res:Proposition-30]{30} shows that with the constant shape \(\lambda\equiv1\), which is quasi-exponential, a realization serving every schedule forces the ordinal loadings to satisfy a linear recurrence, and that harmonic cumulative loadings admit none. So quasi-exponential shapes alone do not suffice under this definition, already for an exponential shape; the obstruction comes from the loadings.

\textbf{Conjecture 2 (quasi-exponential shapes and recurrent loadings).} Under this definition, a realization exists for every schedule and every scale exactly when each \(\lambda_j\) is quasi-exponential and each sequence of ordinal loadings satisfies a linear recurrence (3.141).

For one factor, Proposition \hyperref[res:Proposition-31]{31} gives sufficiency for every bounded Lipschitz function of an autonomous Lipschitz scale state, with any correlation, and Propositions \hyperref[res:Proposition-30]{30} and \hyperref[res:Proposition-33]{33} give necessity of the recurrence for every nonzero quasi-exponential shape with a constant scale; so the one-factor constant-scale case of the conjecture's recurrence condition is settled (Proposition \hyperref[res:Proposition-33]{33}(c)). Necessity of the recurrence under a stochastic scale, necessity of quasi-exponential shapes when meetings are present, several factors, and whether the requirement that one realization serve schedules with arbitrarily many meetings is the right scope for the SOFR use, are open. On a finite horizon, Proposition \hyperref[res:Proposition-51]{51} gives a quantitative substitute: a supplied recurrent approximation with loading error \(\varepsilon\) gives a realized model whose forward-curve and bond-price mean-square errors are of order \(\varepsilon^2\), with schedule-free constants when the loadings are bounded; which loading sequences have accurate recurrent approximations of small order is open. Without meetings, a scale driven by a separate scalar Markov diffusion gives the same quasi-exponential answer: \cite{bjork2004finite} Proposition 5.3 proves it under its smoothness and non-degeneracy assumptions when the curve and scale noises are not perfectly correlated, and gives sufficiency when they are.

Proposition \hyperref[res:Proposition-27]{27} disproves the proposed exclusion of every bounded piecewise exponential-polynomial family when the variance state enters the curve. Its three-exponent family closes because the new \(4\lambda\) coefficient has finite variation, even though the drift-generated \(2\lambda\) coefficient has a martingale part. It does not classify the exponent sets that can close or settle Conjecture 2. \cite{brace2024sofr} §3.5 and §3.6 is the exponential instance relevant to that conjecture. The Wiener-driven criteria of \cite{bjork2001existence} and \cite{tappe2010alternative}, the Lévy-driven ones of \cite{tappe2012existence} and the multi-curve ones of \cite{fontana2024geometry} have no predictable jump times.

\subsection{The splice}\label{sec:5.5}

Two constructions attach a smooth long-end block to a scheduled-step front end. \cite{gellert2021short} §3 (eqs. 4 to 13 and 48 to 53) drives the step block and a Vasiček block by independent Brownian motions, so forward rates and bond prices factorize; independence is assumed there. \cite{brace2024sofr} eq. 36 puts the meeting indicator and the exponential decay in the same factor, so bond prices (their eq. 43) carry exponential integrals restarted at each meeting, and their §5.2 (eqs. 84 to 87) shows the piecewise meeting weights can mimic an exponential mean-reversion profile. \cite{fontana2022term} Example 3.10 gives the matching condition between continuous and jump shapes for a Cheyette block to remain one-dimensional, and \cite{bjork1999interest} Theorem 6.1 condition 78 is the general closure-under-jumps requirement for intensity-form jumps.

\textbf{Conjecture 3.} Beyond the scope of Propositions \hyperref[res:Proposition-34]{34} to 47, 52 and 53, in particular for varying exponents in the correlated splice and for unbounded scales outside the class of Proposition \hyperref[res:Proposition-53]{53}, a fixed-size exponential-polynomial block attached to a scheduled-step front end is consistent if and only if either the block's noise, apart from its level, is uncorrelated with the front end's aggregate maturity-step noise wherever the step volatility jumps, or the family includes the indicator-times-exponential terms.

\emph{What is proved.} In the following scope the statement of Conjecture 3 is proved, the "if" direction by the two constructions and the "only if" direction by the propositions cited: deterministic volatilities (Propositions \hyperref[res:Proposition-34]{34} and \hyperref[res:Proposition-35]{35}), predictable bounded scales on the step factor, the block and the correlation (Proposition \hyperref[res:Proposition-41]{41}), a block volatility driven by the block's own state for an observable block (Proposition \hyperref[res:Proposition-43]{43}), and several step and block factors (Proposition \hyperref[res:Proposition-44]{44}), for quasi-exponential blocks without an exponent zero and, in the correlated splice, for blocks with a full polynomial part, with one step factor (Proposition \hyperref[res:Proposition-47]{47}) and with several (Proposition \hyperref[res:Proposition-52]{52}), the last also for locally square-integrable noises, which include Heston-type scales, conditional on their existence (Proposition \hyperref[res:Proposition-53]{53}). Whether the block's own restrictions survive the splice is answered for the uncorrelated splice: the front end absorbs nothing of a biexponential-polynomial block's equation (Proposition \hyperref[res:Proposition-37]{37}), or of an exponential-polynomial block's with positive varying exponents (Proposition \hyperref[res:Proposition-46]{46}), and only the affine part of a block with constant or linear terms (Propositions \hyperref[res:Proposition-38]{38}, \hyperref[res:Proposition-40]{40} and \hyperref[res:Proposition-45]{45}), so that only a level is freed.

\emph{Conventions.} Three parts of this classification are counting conventions, not results. A block whose noise is orthogonal to the aggregate maturity-step noise of several step factors is counted as the independence splice, although single step factors may be correlated with it (Proposition \hyperref[res:Proposition-44]{44}); a block level correlated with the step noise is counted as part of the front end, since constants lie in both families (Proposition \hyperref[res:Proposition-47]{47}); and a correlation switched off wherever the step factor has maturity steps is counted as the independence splice locally in time (Proposition \hyperref[res:Proposition-41]{41}). A reader who counts any of these separately obtains a further construction, which the cited proposition characterizes, for Proposition \hyperref[res:Proposition-44]{44} only up to the affine condition on the block residual plus the cross term in its part (a); with the conventions there are only the two. Proposition \hyperref[res:Proposition-52]{52} states the correlated splice with several step factors and a full polynomial part without these conventions, naming the covariance and drift restrictions directly.

\emph{What is not done.} The indicator-times-exponential branch is shown consistent here for deterministic and bounded predictable scales (Propositions \hyperref[res:Proposition-34]{34}, \hyperref[res:Proposition-35]{35} and \hyperref[res:Proposition-41]{41}, part (c)(ii)). For the Heston-scale model of \cite{brace2024sofr} eqs. 35 to 41 its realizability is taken from that paper and not verified here: that paper proves neither that its square-root variance equations have solutions nor that its discounted bonds are true martingales (Section \hyperref[sec:4]{4}), and making the model rigorous would need sourced existence and martingale inputs, which are not pursued. Also open are correlated splices with varying exponents, polynomial parts with structurally absent low-degree coefficients, front ends with a diffusive ramp compensated by the block, unbounded scales outside the class of Proposition \hyperref[res:Proposition-53]{53}, and consistency from every initial time rather than from time zero. The use is splice-zone relative value.

\subsection{Empirical programme}\label{sec:5.6}

Not started. Planned: fit the family of Theorem \hyperref[res:Theorem-11]{11} to SOFR futures and options, report residuals against the step-plus-spline construction, and test Conjecture 1 for a specified contract collection. Confounds to handle: quarter-end and month-end SOFR spikes, which are scheduled-date effects and fit the framework; the FF--SOFR basis; futures margin funding; compounded-rate settlement; transaction costs of a quarter to half a basis point.

\appendix

\appsection{Machine verification}\label{sec:A}

The results through Remark 12 have formal statements and proofs in Lean 4 against Mathlib, except Corollary \hyperref[res:Corollary-2]{2}. Propositions \hyperref[res:Proposition-13]{13} and \hyperref[res:Proposition-14]{14} have the conditional formalizations described below; their stochastic derivations are not formalized. Proposition \hyperref[res:Proposition-15]{15} has a full formalization for its stated scalar Gaussian model and ideal European bond calls. Proposition \hyperref[res:Proposition-16]{16} is fully formalized, with the standard stochastic calculus results it uses taken as fields of audited hypothesis structures, and the solutions of its variance equations as hypotheses. Proposition \hyperref[res:Proposition-17]{17} is fully formalized for Proposition \hyperref[res:Proposition-15]{15}'s ideal prices, deterministic interval feasibility, sharp individual bounds and local inverse stability. Proposition \hyperref[res:Proposition-19]{19} is fully formalized for its continuously compounded futures and ideal European cash calls. Corollary \hyperref[res:Corollary-21]{21} is fully formalized for exact discrete bond-implied compounding and its late American cash exercise problem. Corollary \hyperref[res:Corollary-22]{22} is fully formalized for the exercise window crossing exactly one unrevealed meeting. Proposition \hyperref[res:Proposition-20]{20} is fully formalized for its stated payoffs of the post-window state. Proposition \hyperref[res:Proposition-23]{23} is fully formalized for the futures-style contract on the compounded rate. Proposition \hyperref[res:Proposition-24]{24} is fully formalized on its calendar, with the assumptions (A1) to (A4) entering as the calendar, tick sizes and reference sizes. Proposition \hyperref[res:Proposition-25]{25} is fully formalized on the deterministic cone of Proposition \hyperref[res:Proposition-13]{13}. Proposition \hyperref[res:Proposition-26]{26} is fully formalized in the same way, with the predictability result also taken as a field, the solutions of its variance equations as hypotheses, and the given noise and forward-loading integrands assumed in the admissible Itô domain; it covers the conditional source process, its regular conditional laws and the three-factor printed specialization, and existence of the variance-state solutions and calibrated rank three are not proved. Proposition \hyperref[res:Proposition-27]{27} is formalized under the stated two-filtration calculus inputs, from which it derives the stopped strong solution, actual discounted-bond martingales, variance identities and nonzero short-rate jump; existence of the calculus structures on the Brownian basis is not proved. Proposition \hyperref[res:Proposition-28]{28} is formalized under the supplied calculus, except that "semimartingale" for \(M_{j,k}\) is paper-only. Proposition \hyperref[res:Proposition-29]{29} is formalized in parts (a) to (e) under the supplied calculus and the audited strong-existence, predictability and time-dependent Markov inputs; the paragraph after its statement is descriptive. Proposition \hyperref[res:Proposition-30]{30} is formalized under the supplied calculus, with only the parts of the realization notion its proof uses. Proposition \hyperref[res:Proposition-31]{31} is formalized under the supplied calculus and predictability input, for any adapted scale state with continuous paths. Lemma \hyperref[res:Lemma-32]{32} and Proposition \hyperref[res:Proposition-33]{33} are formalized under the supplied calculus; the sufficiency half of Proposition \hyperref[res:Proposition-33]{33}(c) is formalized in the one-time sense of the rank bounds, the full realization being Proposition \hyperref[res:Proposition-31]{31}. Proposition \hyperref[res:Proposition-34]{34} is formalized under the supplied calculus, Proposition \hyperref[res:Proposition-35]{35} under the supplied calculus with controllability and \(c\neq0\) as standing hypotheses, Lemma \hyperref[res:Lemma-36]{36} as a deterministic statement, Propositions \hyperref[res:Proposition-37]{37} and \hyperref[res:Proposition-38]{38} under the supplied calculus, Propositions \hyperref[res:Proposition-39]{39} and \hyperref[res:Proposition-40]{40} as deterministic statements with pathwise parts, Proposition \hyperref[res:Proposition-41]{41} under the supplied calculus with the standing hypotheses of Proposition \hyperref[res:Proposition-35]{35}, Lemma \hyperref[res:Lemma-42]{42} and Propositions \hyperref[res:Proposition-43]{43} and \hyperref[res:Proposition-44]{44} under the supplied calculus, Proposition \hyperref[res:Proposition-45]{45} at a parameter point, Proposition \hyperref[res:Proposition-46]{46} under the supplied calculus and, for (b), the cited exponent restrictions, Proposition \hyperref[res:Proposition-47]{47} under the supplied calculus, Proposition \hyperref[res:Proposition-48]{48} without conditions, Proposition \hyperref[res:Proposition-49]{49} with (a) and (b) under an explicit Gaussian-law hypothesis and (c) to (e) without conditions, Proposition \hyperref[res:Proposition-50]{50} with (a) and (b) under the corresponding hypothesis with a shape and (c) and (d) without conditions, Proposition \hyperref[res:Proposition-51]{51} with its deterministic estimates without conditions, (a) under the supplied calculus and (b) and (c) also assuming a Brownian driver, Proposition \hyperref[res:Proposition-52]{52} with Assumption 2.1 as an explicit hypothesis and no hypothesis structure, and Proposition \hyperref[res:Proposition-53]{53} in the same way. The hypothesis structures behind Propositions \hyperref[res:Proposition-27]{27}, \hyperref[res:Proposition-29]{29} and \hyperref[res:Proposition-46]{46}(b) have vacuous instances, because their fields assume Brownian drivers that the trivial base lacks, with separate checks that their premises and conclusions are inhabited; the other calculus structures have degenerate instances, and \(\texttt{HJMScheduled}\) a nonzero one. Neither kind gives a nontrivial model of its fields, so every result resting on them is conditional on the cited inputs. Each formal statement is a proposition in a module under \(\texttt{lean/Standalone/}\), which imports Mathlib only and can be read without the proofs, and is proved in a module under \(\texttt{lean/Novel/}\). The axiom audit of every Lean proof reports only \(\texttt{propext}\), \(\texttt{Classical.choice}\) and \(\texttt{Quot.sound}\); no proof uses \(\texttt{sorry}\). Corollary \hyperref[res:Corollary-2]{2} uses Fubini's theorem and the vanishing of the Itô integral of a null integrand and is not formalized.

Statements live in the namespace \(\texttt{Standalone}\) and proofs in the namespace \(\texttt{Novel}\). For each module \(\texttt{X}\) in the table, the statement is \(\texttt{X.statement}\) in \(\texttt{lean/Standalone/X.lean}\), which imports Mathlib only and can be read without the proofs, and its proof is \(\texttt{XProof.x}\) in \(\texttt{lean/Novel/XProof.lean}\), where \(\texttt{x}\) is \(\texttt{X}\) with a lower-case first letter. The first module of each row is the main declaration and the others are the targets it builds on. The proof modules \(\texttt{MGF}\) (Lemma \hyperref[res:Lemma-4]{4}) and \(\texttt{MGFUniqueness}\), \(\texttt{CondCalculus}\), \(\texttt{JumpLawCore}\) and \(\texttt{JumpLawExample}\) (Proposition \hyperref[res:Proposition-9]{9}) have no separate statement.

{\footnotesize\setlength{\tabcolsep}{4pt}
\begin{longtable}{lP{384pt}}
\toprule
\textbf{Result} & \textbf{Lean modules} \\
\midrule
\endhead
\bottomrule
\endlastfoot
Lemma \hyperref[res:Lemma-3]{3} & \(\texttt{LemmaA}\) \\
Lemma \hyperref[res:Lemma-4]{4} & \(\texttt{LemmaB}\) \\
Lemma \hyperref[res:Lemma-5]{5} & \(\texttt{StepJumpsVanish}\) \\
Lemma \hyperref[res:Lemma-6]{6} & \(\texttt{StepDriftVanish}\) \\
Theorem \hyperref[res:Theorem-1]{1} & \(\texttt{Theorem1}\) \\
Lemma \hyperref[res:Lemma-7]{7} & \(\texttt{PiecewiseAffineIntegral}\) \\
Proposition \hyperref[res:Proposition-8]{8} & \(\texttt{RampDriftIdentified}\) \\
Proposition \hyperref[res:Proposition-9]{9} & \(\texttt{JumpLawTiltedGaussian}\) \\
Proposition \hyperref[res:Proposition-10]{10} & \(\texttt{D2Instance}\) \\
Theorem \hyperref[res:Theorem-11]{11} & \(\texttt{Theorem2a}\) \\
Remark 12 & \(\texttt{D3EventVariances}\) \\
Proposition \hyperref[res:Proposition-13]{13} & \(\texttt{GaussianMeetingVarianceCone}\) \\
Proposition \hyperref[res:Proposition-14]{14} & \(\texttt{StochasticMeetingVariance}\) \\
Proposition \hyperref[res:Proposition-15]{15} & \(\texttt{BondOptionMeetingVariances}\) \\
Proposition \hyperref[res:Proposition-16]{16} & \(\texttt{ZeroMeanReversionUpstreamBridge}\), \(\texttt{ZeroMeanReversionVarianceSupport}\) \\
Proposition \hyperref[res:Proposition-17]{17} & \(\texttt{BondOptionPriceIntervals}\) \\
Lemma \hyperref[res:Lemma-18]{18} & \(\texttt{CompoundedFuturesIdentification}\) (within Proposition \hyperref[res:Proposition-19]{19}) \\
Proposition \hyperref[res:Proposition-19]{19} & \(\texttt{CompoundedFuturesIdentification}\) \\
Proposition \hyperref[res:Proposition-20]{20} & \(\texttt{PreWindowVarianceAggregates}\) \\
Corollary \hyperref[res:Corollary-21]{21} & \(\texttt{LateAmericanExercise}\) \\
Corollary \hyperref[res:Corollary-22]{22} & \(\texttt{CrossMeetingAmerican}\) \\
Proposition \hyperref[res:Proposition-23]{23} & \(\texttt{FuturesStyleExercise}\) \\
Proposition \hyperref[res:Proposition-24]{24} & \(\texttt{ListedSr3Identification}\) \\
Proposition \hyperref[res:Proposition-25]{25} & \(\texttt{SurvivingScaleRestrictions}\) \\
Proposition \hyperref[res:Proposition-26]{26} & \(\texttt{PositiveMeanReversionSupport}\) \\
Proposition \hyperref[res:Proposition-27]{27} & \(\texttt{BoundedVarianceCompletion}\), \(\texttt{BoundedVarianceBond}\), \(\texttt{BoundedVarianceExistence}\), \(\texttt{BoundedVarianceIndependence}\), \(\texttt{BoundedVarianceIntegralComparison}\), \(\texttt{BoundedVarianceIto}\), \(\texttt{BoundedVarianceMartingale}\), \(\texttt{BoundedVarianceState}\), \(\texttt{BoundedVarianceUniqueness}\) \\
Proposition \hyperref[res:Proposition-28]{28} & \(\texttt{SeparableMeetingAssembly}\), \(\texttt{SeparableMeetingShapes}\), \(\texttt{SeparableMeetingCoefficients}\), \(\texttt{SeparableMeetingIntegrals}\), \(\texttt{SeparableMeetingRepresentation}\) \\
Proposition \hyperref[res:Proposition-29]{29} & \(\texttt{SeparableMeetingMarkovSeveral}\), \(\texttt{SeparableMeetingScaleVersion}\), \(\texttt{SeparableMeetingMarkovCoefficients}\), \(\texttt{SeparableMeetingMarkovSolution}\), \(\texttt{SeparableMeetingMarkovRealization}\), \(\texttt{SeparableMeetingMarkovProperty}\) \\
Proposition \hyperref[res:Proposition-30]{30} & \(\texttt{MeetingLoadingObstruction}\), \(\texttt{MeetingLoadingHankel}\), \(\texttt{MeetingLoadingDimension}\), \(\texttt{MeetingLoadingLaw}\), \(\texttt{MeetingLoadingRank}\), \(\texttt{MeetingLoadingCurve}\) \\
Proposition \hyperref[res:Proposition-31]{31} & \(\texttt{ExternalScaleAssembly}\), \(\texttt{ExternalScaleDrift}\), \(\texttt{ExternalScaleStateP}\), \(\texttt{ExternalScaleStateQ}\), \(\texttt{ExternalScaleConditions}\), \(\texttt{RecurrentLoadingAssembly}\), \(\texttt{RecurrentLoadingAlgebra}\), \(\texttt{RecurrentLoadingDiffusion}\), \(\texttt{RecurrentLoadingDrift}\), \(\texttt{RecurrentLoadingRestartP}\), \(\texttt{RecurrentLoadingRestartQ}\), \(\texttt{RecurrentLoadingStateP}\), \(\texttt{RecurrentLoadingStateQ}\), \(\texttt{RecurrentLoadingXi}\), \(\texttt{RecurrentLoadingXiEquation}\) \\
Proposition \hyperref[res:Proposition-33]{33} & \(\texttt{RecurrentLoadingCharacterization}\), \(\texttt{RecurrenceNecessityAssembly}\), \(\texttt{RecurrenceNecessityAlgebra}\), \(\texttt{RecurrenceNecessityReduction}\), \(\texttt{RecurrenceNecessityCurve}\), \(\texttt{RecurrenceNecessityRank}\) \\
Lemma \hyperref[res:Lemma-32]{32} & \(\texttt{RecurrenceNecessityGaussian}\) \\
Proposition \hyperref[res:Proposition-34]{34} & \(\texttt{SpliceCrossTermAssembly}\), \(\texttt{SpliceCrossTermDrift}\), \(\texttt{SpliceCrossTermCurve}\), \(\texttt{SpliceCrossTermAlpha}\), \(\texttt{SpliceCrossTermAnalytic}\), \(\texttt{SpliceCrossTermNecessity}\), \(\texttt{SpliceCrossTermOpen}\), \(\texttt{SpliceCrossTermUncorrelated}\), \(\texttt{SpliceCrossTermConsistency}\), \(\texttt{SpliceCrossTermSufficiency}\) \\
Proposition \hyperref[res:Proposition-35]{35} & \(\texttt{SpliceQuasiExponentialAssembly}\), \(\texttt{SpliceQuasiExponentialCross}\), \(\texttt{SpliceQuasiExponentialCurve}\), \(\texttt{SpliceQuasiExponentialBlock}\), \(\texttt{SpliceQuasiExponentialAlgebra}\), \(\texttt{SpliceQuasiExponentialKey}\), \(\texttt{SpliceQuasiExponentialConsistency}\) \\
Lemma \hyperref[res:Lemma-36]{36} & \(\texttt{SharefFilipovicMaxFactors}\) \\
Proposition \hyperref[res:Proposition-37]{37} & \(\texttt{SharefFilipovicAssembly}\), \(\texttt{SharefFilipovicIndependence}\), \(\texttt{SharefFilipovicResidual}\), \(\texttt{SharefFilipovicSplit}\), \(\texttt{SharefFilipovicIto}\), \(\texttt{SharefFilipovicPartA}\), \(\texttt{SharefFilipovicPartB}\) \\
Proposition \hyperref[res:Proposition-38]{38} & \(\texttt{SpliceAffineOverlapAssembly}\), \(\texttt{SpliceAffineOverlapAbsorb}\), \(\texttt{SpliceAffineOverlapNS}\), \(\texttt{AnalyticPrimitive}\) \\
Proposition \hyperref[res:Proposition-39]{39} & \(\texttt{CorrelatedFactorsAssembly}\), \(\texttt{CorrelatedFactorsReduction}\), \(\texttt{NonnegPolySOS}\), \(\texttt{CorrelatedFactorsGram}\), \(\texttt{CorrelatedFactorsExists}\), \(\texttt{CorrelatedFactorsHankel}\), \(\texttt{CorrelatedFactorsAlways}\), \(\texttt{CorrelatedFactorsSign}\), \(\texttt{CorrelatedFactorsSmall}\), \(\texttt{CorrelatedFactorsProcess}\) \\
Proposition \hyperref[res:Proposition-40]{40} & \(\texttt{SvenssonSplice}\) \\
Proposition \hyperref[res:Proposition-41]{41} & \(\texttt{SpliceRandomScalesAssembly}\), \(\texttt{SpliceRandomScalesPath}\), \(\texttt{SpliceRandomScalesCurve}\), \(\texttt{SpliceRandomScalesNecessity}\), \(\texttt{SpliceRandomScalesSufficiency}\) \\
Lemma \hyperref[res:Lemma-42]{42} & \(\texttt{SpliceStateBlockObservability}\) \\
Proposition \hyperref[res:Proposition-43]{43} & \(\texttt{SpliceStateBlockAssembly}\), \(\texttt{SpliceStateBlockCross}\), \(\texttt{SpliceStateBlockAX01}\) \\
Proposition \hyperref[res:Proposition-44]{44} & \(\texttt{SpliceSeveralFactorsAssembly}\), \(\texttt{SpliceSeveralFactorsCore}\), \(\texttt{SpliceSeveralFactorsAX01}\), \(\texttt{SpliceSeveralFactorsSufficiency}\), \(\texttt{SpliceSeveralFactorsBlocks}\), \(\texttt{SpliceSeveralFactorsMinimal}\) \\
Proposition \hyperref[res:Proposition-45]{45} & \(\texttt{SplicePolynomialParts}\) \\
Proposition \hyperref[res:Proposition-46]{46} & \(\texttt{SpliceVaryingExponentsAssembly}\), \(\texttt{SpliceVaryingExponents}\), \(\texttt{SpliceVaryingExponentsRestrictions}\) \\
Proposition \hyperref[res:Proposition-47]{47} & \(\texttt{SpliceExponentZeroAssembly}\), \(\texttt{SpliceExponentZeroQuasiPoly}\), \(\texttt{SpliceExponentZeroCross}\), \(\texttt{SpliceExponentZeroAX01}\), \(\texttt{SpliceExponentZeroResidual}\) \\
Proposition \hyperref[res:Proposition-48]{48} & \(\texttt{JumpShapeAssembly}\), \(\texttt{JumpShapeCond}\), \(\texttt{JumpShapeCondB}\), \(\texttt{JumpShapeKernel}\), \(\texttt{JumpShapeProfile}\), \(\texttt{JumpShapeAffine}\), \(\texttt{JumpShapeOrigin}\), \(\texttt{JumpShapeExist}\), \(\texttt{JumpShapeTwoPoint}\), \(\texttt{JumpShapeForward}\), \(\texttt{JumpShapeGeneral}\), \(\texttt{JumpShapeLater}\) \\
Proposition \hyperref[res:Proposition-49]{49} & \(\texttt{DiffusionMeetingAssembly}\), \(\texttt{DiffusionMeetingGauss}\), \(\texttt{DiffusionMeetingConsistency}\), \(\texttt{DiffusionMeetingPricing}\), \(\texttt{DiffusionMeetingIdentities}\), \(\texttt{DiffusionMeetingFutures}\), \(\texttt{DiffusionMeetingRank}\), \(\texttt{DiffusionMeetingRankValue}\), \(\texttt{DiffusionMeetingInversion}\), \(\texttt{DiffusionMeetingCalendar}\), \(\texttt{DiffusionMeetingPrecision}\), \(\texttt{DiffusionMeetingInstance}\) \\
Proposition \hyperref[res:Proposition-50]{50} & \(\texttt{MaturityShapeAssembly}\), \(\texttt{MaturityShapeGauss}\), \(\texttt{MaturityShapeConsistency}\), \(\texttt{MaturityShapeIdentities}\), \(\texttt{MaturityShapePricing}\), \(\texttt{MaturityShapeRank}\), \(\texttt{MaturityShapeCalendar}\), \(\texttt{MaturityShapeS1}\), \(\texttt{MaturityShapeS2}\), \(\texttt{MaturityShapeListing}\), \(\texttt{MaturityShapeInstance}\) \\
Proposition \hyperref[res:Proposition-51]{51} & \(\texttt{RecurrentApproxAssembly}\), \(\texttt{RecurrentApproxEstimates}\), \(\texttt{RecurrentApproxMeanSquare}\), \(\texttt{RecurrentApproxPriceBounds}\), \(\texttt{RecurrentApproxBond}\), \(\texttt{RecurrentApproxRealization}\) \\
Proposition \hyperref[res:Proposition-52]{52} & \(\texttt{UnifiedSpliceAssembly}\), \(\texttt{UnifiedSpliceAlgebra}\), \(\texttt{UnifiedSpliceNecessity}\), \(\texttt{UnifiedSpliceConverse}\), \(\texttt{UnifiedSpliceStep0}\), \(\texttt{UnifiedSpliceQuantifiers}\), \(\texttt{UnifiedSpliceConverseAX01}\), \(\texttt{UnifiedSpliceExpPoly}\), \(\texttt{UnifiedSpliceDiag}\), \(\texttt{UnifiedSpliceSpecial}\) \\
Proposition \hyperref[res:Proposition-53]{53} & \(\texttt{SpliceLocalizationAssembly}\), \(\texttt{SpliceLocalization}\), \(\texttt{SpliceLocalizationConverse}\), \(\texttt{SpliceLocalizationRegularity}\) \\
\end{longtable}}

The declarations for Propositions \hyperref[res:Proposition-13]{13} and \hyperref[res:Proposition-14]{14} are conditional formalizations, described below.

The hypothesis structures imported by proof modules, with the results that use them:

\begin{itemize}
  \item \(\texttt{lean/Upstream/HJMScheduled.lean}\): Theorem \hyperref[res:Theorem-1]{1}, Proposition \hyperref[res:Proposition-10]{10}, Theorem \hyperref[res:Theorem-11]{11} and Remark 12
  \item \(\texttt{lean/Upstream/ItoCalculus.lean}\): Propositions \hyperref[res:Proposition-16]{16} and \hyperref[res:Proposition-26]{26} to 29
  \item \(\texttt{lean/Upstream/Predictability.lean}\): Propositions \hyperref[res:Proposition-26]{26}, \hyperref[res:Proposition-27]{27} and \hyperref[res:Proposition-29]{29}
  \item \(\texttt{lean/Upstream/LipschitzSDE.lean}\): Propositions \hyperref[res:Proposition-27]{27} and \hyperref[res:Proposition-29]{29}, and Proposition \hyperref[res:Proposition-46]{46} (Brownian premise only)
  \item \(\texttt{lean/Upstream/LipschitzMarkov.lean}\): Proposition \hyperref[res:Proposition-29]{29}, for its test-function predicate only
  \item \(\texttt{lean/Upstream/LipschitzMarkovTime.lean}\): Proposition \hyperref[res:Proposition-29]{29}
  \item \(\texttt{lean/Upstream/ExpPolyConsistency.lean}\): Proposition \hyperref[res:Proposition-46]{46}
  \item \(\texttt{lean/Upstream/ExponentialMartingale.lean}\): Proposition \hyperref[res:Proposition-27]{27}
  \item \(\texttt{lean/Upstream/IntegralApproximation.lean}\): Proposition \hyperref[res:Proposition-27]{27}
\end{itemize}

\emph{Conventions of the formal statements.} The dates are a function \(\tau:\mathbb N\to\mathbb R\) with \(\tau(0)=0\), strictly increasing on \(\{0,\dots,N\}\), with \(I_k=[\tau(k),\tau(k+1))\) and \(I_N=[\tau(N),\infty)\), so no extended real enters; a step or piecewise-affine function is any function with the prescribed values on each \(I_k\); \(\mathbb R^d\) is \(\texttt{EuclideanSpace}\ \mathbb R\ (\texttt{Fin}\ d)\); integrals are Mathlib's interval integrals against Lebesgue measure; conditional expectations are Mathlib's, which vanish on non-integrable functions, so a conditional-expectation identity forces integrability. From Section \hyperref[sec:3.25]{3.25} on, the time interval \(I_k\) is often represented by \((T_k,T_{k+1}]\), which changes no integral, and the last date by the horizon.

\emph{Theorem \hyperref[res:Theorem-1]{1} to Remark 12.} Lemmas \hyperref[res:Lemma-3]{3} to 7 are stated as in the text, with their parts bundled; the proofs of Lemmas \hyperref[res:Lemma-5]{5} and \hyperref[res:Lemma-6]{6} avoid the constants \(c_k\), \(A_k\), \(C_k\). Theorem \hyperref[res:Theorem-1]{1} restates the regularity conditions and Assumptions 2.1 and 2.2 as the fields of \(\texttt{Upstream.HJMScheduled}\), with an abstract sub-sigma-algebra in place of \(\mathcal F_{T_n-}\); a second declaration, \(\texttt{theorem1\_of\_HJMScheduled}\), applies it to an instance of the structure. That structure has a nonzero deterministic instance, \(\texttt{HJMScheduled.nonvacuityInstance}\), outside the step family. Corollary \hyperref[res:Corollary-2]{2} uses Fubini's theorem and the vanishing of the Itô integral of a null integrand and is not formalized. Proposition \hyperref[res:Proposition-8]{8} asserts the cycle (A) \(\Rightarrow\) (A\('\)) \(\Rightarrow\) (B) \(\Rightarrow\) (A) with (B) \(\Leftrightarrow\) (B\('\)); the computation on the parametrization (2.5) after it is not formalized. Proposition \hyperref[res:Proposition-9]{9} encodes the \([0,\infty]\)-valued conditional identities by integrals over the sets of \(\mathcal G\) and uses no regular conditional distributions; (3.28) is proved through a one-sided uniqueness theorem for moment generating functions, \(\texttt{MGFUniqueness.map\_eq\_of\_mgf\_eqOn\_Ico}\). Proposition \hyperref[res:Proposition-10]{10} verifies every regularity field and both assumptions for the product of Gaussian coordinate laws; the reconstruction of the curve on a Brownian extension at the end of Section \hyperref[sec:3.7]{3.7} is not formalized. Theorem \hyperref[res:Theorem-11]{11} formalizes the three equivalences, and a second declaration applies necessity to the standing-hypothesis structure; the coefficient calculation after it is not formalized. Remark 12 constructs the standing-hypothesis instance on the product Borel measure with the uncompleted coordinate filtration; completion, the path-space Markov assertion, the bank-account state assertion, the option factorization and the multivariate Gaussian packaging are proved in the text only.

\emph{Propositions \hyperref[res:Proposition-13]{13} to 15.} Proposition \hyperref[res:Proposition-13]{13} is formalized after assuming its variance-integral representation; the Brownian construction, the maturity limits, the conditional-variance derivation and the Gaussian law are not formalized. Proposition \hyperref[res:Proposition-14]{14} takes its stochastic ingredients (the centered representation, conditional isometry, cross-factor orthogonality and conditional first-moment identity) as explicit hypotheses and derives Mathlib's conditional variance formula and the deterministic parts; it constructs no Brownian motion or solution. Proposition \hyperref[res:Proposition-15]{15} is fully formalized on the finite product Gaussian space of Remark 12 with the actual bank account and bonds, with no pricing premise.

\emph{Proposition \hyperref[res:Proposition-16]{16}.} The first layer proves the distributional consequences from explicit premises: the transition law, product kernel, Laplace uniqueness, image support, affine hull, covariance, rank and atom. The second layer discharges those premises from \(\texttt{Upstream.ItoCalculus}\), for an adapted nonnegative solution of (3.62) with almost surely continuous paths, taken as a hypothesis; it derives the Itô representation, localization, conditional transforms, isometry and orthogonality, the centered representation and the regular conditional law, and identifies the meeting jumps with the actual short-rate jumps. Not covered: existence of the solutions, the second-moment bound of Section \hyperref[sec:3.12]{3.12}, and anything about calibration or listed options.

\emph{Propositions \hyperref[res:Proposition-17]{17}, \hyperref[res:Proposition-19]{19}, \hyperref[res:Proposition-20]{20} and \hyperref[res:Proposition-23]{23} to 25 and Corollaries \hyperref[res:Corollary-21]{21} and \hyperref[res:Corollary-22]{22}.} Each is formalized in its stated ideal model on the finite Gaussian product space with its completed reveal filtration, from the actual accrual, discounting and exercise definitions, with no pricing, measurability, stopping-time, optimizer or Gaussian-law premise; Lemma \hyperref[res:Lemma-18]{18} is part of the declaration of Proposition \hyperref[res:Proposition-19]{19}, and Proposition \hyperref[res:Proposition-23]{23} reuses those of Propositions \hyperref[res:Proposition-19]{19} and \hyperref[res:Proposition-20]{20} and Corollary \hyperref[res:Corollary-21]{21}. American values are suprema over all stopping times of the completed filtration valued in the exercise window. Corollary \hyperref[res:Corollary-22]{22} is stated with the number of earlier meetings as its index, and on the single adapted version (3.92). In Proposition \hyperref[res:Proposition-20]{20} the inclusion of the cash calls of Corollaries \hyperref[res:Corollary-21]{21} and \hyperref[res:Corollary-22]{22} in (c) and (d) is argued in the text only. Proposition \hyperref[res:Proposition-24]{24} checks the calendar by decision, as day numbers from 1 January 2026 with actual-over-360 accrual, and takes (A1) to (A4) as the calendar, tick sizes and reference sizes; the forward is taken as one in the exact evaluation. Proposition \hyperref[res:Proposition-25]{25} reuses Proposition \hyperref[res:Proposition-13]{13}'s matrix and Proposition \hyperref[res:Proposition-16]{16}'s closedness lemma; the sentence of (d) on stochastic volatility claims nothing and is not formalized. Listed settlement, margin, exercise into futures, unknown initial curves, noisy prices, calibration and arbitrage conclusions are outside every one of these statements.

\emph{Proposition \hyperref[res:Proposition-26]{26}.} Its inputs are \(\texttt{Upstream.ItoCalculus}\), \(\texttt{Upstream.Predictability}\), the driver covariations, nonnegative adapted solutions of (3.54) with almost surely continuous paths, initial state one, the supremum-square moment bound and admissibility of the original integrands. It derives the transforms, laws, conditional independence and the exact support, affine hull, mean, covariance, rank and vertex mass. It does not construct the solutions, prove rank three for the incomplete calibration, or cover time-dependent mean reversion.

\emph{Proposition \hyperref[res:Proposition-27]{27}.} Its inputs are the two calculus structures on one Brownian basis with their filtrations and covariations and the \(\texttt{ItoCalculus}\), \(\texttt{LipschitzSDE}\), \(\texttt{Predictability}\), \(\texttt{ExponentialMartingale}\) and \(\texttt{IntegralApproximation}\) interfaces. It derives the stopped strong solution and its uniqueness, the identification of the integrals between filtrations, the variance identities, the true martingale property of the discounted bonds and the nonzero short-rate jump. The \(\texttt{LipschitzSDE}\) instance is vacuous, and no instance constructs the two structures on a Brownian basis.

\emph{Proposition \hyperref[res:Proposition-28]{28}.} The drivers are all coordinates of an \(\texttt{ItoCalculus}\) instance, and the shapes are continuous on \(\mathbb R\). The statement proves the coefficients (3.125) on each interval, (2.3) on every path, (3.128) for each fixed maturity on one event for all times without a stochastic Fubini theorem, the start and freeze properties, (3.129) and the regularity of \(A_{j,k}\) and \(V_{j,m}\). The word "semimartingale" for \(M_{j,k}\) is paper-only: the local-martingale property of the integral of an integrand with finite path energy is not a field of the calculus.

\emph{Proposition \hyperref[res:Proposition-29]{29}.} Its inputs are \(\texttt{ItoCalculus}\) with the Brownian premise and solution class of \(\texttt{LipschitzSDE}\), \(\texttt{Predictability}\), and \(\texttt{LipschitzMarkovTime}\), whose field is the Markov property at deterministic times for one solution; the instances of \(\texttt{LipschitzSDE}\) and \(\texttt{LipschitzMarkovTime}\) are vacuous, with separate checks. The several-time property is proved for any adapted process with the one-time property, by products and a \(\pi\)-\(\lambda\) argument, since Mathlib's conditional-independence lemmas need a standard Borel space. The paragraph after the statement has no formal counterpart.

\emph{Propositions \hyperref[res:Proposition-30]{30} to 33.} Only what the proofs use of a realization is assumed: the state at one time in a set \(Z\subseteq\mathbb R^q\), not necessarily measurable, and the curve map at finitely many maturities, locally Lipschitz on \(Z\). The driver is a Brownian motion of Mathlib with its Gaussian law and independent increments, and the calculus enters through the standalone restatement of its fields. The dimension step uses the Hausdorff dimension of a Lipschitz image. Proposition \hyperref[res:Proposition-31]{31} is stated for any adapted scale state with continuous paths, with continuous-adapted predictability as a hypothesis, and (3.146) in the version form of (3.139); its constant-scale case is also formalized separately (\(\texttt{RecurrentLoadingAssembly}\)), with the characterization for \(\lambda\equiv1\) stated in Proposition \hyperref[res:Proposition-33]{33}(c) (\(\texttt{RecurrentLoadingCharacterization}\)). Lemma \hyperref[res:Lemma-32]{32} is proved for bounded Borel integrands continuous almost everywhere. In Proposition \hyperref[res:Proposition-33]{33} the reduction uses the Krylov matrix and the companion matrix of the characteristic polynomial, and the sufficiency half of (c) is stated in the one-time sense of the rank bounds.

\emph{Propositions \hyperref[res:Proposition-34]{34}, \hyperref[res:Proposition-35]{35}, \hyperref[res:Proposition-41]{41}, \hyperref[res:Proposition-43]{43} and \hyperref[res:Proposition-44]{44}.} Consistency is in the version form: a random element with measurable evaluations agreeing with the curve almost surely at every maturity. The proofs are pathwise in the drift, through the linear independence of polynomials times exponentials with distinct exponents and analytic continuation across the meetings. Propositions \hyperref[res:Proposition-35]{35} and \hyperref[res:Proposition-41]{41} take controllability, \(c\neq0\) and invertibility of \(A\) as standing hypotheses; the reduction to that case is not formalized. Propositions \hyperref[res:Proposition-43]{43} and \hyperref[res:Proposition-44]{44} are stated on one path, with Assumption 2.1 for almost every time and every maturity; Proposition \hyperref[res:Proposition-44]{44}(b) is proved with the Cayley--Hamilton theorem and a Bézout identity. Not formalized, as remarks: that the enlarged family of Proposition \hyperref[res:Proposition-34]{34}(c)(ii) does not depend on \(t\), the sentences on what an enlargement absorbing the jump must contain, and, for Proposition \hyperref[res:Proposition-35]{35}, the open set of times.

\emph{Lemma \hyperref[res:Lemma-36]{36} and Propositions \hyperref[res:Proposition-37]{37} to 40 and 45 to 47.} Lemma \hyperref[res:Lemma-36]{36} is the source's proposition proved as a lemma of this paper, with no hypothesis structure. Propositions \hyperref[res:Proposition-37]{37}, \hyperref[res:Proposition-38]{38} and \hyperref[res:Proposition-46]{46} take Assumption 2.1 for the drift that Itô's formula gives for the representation (3.161). Proposition \hyperref[res:Proposition-38]{38} uses the lemma that the primitive of a real-analytic function is real-analytic, and Proposition \hyperref[res:Proposition-39]{39} the lemma that nonnegative polynomials in one variable are sums of squares, both proved here since Mathlib lacks them. Propositions \hyperref[res:Proposition-39]{39}, \hyperref[res:Proposition-40]{40} and \hyperref[res:Proposition-45]{45} are statements at a parameter point, with the pathwise parts of 39(e) and 40(a) for paths right-continuous at the point. Proposition \hyperref[res:Proposition-46]{46}(b) restates \(\texttt{Upstream.ExpPolyConsistency}\), whose fields are the cited conclusions of \cite{filipovic2000exponential} Theorem 3.2 and Corollaries 3.4 and 3.5, under the Brownian premise, an adapted driver-form Itô process with \(\mathcal F_0\)-measurable initial value, boundedness of the family and the consistency equation almost everywhere; its instance is vacuous, with separate checks. Proposition \hyperref[res:Proposition-47]{47} is stated on one path.

\emph{Proposition \hyperref[res:Proposition-48]{48}.} The formal statement takes a regular conditional distribution of \(X\) given \(\mathcal G\) as a hypothesis with its defining property, states (H) through set integrals over \(\mathcal G\), so that no integrability is assumed, and allows \(L=\infty\). It proves (a) with one \(\mathcal G\)-measurable null set, (b) with the derivatives of Mathlib's cumulant generating function at \(-\tau\), the uniqueness of the law through the moment-generating-function uniqueness lemma used for Proposition \hyperref[res:Proposition-9]{9}, and existence with joint measurability through parametric kernel integrals; (c)(i) and (c)(ii) with (3.184) and (3.185); (d) with its equivalence, the agreement of the conditional laws, the case of a right-continuous curve and the bond formula, and the general-curve form with its equivalence; and (e). The numerical values of the gap in the discussion after the proposition are not formalized; the bound (3.185) is.

\emph{Proposition \hyperref[res:Proposition-49]{49}.} The probabilistic parts of (a) and (b) are proved under an explicit hypothesis, not a field of the calculus structures: the Wiener integrals of deterministic integrands are linear, jointly Gaussian with the stated covariances and independent of the jumps, \(Y\) is continuous and satisfies the stochastic Fubini identity, and the filtration has the stated measurability and independence properties. The formal model allows any bounded measurable \(\sigma^2\) and initial curve. The hypothesis is shown satisfiable, with nondegenerate Gaussian jumps, only by the model of Section \hyperref[sec:3.10]{3.10}, which has no diffusion; no Brownian motion is constructed, so with \(\sigma\neq0\) the parts (a) and (b) are conditional. The deterministic identities of (b), including (3.189), the rank criterion (3.190) with the inversion formulas and the strictly positive pairs, the calendar facts of (d) with both non-identifying pairs and the rank of the stacked matrices, and the error formula (3.191) with the sharp bound (3.192) are proved without conditions. The table of practically identified meetings, the linear-programming comparison and the size \(1.9\times10^{-9}\) of the \(z\)-difference are not formalized.

\emph{Proposition \hyperref[res:Proposition-50]{50}.} As for Proposition \hyperref[res:Proposition-49]{49}, the probabilistic parts of (a) and (b) are proved under an explicit hypothesis, the Gaussian law of Proposition \hyperref[res:Proposition-49]{49} with the process \(\int_0^t\sigma(s)\phi(s,x)\,dW_s\) in place of \(Y\): it is the Wiener integral of \(\mathbf 1_{[0,t]}\phi(\cdot,x)\), integrable along the diagonal and in \(x\), and satisfies stochastic Fubini in the two forms the proof uses. It is shown satisfiable, for every bounded Borel shape, only by the model of Section \hyperref[sec:3.10]{3.10} without diffusion; no Brownian motion is constructed. The formal model allows any bounded measurable \(\sigma^2\) and initial curve and any bounded Borel shape. The deterministic identities of (a) and (b), including (3.195), the criterion (3.196) with its consequences, and (d1) to (d5) are proved without conditions: (d2) for every bounded Borel shape, (d3) in exact arithmetic with both pairs, which prove non-identification without stating the ranks, (d4) with the uniqueness of the zero of (3.197) and its bracket through exact Taylor bounds on the exponential, and (d5) with its gap counts, the span and rank for every diffusion column with \(D_{1,1}\neq0\), the identified meetings, and, in day numbers, the weekdays and the unchanged gap assignments under effective dates. That the business day after each decision is not a holiday is a check, since the holiday calendar is not modelled. The condition for \(\tilde z/\delta=\int_0^SQ\), the decimal value of \(\kappa^*\), the table, the standard-deviation errors, the linear-programming comparison and the condition numbers are not formalized.

\emph{Proposition \hyperref[res:Proposition-51]{51}.} The deterministic estimates are kept apart from the probability and proved without conditions: part (a)'s bounds on \(\sigma-\tilde\sigma\) and \(\alpha-\tilde\alpha\), the bounds on \(m\) and \(\int(\sigma-\tilde\sigma)^2\), the bounds on the mean and variance of the log-prices, the Gaussian fourth moment, which is derived from the moment-generating function, and the price inequality for jointly Gaussian log-prices. Part (a)'s mean-square identity and (3.200) and (3.201) assume a standard driver of the supplied calculus, for any bounded Borel initial curve. Parts (b) and (c) assume in addition that the driver is a Brownian motion, the hypothesis of Lemma \hyperref[res:Lemma-32]{32} that gives the Gaussian laws; no model with such a driver is constructed. The version of (b) is constructed, with \([0,t]\) split by the number of meetings in \((s,t]\), and stochastic Fubini for that split kernel is proved from finite linearity, so (3.202) holds with the stated constants. Part (c) identifies the approximating model with the realization of Proposition \hyperref[res:Proposition-31]{31} at \(\psi\equiv1\), whose volatility and drift are \(\tilde\sigma\) and \(\tilde\alpha\) by definition, and proves by Tonelli's theorem that the realized bond price equals \(\tilde P\) almost surely; the state dimension is that of Proposition \hyperref[res:Proposition-31]{31}. The exact mean-square error in the example of (a) is formalized. Not formalized: the characterization of bounded recurrences, the growth of the constants for growing recurrences, the remark on shapes that are not quasi-exponential, and the numerical check.

\emph{Proposition \hyperref[res:Proposition-52]{52}.} The result is pathwise algebra together with the passage from Assumption 2.1, which enters as an explicit hypothesis in its integrated form; no hypothesis structure is used. Because the block is linear in its state, the drift and volatility of the representation are written out, as the product rule gives them, rather than derived from the calculus structure, and the existence of \(Z\), \(L_m\) and \(C_m\) is not constructed. Formalized: the identity (3.206), the general jump and (3.205); (a) with the level jump, through Lemma \hyperref[res:Lemma-42]{42} and the descending-degree argument; (b)(iv) at a point and, through the rational maturities, continuity and Fubini, with the proposition's quantifiers; the converse back to Assumption 2.1 with explicit front-end drifts, whose measurability as processes is not addressed and is not needed by the formal assumption; (3.207) for \(\mathcal A\) diagonalizable over \(\mathbb C\), through a proved linear independence of exponential polynomials, under every resonance; the counterexample (d); the aggregation example and the reduction given a minimal realization; and in (e), the sharpening of Proposition \hyperref[res:Proposition-44]{44}, \(K=0\) in the uncorrelated splice and the entries of \(a_\sharp\). Not formalized: the split in (c) of "\(R_\sharp\) affine" into the two conditions; the existence of the minimal realization and the block-by-block form of Proposition \hyperref[res:Proposition-44]{44}(b) and (c); the carry-over of Proposition \hyperref[res:Proposition-45]{45} in (e) beyond the entries of \(a_\sharp\), including the drift shifts; the remark on non-diagonalizable \(\mathcal A\); and the symbolic check.

\emph{Proposition \hyperref[res:Proposition-53]{53}.} As for Proposition \hyperref[res:Proposition-52]{52}, Assumption 2.1 is an explicit hypothesis and no hypothesis structure is used. The formal statements of Proposition \hyperref[res:Proposition-52]{52}, including the passage from Assumption 2.1 and the converse back to it, put no bound, measurability or integrability hypothesis on the coefficients at any \((u,\omega)\), so (a) and (b) hold verbatim for the class (L), with the same quantifiers; the class appears in Lean only in (c). There the drifts of the converse are written out, with the front end's primitive proved by induction over the meetings, Assumption 2.1 holds with them, and they are integrable on \([0,H]\) on every path whose noises \(V_m\) and block rows are square integrable and whose block drifts and state are integrable. Maturity-uniform envelopes for \(\sigma\) and \(\alpha\), square integrable and integrable on \([0,H]\), give the integrability part of the regularity of Section \hyperref[sec:2]{2}. The product counterexample, the stability under a bounded factor, the integrability of products of square-integrable rows and the shift invariance with the frozen curve are formalized. Not formalized: the construction of \(L_m\) and \(C_m\) as stochastic integrals in (c), whose converse is proved as Assumption 2.1 at each \((u,\omega)\) with the explicit drifts; progressive measurability in \((t,T)\); the localizing times; the Heston-type remark; and the scripted checks.

\bibliographystyle{plainnat}
\bibliography{refs}
\end{document}
